\documentclass[11pt,a4paper]{article}
\usepackage[T1]{fontenc}
\usepackage[utf8]{inputenc}
\usepackage{lmodern}
\DeclareUnicodeCharacter{0393}{\ensuremath{\Gamma}}
\pdfmapfile{+lm.map}
\pdfmapfile{+cm.map}
\pdfmapfile{+symbols.map}
\pdfmapfile{+euler.map}
\pdfmapfile{+rsfs.map}
\usepackage[margin=27mm]{geometry}
\usepackage{amsmath,amssymb,amsthm,mathrsfs}
\usepackage{booktabs,longtable,array}
\usepackage{graphicx}
\usepackage{hyperref}
\usepackage{xurl}
\usepackage{microtype}
\hypersetup{colorlinks=true,linkcolor=blue,urlcolor=blue}
\setlength{\emergencystretch}{3em}
\providecommand{\tightlist}{\setlength{\itemsep}{0pt}\setlength{\parskip}{0pt}}
\providecommand{\pandocbounded}[1]{#1}
\newtheorem{theorem}{Theorem}[section]
\newtheorem{proposition}[theorem]{Proposition}
\newtheorem{lemma}[theorem]{Lemma}
\newtheorem{corollary}[theorem]{Corollary}
\theoremstyle{definition}
\newtheorem{definition}[theorem]{Definition}
\newtheorem{remark}[theorem]{Remark}
\title{DEG\_L: Degrees, theta norms, and semiregularity\\
in the 032 construction}
\author{DEG\_L research run\\GPT-6 Astra Pro}
\date{21 November 2026}
\begin{document}
\maketitle
\begin{abstract}
This working paper separates the effective input surface, its nonzero mixed
period, and the signed codimension-two output of the 032 construction.
It proves exact Chow cancellations, integral norm-and-trace division,
fixed-field Weil projectors, and quantitative extraction of any fixed Weil
class from an effective witness. The theta spanning theorem supplies an
exact finite character-theta formula for the energy, and an explicit
surface-degree inequality. The exact quartic spectral formula is a finite rank-two standard-$L(1)$
norm sum, with source representations, multiplicity copies and measures
retained. Sharp local estimates distinguish Schwartz support from level.
The actual CM Hom covolume is computed in a fixed simple CM isogeny
class, giving a conductor-dependent bound and a sharp limitation for
fixed-reference witness transport even after passing to reduced images.
The original construction does not specify a smooth embedding; universal
semiregularity of its deeper-level morphism presentations is false. The
remaining deformation problem is reduced to a gluable weakly equivariant
replacement object with its full Hodge-compatible character. An Adams
operation on its finite twisting gerbe then constructs a signed-cycle
bound even in virtual rank zero. A proper Hilbert parameter-space
argument carries a uniform effective witness bound on dense fibres to
bounded division on the entire family. The general uniform witness bound
in the open Weil families is not asserted. This is a live checkpoint
manuscript; the session has not ended.
\end{abstract}
\tableofcontents
\clearpage
\section{Scope, assumptions, and check protocol}
The prompt and the attached bridge sketch determine the questions. The
sketch is treated as a map of questions, not as a source of proved facts.
As requested, the attached 032 CM Hodge theorem and the stated Alp\"oge
Dirichlet GRH theorem are admitted inputs. The full Lean dependency archive
was not supplied or compiled. General Hecke GRH is a conditional input
attributed by the prompt to Masullo, \S1.2; the attached index does not
contain a complete numbered Hecke-GRH statement. No trace-formula number
is silently substituted for it.

Every finding has a stable run identifier: T6.x, Bx, S-xx, DIVx, or DERx within
DEG\_L\_20261121. Findings are formulations proved in this run, without a
claim that all elementary ingredients are new to the literature. A
conditional implication is proved only with its displayed hypotheses.
The detailed event ledger distinguishes whole-proof reading, targeted
checks, background use, and computations. Exact reasoning effort is not
exposed by the execution environment and is not invented.

The cycle equality $Nt=[Z^+]-[Z^-]$ means equality in cohomology under the
cycle-class map, unless a Chow-group equality is explicitly written. The
word degree always includes the chosen polarization and distinguishes a
pushforward cycle with multiplicity from its reduced image.

\part{Checked findings and explicitly conditional reductions}

\section{The surface degree and the Rallis comparison}

% DEG_L T6 proof module. Generated from rallis_t6.md after source and proof checks.
% Stable theorem identifiers are \label{DEG-L.T6.1}, ..., \label{DEG-L.T6.22}.
% Requires amsmath, amssymb, hyperref; this is an \input-able module, not a document.
% Imported background is identified in the text. No general RH audit is claimed.
\providecommand{\tightlist}{\setlength{\itemsep}{0pt}\setlength{\parskip}{0pt}}
\providecommand{\passthrough}[1]{#1}
\section{T6 research report --- DEG\_L / rallis\_t6}\label{t6-research-report-deg_l-rallis_t6}

Date: 2026-11-21. Actor: GPT-6 Astra Pro, effort not exposed, child run \texttt{/root/rallis\_t6}. This report records checked reductions, not a claimed proof of T6 as originally phrased. The later T6.9 strengthens the early analytic gate: all needed one-form theta expansions do exist by the supplied human proof and checked literature; their coefficients and integral normalization remain to be bounded.

\subsection{1. What the attached 032 construction actually supplies}\label{what-the-attached-032-construction-actually-supplies}

Source: \texttt{OPENAI\_032\_The-rational-Hodge-conjecture-for-CM-abelian-varieties.pdf}, Proposition 3.1, printed pp.~11--12, equations (3.1)--(3.5); Corollary 3.2, p.~12; Proposition 4.6, p.~23; Corollary 4.7, pp.~24--25, equations (4.18)--(4.19); Corollary 6.7, p.~45; proof of Theorem 2.5, p.~46, equation (7.1).

Write \(A=\prod_{i=1}^4B_i\), with each \(B_i=B(v_i)\) of dimension \(s\), and \(n=4s\). In the nontrivial surface construction \(s\ge2\), so \(n\ge8\). The surface \(S\) is a smooth projective ball quotient. The paper obtains rational Hodge maps \(h_i:H^1(B_i,\mathbf Q)\to H^1(S,\mathbf Q)\) and integers \(m_i>0\) so that \(f_i:S\to B_i\) has \(f_i^*=m_i h_i\). It defines \[
f=(f_i),\quad z=f_*[S]\in CH^{n-2}(A),\quad
k=\int_S f_1^*\alpha_1\,f_2^*\alpha_2\,f_3^*\alpha_3\,f_4^*\alpha_4\ne0.
\] The generators \(\alpha_1,\alpha_2\) are holomorphic and \(\alpha_3,\alpha_4\) antiholomorphic; their Betti coefficients are algebraic. With copied factors \(B_i'\) and divisor kernels \(D_i\) having pure class \(\alpha_i\otimes\beta_i\), the correspondence \(K=\prod_i q_i^*D_i\) gives \[
\operatorname{cl}(K_*z)=k\,t_\beta,
\qquad t_\beta=\beta_1\otimes\beta_2\otimes\beta_3\otimes\beta_4.
\] Thus three different quantities must be distinguished: the degree of the effective pushforward surface \(z\); a mixed period \(k\) detecting a switch line; and the codimension-two output \(K_*z/k\) with algebraic coefficients. They have different degrees, codimensions, and positivity properties. The paper itself explicitly says that \(z\) is a functional on the line, not a vector on it. It only proves \(\dim f(S)=2\), not that \(f\) embeds \(S\).

The theta construction first supplies nonzero total two-forms and their mixed pairing. Section 7 then expands finite wedges, chooses a nonzero summand, expands each theta class in rational CM-map images, and chooses another nonzero summand. The final \(k\) is not stated to be the total theta norm. No Rallis formula is stated in the attached 032 paper.

\subsection{2. Checked findings}\label{checked-findings}

\subsubsection{DEG-L.T6.1 --- Exact geometric degree and its positive Gram expression}\label{DEG-L.T6.1}

Let \(H\) be an ample line bundle on a complex abelian variety \(A\) of dimension \(n\), and let \(f:S\to A\) be a morphism from a smooth projective surface with two-dimensional image. Put \(z=f_*[S]\). Choose translation-invariant holomorphic one-forms \(\omega_1,\ldots,\omega_n\) such that the invariant positive representative of \(c_1(H)\) is \[
\Omega_H={i\over2}\sum_{a=1}^n\omega_a\wedge\overline{\omega_a}.
\] For \(\eta_a=f^*\omega_a\) and \(\xi_{ab}=\eta_a\wedge\eta_b\), one has \[
\deg_H(z)=\int_S(f^*\Omega_H)^2
={1\over2}\sum_{a<b}\int_S\xi_{ab}\wedge\overline{\xi_{ab}}.
\tag{G1}
\] If \(e=[\mathbf C(S):\mathbf C(f(S))]\), then \(z=e[f(S)]\) and \(\deg_H(f(S))=\deg_H(z)/e\).

\textbf{Proof.} The projection formula gives the first equality. A positive definite Hermitian form can be diagonalized, giving the displayed expression for \(\Omega_H\). Expanding its square and moving each antiholomorphic factor past the next holomorphic factor yields \[
(f^*\Omega_H)^2={1\over4}\sum_{a,b}\eta_a\wedge\eta_b\wedge\overline{\eta_a}\wedge\overline{\eta_b}.
\] The terms with \(a=b\) vanish; the terms with \((a,b)\) and \((b,a)\) agree. This proves (G1). In local holomorphic coordinates, if \(\xi=u\,dz_1\wedge dz_2\), then \(\xi\wedge\bar\xi=4|u|^2dx_1dy_1dx_2dy_2\), so every summand is nonnegative. The assertion about \(e\) is the definition of proper pushforward of an integral cycle. QED.

\textbf{Check.} Read the entire proof, independently recomputed both exterior-algebra signs and the factor \(1/2\), and compared \(A=S\) of dimension two with the standard Euclidean Kähler form. In that test, \(\Omega_H^2=2\,dV\) and \(\xi_{12}\bar\xi_{12}=4\,dV\).

\textbf{Meaning.} A degree is a positive sum of norms of all polarization-coordinate wedges. It is not the selected mixed period of four CM lines. Even when Rallis computes certain wedge norms, it takes additional data to identify every wedge in (G1), its automorphic spectral projections, and the conversion from invariant forms to algebraic Betti generators.

\subsubsection{DEG-L.T6.2 --- Zero degree of the pure switch class}\label{DEG-L.T6.2}

Give \(A=\prod_{i=1}^4B_i\) the product polarization \(H=\bigotimes_i\operatorname{pr}_i^*H_i\), where each \(H_i\) is ample. Every class \[
w\in H^1(B_1)\otimes H^1(B_2)\otimes H^1(B_3)\otimes H^1(B_4)\subset H^4(A)
\] has \[
\int_A w\,c_1(H)^{n-2}=0.
\tag{G2}
\] Consequently a nonzero rational switch class cannot be the cohomology class of an effective codimension-two cycle for this polarization. A nonzero pure tensor \(t_\beta\) also satisfies \(t_\beta^2=0\).

\textbf{Proof.} Expand \(c_1(H)^{n-2}\) into products \(\prod_i\operatorname{pr}_i^*c_1(H_i)^{a_i}\). In factor \(i\), the product with \(w\) has degree \(1+2a_i\), which is odd. Integration on \(B_i\), of even real dimension \(2\dim B_i\), therefore vanishes. Every term vanishes. Every nonzero effective cycle has strictly positive degree against an ample polarization, so (G2) rules out effectivity. For the last assertion, squaring a pure tensor repeats the same degree-one class in every factor, and the square of a degree-one class vanishes in characteristic zero. QED.

\textbf{Check.} Entire proof reread; parity argument checked separately for unequal factor dimensions and arbitrary sums of pure tensors. It does not say that a rational combination of conjugate switch lines has zero polarized self-pairing; that would be false. The statement about \(t_\beta^2\) concerns one pure algebraic-coefficient tensor only.

\textbf{Meaning.} The output of 032, Proposition 3.1, has degree zero for a natural ample polarization. Therefore an equality between its degree and a nonzero positive theta norm cannot hold. The complexity \(D\) in Route 2 is the separate degree of positive and negative effective representatives; (G2) does not bound it. For the input surface \(z\), the ordinary square is already zero by codimension when \(n\ge8\), since \(2(n-2)>n\). A numerical self-intersection would require a different pairing and an explicit dualizing correspondence.

\subsubsection{DEG-L.T6.3 --- Level alone does not determine the raw surface degree}\label{DEG-L.T6.3}

Fix every datum through the rational maps \(h_i\) and the source surface \(S\) in 032, Proposition 3.1. Replacing all clearing integers \(m_i\) by \(Mm_i\) and \(f_i\) by \([M]\circ f_i\) is still an allowed construction on exactly the same source level and with the same initial theta inputs. If \(H\) is symmetric, then \[
\deg_H(z_M)=M^4\deg_H(z),\qquad k_M=M^4k.
\tag{G3}
\] Thus there is no degree of the raw pushforward surface determined by the level alone under the choices specified in 032.

\textbf{Proof.} Multiplication by \(M\) acts as \(M\) on \(H^1(A)\) and as \(M^2\) on \(H^2(A)\), so \([M]^*c_1(H)=M^2c_1(H)\). Applying the projection formula twice gives \(\deg_H([M]_*z)=\int_S f^*[M]^*c_1(H)^2=M^4\deg_H(z)\). Each of the four factors in the period acquires \(M\). Positivity in T6.1 and the nonzero period guarantee \(\deg_H(z)>0\). QED.

\textbf{Check.} Entire proof reread. The assertion concerns \(f_*[S]\), the cycle defined by 032, not the reduced image \([f(S)]\); the generic mapping degree can change. Even a symmetric assumption is unnecessary for the cohomological degree identity, but retained for a standard polarization normalization.

\subsubsection{DEG-L.T6.4 --- Auxiliary-cover division cancels in Chow}\label{DEG-L.T6.4}

Keep the target \(A\), kernels \(K\), and algebraic generators fixed. For a finite surjective morphism \(\pi:S'\to S\) of degree \(d\) between smooth projective surfaces, use \(f'=f\pi\) and the pulled one-forms. Then \[
z'=d z,\quad k'=d k,\quad {K_*z'\over k'}={K_*z\over k}
\quad\hbox{in }CH^2(A')\otimes\overline{\mathbf Q}.
\tag{G4}
\] In particular this applies to every deeper principal cover used in 032. Every fixed rational coefficient-descent combination gives one identical rational Chow class throughout the tower. A single rational representative therefore gives bounded \(N,D\) throughout this auxiliary tower.

\textbf{Proof.} Proper pushforward gives \(\pi_*[S']=d[S]\), and functoriality gives \(z'=f_*\pi_*[S']=dz\). The integration formula for a finite map gives \(k'=dk\); alternatively use its projection formula with the degree-four class on \(S\). The correspondence action is linear, so the quotient is identical before taking any cycle-class map. Coefficient trace and rational linear combination commute with this identity. Once such a fixed rational cycle is written as a difference of effective cycles after clearing denominators, the same \(N,D\) work at every cover. QED.

\textbf{Check.} Entire proof reread; arbitrary finite surjective maps work, although the 032 covers are finite étale. This proves division for auxiliary covers at a fixed abelian target. It does not transport a cycle to new targets in a dense Hecke orbit.

\subsubsection{DEG-L.T6.5 --- Clearing-map division also cancels with biextension kernels}\label{DEG-L.T6.5}

The divisor kernels in Proposition 3.1 may be chosen as algebraic linear combinations of birigidified Poincaré-type line bundles, without changing their Betti classes. For these choices, write \(u=\prod_i[M_i]\) on \(A\), with \(M_i>0\). Then \[
(u\times1)^*K=\Bigl(\prod_i M_i\Bigr)K
\quad\hbox{in }CH^4(A\times A')\otimes\overline{\mathbf Q}.
\] Replacing \(f_i\) by \([M_i]f_i\) leaves \(K_*z/k\) unchanged in Chow.

\textbf{Proof.} For a divisor class \(D\) on \(B\times B'\), replace any representative line bundle \(L\) by its birigidified part \[
L^{\mathrm b}=L\otimes\operatorname{pr}_1^*(L|_{B\times0})^{-1}
\otimes\operatorname{pr}_2^*(L|_{0\times B'})^{-1},
\] with the evident normalization of the fiber at \((0,0)\). Extend this operation linearly to the algebraic-coefficient divisor kernel. Since its prescribed Betti class is purely in \(H^1(B)\otimes H^1(B')\), this operation leaves that class unchanged. A birigidified line bundle corresponds to a morphism \(B'\to\widehat B\) sending zero to zero, hence to a homomorphism, and is the pullback of the Poincaré bundle. Its biextension property gives \(([M]\times1)^*D=M D\) in \(CH^1\). Multiplying the four kernel identities yields the displayed formula.

The projection formula for \(u\times1\) then gives \[
K_*u_*z=p'_*\bigl(p^*z\cdot(u\times1)^*K\bigr)
=\Bigl(\prod_iM_i\Bigr)K_*z.
\] The period is multiplied by exactly the same product, because the four pullbacks on \(H^1\) multiply by \(M_i\). Division cancels in Chow. QED.

\textbf{Check.} Entire proof reread. The factor is \(\prod_iM_i\), not \(\prod_iM_i^2\). For common \(M\) this is \(M^4\), consistent with T6.3. The permitted change of representatives may alter the original unnormalized Chow class, but preserves its required Betti class and gives a valid 032 construction. The claim is not that an arbitrary representative of a pure cohomology divisor automatically has the exact biextension scaling in Chow.

\textbf{Grothendieck locator.} SGA 7 I, Exposé VII, §2.4, printed pp.164--165, states additivity in the three arguments; equation (2.4.2) specifically states additivity in the third argument. Example 2.9.5, printed p.175 (PDF p.180; internal VII43), identifies birigidified line bundles on a product of abelian schemes with \(\mathbf G_m\)-biextensions. Remark 2.9.6, printed p.176 (PDF p.181; internal VII44), gives \(\operatorname{Corr}(P,Q)\simeq\operatorname{Hom}(P,Q^\vee)\simeq\operatorname{Hom}(Q,P^\vee)\). Root visually verified these original scans in this run; this agent checked the transmitted locators against the claims used here and does not claim a separate scan inspection. Primary scan: \url{https://library.slmath.org/nonmsri/sga/sga/pdf/sga7-1.pdf}.

\subsubsection{DEG-L.T6.6 --- A mixed period gives a lower, not upper, surface-degree estimate}\label{DEG-L.T6.6}

Use a product polarization on \(A\). Let the four chosen forms in the 032 period be supported on their four distinct factors, and let \(r_i\) be their norms in the invariant Hermitian cotangent metric for \(H\) (use \(\bar\alpha_i\) for the last two holomorphic cotangent vectors). Then \[
|k|\le (r_1r_2r_3r_4)\deg_H(z).
\tag{G5}
\] There is no upper degree estimate from this single linear period over the cone of effective surface cycles.

\textbf{Proof.} The four unit forms \(\alpha_1/r_1,\alpha_2/r_2,\bar\alpha_3/r_3,\bar\alpha_4/r_4\) are orthonormal because they lie on different polarized factors. Extend them to an orthonormal cotangent basis in T6.1. Cauchy--Schwarz and \(2ab\le a^2+b^2\) give \[
{|k|\over\prod_i r_i}\le {1\over2}(\|\xi_{12}\|^2+\|\xi_{34}\|^2)
\le\deg_H(z).
\] For the final assertion, take a nonzero effective complete-intersection surface \(Y\) with class a positive multiple of \(c_1(H)^{n-2}\). Its pairing with the switch tensor is zero by T6.2, while its degree is positive. The effective cycles \(z+mY\) all give the same mixed period and have degrees tending to infinity. QED.

\textbf{Check.} Entire proof reread; this is a statement about all effective cycles, not necessarily individual connected 032 ball-quotient images. The common-source scaling counterexample T6.3 addresses the actual construction's nonuniqueness separately. The lower bound depends on metric norms/CM periods of the generators; an L-value estimate for \(k\) alone does not eliminate them.

\subsection{3. What Rallis can actually compute, and the remaining explicit gate}\label{what-rallis-can-actually-compute-and-the-remaining-explicit-gate}

Primary source: W. T. Gan, Y. Qiu, S. Takeda, \emph{The regularized Siegel--Weil formula (the second term identity) and the Rallis inner product formula}, arXiv:1207.4709v3, §1.1, §11.6, §11.7 Theorem 11.4, printed pp.~54--55. URL: \url{https://arxiv.org/pdf/1207.4709}.

The theorem is about theta \textbf{integrals of specified automorphic representations}, not a kernel evaluated at an unspecified source point. For a unitary source of dimension \(n_0\) and target of dimension \(m\), its normalization uses \(s_{m,n_0}=(m-n_0)/2\) and the standard value at \(s_{m,n_0}+1/2\), retaining the splitting-character twist and normalized doubling integrals. In the rank-two source / rank-three target case it gives \[
\langle\theta(\phi_1,f_1),\theta(\phi_2,f_2)\rangle
=2L(1,\pi\times\chi_V)\,Z^*(1/2,\phi_1\otimes\bar\phi_2,f_1,f_2),
\tag{R1}
\] when the theorem's hypotheses and conventions hold. Here the target is globally anisotropic in 032, hence has Witt index zero, so the lower-tower vanishing condition is empty. The scalar automorphic inner product is converted exactly to the differential-form inner product in T6.10 below, including component/level volume.

For the rank-one source / rank-three target, the formal point is \(3/2\) in unitary normalization; it lies outside the displayed dimensional range of Theorem 11.4. A convergent-range theorem must be used instead. The bounds agent located, and this agent independently opened and checked, J. Haan, \emph{The Bessel period of U(3) and U(2) involving a non-tempered representation}, arXiv:1403.5061v7, §3.3, Theorem 3.4, printed p.~8. In its conventions the character-lift norm is \[
\langle\theta(\bar f_1,\phi_1),\theta(\bar f_2,\phi_2)\rangle
={L_E(3/2,\mathrm{BC}(\tau)\otimes\gamma^3)\over L_F(3,\epsilon_{E/F})}
\prod_v Z_v^\#(1,f_{1,v},f_{2,v},\Phi_{1,v}).
\] Its §3.6.1, p.~9, states absolute convergence and nonvanishing of the numerator Euler product at \(3/2\). Thus this particular global factor does not require GRH. The version is identified in the PDF itself; use the unversioned URL \url{https://arxiv.org/pdf/1403.5061}. Do not cite GQT §11.7, Theorem 11.4 for the rank-one/rank-three case.

The total theta tensor in 032, Proposition 4.6, uses a Hermitian plane \(W\). Its dual partner is \(U(W)\), not \(U(1)\). Spectral projection can therefore require a representation of a unitary group of rank two. Its degree-two standard \(L\)-function is not automatically a Hecke-character \(L\)-function. No claim is made here that a particular nonabelian constituent must occur in the paper's chosen tensor: rather, character-only decomposition is a specific additional theorem to be proved, not a consequence of the mere existence of the theta kernel. All local zeta factors, level volume, input scalars and CM-period comparisons remain part of the equality.

The further attacks below complete the existential analytic comparison: T6.9 supplies every one-form theta expansion, T6.10 fixes the scalar/geometric normalization, T6.11 gives an exact finite Hecke-\(L(3/2)\) energy formula and degree upper bound, and T6.12--T6.13 give an exact finite standard-\(L(1)\) degree expression with all comparison/input data retained. What is still not supplied is an effective bound for those actual integral realization/input data on a varying target orbit. T6.7--T6.8 separately control period division and signed correspondence degrees once the geometric and kernel data are bounded.

The exact-cancellation findings T6.4 and T6.5 remove two avoidable sources of apparent denominator/level growth \textbf{before} pursuing that harder gate. They leave the meaningful quantities: primitive integral CM maps, their polarizations and periods, the genuinely new choices of source/correspondence, and their behavior when the abelian target varies on a Hecke orbit.

\subsection{4. Approaches tried and their disposition}\label{approaches-tried-and-their-disposition}

\begin{enumerate}
\def\labelenumi{\arabic{enumi}.}
\tightlist
\item
  Direct identification of the final 032 class with a positive Rallis norm: refuted for product polarization by T6.2.
\item
  Identifying the selected four-form period with geometric degree: refuted in that form; exact replacement is T6.1, and T6.6 shows the selected period is only a lower-bound witness.
\item
  Defining raw degree as a function only of principal level: refuted without further normalization by T6.3.
\item
  Treating source-cover degree growth as a genuine obstruction to division: removed by the exact Chow cancellation T6.4.
\item
  Treating enlarged denominator-clearing maps as a genuine obstruction: removed after a legitimate biextension choice by T6.5.
\item
  Applying a rank-one Rallis formula directly to a rank-two wedge tensor: not justified. T6.12--T6.13 instead prove a finite cuspidal rank-two spectral expansion and standard-\(L(1)\) identity. No character-only or single-value equality has been proved.
\item
  Bounding the surface degree using only one-form norms: positive alternative sent to the bounds agent. Subsequently strengthened by T6.9 below: existence of every needed theta expansion is established from the human proof, and T6.10 fixes the scalar/geometric measure conversion. Effective coefficient/integral and test-input bounds remain. Hodge index gives \(\deg_H(z)\le E_f^2/\operatorname{vol}(S)\), where \(E_f=\int_Sf^*\Omega_H\wedge\omega_S\). This could use rank-one theta decompositions, but the rational CM extraction does not by itself identify all polarization-coordinate pullbacks as theta lifts with controlled coefficients.
\end{enumerate}

\subsection{5. Source-check ledger}\label{source-check-ledger}

\begin{itemize}
\tightlist
\item
  032 Proposition 3.1 and Corollary 3.2, pp.11--12: read entire statements and proofs, equations and codimension checked in this run. Imported CM theorem itself is background, not independently reproved.
\item
  032 Proposition 4.6 and Corollary 4.7, pp.23--25: read entire proofs to check fixed input, finite decomposition, dependence on arbitrary constant-line trivializations, and covering formula. No independent proof of analytic theta construction claimed.
\item
  032 Theorem 5.1 proof and Corollary 6.7, pp.44--45; Theorem 2.5 proof, p.46: read to track extraction from complex theta classes into algebraic Betti vectors and then rational maps. Extraction supplies spans, not a normalized unique cycle or coefficient bound.
\item
  Gan--Qiu--Takeda §11.6 and Theorem 11.4: primary arXiv PDF read via web, with theorem range and normalized factors checked. This is use-and-recheck, not a reproof.
\item
  SGA 7 I original scan: root visually verified Exposé VII §2.4 pp.164--165, Example 2.9.5 p.175 and Remark 2.9.6 p.176. This agent records that communicated check; no separate scan inspection or transcription correction is claimed.
\end{itemize}

\subsection{6. Further checked attack: algebraic norm clears the period without a lower bound}\label{further-checked-attack-algebraic-norm-clears-the-period-without-a-lower-bound}

\subsubsection{DEG-L.T6.7 --- Integral coefficient normalization and trace division}\label{DEG-L.T6.7}

Retain the 032 hypothesis that \textbf{every scalar conjugate} of each tensor \(\alpha_i\otimes\beta_i\) is of type \((1,1)\); one algebraic-coefficient Hodge class without this hypothesis would not suffice. After rescaling the chosen algebraic Betti generators, fix a finite Galois field \(L/\mathbf Q\), of degree \(d\), containing all their coefficients. They may be taken in \(H^1(B_i,\mathbf Z)\otimes\mathcal O_L\) and \(H^1(B_i',\mathbf Z)\otimes\mathcal O_L\). The divisor kernels can then be chosen in \(CH^1(B_i\times B_i')\otimes\mathcal O_L\), their product \(K\) in \(CH^4(A\times A')\otimes\mathcal O_L\), and the nonzero period \(k\) lies in \(\mathcal O_L\).

Let \(N=|N_{L/\mathbf Q}(k)|\). For every \(a\in\mathcal O_L\) the cycle \[
C_a=\operatorname{Tr}_{L/\mathbf Q}\bigl(a(N/k)K_*z\bigr)\in CH^2(A')
\] is integral and satisfies \[
\operatorname{cl}(C_a)=N\,\operatorname{Tr}_{L/\mathbf Q}(a t_\beta).
\tag{G6}
\] Thus clearing the algebraic period needs \textbf{no analytic lower bound} for \(|k|\).

More explicitly, write \(K=\sum_{j=1}^J a_j\Gamma_j\) with \(a_j\in\mathcal O_L\) and integral Chow correspondences \(\Gamma_j\). Put \[
M=\max_\sigma|\sigma k|,\qquad B_j=\max_\sigma|\sigma(a a_j)|.
\] Then \(M\ge1\), \(N\le M^d\), and the integral coefficient \[
c_j=\operatorname{Tr}_{L/\mathbf Q}\bigl(a a_j N/k\bigr)
\] satisfies \(|c_j|\le d B_jM^{d-1}\). If \(\Gamma_{j*}z=W_j^+-W_j^-\) with effective cycles of degrees at most \(D_j\), then \(C_a\) admits effective positive and negative parts of degrees at most \[
dM^{d-1}\sum_j B_jD_j.
\tag{G7}
\]

\textbf{Proof.} The finitely many algebraic coordinates of the generators have a common integer denominator, so multiplying the generators clears it. For each \(i\) the class \(\alpha_i\otimes\beta_i\) is then an integral \(\mathcal O_L\)-linear combination of Betti classes. Expand it in an integral basis of \(\mathcal O_L\). Its coefficient classes in integral cohomology are all of type \((1,1)\): scalar conjugation preserves the Hodge property of the whole tensor, and the embedding matrix of the integral basis is invertible, so each coefficient is a complex linear combination of these conjugate \((1,1)\) classes. The integral Lefschetz \((1,1)\) theorem supplies line bundles for the integral coefficient classes. Choosing their birigidified parts as in T6.5 preserves the required class. This produces the integral-coefficient kernels. Pullback by actual morphisms, cup product and evaluation on the integral fundamental class show \(k\in\mathcal O_L\).

The monic characteristic polynomial of the algebraic integer \(k\) shows \(N/k\in\mathcal O_L\) (equivalently, up to sign it is the product of all conjugates except the identity embedding). Hence every \(c_j\) is an integer, and taking traces of \(\operatorname{cl}(K_*z)=kt_\beta\) proves (G6). Since \(N\) is a nonzero integer and \(N\le M^d\), one has \(M\ge1\). For each embedding \(\sigma\), \[
|\sigma(N/k)|=\prod_{\tau\ne\sigma}|\tau(k)|\le M^{d-1}.
\] Triangle inequality for the trace gives the coefficient estimate. Replace \(c_jW_j\) by \(|c_j|W_j^+-|c_j|W_j^-\) if \(c_j\ge0\), and interchange signs if \(c_j<0\). Adding all positive and negative parts proves (G7). QED.

\textbf{Check.} Entire proof reread, including integral coefficient descent, the distinction between an integral basis of \(\mathcal O_L\) and a basis of eigenvectors, and the trace bound when \(k\) lies in a proper subfield. Repeated conjugates cause no problem. No lower bound for any individual \(\sigma k\) was used.

\textbf{Scope.} The targets \(\operatorname{Tr}(a t_\beta)\) are integral Hodge classes. An arbitrary prescribed rational class in their rational span may require a further rational coefficient \(a\) and a clearing integer. Bounding that integral spanning index is a remaining lattice problem; it is not suppressed by this theorem. Rescaling the generators changes their metric sizes and the kernels, so it is not a uniformity argument by itself.

\subsubsection{DEG-L.T6.8 --- Explicit signed-degree bound for the four divisor kernels}\label{DEG-L.T6.8}

Let \(A,A'\) have dimension \(n\), with polarizations \(H,H'\), and put \(\mathscr L=p^*H\otimes p'^*H'\) on \(A\times A'\). Let \(D_1,\ldots,D_4\) be integral divisor classes. Choose expressions \(D_i=D_i^+-D_i^-\) with \(D_i^\pm\) represented by basepoint-free effective divisors, and positive integers \(b_i\) such that both \(D_i^\pm\) and \(b_i c_1(\mathscr L)-D_i^\pm\) are nef. Such choices exist for every fixed divisor class after tensoring with a sufficiently high power of an ample bundle. Let \(z\) be an effective surface cycle on \(A\). Then \[
(D_1D_2D_3D_4)_*z=W^+-W^-,
\] with effective codimension-two cycles on \(A'\) satisfying \[
\deg_{H'}W^\pm\le48\Bigl(\prod_i b_i\Bigr)\,(H'^n)\,\deg_H z.
\tag{G8}
\]

\textbf{Proof.} Expand the product of four differences into sixteen signed products. For each product choose the four divisors successively in their basepoint-free linear systems so that they meet each component of \(z\times A'\) properly. The resulting intersection cycle is effective; proper pushforward is effective, with components whose dimension drops contributing zero. Its degree is the intersection of that product with \(p'^*c_1(H')^{n-2}\) on \(z\times A'\).

Replacing one nef divisor \(D_i^\pm\) by \(b_i c_1(\mathscr L)\) increases or preserves this intersection, since the difference and the other divisors are nef. Iterating bounds the degree by \[
\Bigl(\prod_i b_i\Bigr)
\int_{z\times A'}(p^*c_1(H)+p'^*c_1(H'))^4p'^*c_1(H')^{n-2}.
\] Only the term with two factors from \(A\) and two from \(A'\) contributes. Its binomial coefficient is \(\binom42=6\), and its integral is \(6\deg_H(z)(H'^n)\). There are eight terms of each sign. Summing them proves (G8). QED.

\textbf{Check.} Entire proof reread; dimensions are \(\dim(z\times A')=n+2\), four divisor intersections leave dimension \(n-2\), and the final degree uses exactly \(n-2\) target hyperplanes. Coefficient \(48=8\binom42\) checked. Nef intersection monotonicity is used only on an effective cycle with all remaining factors nef. The assertion is a bound for a selected signed representative, not for every rationally equivalent cycle.

Combining T6.7 and T6.8 gives a completely explicit denominator/degree bound from the source surface degree, the finitely many integral divisor-kernel choices, their coefficient embeddings, the coefficient-field degree, and the Galois conjugate periods. By T6.6 applied to each conjugate (reordering the two holomorphic and two antiholomorphic factors as necessary), \[
M\le G\deg_H z,\qquad
G=\max_\sigma\prod_{i=1}^4\|\sigma\alpha_i\|_H.
\] For fixed \(L\), generators, and kernel data, the displayed bounds are polynomial of degree at most \(d\) in \(\deg_H z\). This is an actual sufficient criterion for bounded division if those data are uniformly bounded on the chosen target orbit. It does not prove their uniform boundedness. It removes the separate purported need for a lower bound on the extracted algebraic period; the remaining nonvanishing input is exactly the already established \(k\ne0\).

\subsection{7. Further checked attack: the human proof supplies every one-form theta expansion}\label{further-checked-attack-the-human-proof-supplies-every-one-form-theta-expansion}

\subsubsection{DEG-L.T6.9 --- Finite-level theta spanning removes the existence gate}\label{DEG-L.T6.9}

For a sufficiently deep principal ball-quotient surface in the 032 construction, \textbf{every} holomorphic one-form, hence every pullback of a polarization coordinate under its actual map to an abelian variety, is a finite complex linear combination of genuine integrated theta lifts from automorphic characters of one-dimensional unitary groups. The finite combination can be made at the chosen common principal level. This is an existence statement; it does not give bounds on its coefficients or on its comparison with the integral Betti/Hom lattice.

\textbf{Sources.} The supplied van Hoften--Sempliner paper, Proposition 3.1.7 and Theorem 3.1.11, printed p.11; proof of Proposition 4.1.8, printed p.15, especially its explicit all-\(H^{1,0}\) spanning assertion. Primary literature checked: Bergeron--Millson--Mœglin, \emph{The Hodge conjecture and arithmetic quotients of complex balls}, Acta Mathematica 216 (2016), Theorem 9.1, pp.73--75, with the range stated in §9.1, p.72; \(3(a+b)+|a-b|<2m\) holds for \(m=3\), \((a,b)=(0,1)\) because \(4<6\). The older arXiv:1306.1515v2 has different numbering, Theorem 9.3, pp.57--59. Also Y. Liu, arXiv:2102.11518v2, Proposition 4.13, printed p.29, and Theorem 4.18, pp.32--33. The actual printed hypothesis in Proposition 4.13 is \(n\ge3\), confirmed by rendering the page: web extraction incorrectly displays \(n>3\).

\textbf{Proof of the finite-level form needed here.} First pass, if necessary, to a deeper principal level. For an integral isometry of the Hermitian lattice, its determinant is an algebraic unit all of whose complex absolute values are one. Hence it is a root of unity. The group of roots of unity of \(E\) is finite, so a sufficiently deep principal congruence condition kills every nontrivial determinant. At that level the principal arithmetic groups for \(U(V)\) and \(SU(V)\) coincide. The ball quotient in question is therefore a connected component of the compact unitary Shimura variety at that level. The extra cover is harmless for the normalized Chow construction by T6.4.

Extend a holomorphic one-form on this connected component by zero on the other components. The cited spanning theorem on the direct-limit cohomology writes its class as a finite sum of integrated cohomological theta classes; in degree one the source has dimension one, so the source representations are automorphic characters. Take a common deeper level on which all finitely many inputs are fixed, and average over the original compact open level subgroup. The class to be expressed is already invariant, and equivariance of the theta lift moves this averaging to its finite Schwartz input. Each such average is a finite sum, since a smooth finite-adelic input has an open stabilizer. Thus every summand can be taken at the original level. Restrict to the chosen component. Equality of classes gives equality of holomorphic forms because \(H^0(S,\Omega^1_S)\to H^1(S,\mathbf C)\) is injective on a compact Kähler variety. QED.

\textbf{Check.} Entire argument reread. The pro-level span is used only for a finite collection of vectors; compact averaging is finite on each smooth input. The determinant argument uses norm-one \textbf{units}, not arbitrary norm-one elements of \(E^\times\), which need not be roots of unity. No claim that arbitrary arithmetic finite covers are principal was used. It suffices to deepen within the principal tower already authorized by 032. Published BMM theorem and its whole proof were read; the classification theorem itself is cited background, not independently reproved.

\textbf{Concrete energy consequence.} Let \(\omega_S\) be an invariant Kähler form and let \(\eta_j=f^*\omega_j=\sum_a c_{ja}\theta_a\) be one of these finite expansions. Define the geometric Gram form \[
G_{ab}=\int_S {i\over2}\theta_a\wedge\bar\theta_b\wedge\omega_S.
\] Then the energy is exactly \(E_f=\sum_j\sum_{a,b}c_{ja}\bar c_{jb}G_{ab}\). Each full-Shimura-variety diagonal norm can be expressed by the rank-one Rallis formula (Haan, Theorem 3.4), after its explicitly chosen scalar-cohomological and measure normalization. Restricting to one component only lowers that nonnegative norm, so with those full norms denoted by \(Q_a\) one has \[
E_f\le\sum_j\left(\sum_a|c_{ja}|\sqrt{Q_a}\right)^2,
\qquad
\deg_H f_*[S]\le{1\over\int_S\omega_S^2}
\left\{\sum_j\left(\sum_a|c_{ja}|\sqrt{Q_a}\right)^2\right\}^{\!2}.
\tag{G9}
\] The last inequality is the Hodge-index bound proved in the bounds report. This is an unconditional finite-character-\(L(3/2)\) comparison after specifying the geometric realization coefficients and all local/measure factors. It is not an effective degree bound in the level alone until those quantities are controlled. The actual existence of the expansions, formerly left as a gate in the earlier analysis, is now resolved.

A further exact identity locates what prevents one-form norms from giving the degree itself. Put \(\beta=f^*c_1(H)\), \(V=\int_S\omega_S^2\), and \(\beta_0=\beta-(E_f/V)\omega_S\). The class \(\beta_0\) is primitive of type \((1,1)\), so with its Hodge norm \(\|\beta_0\|^2=-\int_S\beta_0^2\), \[
\deg_H f_*[S]={E_f^2\over V}-\|\beta_0\|^2.
\tag{G10}
\] This follows by expanding the square and using \(\beta_0\cdot\omega_S=0\). Thus the precise missing datum for an equality is the primitive degree-two component of the pulled polarization. The rank-two theta/see-saw approach is relevant to this remaining fourth-moment term; the rank-one Rallis norms already suffice for the upper comparison (G9).

\subsection{8. Exact scalar normalization and the finite-character energy formula}\label{exact-scalar-normalization-and-the-finite-character-energy-formula}

\subsubsection{DEG-L.T6.10 --- The scalar/differential-form norm conversion is explicit}\label{DEG-L.T6.10}

Let \(G=U(V)\) be the compact-quotient unitary group used above, let \(K_f\) be neat, and put \[
X_K=G(F)\backslash(D\times G(\mathbf A_f)/K_f).
\] The notation includes every connected component. Choose the invariant metric on the two-dimensional ball so that an orthonormal holomorphic coframe satisfies \[
\omega={i\over2}(e_1\wedge\bar e_1+e_2\wedge\bar e_2),\qquad dV={\omega^2\over2}.
\] Normalize Haar measure on \(K_\infty\) to have mass one. Write the archimedean Haar factor in the chosen Tamagawa factorization as \(dg_\infty=c_\infty\,dV\,dk\), and set \[
a_K=c_\infty\operatorname{vol}_{dg_f}(K_f).
\] For a holomorphic one-form \(\eta\), let \(F_\eta(g)\in E=(\mathfrak p^{1,0})^*\) be its moving-frame automorphic realization and \(F_{\eta,e}(g)\) its coefficient along any fixed unit vector \(e\in E\). Then \[
\int_{X_K}{i\over2}\eta\wedge\bar\eta\wedge\omega
={2\over a_K}\int_{[G]}|F_{\eta,e}(g)|^2\,dg.
\tag{G11}
\] For a holomorphic two-form \(\xi\), let \(F_\xi\) denote its scalar coefficient along the unit frame \(e_1\wedge e_2\). Then \[
\int_{X_K}\xi\wedge\bar\xi
={4\over a_K}\int_{[G]}|F_\xi(g)|^2\,dg.
\tag{G12}
\] Consequently, extend each \(\eta_j=f^*\omega_j\) by zero from the chosen component \(S\) to \(X_K\). Its wedge extension is also zero on the other components, and the degree formula becomes \[
\deg_H f_*[S]={2\over a_K}\sum_{j<k}\|F_{\eta_j\wedge\eta_k}\|_{L^2([G],dg)}^2.
\tag{G13}
\]

\textbf{Proof.} More generally, let \(E\) be an irreducible unitary representation of \(K_\infty\) of dimension \(r\), and let \(F(gk)=\rho(k)^{-1}F(g)\) realize a section of the corresponding homogeneous bundle. For a unit vector \(e\), Schur averaging gives \[
\int_{K_\infty}|\langle\rho(k)^{-1}v,e\rangle|^2dk={\|v\|^2\over r}.
\] Indeed, the average of the rank-one orthogonal projector onto \(\mathbf Ce\) commutes with \(\rho\); its trace is one, so it is \(r^{-1}\operatorname{id}_E\). Right invariance of Haar measure and the quotient measure formula now give \[
\int_{[G]}|F_e|^2dg={a_K\over r}\int_{X_K}\|F\|^2dV.
\] Neatness removes finite stabilizer weights; alternatively the same equality holds with the usual orbifold weights consistently on both sides. The cotangent representation of \(U(2)\times U(1)\) is irreducible of dimension two, while its determinant has dimension one. In the stated coframe, \[
{i\over2}\eta\wedge\bar\eta\wedge\omega=\|\eta\|^2dV,
\qquad
\xi\wedge\bar\xi=4\|\xi\|^2dV.
\] These pointwise identities prove (G11) and (G12). Insert (G12) into (G1) to obtain (G13). QED.

\textbf{Normalization consequence.} If a cohomological theta cocycle is written \(\varphi=\varphi_1\otimes e_1+\varphi_2\otimes e_2\), the scalar theta lift with Schwartz input \(\varphi_1\) is exactly \(F_{\eta,e_1}\). Thus no unspecified scalar/cohomological comparison remains once the unit coframe, actual cocycle, Haar factorization and \(K_f\) are specified. The lemma does not require the entire form to belong to a single irreducible automorphic representation. The supplied human paper uses probability measures on adelic class sets (§4.1, printed p.12); Haan and GQT use their stated Tamagawa normalizations. These choices must be converted, not silently equated.

\textbf{Check.} Entire proof reread. The Schur argument applies to the fiber representation, not to an assumed irreducible global representation. Both pointwise constants were recomputed in a unit coframe; \(4\) in (G12), \(2\) in (G13), and the factor \(a_K^{-1}\) agree with T6.1. Finite stabilizers and all connected components are included explicitly.

\subsubsection{\texorpdfstring{DEG-L.T6.11 --- An exact finite Hecke-\(L(3/2)\) energy formula}{DEG-L.T6.11 --- An exact finite Hecke-L(3/2) energy formula}}\label{DEG-L.T6.11}

Use the actual morphism \(f:S\to A\), its polarized coordinate forms, and their zero extensions to \(X_K\). The oscillator decomposition in van Hoften--Sempliner, Proposition 3.1.7 and §4.1.7--Proposition 4.1.8 (equivalently the cited Liu and BMM results), gives an orthogonal finite decomposition of these forms into their actual automorphic summands. For each occurring summand \(\sigma\), choose its rank-one theta realization, a nonzero source character vector \(f_\sigma\in\tau_\sigma\), and scalar Schwartz inputs \(\phi_{j\sigma}\) such that the coefficient of the \(\sigma\)-part of \(\eta_j\) along the unit cotangent vector is \[
F_{j\sigma}=\theta(\bar f_\sigma,\phi_{j\sigma}).
\] Finite sums of inputs are allowed. The multiplicity-free oscillator decomposition and the uniqueness of \((W,\pi)\) with the auxiliary source splitting fixed are exactly the facts used here; a general multiplicity-one assertion for all automorphic representations is not needed.

To use Haan's notation, take \(\gamma_\sigma\) to be the splitting character on the \(U(V)\) side. Changing the source-side splitting to \(\gamma_\sigma^3\) twists \(\tau_\sigma\) by the corresponding determinant character and preserves the displayed actual theta lift. This follows from the splitting-change formula, Wu §2.2, equation (2.9), printed p.5. Thus each summand can be put into Haan's convention without changing the form. Define \[
\mathcal Q_{j\sigma}=
{L_E(3/2,\operatorname{BC}(\tau_\sigma)\otimes\gamma_\sigma^3)
 \over L_F(3,\epsilon_{E/F})}
\,\mathcal Z^\#_\sigma(1;\phi_{j\sigma},\phi_{j\sigma};f_\sigma,f_\sigma),
\tag{G14}
\] where \(\mathcal Z^\#\) is Haan's product of normalized local integrals for pure tensors and its sesquilinear extension to finite sums. The source vector's norm is included in this expression. Haan, Theorem 3.4, identifies \(\mathcal Q_{j\sigma}\) with \(\|F_{j\sigma}\|^2_{L^2([G])}\); in particular the whole expression is real and nonnegative, without any claim that each displayed factor separately is positive.

\textbf{Inner forms and archimedean conventions.} Haan's §2.1, printed pp.3--4, starts with arbitrary Hermitian and skew-Hermitian spaces; §3.3, Theorem 3.4 imposes dimensions one and three, without a signature or quasi-split condition. Its use here is with the actual inner form \(G\): signature \((2,1)\) at the active place and definite at every other real place. The line source is compact \(U(1)\) at every real place; its active skew-Hermitian signature is \((0,1)\), and its other signs are those of the chosen CM theta datum. To match Haan's convention in which the source is Hermitian and the target skew-Hermitian, choose \(\delta\in E^\times\) with \(\bar\delta=-\delta\) and multiply the source and target forms by \(\delta^{-1}\) and \(\delta\), respectively. Under the tensor-factor exchange this preserves the symplectic space; the same additive character and transported Schrödinger/Fock polarization are used. Equivalently one may use the explicit role exchange of Wu §2.2, printed pp.5--6. The groups themselves are unchanged.

At the active place take the actual BMM cohomological cocycle \(\varphi^{1,0}\) and its coefficient along a unit cotangent vector, rather than replacing it by a spherical Gaussian. At a compact target place take the prescribed invariant Gaussian factor of that cocycle, with its matching splitting character. All these archimedean inputs and their local zeta factors remain in \(\mathcal Z^\#\). The already nonzero global theta summand ensures that its local realization exists; no unsupported assertion that arbitrary character/sign data give a nonzero lift is used. The dyadic restriction in Haan's main Bessel-period Theorem 1.2 is not a hypothesis of the Rallis statement in §3.3: the latter is stated for the arbitrary spaces of §§2.1 and3.1 and follows there from Li's convergent-range theorem. In the present case the target is globally anisotropic, so the scalar lifts are square integrable and cuspidal automatically. The Weil symplectic form convention includes the trace factor and additive character; if these are changed, the Schwartz input and local factors must be transported as well.

\textbf{Exact infinity type and the all-place ratio.} For an actual oscillator summand, van Hoften--Sempliner §4.1.7 permits the choice \(\gamma_\sigma=\mu_\sigma\) from §3.1.3, whose angular weight at each real place is \(\varepsilon_v\in\{1,-1\}\). After the preceding compatible splitting change, \(\tau_{\sigma,v}=1\) at every archimedean place. Indeed, the center of the target acts trivially on its symmetric domain and cotangent representation (and the compact target places have the trivial coefficient representation); Haan, Remark 2.1, printed p.4, identifies the image of the line member with that central image under his power-compatible splittings. Hence its character is trivial, independently of the convention that inverts the source character. Thus the Hecke character in (G14) has infinity type \((z/|z|)^{3\varepsilon_v}\). With the explicit standard local factors \(\Gamma_{\mathbf C}(s)=2(2\pi)^{-s}\Gamma(s)\) and \(\Gamma_{\mathbf R}(s)=\pi^{-s/2}\Gamma(s/2)\), the numerator factor at \(s=3/2\) is \(\Gamma_{\mathbf C}(3)\) and the denominator factor is \(\Gamma_{\mathbf R}(4)\). Their ratio is \(1/(2\pi)\) at each real place. Consequently the all-place local-product ratio is \((2\pi)^{-[F:\mathbf Q]}\) times its finite Euler-product ratio. This fixes the archimedean \(L\) factors; it does not remove the archimedean zeta inputs. A conductor-completed functional-equation symbol \(\Lambda\) is a different normalization and is not silently substituted here. Writing \(\gamma=\psi\|\cdot\|_E^{1/2}\) as in van Hoften--Sempliner Lemma 3.1.5 gives the arithmetic shift \(L(s,\operatorname{BC}(\tau)\gamma^3)=L(s+3/2,\operatorname{BC}(\tau)\psi^3)\): the value is at arithmetic argument \(3\), with weight-three center \(2\). The complete normalization proof and finite Euler-product bounds are also recorded as bounds B1--B1a, independently checked in this run.

One then has the \textbf{exact} finite formula \[
E_f={2\over a_K}\sum_{j,\sigma}\mathcal Q_{j\sigma},
\qquad
\deg_H f_*[S]\le
{4\over a_K^2\int_S\omega^2}
\left(\sum_{j,\sigma}\mathcal Q_{j\sigma}\right)^2.
\tag{G15}
\] The corresponding equality with the primitive defect is \[
\deg_H f_*[S]=
{4\over a_K^2\int_S\omega^2}
\left(\sum_{j,\sigma}\mathcal Q_{j\sigma}\right)^2
-\|\beta_0\|^2.
\tag{G16}
\]

\textbf{Proof.} T6.9 supplies actual finite expansions of the zero extensions. Project them to the oscillator summands in the cited direct-sum decomposition, and choose one rank-one model for each summand. The theta image spans that summand; finite linearity of the Schwartz input absorbs the finite expansion, with the one-dimensional source vector fixed. Distinct actual summands are orthogonal for the adelic invariant inner product. Therefore the norm of each scalar coefficient is the sum of the norms of its summands. Apply Haan's bilinear formula to each norm and (G11) to the differential form. Zero extension makes the full geometric norm equal to the norm on \(S\), so there is no component restriction inequality at this stage. Finally apply the Hodge-index identity (G10). QED.

\textbf{Check and scope.} Entire proof reread, including zero extension, representation orthogonality, finite linearity, source twists, and the constants in (G15)--(G16). Haan's actual theorem is bilinear, so it applies to Gram entries before finite inputs are combined. Its global character factor is absolutely convergent and nonzero at \(3/2\) (§3.6.1, printed p.9). This proves an exact energy expression and a degree inequality in terms of finitely many Hecke values, rather than an asserted equality of the degree with one Hecke value. It is not yet an effective level bound: the finite geometric/theta realization coefficients and the finite and archimedean Schwartz inputs are actual additional data whose size must be bounded. Their existence and the scalar norm conversion are no longer missing hypotheses.

\subsection{9. Rank-two attack: the wedge has a genuinely finite cuspidal spectrum}\label{rank-two-attack-the-wedge-has-a-genuinely-finite-cuspidal-spectrum}

\subsubsection{DEG-L.T6.12 --- Product, toric period, and finite Rallis expansion}\label{DEG-L.T6.12}

Let \(\Theta_i\) be two genuine integrated holomorphic one-form theta lifts from characters \(\chi_i\) of \(T_i=U(W_i)\), with compatible Weil splittings. Both \(W_i\) have signature \((0,1)\) at the active real place. Put \[
W=W_1\perp W_2,\quad H=U(W),\quad T=T_1\times T_2\subset H,\quad \chi=\chi_1\boxtimes\chi_2.
\] Choose the rank-two cocycle \(\varphi\) to be the tensor/wedge product of the actual rank-one cocycles, under the direct-sum identification of Weil representations. This definition includes its normalization. Let \(\mathcal K_\varphi(g,h)\) be the resulting holomorphic two-form theta kernel. Then \[
\Theta_1\wedge\Theta_2
=\int_{[T]}\mathcal K_\varphi(-,t)\,\overline{\chi(t)}\,dt.
\tag{G17}
\] This is the theta transform of a toric-period functional on \(H\), rather than the theta lift of a single character of \(H\).

Assume the finite input is fixed by an open compact subgroup of \(H(\mathbf A_f)\) and the archimedean input is \(K_{H,\infty}\)-finite, as holds for the polynomial-Gaussian cocycles just specified. Fix a common neat \(G\)-level \(K_f\). There are finitely many irreducible cuspidal \textbf{realizations} \(\pi_\lambda\subset L^2([H])\) and finitely many orthonormal vectors \(b_{\lambda r}\in\pi_\lambda\) for which \[
\mathcal K_\varphi(g,h)=
\sum_{\lambda,r}\Theta_\varphi(b_{\lambda r})(g)b_{\lambda r}(h),
\qquad
\Theta_\varphi(b)(g)=\int_{[H]}\mathcal K_\varphi(g,h)\bar b(h)\,dh.
\tag{G18}
\] The equality is an equality of smooth two-form-valued functions, and the sum is genuinely finite. Write \[
p_{\lambda r}=\int_{[T]}b_{\lambda r}(t)\overline{\chi(t)}\,dt.
\] Thus the wedge is the finite sum \(\sum_{\lambda,r}p_{\lambda r}\Theta_\varphi(b_{\lambda r})\).

Let \(\phi\) denote the scalar coefficient of \(\varphi\) in the unit canonical frame from T6.10. With inner products linear in the first variable, its exact norm on the full Shimura variety is \[
\int_{X_K}(\Theta_1\wedge\Theta_2)\wedge
                 \overline{(\Theta_1\wedge\Theta_2)}
={8\over a_K}\sum_\lambda L(1,\pi_\lambda\times\chi_V)
\sum_{r,s}p_{\lambda r}\bar p_{\lambda s}
\mathcal Z^*_{\lambda}(1/2;\phi,\phi;b_{\lambda r},b_{\lambda s}).
\tag{G19}
\] Here \(\chi_V\) is the retained source-side splitting character in GQT's notation; \(\mathcal Z^*\) is precisely its normalized doubling form, extended to finite sums of pure tensors. Its unramified denominators are retained. The standard \(L\)-function has degree two. It is not in general a Hecke-character \(L\)-function. No lower-tower, continuous-spectrum, or residual-spectrum term has been omitted.

\textbf{Proof.} The direct-sum Schrödinger model is the tensor product of the two models. At a point of \(T\), its theta sum is the product of the two theta sums; applying the wedge map to the differential-form factors gives the identity of kernels. Absolute convergence and compactness of the two torus quotients allow their integrations to be combined. This proves (G17). Compatibility of the splitting is substantive: the character on the \(G\) side is the product of the two rank-one characters, and the character on the \(H\) side is common to both lines. More explicitly, if the original line models use source-side characters \(\chi_{V,1},\chi_{V,2}\) and target-side characters \(\gamma_1,\gamma_2\), choose one \(\chi_{V,0}\) restricting to \(\epsilon_{E/F}^{3}\) on \(\mathbf A_F^\times\). The ratios \(\chi_{V,0}/\chi_{V,i}\) are trivial there and descend through \(x\mapsto x/\bar x\) to automorphic characters of \(U(1)\). Changing only the source-side splittings and retwisting the two source vectors by these characters preserves each theta integral: the unitary character multiplying its kernel cancels the conjugate character multiplying its source vector. Keep \(\gamma_1,\gamma_2\) unchanged. The direct-sum model has source-side character \(\chi_{V,0}\) and target-side character \(\gamma_1\gamma_2\). This works when their infinity types differ. Forcing a common target-side character instead would introduce a target determinant twist and would require tracking a changed automorphic form; that operation is not being used. See Wu, Proposition 2.2 and §2.2 equations (2.8)--(2.9), printed pp.4--5; the corresponding common auxiliary choice is explicit in the human paper §4.1.7, printed p.15. The direct-sum theta-product argument also appears in 032, Remark 4.3 and Proposition 4.6.

The form \(W\) is definite at the active real place, so it has no nonzero rational isotropic vector. Thus \(H\) is anisotropic over \(F\). The target \(V\) is definite at a nonactive real place, so \(G\) is anisotropic as well. Their centers are norm-one tori. Hence both adelic quotients are compact and their automorphic \(L^2\) spectra are discrete and cuspidal; cuspidality here is vacuous because there is no proper \(F\)-parabolic. In particular there is no continuous or residual contribution from an Eisenstein series on \(H\). Use the standard discrete admissible decomposition of this compact automorphic quotient, with finite multiplicities.

We give the orthogonality step explicitly, including distinct copies of an isomorphic representation. For two orthogonal irreducible subrepresentations \(\mathcal V_1,\mathcal V_2\subset L^2([H])\), the mixed convergent-range doubling identity gives the inner product of theta lifts from \(f_i\in\mathcal V_i\) as the value of a mixed doubling zeta integral. The identity is the kernel/see-saw identity underlying GQT, Proposition 11.1(ii) and Proposition 11.2, before an irreducible representation is fixed. For large \(\operatorname{Re}s\), unfold that integral. Every rational maximal isotropic plane in \(W\oplus(-W)\) is the graph of an isometry of \(W\): its projection to either copy of \(W\) is injective, since a kernel vector would be isotropic in \(W\), and dimensions then make it an isomorphism. There is therefore only the graph orbit. The unfolded integral is an integral of a Siegel section against a matrix coefficient \[
\langle R(h)f_1,f_2\rangle_{L^2([H])}
\] (up to the fixed splitting characters, which cancel in the usual doubling normalization). This coefficient is zero because \(\mathcal V_1\) is invariant and orthogonal to \(\mathcal V_2\). The zeta integral is identically zero and stays zero on meromorphic continuation. The mixed theta inner product is zero. This argument works for distinct orthogonal realization copies as well as inequivalent representations; it does not assume global multiplicity one.

For each \(h\), the two-form \(\mathcal K_\varphi(-,h)\) belongs to the fixed finite-dimensional space \(\mathcal F=H^0(X_K,\Omega^2)\). The \(H\) action commutes with the \(G\) action and the cohomological differential, so varying \(h\) preserves this property and the chosen \(G\)-level. The spectral coefficients \(\Theta_\varphi(b)\) also lie in \(\mathcal F\). By the preceding orthogonality, the nonzero images belonging to distinct orthogonal irreducible source realizations are mutually orthogonal nonzero subspaces of \(\mathcal F\). There can therefore be only finitely many such realizations. Within each one, the fixed finite level and the finitely many archimedean compact types of \(\varphi\) select a finite-dimensional subspace by admissibility. All other spectral coefficients vanish. The \(L^2\) expansion is consequently finite; both sides are smooth, so the equality holds everywhere, proving (G18). Applying the toric functional to this finite equality is legitimate and introduces no Sobolev convergence issue.

For every retained \(\pi_\lambda\), the global theta lift is nonzero, hence all its local theta lifts are nonzero. In GQT, the dimensions are \(n=2,m=3\), \(\epsilon_0=0\), and the target Witt index is \(r=0\). Thus \(2<3\le4\), \(r\le n\), and the lower-tower hypothesis is empty. Theorem 11.4(ii), printed pp.54--55, applies at \(s_{m,n}=1/2\) and gives \[
\langle\Theta_\phi(b_{\lambda r}),\Theta_\phi(b_{\lambda s})\rangle
=2L(1,\pi_\lambda\times\chi_V)
\mathcal Z^*_{\lambda}(1/2;\phi,\phi;b_{\lambda r},b_{\lambda s}).
\] The \(L\)-value here is holomorphic; a residue is not required. Expand the finite squared norm, use the orthogonality already proved, and apply (G12). This proves (G19). QED.

\textbf{Check.} Entire proof reread. Definite-at-one-place anisotropy, the graph-orbit argument, orthogonal-copy issue, and the finite-dimensional-image argument were checked separately. Finite image alone would not imply finite spectral support: orthogonality and admissibility are essential. The toric functional is applied only after finiteness has been proved. The conjugate in the definition of \(\Theta_\varphi\) agrees with the rendered GQT §11.2, printed p.52; web extraction omits that overbar. The factor \(8/a_K=4/a_K\times2\) in (G19) was checked. This proof uses the standard discrete admissible spectral decomposition and the cited convergent Siegel--Weil/doubling identity as background theorems, rather than claiming to reprove them.

\subsubsection{\texorpdfstring{DEG-L.T6.13 --- General degree as a finite standard-\(L(1)\) norm sum}{DEG-L.T6.13 --- General degree as a finite standard-L(1) norm sum}}\label{DEG-L.T6.13}

The preceding statement extends to every actual map \(f\) in the 032 construction in the following precise, non-effective form. There exist finitely many rank-two skew-Hermitian spaces \(W_\lambda\), irreducible cuspidal representations \(\pi_\lambda\) of \(U(W_\lambda)\) with nonzero lift to \(G\), and finite theta input tensors \(q_{jk,\lambda}\) such that \[
\deg_H f_*[S]
={4\over a_K}\sum_{j<k}\sum_\lambda
 L(1,\pi_\lambda\times\chi_{V,\lambda})
 \mathcal Z^*_{\lambda}(1/2;q_{jk,\lambda},q_{jk,\lambda}).
\tag{G20}
\] To make this a defined equality, the data include unitary \(G(\mathbf A)\)-intertwiners when isomorphic target representations occur through different theta models or in several copies. The input tensors encode the actual zero-extended wedge forms through these identifications. They are not normalized to have size one. The formula retains the local factors and all input coefficients; it does not assert that these quantities can be recovered or bounded from the source level alone.

\textbf{Proof.} Expand the zero extension of each \(\eta_j\) by T6.9. Each \(\eta_j\wedge\eta_k\) is a finite sum of products of rank-one theta forms. Apply T6.12 separately to every such product, taking a common \(G\)-level; all relevant rank-two source groups are anisotropic at the active place. This places all the scalar wedge coefficients in the sum of finitely many global rank-two theta images. Every nonzero such image is an irreducible cuspidal representation by Wu, Theorem 5.3, printed p.15, since the target is anisotropic. Their sum is a finite-length unitary representation. Decompose it into orthogonal irreducible actual copies.

For each resulting copy choose one theta image isomorphic to it and a unitary intertwiner from that image to the chosen copy. Transport the orthogonal component of every scalar wedge coefficient back by this isometry. It is a smooth vector of finite compact type in that theta image, and hence a finite sum of theta integrals by the definition of the global theta image (equivalently its irreducibility as an admissible module). Denote the resulting finite input tensor by \(q_{jk,\lambda}\). Unitary transport and orthogonality express the squared norm of the scalar wedge as the sum of squared norms of these theta integrals. Apply GQT §11.7, Theorem 11.4(ii) to each sum, extending its bilinear identity by finite linearity, and insert the result into (G13). The factor is \(2/a_K\) from (G13) times \(2\) from Rallis. QED.

\textbf{Check and practical distinction.} Entire proof reread, including the finite-length span and the possible repetition of an abstract target representation. The unitary identifications are explicitly retained; simply declaring different source groups or source characters orthogonal would be invalid. In the special case of one common rank-two source and splitting, (G20) has the concrete toric-period/input Gram coefficients of (G19). With different source models, the isotypic comparison data are additional data. The theorem gives a finite spectral identity and locates the exact analytic factors; it does not give a canonical single Hecke value, an effective closed formula in level/conductor alone, or a bound for the integral realization coefficients. For the desired upper degree estimate, T6.11 is stronger operationally: it already avoids rank-two comparison data and uses only Hecke-character values at \(3/2\).

\subsubsection{DEG-L.T6.14 --- Exact level-index dependence for fixed adelic data}\label{DEG-L.T6.14}

Keep the group, invariant metric, additive character, splittings and all automorphic/Schwartz inputs fixed. Let \(K_f'\subset K_f\) be neat compact open subgroups fixing those inputs, of index \(d=[K_f:K_f']\). For the full Shimura varieties, the finite étale map \(X_{K'}\to X_K\) has degree \(d\), and \[
a_{K'}={a_K\over d},\qquad
E_{K'}=dE_K,\qquad
\int_{X_{K'}}\xi\wedge\bar\xi
=d\int_{X_K}\xi\wedge\bar\xi.
\tag{G21}
\] The scalar automorphic functions on \([G]\), their Rallis norms, their global \(L\)-values and their local integrals are unchanged. Thus (G11)--(G20) have exactly the required factor \(d\) from \(a_K^{-1}\). If the maps to the abelian variety are pulled back through this cover, the raw pushforward surface degree also multiplies by \(d\). On a selected connected component the analogous factor is its actual covering degree, which need not equal the index for the full disconnected Shimura variety.

\textbf{Proof.} Haar measure gives \(\operatorname{vol}(K_f')=\operatorname{vol}(K_f)/d\), with \(c_\infty\) unchanged. A fixed right-\(K_f\)-invariant automorphic function is the same function when regarded as right-\(K_f'\)-invariant, so its adelic norm does not change. Apply T6.10, or integrate pulled-back forms over the finite étale cover, to obtain the displayed identities. Proper pushforward of the fundamental cycle gives the degree assertion. The statement for one component uses its own finite covering degree. QED.

\textbf{Check and limit.} Entire proof reread. The source-volume factor in the Hodge-index bound also scales by \(d\), so \(E^2/\int\omega^2\) scales by \(d\), not \(d^2\). This is the same cover growth that cancels exactly in the normalized Chow output by T6.4. It describes fixed-data level change; allowing new Schwartz vectors, primitive integral maps or target polarizations changes additional data and is not covered by this invariance assertion.

\subsection{10. A nonmonomial rank-two constituent can actually occur in a holomorphic wedge}\label{a-nonmonomial-rank-two-constituent-can-actually-occur-in-a-holomorphic-wedge}

\subsubsection{DEG-L.T6.15 --- Existence of a Steinberg constituent in an admissible balanced wedge}\label{DEG-L.T6.15}

There exist a CM extension \(E/F\) with \([F:\mathbf Q]=2\), a Hermitian space \(V\) of dimension three with signatures \((2,1)\) and \((3,0)\), a neat compact unitary Shimura surface \(X_K\) attached to \(G=U(V)\), and two integrated holomorphic one-form theta lifts \(\eta_1,\eta_2\) from unitary characters of \(U(1)\), with the following property. In the finite rank-two Rallis decomposition of \[
\eta_1\wedge\eta_2\ne0
\] from T6.12, an actual source realization \(\pi\) occurs with nonzero contribution, and its stable base change \(\Pi\) to \(\mathrm{GL}_2(\mathbf A_E)\) is cuspidal and has a twist of the Steinberg representation at a finite place of \(E\). In particular the Weil--Deligne monodromy operator there is nonzero. Twisting by the fixed character in the Rallis standard \(L\)-function does not remove this monodromy.

Consequently its local standard parameter at this place is neither a direct sum of one-dimensional parameters nor the induction of a one-dimensional parameter from a quadratic extension. Thus this contributing source cannot be replaced, at the level of compatible local parameters, by those Hecke-character data. This gives an actual counterexample to a universal \textbf{character-only spectral reduction} for holomorphic rank-one theta wedges of this kind. No equality or inequality of isolated numerical \(L\)-values is being inferred from monodromy. The stronger refinement concerning equality of almost-everywhere Euler factors is not claimed here; it would additionally use a precisely sourced strong multiplicity-one and automorphic-induction theorem. Nor is a statement being made about one preselected indecomposable four-factor CM relation.

The pair and its complex conjugate give a nonzero balanced four-factor period on a connected component of \(X_K\). Thus the example lies in the geometric class of compact unitary surfaces and CM one-forms used by 032. The four-factor relation obtained this way is decomposable; a claim that this construction produces a distinct indecomposable switch relation would require an additional argument and is not made.

\textbf{Imported results and precise scope.} The proof uses the following established results as background, together with the previously checked theta-product and mixed-orthogonality identities. It does not purport to reprove the trace formula or local Langlands correspondence.

\begin{itemize}
\tightlist
\item
  Gan--Gross--Prasad, \emph{Restrictions of representations of classical groups: examples}, Astérisque 346 (2012), §15, printed p.161 and Lemma 15.4 p.162: the local/global classification of Hermitian spaces, and existence of a dimension-two definite space whose unitary group is quasi-split at every finite place when the totally real degree is even.
\item
  Shin, \emph{Automorphic Plancherel density theorem}, author version §5.2, Example 5.6 and Theorem 5.7, printed pp.26--27: existence with a prescribed positive-Plancherel family of local discrete-series representations, one auxiliary variable finite level, fixed algebraic coefficient representation at infinity, and unramifiedness elsewhere.
\item
  Labesse, \emph{Changement de base CM et séries discrètes}, author version Remark 5.2 p.42, Corollary 5.3 pp.43--44, and Theorem 5.9 p.49: weak base change, compatibility at every split place, and strong local compatibility at finite quasi-split places once the base change is cuspidal. The extra unramifiedness restrictions in §5.2 concern Theorems 5.4--5.5 and are not hypotheses of these statements.
\item
  Atobe--Gan, arXiv:1602.01299v1, Appendix D, Theorem D.2(i), printed p.43, citing Gan--Ichino for the unitary case: for \(m=n+1\), if the source parameter does not contain the one-dimensional summand \(\chi_V\), both target pure inner forms have nonzero local theta lift. Mínguez, Ann. Sci. Éc. Norm. Supér. 41 (2008), Theorem 1(2), p.718, gives nonzero local theta lift for every irreducible representation in the split \(\mathrm{GL}_2\)-to-\(\mathrm{GL}_3\) case.
\item
  BMM, \emph{The Hodge conjecture and arithmetic quotients of complex balls}, Acta Mathematica 216 (2016), Proposition 5.17 and its proof, Lemma 5.18, and Proposition 5.19, printed pp.48--50: the product cocycle is closed and holomorphic, its definite source type is a determinant character, and its top polynomial is a nonzero determinant. Haan, Remark 2.1, p.4, and Wu, Proposition 2.2 and equations (2.8)--(2.9), pp.4--5, supply the central and direct-sum splitting compatibility.
\item
  GQT §11.6, p.54, gives the unramified doubling integral; §11.7, Theorem 11.4, pp.54--55, and §11.8, Proposition 11.6, pp.55--56, give the required inner product and local nonvanishing statements. The resulting local/global criterion is Theorem 11.7(ii), pp.56--57; its additional condition \(m=d(n)+1\) holds here.
\item
  The standard \(L\)-function of a unitary cuspidal \(\mathrm{GL}_2\) representation, twisted by a unitary Hecke character, is nonzero at \(1\). A primary accessible formulation is Shahidi, \emph{On nonvanishing of L-functions}, Bull. Amer. Math. Soc. 2 (1980), the Theorem on p.462, with proof outline pp.462--463, applied to \(\mathrm{GL}_2\times\mathrm{GL}_1\). The finite omitted local factors are nonzero and finite at \(1\) for these generic unitary local representations. No GRH is used here.
\end{itemize}

\textbf{Proof.} We separate the choices, the exact archimedean vector, and the passage from a rank-two lift to a rank-one product.

\textbf{1. The global source and its finite components.} Choose a real quadratic field \(F\) and a CM quadratic extension \(E/F\); one may take a biquadratic CM field \(E=F(i)\). The cited Hermitian classification gives a dimension-two space that is negative definite at both real places and quasi-split at every finite place. Multiply its form by a fixed trace-zero element when converting it to the skew-Hermitian convention for a theta source. Denote the resulting space by \(W\) and its group by \(H=U(W)\). Orthogonal diagonalization over \(E\) gives \[
W=W_1\perp W_2,
\] where both lines are negative at every real place. The definiteness implies that \(H\) is anisotropic over \(F\), with compact adelic quotient and compact real groups.

Choose a unitary splitting character \(\gamma\) of \(\mathbf A_E^\times/E^\times\) restricting to \(\epsilon_{E/F}\), with its weight-one CM infinity type chosen for these negative lines. At the active real place use the actual holomorphic one-form cocycle, and at the other real place use the invariant Gaussian for the compact target. The character is chosen to make these genuine relative cochains; an arbitrary previously fixed infinity type would not suffice. Existence of the required CM character is the construction in van Hoften--Sempliner, §3.1.3 and Lemma 3.1.5. Each line uses source-side character \(\gamma^3\) and target-side character \(\gamma\); the sum uses source-side character \(\gamma^3\) and target-side character \(\gamma^2\).

Choose a finite split place \(v_s\) of \(F\), a finite nonsplit place \(v_i\), and another split place \(v_0\). Take a finite set \(S\), containing \(v_s,v_i\) but not \(v_0\), that includes every place where the extension, source lattices, splitting or additive character is ramified. At every place in \(S\), prescribe the unitary unramified-twist orbit of the Steinberg representation of the quasi-split group \(H(F_v)\). At a split place this is the usual \(\mathrm{GL}_2(F_v)\) orbit. These families satisfy Shin's positive-Plancherel and regularity hypotheses by Example 5.6, taking the Levi subgroup to be the group itself. Apply Theorem 5.7 with auxiliary place \(v_0\) and trivial algebraic coefficient representation \(\xi=1\).

The real groups are compact, so being \(\xi\)-cohomological means exactly that the real representation is trivial: the relative Lie algebra quotient is zero, and degree-zero relative cohomology is the invariant subspace. Thus we obtain an actual cuspidal automorphic realization \(\pi\subset L^2([H])\) with \[
\begin{aligned}
\pi_\infty&=1,\\
\pi_v&\text{ is a unitary Steinberg twist for }v\in S,\\
\pi_v&\text{ is unramified for }v\notin S\cup\{v_0\}.
\end{aligned}
\] Shin's center condition holds because both the global and real split ranks of this norm-one center are zero. The remaining group-unramifiedness hypothesis holds by the choice of \(S\). Compactness also makes cuspidality automatic in this case.

Labesse applies because \([F:\mathbf Q]\ge2\), the real representations are discrete series of compact groups, and his hypothesis \((*)\) holds automatically for a compact real group. Let \(\Pi\) be the resulting isobaric base change. At either place \(w_s\mid v_s\) of \(E\), local compatibility gives a unitary Steinberg twist. It follows that \(\Pi\) is cuspidal: in rank two the other possibilities in the discrete/isobaric classification are an isobaric sum of two characters, with a sum of one-dimensional local parameters, or the residual representation obtained from a repeated character. The latter has local Langlands quotient of the form \(\nu\circ\det\), again with zero monodromy. Neither can have the matched Steinberg parameter. This uses the discrete-spectrum classification already invoked in the proof of Labesse's Corollary 5.3. Labesse's Theorem 5.9 now gives local base-change compatibility at every finite place, since \(H\) is quasi-split there. This also identifies the source standard \(L\)-function used by GQT with the standard \(\mathrm{GL}_2\) function of \(\Pi\), retaining its specified splitting-character twist.

\textbf{2. A coherent target with all finite lifts nonzero.} At every nonsplit place in \(S\), the source parameter is a twist of \(S_2\). It contains no \(\chi_V\otimes S_1\) summand as a Weil--Deligne parameter, even when its semisimplified Weil characters happen to have a related twist. Theorem D.2(i) therefore makes the lift nonzero for both dimension-three Hermitian spaces. At a split place, including \(v_0\), Mínguez gives nonvanishing without a restriction on \(\pi_v\).

Outside \(S\cup\{v_0\}\), take the unramified target Hermitian space and the spherical Schwartz data. The unramified local doubling calculation gives a nonzero normalized integral at \(s_{3,2}=1/2\); hence the local theta lift is nonzero. The normalizing denominator has no pole there. Equivalently, this is the unramified theta correspondence. All the relevant characters and lattices are unramified at these places by construction.

Prescribe signature \((2,1)\) at one real place and \((3,0)\) at the other. Choose one of the allowed finite target forms at each remaining place. The local/global Hermitian product condition can be corrected at \(v_i\), where both choices give nonzero lift. The exact norm-class sequence recalled by Gan--Gross--Prasad on p.161, together with the Hermitian classification there, then gives a global target \(V\) with all these localizations. It is anisotropic over \(F\), since it is definite at one real place. Consequently its Shimura varieties at neat level are compact smooth surfaces. Notice that the two odd-dimensional local Hermitian spaces may have isomorphic unitary groups; the choice is a choice of Hermitian space and Weil model and must still be retained.

\textbf{3. The prescribed holomorphic cocycle survives at infinity.} At the active place take the product of the two rank-one holomorphic cocycles, and at the other place their Gaussian product. The resulting rank-two relative cocycle \(\varphi_\infty\) is closed and has type \((2,0)\). BMM's Proposition 5.19 evaluates its scalar top component as a nonzero two-by-two determinant polynomial times the Gaussian. Proposition 5.17(8) says that its source \(U(2)\) type is one dimensional; changing from the polynomial normalization to the actual Weil splitting only changes this type by a determinant character. At the compact target place the Gaussian likewise has a one-dimensional source type.

Under the compatible power splittings, Haan's central identity on each line, followed by Wu's direct-sum identity, identifies the action of \(zI_2\) in the source with that of \(zI_3\) in the target. The properly twisted relative target cochain has trivial target central action. Hence this one-dimensional source type has trivial central action. Every genuine one-dimensional representation of compact \(U(2)\) is a power of its determinant, and its value on \(zI_2\) is \(z^{2k}\); triviality for all \(z\) forces \(k=0\). Therefore \(\varphi_\infty\) is invariant under all of \(H(F_\infty)\), rather than merely under a subgroup. Source action is on the Schwartz factor, so its nonzero scalar canonical-frame component is invariant as well.

For \(\pi_\infty=1\), the corresponding raw compact local doubling pairing of this scalar component with itself is \[
\int_{H(F_v)}\langle\omega_v(h)\phi_v,\phi_v\rangle\,dh
=\operatorname{vol}(H(F_v))\,\|\phi_v\|^2>0.
\] For precision, GQT §11.3, printed p.52, normalizes the doubling intertwiner by \(\delta(\phi_1\otimes\bar\phi_2)(0)=\langle\phi_1,\phi_2\rangle\). Its extra \(\chi_V\circ\det\) character is on the second source factor. At \((h,1)\) that character is one, so the Siegel--Weil section at \(s_{3,2}\) is exactly \(\langle\omega_v(h)\phi_v,\phi_v\rangle\). GQT §11.6, p.54, integrates it against the source matrix coefficient; take the trivial source vector to have norm one. This verifies the displayed compact integral without an unspecified scalar or character. Division by its nonzero local \(L\) factor does not change nonvanishing. Thus it is this exact holomorphic cocycle, not just an unspecified local theta vector, that survives.

At the finitely many finite places requiring a choice, GQT Proposition 11.6 supplies nonzero local doubled matrix entries from the actual target Weil representation. At almost all places use the spherical entries. Pair the same fixed scalar archimedean component on both sides. The source standard value \(L(1,\pi\times\chi_V)\) is nonzero by the cuspidal \(\mathrm{GL}_2\) nonvanishing theorem; its holomorphy also follows from GQT once all the local lifts are nonzero. In GQT the dimensions are \(n=2,m=3\) and the global target Witt index is \(r=0\), so the dimensional range holds and the lower-tower condition is empty. Theorem 11.4 makes the resulting global theta inner product nonzero. Some global lift with the fixed \(\varphi_\infty\) and suitable finite Schwartz input is therefore a nonzero holomorphic two-form. This argument only needs nonzero bilinear finite entries; it does not assume a positivity assertion about each normalized finite local factor.

\textbf{4. A toric character detects this source image.} Put \(T=U(W_1)\times U(W_2)\). The quotient \([T]\) is a compact abelian group. For any nonzero smooth \(b\in\pi\), there is a finite-adelic translate \(b'\) with \(b'(1)\ne0\): choose a point where \(b\) is nonzero and use \(\pi_\infty=1\) to remove its real component. The restriction \(b'|_{[T]}\) is nonzero. It is invariant under \(T(F_\infty)\) and an open compact subgroup at the finite places, so it factors through a finite abelian quotient. Fourier completeness on that quotient gives a finite-order automorphic character \[
\chi=\chi_1\boxtimes\chi_2,\qquad \chi_\infty=1,
\] whose toric functional \(\ell_\chi(b)=\int_{[T]}b(t)\overline{\chi(t)}\,dt\) is nonzero on \(\pi\). There is no need to prescribe this character before constructing \(\pi\), and no distinguished-globalization theorem is being assumed.

Let \(P_\pi\) be the orthogonal projection to this chosen actual source realization. Form the span \[
\mathcal I=\operatorname{span}\{P_\pi\mathcal K_{\varphi_\infty\otimes\phi_f}(g,-):
                  g\in G(\mathbf A),\ \phi_f\in\mathcal S_f\}.
\] Here scalar canonical-frame coefficients are understood when evaluating the two-form kernel. The nonzero theta lift just proved shows that \(\mathcal I\ne0\). Right translation by \(h_f\in H(\mathbf A_f)\) replaces \(\phi_f\) by \(\omega(h_f)\phi_f\); hence \(\mathcal I\) is an invariant subspace of \(\pi\). Because \(\pi_\infty=1\), its smooth \(H(\mathbf A_f)\)-module is irreducible. Thus \(\mathcal I=\pi^\infty\). Some projected kernel therefore has a nonzero \(\ell_\chi\)-period. T6.12's orthogonality for distinct actual source realizations implies that its nonzero target theta contribution cannot cancel against the other source terms after this toric integration.

\textbf{5. Extract a product of two rank-one forms.} Finite-adelic Schwartz functions on a direct sum are algebraic finite sums of products of Schwartz functions on the two summands. To see finiteness directly at one place, a locally constant compactly supported function is constant on finitely many product cosets inside a compact support; their indicator functions are products. At almost all places use the factorized standard lattice indicator. Therefore the chosen \(\phi_f\) is a finite sum of decomposable inputs. The archimedean cocycle was already defined as the product/wedge of the line cocycles. Apply the exact theta-product identity (G17) with \(\chi=\chi_1\boxtimes\chi_2\). The resulting toric period is a finite sum of wedges of actual integrated holomorphic line theta forms. Since its \(\pi\)-contribution is nonzero, at least one such wedge has nonzero \(\pi\)-contribution. Call its factors \(\eta_1,\eta_2\). In particular neither line lift is zero; this conclusion is obtained from the nonzero product and is not an unverified separate local/global assertion about the two chosen torus characters.

Choose a common sufficiently small neat target level fixing these two inputs. A nonzero holomorphic top form has \[
\int_{X_K}(\eta_1\wedge\eta_2)\wedge
                  \overline{(\eta_1\wedge\eta_2)}>0,
\] so some connected component has a nonzero balanced four-factor period. These are CM one-forms by the oscillator description used in T6.9. Their span is contained in the CM part of the Albanese cohomology, exactly the Hodge-theoretic input of the 032 map construction. The claim here concerns these actual theta forms and their wedge. It does not assert that their initially chosen analytic normalizations are integral Betti generators, nor does it remove the integral-realization gate in the quantitative problem.

Finally \(\Pi_{w_s}\) has parameter \(\nu\otimes S_2\), with nonzero monodromy. A one-dimensional Weil--Deligne representation has zero nilpotent operator; a direct sum of such characters, or induction of a character through a quadratic extension, also has zero operator. Twisting preserves whether this operator is zero. Thus the relevant rank-two parameter is neither a sum nor a quadratic induction of characters. The nonzero \(\pi\)-component contributes a nonzero norm term in (G19), so it is a genuinely occurring standard factor. QED.

\textbf{Albanese degree corollary.} The same construction yields an actual morphism \(f:S\to A\) to a polarized CM abelian variety, with surface image, for which the finite exact degree identity T6.13 has a nonzero contribution from a nonmonomial rank-two source.

\textbf{Proof of the corollary.} Fix the ambient full Shimura variety and the actual irreducible target theta realization \(\sigma\) coming from the chosen source \(\pi\). It is irreducible by the global theta theorem used in T6.13. Write the original scalar wedge coefficient as \(F_\xi\), and let \(P_\sigma\) denote orthogonal projection in this full automorphic Hilbert space. The construction proves \(P_\sigma F_\xi\ne0\). Decompose the geometric form into its zero extensions from connected components, \[
F_\xi=\sum_S F_{\xi,S}.
\] There are finitely many components, so at least one satisfies \(P_\sigma F_{\xi,S}\ne0\). Fix this \(S\). This is an assertion about a specified ambient projection; it does not assign an intrinsic automorphic action to one isolated connected component. In particular the restricted wedge on \(S\) is nonzero.

If the full CM description of \(H^1\) in T6.9 is used, take \(A=\operatorname{Alb}(S)\). For the weaker input that only the two chosen forms lie in its CM part, take the corresponding CM isogeny quotient of \(\operatorname{Alb}(S)\). This quotient exists because polarizable rational Hodge structures of weight one are semisimple and correspond to abelian varieties up to isogeny. Clear one isogeny denominator to obtain an actual morphism \(f:S\to A\). Its pullback on holomorphic forms contains the plane spanned by \(\eta_1|_S,\eta_2|_S\). Consequently \(\dim f(S)=2\).

Choose any ample polarization on \(A\). Gram--Schmidt applied to the preimage of this two-dimensional plane, using the Hermitian form of its invariant polarization, gives the first two members \(\omega_1,\omega_2\) of a unit invariant coframe. Their pulled-back wedge is \(c\,\xi|_S\) for some \(c\ne0\). Extend to a full unit coframe. Its term in T6.1 therefore has a nonzero \(\sigma\)-projection after zero extension, and T6.13 may use the original \(\pi\)-theta model for that projection. The associated norm term is nonzero.

At the split Steinberg place, Mínguez's parameter formula identifies the local target parameter, up to the fixed characters, with the contragredient rank-two parameter plus a one-dimensional summand. Its monodromy is therefore nonzero. A rank-two theta model with zero-monodromy source could not produce this same target representation at that place. Thus merely changing the rank-two model for this actual target summand does not turn it into an induction or a sum of characters. This corollary is an actual-map degree statement; the target is not asserted to be the four particular \(B_i\) of a prescribed indecomposable switch relation. QED.

\textbf{CM-motive consistency check.} The conclusion concerns the standard \(L\)-function of the \textbf{source in a theta norm formula}. It does not identify that degree-two source parameter with the Galois representation carried by the image of \(\Lambda^2H^1\) in the surface's \(H^2\). The target transfer is endoscopic: its parameter has a rank-two part and a character part. An individual endoscopic cohomological packet contribution need not contain the full standard transfer. Chenevier--Harris, \emph{Construction of automorphic Galois representations, II}, Introduction, printed p.2, gives explicit even-degree examples in which the \(U(2)\times U(1)\)-to-\(U(3)\) transfer contributes only the one-dimensional Galois representation from \(U(1)\). This directly rules out treating the full standard transfer as automatically equal to each geometric cohomology contribution. No packet-specific Galois calculation for the present Steinberg example is asserted here. The CM nature of the cup image therefore does not negate any step of the analytic proof; it remains compatible with a different, nonmonomial source factor in its norm expression.

\textbf{Check and limitation.} The entire proof has been read as one chain, with separate checks of Shin's real/center hypotheses, base-change cuspidality at the split Steinberg place, global Hermitian coherence at the nonsplit Steinberg place, exact archimedean invariance, nonzero bilinear local/global theta pairing, toric Fourier detection, irreducible-source spanning, and the finite tensor extraction. Bounds independently checked the archimedean argument and the conditional toric/kernel argument, including the mixed-orthogonality requirement. The global existence step imports the cited theorems with their hypotheses as stated. The example proves a spectral obstruction for an admissible balanced wedge; it does not prove a corresponding obstruction for every indecomposable 032 switch relation, or disprove a numerical identity with unrestricted auxiliary coefficients.

\subsection{11. The distinct indecomposable switch gate is closed for an explicit relation}\label{the-distinct-indecomposable-switch-gate-is-closed-for-an-explicit-relation}

\subsubsection{DEG-L.T6.16 --- A distinct indecomposable switch with a Steinberg source}\label{DEG-L.T6.16}

There exist a Galois CM field \(E\) of degree eight, four \textbf{distinct} odd CM sign functions \(v_i\), and an actual four-map surface witness in the framework of 032, Proposition 3.1, with the following properties.

\begin{enumerate}
\def\labelenumi{\arabic{enumi}.}
\tightlist
\item
  \(v_1+v_2=v_3+v_4\), and the selected balanced list \(v_1,v_2,-v_3,-v_4\) has no nonempty proper balanced sublist. Thus this is an indecomposable four-factor switch in the literal tensor sense, rather than a wedge paired with its own conjugate. Distinctness here concerns the labeled CM types; it does not assert that the four abelian varieties are pairwise nonisogenous or simple.
\item
  There are a smooth connected projective compact ball quotient \(Y\), actual morphisms \(f_i:Y\to B(v_i)\), and nonzero eigenforms \(\alpha_i\in U(v_i)_1\), chosen with algebraic Betti coefficients, such that \[
  \int_Y f_1^*\alpha_1\wedge f_2^*\alpha_2
         \wedge\overline{f_3^*\alpha_3\wedge f_4^*\alpha_4}\ne0.
  \tag{T16.1}
  \]
\item
  In a specified full unitary Shimura variety containing \(Y\) as a connected component at a sufficiently deep level, the zero extension of the first wedge in (T16.1) has nonzero projection to an actual target automorphic isotypic space whose dimension-two theta source has cuspidal base change with a Steinberg local component. In particular its Weil--Deligne monodromy is nonzero at a fixed finite split place.
\item
  For the actual product morphism \(f=(f_1,f_2,f_3,f_4)\), some nonzero term in the exact degree formula T6.13 has such a rank-two source. The image of \(f\) is a surface. A product polarization can be used. The source parameter cannot be a sum of characters or a quadratic induction of a character.
\end{enumerate}

The assertion is existential for this explicit switch, not for every prescribed indecomposable relation. It concerns the actual source in an exact theta-norm expansion of the degree; it does not identify that source standard parameter with a geometric Galois constituent of the CM cup image.

\textbf{Proof.}

\textbf{1. The four sign functions and the rational spaces}

Take a totally real biquadratic field \(F\), for example \(\mathbf Q(\sqrt2,\sqrt3)\), and put \(E=F(i)\). Choose the four representatives of \(\operatorname{Gal}(E/\mathbf Q)/\langle c\rangle\) that fix \(i\), and label them \(\mu,a,b,c_0\), with \(\mu=1\). (The label \(c_0\) is a real place and is not complex conjugation.) On these representatives set \[
\begin{array}{c|rrrr}
 &\mu&a&b&c_0\\\hline
v_1&+1&+1&+1&+1\\
v_2&+1&-1&-1&+1\\
v_3&+1&+1&-1&+1\\
v_4&+1&-1&+1&+1.
\end{array}
\tag{T16.2}
\] Extend each function oddly under complex conjugation. The equality of the two pair sums is immediate. The four functions are distinct. A proper balanced sublist cannot have one element. It cannot have three, since the sum of all four is zero and the omitted sign function is nonzero. It cannot have two: the two positive functions cannot cancel, since both take value \(+1\) at \(\mu\); the two negative functions cannot cancel; and a positive and a negative could cancel only if two of the four displayed functions were equal. This proves indecomposability. Since every element of this Galois group has order two, the inversion in 032's convention \(d_i(\tau)=v_i(\tau^{-1})\) makes no change.

Use \(\delta=i\) in 032 §4.2, so its imaginary coefficient is positive at all four chosen representatives. Accordingly \(d_i=+1\) means a negative Hermitian multiplicity line in the Fock convention used here. Choose a dimension-two Hermitian space \(W\) which is negative definite at \(\mu,c_0\), has signature \((1,1)\) at \(a,b\), and is quasi-split at every finite place. Such a space exists by the Hermitian classification and the norm-class product sequence in Gan--Gross--Prasad §15, printed p.161: a quasi-split plane has determinant class \(-1\), and the number of definite real places in this prescription is even. This uses the general classification there, rather than applying their all-definite Lemma 15.4 to the mixed signature prescription.

Choose by weak approximation a vector in \(W(E)\) of negative norm at all real places. Its line \(W_1\) and orthogonal complement \(W_2\) have precisely the signs \(d_1,d_2\) of (T16.2). Thus \(W=W_1\perp W_2\). Convert to a skew-Hermitian theta source by the same fixed trace-zero scalar as in T6.15. Write \(H=U(W)\). It is anisotropic over \(F\), because it is definite at \(\mu\); hence \([H]\) is compact, although two of its real factors are noncompact.

Choose weight-one unitary splitting characters \(\gamma_1,\gamma_2,\gamma_3\) for the first three line types, all restricting to \(\epsilon_{E/F}\), by the CM-character construction used in T6.15. Set \[
\gamma_4=\gamma_1\gamma_2\gamma_3^{-1},\qquad
\chi_W=\gamma_1\gamma_2=\gamma_3\gamma_4,\qquad
\chi_V=\gamma_1^3.
\tag{T16.3}
\] The infinity type of \(\gamma_4\) is the required fourth type, and its restriction is again \(\epsilon_{E/F}\). Use the common source splitting \(\chi_V\) for every line and the target splitting \(\gamma_i\) for that line. The source-splitting retwist in T6.12 changes only the source vector and preserves the actual rank-one theta form. For unintegrated kernels evaluated at the source identity, the retwisting character has value one. Crucially the two pairs have exactly the same total target splitting \(\chi_W\), not merely the same infinity type.

\textbf{2. The required real representation, including its exact Gaussian vector}

At \(\mu\) take the product of the two holomorphic one-form cocycles, with its nonzero determinant polynomial. At \(c_0\) take the product Gaussian. The same BMM and Haan central-action calculation as in T6.15 makes their source \(U(2)\) representations trivial.

At each of \(a,b\), the target is compact \(U(3)\) and the source is \(U(1,1)\). The total target splitting has zero angular exponent there. The common source splitting is calibrated by the all-negative choice \(\gamma_1^3\): in BMM's notation the twist is \(k=0,\ell=-3\). One can fix this calibration without any ambiguity in the notation for an angular character: at a compact target with a negative definite source plane, the same common splitting makes the Gaussian source invariant, and Proposition 4.15 gives its exponent \((\ell+3)/2\), hence \(\ell=-3\). At a mixed source the Gaussian therefore has maximal-compact character \[
(z_+,z_-)\longmapsto z_+^{-3},
\tag{T16.4}
\] where the positive source line is written first. Its character on the ordered pair \((W_1,W_2)\) is consequently \((0,-3)\).

Here is an elementary verification of the entire source representation generated by this vector. In the Fock space write \(x_1,x_2,x_3\) and \(y_1,y_2,y_3\) for the positive and negative variables, and put \[
R=\sum_{j=1}^3x_jy_j,
\quad E_+=R,
\quad E_-=-\sum_{j=1}^3\partial_{x_j}\partial_{y_j},
\quad H_0=\operatorname{Euler}+3.
\tag{T16.5}
\] After the usual Fock-unit normalization these are the source \(\mathfrak{sl}_2\) generators; BMM Proposition 5.1, printed p.40, gives the lowering operator, and the raising operator is its negative adjoint. Direct differentiation gives \[
[E_+,E_-]=H_0,\quad [H_0,E_\pm]=\pm2E_\pm,
\qquad E_-R^j=-j(j+2)R^{j-1}.
\tag{T16.6}
\] Thus the Gaussian generates \(\mathbf C[R]\), whose genuine \(U(1)\times U(1)\) weights are \((-3-j,j)\), \(j\ge0\), and whose squared monomial norms are \(j!(3)_j\) times the squared vacuum norm. The lowering coefficients are nonzero for every \(j>0\), so this is an irreducible highest-weight module. It is the discrete series \(D_3^-\) specified by these weights. For example, its disk model has reproducing kernel \((1-z\bar w)^{-3}\), or equivalently weighted area norm \(\int_{|z|<1}|f(z)|^2(1-|z|^2)\,dx\,dy\). The automorphy factor of exponent three gives the vacuum coefficient of modulus \((\cosh t)^{-3}\) on a Cartan element with entries \(\cosh t,\sinh t\). Its square is integrable against the Cartan density proportional to \(\sinh(2t)\,dt\). This also verifies square integrability directly.

The real source representation to be globalized is therefore \[
\pi_\infty=1_{\mu}\boxtimes D_{3,a}^-\boxtimes D_{3,b}^-\boxtimes1_{c_0}.
\tag{T16.7}
\] Its chosen archimedean vector is the product of the invariant vectors and the two lowest compact-type vectors. That compact-type space is one dimensional.

We also check the extra base-change hypothesis which was automatic in T6.15. For the derived \(\mathfrak{sl}_2\) at a mixed place, use the Casimir \(C=H_0^2-2H_0+4E_+E_-\). It acts by \(3\). Since \(E_-=-E_+^*\), a unit vector of \(H_0\)-weight \(m\) in any unitary representation with this infinitesimal character satisfies \[
\|E_-v_m\|^2=\frac{(m-3)(m+1)}4,
\qquad
\|E_+v_m\|^2=\frac{(m-1)(m+3)}4.
\tag{T16.8}
\] The fixed central character imposes odd parity. Weights \(1,-1\) are impossible by nonnegativity. Starting from any nonzero compact-weight vector, the ladder identities and irreducibility force either the weights \(3,5,\ldots\), or \(-3,-5,\ldots\). They are exactly the two discrete-series members with this central and infinitesimal character. At the compact places the specified infinitesimal character permits only the trivial representation. Labesse's criterion in Remark 5.2, printed p.42, therefore gives hypothesis \((*)\) for (T16.7). This is an explicit verification; no assertion that arbitrary noncompact cohomological infinity types satisfy \((*)\) is used.

Finally the local doubling integral detects the \textbf{same} Gaussian vector. GQT §11.3, p.52, normalizes the Siegel--Weil section at \((h,1)\) as the oscillator matrix coefficient. In §11.6, p.54, its other factor is the \textbf{complex conjugate} of the source matrix coefficient (the overbar is visible in the original PDF). Identify the cyclic Gaussian module unitarily with \(D_3^-\) and choose its corresponding source vector of norm one. Schur orthogonality gives \[
Z_v(1/2)=\|\phi_v\|^2\int_{H(F_v)}
       |\langle\pi_v(h)f_v,f_v\rangle|^2\,dh
       =d(\pi_v)^{-1}\|\phi_v\|^2>0.
\tag{T16.9}
\] The compact-place version is the positive compact integral of T6.15. Thus every factor of this exact archimedean cocycle has nonzero local pairing.

\textbf{3. Globalization, base change, and a nonzero holomorphic rank-two lift}

Choose a finite split Steinberg place \(v_s\), a finite nonsplit Steinberg place \(v_i\), and two additional split places \(u_1,u_2\). Let \(S\) contain all real places, \(v_s,v_i\), and all finite bad places for the source, extension, characters and additive character, but exclude \(u_1,u_2\). At the finite places of \(S\), prescribe the unitary unramified-twist Steinberg orbits, as in T6.15. At infinity prescribe (T16.7).

Apply Shin Theorem 5.13, printed p.28, based on §4.3 and Theorem 4.8, pp.19--21. That version explicitly permits \(S\) to contain \textbf{all} infinite places and prescribes actual representations there. Its two auxiliary places can be chosen split. To make the local test precise, the characteristic function of the desired archimedean discrete-series tuple belongs to Shin's space \(\mathcal F\): its infinitesimal-character support is a point; twisting a representation of the full real group by an unramified unitary character does not change it, because its center is compact; and unitary parabolic induction from a proper real Levi cannot contain these strict discrete series. The latter assertion also follows from (T16.8) and the distinct infinitesimal characters of tempered principal series. Products with the finite discrete-series-orbit indicators of Example 5.6 remain in \(\mathcal F\). Their Plancherel integral is positive. Thus one may apply Theorem 4.8 directly to this indicator and use the proof of Theorem 5.13. This avoids an unproved selection of a member merely from a cohomological packet.

We obtain a cuspidal automorphic \(\pi\) with exactly (T16.7), Steinberg twists at all finite places of \(S\), unramified outside \(S\cup\{u_1,u_2\}\), square-integrable at \(u_1\), and supercuspidal at \(u_2\). Labesse Corollary 5.3 applies by the preceding verification of \((*)\). Its local compatibility at the split place \(v_s\) implies that its rank-two isobaric base change \(\Pi\) is cuspidal, by exactly the rank-two argument in T6.15. Theorem 5.9 then identifies all finite local base changes, since the source is quasi-split at every finite place.

Choose the dimension-three target with signature \((2,1)\) at \(\mu\), and positive definite at \(a,b,c_0\). At the nonsplit bad places Atobe--Gan Theorem D.2(i) permits both target spaces for a Steinberg source: its Weil--Deligne parameter has no \(\chi_V\otimes S_1\) summand. At every split place, including both auxiliary places, Mínguez's theorem makes the lift nonzero for every source representation. At the remaining places use the unramified theta lift. Correct global Hermitian coherence at \(v_i\), where both choices are allowed. The resulting \(V\) is anisotropic over \(F\), because it is definite at three real places.

The fixed real cocycle survives by (T16.9). Choose the finite bilinear local entries by GQT Proposition 11.6. The twisted unitary cuspidal \(\mathrm{GL}_2\) value at \(1\) is nonzero by the theorem cited in T6.15. GQT Theorem 11.4, with \(n=2,m=3,r=0\), therefore gives a nonzero global theta lift with this \textbf{fixed} real cocycle. Fix an actual source realization and a corresponding actual irreducible target realization \(\sigma\). The latter is irreducible by Wu's theorem as used in T6.13. At \(v_s\), its local parameter has nonzero monodromy by Mínguez's formula. These are precisely the compatibility conditions checked in T6.15, with the real step replaced by (T16.7)--(T16.9).

\textbf{4. A total 032 theta input retaining this constituent}

We first obtain a total unintegrated theta form with nonzero \(\sigma\) projection. Let \(\mathcal K_{\phi}(g,h)\) denote the rank-two kernel with the fixed real cochain and variable finite Schwartz input. Project in the source variable to the chosen actual realization of \(\pi\). The span of all scalar canonical-frame coefficients \[
\mathcal I=\operatorname{span}_{\phi_f,g}
 P_\pi\mathcal K_{\phi_\infty\otimes\phi_f}(g,-)
\tag{T16.10}
\] is nonzero by the nonzero global lift. It is stable under \(H(\mathbf A_f)\), by replacing \(\phi_f\) with its source translate. Its archimedean compact-type space is the one-dimensional space just specified. Thus irreducibility of the smooth finite representation gives \(\mathcal I=\pi_f^\infty\otimes\ell_\infty\).

Evaluation at the source identity is nonzero on this space after a suitable finite translate. Here is the needed density argument, since \(\pi_\infty\) is no longer trivial. The Cayley transform from skew-adjoint endomorphisms to \(U(W)\) gives a rational open chart. Weak approximation in its affine coordinates shows that \(H(F)\) is dense in \(H(F_\infty)\). For a point represented by \((h_\infty,h_f)\), choose \(q_n\in H(F)\) with \(q_{n,\infty}\to h_\infty\). The classes of \((1,q_{n,f}^{-1}h_f)\), represented also by \((q_{n,\infty},h_f)\), approach that point. Consequently the image of \(H(\mathbf A_f)\) is dense in \([H]\). A nonzero continuous automorphic vector in \(\pi_f^\infty\otimes\ell_\infty\) is nonzero at some such finite point; a finite right translate then has nonzero value at the identity and still lies in the same archimedean line.

It follows from (T16.10) that some projected kernel has nonzero value at \(h=1\). The mixed-source orthogonality and finite holomorphic spectral support established in T6.12 imply that this yields a nonzero target \(\sigma\) contribution in \(\mathcal K_\phi(-,1)\); the other source realization terms cannot cancel it. At infinity this unintegrated form is exactly the double-alternating degree-two theta tensor of 032, up to its allowed nonzero constant-line trivialization. Indeed the compact-place factors are its Gaussians and the active factor is the determinant obtained from the product of the two first tensors; this is the explicit coefficient calculation in 032, Proposition 4.6(i), pp.23--24, and BMM Proposition 5.19. Thus choose a finite total input for which \[
\xi=\Xi_W(\phi;t_0),\qquad P_\sigma\xi\ne0.
\tag{T16.11}
\] Take a sufficiently deep neat finite level fixing this input. One may work with the finite disjoint union of connected components of the full unitary Shimura variety. Each component admits a further principal special-unitary cover of the kind used in 032: the integral determinant subgroup is finite and is killed by a sufficiently deep neat level. Only finitely many components and inputs will be needed.

\textbf{5. A Hecke projector preserves the two pair-type spans}

At the fixed finite level let \(\mathcal H\) be the finite-dimensional space of holomorphic top forms on the full Shimura variety. The image of the finite-level Hecke algebra on \(\mathcal H\) is a finite-dimensional star algebra, with adjoints for the Petersson pairing. Its central orthogonal idempotents project to its irreducible isotypic spaces. In particular there is a finite Hecke operator \(Q\) equal to the orthogonal projection to the \(\sigma_f\)-isotypic space. This includes all its actual multiplicity copies, which is sufficient.

For precision, an irreducible smooth finite-adelic representation with nonzero level invariants is determined by its irreducible module for that level Hecke algebra: the full representation is generated by those invariants, and a module isomorphism extends through compact induction and its irreducible quotient. All target representations occurring in \(\mathcal H\) have the same holomorphic top-degree real representation, with trivial compact factors. Equivalently their scalar canonical coefficients are annihilated by the antiholomorphic lowering algebra and have the prescribed canonical compact type, whose irreducible highest-weight quotient is unique. Hence the relevant Hecke isotypic space consists of copies of the same full target representation \(\sigma\). At the Steinberg place every one of these copies has nonzero local monodromy. In particular \(Q\xi\ne0\).

Let \(\mathcal U_{12}\) be the span of products of theta one-forms of the first two line types, allowing all finite Schwartz inputs, and define \(\mathcal U_{34}\) similarly. A finite target translate of a product is the product of the same simultaneous translates of its two factors. Convolution by a compactly supported bi-level-invariant Hecke function is a finite sum of these translates on a fixed smooth vector. It follows that every finite Hecke operator, in particular \(Q\), preserves each of these two pair-type spans. This statement does not assert that it preserves one individual product. The translated individual factors may require a deeper common level than the final Hecke output; the finite expansions are first understood adelically, and all their factors descend together at the principal level chosen in the next step.

Apply 032, Proposition 4.6(iii) and Corollary 4.7, to the single input chosen in (T16.11). The rational decompositions \(W=W_3\perp W_4\) with signs \(d_3,d_4\) approach the first decomposition in total compact-place period data, with the active data exactly fixed. The global splitting choices agree by (T16.3). On the fixed quotient the resulting \(\xi'\) approach \(\xi\) smoothly. The operator \(Q\) is bounded, so \[
\langle Q\xi,Q\xi'\rangle\longrightarrow\|Q\xi\|^2>0.
\tag{T16.12}
\] Choose a sufficiently close rational endpoint. Both endpoint forms are holomorphic, and the finite wedge expansions of 032 give \[
x=Q\xi\in\mathcal U_{12},\qquad
y=Q\xi'\in\mathcal U_{34},\qquad
\langle x,y\rangle\ne0,
\tag{T16.13}
\] with \(x,y\) in the full \(\sigma\)-isotypic space. Continuity can be applied on the finitely many components at once; the same Gaussian estimates prove it on each compact component, so their finite sum has (T16.12).

\textbf{6. Extract actual CM maps without losing the first constituent}

Choose a connected component \(S\) on which the pairing of \(x|_S,y|_S\) is nonzero. There are only finitely many terms in their translated line-theta expansions. Apply 032, Corollary 6.7, printed p.45, to these finite families for all four sign functions and pass to one connected common principal cover \(Y\to S\). The pairing remains nonzero because it is multiplied by the degree of the finite étale cover. At the new full finite level the pullback of the global form \(y\) is still in the same full automorphic \(\sigma\)-isotypic space; pulling back has not changed the scalar adelic function.

Fix nonzero generators of the four selected eigenlines. The finite-family CM extraction and equality of harmonic representatives on a compact Kähler variety write \[
x|_Y=\sum_r a_r A_r,\quad
A_r=h_{1r}^*\alpha_1\wedge h_{2r}^*\alpha_2;
\qquad
y|_Y=\sum_s b_s B_s,\quad
B_s=h_{3s}^*\alpha_3\wedge h_{4s}^*\alpha_4,
\tag{T16.14}
\] where the \(h_{ir}\) denote rational Hodge maps, or corresponding rational homomorphisms from the Albanese. Since \(\langle x|_Y,y|_Y\rangle\ne0\), some \(r\) satisfies \(\langle A_r,y|_Y\rangle\ne0\). The zero extension \(A_r^0\) to the new full Shimura variety then has nonzero \(\sigma\)-projection, because its geometric pairing with the global \(\sigma\)-form \(y\) is that same nonzero number. If scalar Petersson conventions are used instead, T6.10 multiplies this pairing by the fixed positive Haar/geometric conversion factor, so nonvanishing is unchanged. Next expand only \(y|_Y\) in (T16.14); some \(s\) satisfies \(\langle A_r,B_s\rangle\ne0\).

This sequential choice is necessary: a nonzero projected pairing alone would not justify choosing an individual pair whose unprojected period and projected first wedge are both nonzero. Here the first choice guarantees survival of its constituent before the second choice supplies the full period.

Clear the four rational-homomorphism denominators. By 032, Lemma 2.2 and Proposition 3.1, this gives actual morphisms \(f_i:Y\to B(v_i)\). Multiplication by the four nonzero integers scales each selected wedge and period by nonzero factors. It preserves the first wedge's nonzero \(\sigma\)-projection and proves (T16.1).

Finally choose a product ample polarization on \(\prod_i B(v_i)\). In each factor its chosen eigenform can be rescaled to be the first vector of a unit invariant coframe and then extended to such a coframe. The first two factors yield a coordinate wedge proportional to \(A_r\). The exact degree identity T6.1 therefore has a term with nonzero \(\sigma\)-projection; T6.13 can use the rank-two theta model already constructed for that projection. Mínguez's split-place formula prevents any rank-two source with zero monodromy from giving this same target local representation. A one-dimensional Weil--Deligne parameter, its direct sum with another, or its quadratic induction has zero monodromy, while the actual source has Steinberg monodromy. The first wedge is nonzero, so \(f\) has surface image. This proves all four assertions. QED.

\textbf{Exact additional source locators.}

\begin{itemize}
\tightlist
\item
  Shin, \url{https://math.berkeley.edu/~swshin/Plancherel.pdf}: Definition 2.1 and Example 2.3, pp.8--9; §4.3 and Theorem 4.8, pp.19--21; Example 5.6, p.26; §5.3 and Theorem 5.13, p.28. The crucial version includes infinity in the prescribed set. Its full short proof and the scope of §4.3 have been read.
\item
  Labesse, \url{https://webusers.imj-prg.fr/~michael.harris/CBCM.pdf}: Remark 5.2, p.42, Corollary 5.3, pp.43--44, Theorem 5.9, p.49. The new elementary Casimir computation supplies the specific sufficient condition in Remark 5.2.
\item
  BMM, \url{https://archive.ymsc.tsinghua.edu.cn/pacm_download/117/6812-11511_2016_Article_136.pdf}: Proposition 4.15, p.38, and its adjoining proof; Proposition 5.1, p.40; Proposition 5.19, pp.49--50. The local real proof above imports the Fock action, then computes its cyclic module explicitly.
\item
  GQT, \url{https://arxiv.org/pdf/1207.4709}: §11.3, p.52, and §11.6, p.54. Original p.54 visually checked for the overbar on the source matrix coefficient. The remaining local/global source locators are those already in T6.15.
\item
  Gan--Gross--Prasad, \url{https://www.numdam.org/item/AST_2012__346__111_0.pdf}: §15, printed p.161, general Hermitian classification and product exact sequence. The new signature prescription uses two definite and two indefinite real places.
\item
  Supplied 032: §4.1--4.2, pp.13--14; Proposition 4.6 and Corollary 4.7, pp.23--25, entire proofs reread; Corollary 6.7, p.45, entire proof reread. Their fixed-input continuity and finite CM extraction are imported with their stated hypotheses.
\end{itemize}

The entire proof was read and checked by this agent and independently by root and the bounds agent. The checks covered the real splitting calibration, direct Shin indicator application, fixed archimedean-line source spanning, identification with the 032 total tensor, finite Hecke projector and preservation of pair-type spans, and the component/cover/zero-extension order in the last extraction. The listed primary inputs were read at their exact locators. These checks import the cited background theorems; they are not reproofs of the trace formula or theta-correspondence theorems. T6.15 is retained with its original narrower scope.

\subsection{12. Positive analytic normalization in a fixed CM isogeny construction}\label{positive-analytic-normalization-in-a-fixed-cm-isogeny-construction}

\subsubsection{DEG-L.T6.17 --- A fixed CM tensor admits bounded analytic normalization}\label{DEG-L.T6.17}

For \(1\le i\le4\), fix a CM field \(E_i\) of degree \(2g_i\), a CM type \(\Phi_i\) (not necessarily primitive), an invertible fractional \(\mathcal O_{E_i}\)-ideal \(I_{i,0}\), and a polarized abelian variety \[
A_{i,0}=\mathbf C^{g_i}/\Phi_i(I_{i,0}),\qquad L_{i,0},\qquad
V_{i,0}=\chi(L_{i,0}).
\tag{T17.1}
\] Assume the \(E_i\) action has rank one on rational homology and the Rosati involution restricts to complex conjugation. Fix a smooth projective surface \(Y\), a Kähler form \(\omega_Y\), actual morphisms \(f_{i,0}:Y\to A_{i,0}\), and one holomorphic CM eigenline in each factor. Choose polarization-unit invariant eigenforms \(\beta_{i,0}\) in these four lines and suppose that \[
p_0=\int_Y f_{1,0}^*\beta_{1,0}\wedge f_{2,0}^*\beta_{2,0}
            \wedge\overline{f_{3,0}^*\beta_{3,0}\wedge f_{4,0}^*\beta_{4,0}}
\ne0.
\tag{T17.2}
\] The metric normalization and the chosen lines are fixed starting data.

Let the targets \((A_i,L_i)\) vary through the corresponding \(E_i\)-linear CM isogeny classes, with the same Rosati condition. In compatible CM coordinates write \(A_i=\mathbf C^{g_i}/\Phi_i(\Lambda_i)\), put \(V_i=\chi(L_i)\), and set \[
\Lambda_i^- =\{x\in\Lambda_i:\mathcal O_{E_i}x\subset\Lambda_i\},\quad
J_i=[\Lambda_i:\Lambda_i^-],\quad
x_i=\frac{V_i}{V_{i,0}}J_i,\quad t_i=x_i^{1/g_i}.
\tag{T17.3}
\] Define field constants \[
K_i=\frac2\pi\bigl(g_i!\sqrt{|D_{E_i}|}\bigr)^{1/g_i},\qquad
c_i=
\begin{cases}
1,&g_i=1,\\
\bigl((g_i-1)/K_i\bigr)^{g_i-1},&g_i\ge2.
\end{cases}
\tag{T17.4}
\] There exist actual \(E_i\)-linear isogenies \(h_i:A_{i,0}\to A_i\) and positive integers \(M_i\) such that the actual maps \[
f_i=[M_i]\circ h_i\circ f_{i,0},\qquad f=(f_i):Y\to\prod_iA_i
\tag{T17.5}
\] have a nonzero selected four-factor period and uniformly bounded period-normalized energy. Precisely, write \(\omega_i,\omega_{i,0}\) for invariant polarization representatives, let \(\beta_i\) be any polarization-unit eigenform in the corresponding target line, and put \[
\begin{aligned}
p&=\int_Yf_1^*\beta_1\wedge f_2^*\beta_2
                   \wedge\overline{f_3^*\beta_3\wedge f_4^*\beta_4},\\
E&=\int_Y\Bigl(\sum_i f_i^*\omega_i\Bigr)\wedge\omega_Y,\qquad
E_{i,0}=\int_Y f_{i,0}^*\omega_{i,0}\wedge\omega_Y,\\
C_E&=\frac{4\sum_iK_iE_{i,0}}
               {|p_0|^{1/2}\bigl(\prod_i c_i\bigr)^{1/4}},\qquad
V_Y=\int_Y\omega_Y^2.
\end{aligned}
\tag{T17.6}
\] Then \[
p\ne0,\qquad \frac{E}{|p|^{1/2}}\le C_E,
\qquad
\frac{\deg_{\boxtimes_iL_i}(f_*[Y])}{|p|}
\le\frac{C_E^2}{V_Y}.
\tag{T17.7}
\] The constants depend only on the fixed fields, polarized source factors, source surface and its four maps; they are independent of the target lattices, their endomorphism orders, and their polarization volumes. The denominator in (T17.7) is the specified metric period \(p\). These inequalities do not bound the unnormalized integral degree.

If the finitely many source coordinate pullbacks lie in the span of a fixed finite basis of genuine rank-one theta forms on \(Y\), then the common normalization \(|p|^{-1/4}\) gives theta presentations of all target polarization coordinates with uniformly bounded coefficients and every fixed Schwartz seminorm. Any one finite cover needed for theta spanning is taken at the outset and included in the fixed data.

\textbf{Proof.}

\textbf{1. Every CM metric eigenvalue is controlled by a short integral isogeny}

For each factor choose a nonzero integral \(E_i\)-linear homomorphism minimizing the squared complex Hilbert--Schmidt norm \[
Q_i(h)=\operatorname{tr}_{\omega_{i,0}}(h^*\omega_i).
\] The ideal-lattice argument of B19.1--B19.2 applies without primitivity. Indeed rank-one \(E_i\)-homology identifies every \(E_i\)-linear rational homomorphism with multiplication by an element of \(E_i\). The integral Hom lattice is \((\Lambda_i:I_{i,0})\), and its product with \(I_{i,0}\) is \(\Lambda_i^-\). The Rosati condition diagonalizes the polarization metrics. None of these statements uses simplicity. Every nonzero such field scalar is an isogeny. The exact covolume and minimum bound are therefore \[
\operatorname{covol}_{Q_i}\operatorname{Hom}_{E_i}(A_{i,0},A_i)
=2^{-g_i}\sqrt{|D_{E_i}|}\,x_i,
\qquad Q_i(h_i)\le K_i t_i,
\qquad \deg h_i\ge J_i.
\tag{T17.8}
\] Here the last inequality is the elementary inclusion \(h_i I_{i,0}\subset\Lambda_i^-\). The constant in the middle inequality is the Minkowski ball constant computed in B19.2.

Let \(\lambda_{i,\sigma}\) be the positive eigenvalues of \(h_i^*\omega_i\) relative to \(\omega_{i,0}\) in the CM eigenlines. They satisfy \[
\sum_\sigma\lambda_{i,\sigma}\le K_i t_i,
\qquad
\prod_\sigma\lambda_{i,\sigma}
=\deg h_i\frac{V_i}{V_{i,0}}\ge x_i=t_i^{g_i}.
\tag{T17.9}
\] For \(g_i\ge2\), arithmetic--geometric mean on the other \(g_i-1\) eigenvalues gives \[
\lambda_{i,\sigma}
\ge\frac{t_i^{g_i}}
           {(K_i t_i/(g_i-1))^{g_i-1}}
=c_i t_i.
\] For \(g_i=1\) the product inequality gives the same bound with \(c_i=1\). The trace gives the upper bound. Thus \[
c_i t_i\,\omega_{i,0}
\le h_i^*\omega_i\le K_i t_i\,\omega_{i,0}.
\tag{T17.10}
\] This is a bound in every CM eigenline, not merely a trace estimate.

\textbf{2. Integral multiplication balances the four factors}

Put \[
T=\max_i t_i,\qquad
M_i=\left\lceil\sqrt{T/t_i}\right\rceil.
\tag{T17.11}
\] Since \(T/t_i\ge1\), \[
T\le M_i^2t_i\le4T.
\] Multiplication by \(M_i\) is an actual endomorphism, so every map in (T17.5) is integral. The corresponding metric eigenvalues satisfy \[
c_iT\le M_i^2\lambda_{i,\sigma}\le4K_iT.
\tag{T17.12}
\] Pullback preserves the inequalities of positive forms, and integration on \(Y\) gives \[
T\sum_i c_i E_{i,0}\le E\le4T\sum_iK_iE_{i,0}.
\tag{T17.13}
\]

Because the chosen period is a \textbf{fixed eigenmonomial}, there is no sum of potentially canceling scalars under this transport. Write \[
([M_i]h_i)^*\beta_i=b_i\beta_{i,0}.
\] The rank-one eigenline and the polarization-unit convention give \[
|b_i|^2=M_i^2\lambda_{i,\sigma_i},\qquad
p=b_1b_2\overline{b_3b_4}\,p_0.
\tag{T17.14}
\] Consequently \[
|p_0|\bigl(\prod_i c_i\bigr)^{1/2}T^2
\le |p|
\le16|p_0|\bigl(\prod_iK_i\bigr)^{1/2}T^2.
\tag{T17.15}
\] The lower bound proves nonvanishing and, together with the upper energy bound in (T17.13), gives the first inequality in (T17.7). Set \(\beta=\sum_i f_i^*\omega_i\). The Hodge index inequality on the fixed Kähler surface gives \[
\int_Y\beta^2\le\frac{(\int_Y\beta\wedge\omega_Y)^2}{V_Y}.
\] The left side is the degree of the effective pushforward, with its actual generic multiplicity, by T6.1. Dividing by \(|p|\) proves the second inequality in (T17.7). No division by an unverified generic degree has occurred. Also \(p\ne0\) implies that the first pulled-back wedge is nonzero, so \(f(Y)\) is a surface.

\textbf{3. A single common analytic normalization bounds every theta coordinate}

Choose a polarization-unit invariant CM coframe in every target. Multiply \textbf{all} its pulled-back one-forms by \[
r=|p|^{-1/4}.
\] The selected four-factor period then has absolute value one. The sum of the coordinate energies is exactly \(r^2E\), hence at most \(C_E\). This is a common scalar normalization of the four actual maps constructed above; the preceding balancing was by integral endomorphisms.

Now assume the fixed source coordinate pullbacks have a finite theta expansion, as supplied by T6.9 in its stated setting. Since each \(h_i\) is diagonal in the CM eigenlines, all varying target coordinate pullbacks lie in the same fixed finite-dimensional space. Take the span of the finitely many actual theta forms used in expansions of its fixed generators; this is a possibly larger finite-dimensional space containing all the pullbacks. Select a basis \(\theta_1,\ldots,\theta_q\) from those actual theta generators. Let \(G\) be its positive energy Gram matrix, using the convention \[
\mathcal E(\eta)=\int_Y\frac i2\eta\wedge\bar\eta\wedge\omega_Y.
\] If the normalized coordinate forms are \(\eta_j=\sum_k a_{jk}\theta_k\), then \[
\lambda_{\min}(G)\sum_{j,k}|a_{jk}|^2
\le\sum_j\mathcal E(\eta_j)
=r^2E\le C_E.
\tag{T17.16}
\] Fix once and for all one theta presentation for each basis vector. Linear combination gives a linear section from this finite coefficient space to a finite collection of Schwartz/source-vector spaces. Every fixed seminorm on its finite-dimensional image is bounded by a constant times the coefficient norm. Thus the normalized inputs have uniformly bounded seminorms, including the weighted local shell seminorms of B14a at the finitely many nonstandard places. The standard tails, global characters, and levels come from a fixed finite list. This proves the analytic-presentation assertion. QED.

\textbf{Exact scope and relation to the arithmetic construction.}

The initial data required here exist for the explicit T6.16 witness. Each of its four CM factors can first be replaced, once and for all, by an \(E\)-linearly isogenous factor with maximal-order action; composing the initial maps with those fixed isogenies preserves the nonzero selected period and the first wedge's automorphic constituent. The two imprimitive types in (T16.2) cause no difficulty because the proof above only uses rank-one \(E\)-homology, not simplicity. A single further fixed cover may supply the theta spanning statement. Thereafter the target lattices and their conductors may vary arbitrarily inside these four fixed \(E\)-linear isogeny classes.

The multiplication factors \(M_i\) do not introduce a new normalized divisor-kernel Chow class. Under the exact birigidified-kernel hypotheses of T6.5, they multiply both the kernel output and its selected \textbf{algebraically normalized Betti period} by \(\prod_iM_i\), so the quotient cancels exactly. This assertion only concerns the integral multiplications at a fixed target; it does not claim cancellation of the varying isogenies \(h_i\).

The period \(p\) in (T17.7) is instead defined by polarization-unit eigenforms. It need not equal the Betti-algebraic coefficient used to divide a Chow correspondence. The final common scalar \(|p|^{-1/4}\) generally changes the integral Betti lattice. Accordingly the theorem supplies bounded analytic theta presentations in a fixed-source construction and a uniform ratio of degree to metric period. It does not supply a conductor-independent effective integral degree, a bounded denominator for the Betti-algebraic normalization, or bounds uniform over varying CM fields or initial witnesses. In particular T6.16's actual nonmonomial rank-two constituent remains compatible with the finite rank-one Hecke-energy upper bound: the latter is an inequality using all normalized coordinates, not an identity that replaces the rank-two spectral factor by a character.

\textbf{Whole-proof check.}

The proof imports the checked exact CM Hom covolume and Minkowski bound of B19.1--B19.2, and the already checked geometric degree/Hodge-index and theta-spanning statements T6.1/T6.9. The entire proof was read by this agent and independently by root and the bounds agent. These checks recomputed the lower bound for each metric eigenvalue, the integer balancing inequalities, the noncanceling eigenmonomial period, the common fourth-root normalization, and the distinction between metric and Betti period. All three reviews passed before promotion to this module.

\subsection{13. Exact integral coefficient content of the four-factor switch}\label{exact-integral-coefficient-content-of-the-four-factor-switch}

\subsubsection{DEG-L.T6.18 --- Exact integral eigenline content}\label{DEG-L.T6.18}

Use a fixed Galois CM field \(E\), \(\mathcal O=\mathcal O_E\), and four CM types as in T6.16. Fix reference factors \[
A_{i,0}=\mathbf C^s/\Phi_i(I_{i,0})
\] with invertible fractional \(\mathcal O\)-ideals \(I_{i,0}\), and actual maps \(f_{i,0}:Y\to A_{i,0}\) from a fixed smooth projective surface. Assume their selected \((1,1,c,c)\) four-factor Betti period is nonzero. Let \[
A_i=\mathbf C^s/\Phi_i(\Lambda_i),\qquad
\Lambda_i^+=\mathcal O\Lambda_i,\qquad
\Lambda_i^-=\{x\in\Lambda_i:\mathcal O x\subset\Lambda_i\},
\] where the full target lattices need not be \(\mathcal O\)-stable. Put \[
I_i=(\Lambda_i:I_{i,0})=\Lambda_i^-I_{i,0}^{-1},\qquad
\mathfrak r_i=\Lambda_i^-(\Lambda_i^+)^{-1}\subseteq\mathcal O,
\quad J_i=[\Lambda_i:\Lambda_i^-],\quad H_i=[\Lambda_i^+:\Lambda_i].
\tag{T18.1}
\]

Identify \(H^1(A_i,E)\) with \(\operatorname{Hom}_{\mathbf Q}(E,E)\). For \(\tau\in\operatorname{Gal}(E/\mathbf Q)\), write \(e_{i,\tau}(x)=\tau(x)\), and define the \textbf{actual integral coefficient eigenline} \[
\mathcal L_{i,\tau}
=Ee_{i,\tau}\cap\bigl(H^1(A_i,\mathbf Z)\otimes_{\mathbf Z}\mathcal O\bigr).
\tag{T18.2}
\] Use \(\mathcal L_{i,0,\tau}\) for the corresponding reference lattice. Then \[
\begin{aligned}
\mathcal L_{i,\tau}
&=\tau(\Lambda_i^+)^{-1}e_{i,\tau},\\
\sum_{h_i\in\operatorname{Hom}_E(A_{i,0},A_i)}h_i^*\mathcal L_{i,\tau}
&=\tau(\mathfrak r_i)\mathcal L_{i,0,\tau},\\
[\mathcal L_{i,0,\tau}:\sum_{h_i}h_i^*\mathcal L_{i,\tau}]
&=N_{E/\mathbf Q}(\mathfrak r_i)=J_iH_i.
\end{aligned}
\tag{T18.3}
\] The sums are additive subgroup sums; each image is already an \(\mathcal O\)-submodule. If the target polarization is principal and its Rosati involution restricts to conjugation on \(E\), then \(H_i=J_i\), so the last index is \(J_i^2\).

Let \(\mathfrak P_0\subseteq\mathcal O\) be the additive group generated by the selected four-factor periods of the fixed maps \(f_{i,0}\) as the four eigenvectors range through \[
\mathcal L_{1,0,1},\quad\mathcal L_{2,0,1},\quad
\mathcal L_{3,0,c},\quad\mathcal L_{4,0,c}.
\] It is a nonzero integral ideal. Let \(\mathfrak P_A\) be the additive group of periods obtained by varying all actual integral \(E\)-linear maps \(h_i:A_{i,0}\to A_i\) and all eigenvectors in the four corresponding target integral coefficient eigenlines, with maps \(h_if_{i,0}\). Then \[
\mathfrak P_A
=\mathfrak r_1\mathfrak r_2\overline{\mathfrak r_3}\,
             \overline{\mathfrak r_4}\,\mathfrak P_0,
\qquad
[\mathfrak P_0:\mathfrak P_A]=\prod_{i=1}^4J_iH_i.
\tag{T18.4}
\] For principal compatible targets this index is \(\prod_iJ_i^2\). These are exact ideal and additive-index identities, not just bounds on absolute values under one complex embedding.

\textbf{Proof.}

Integral first cohomology is \(\operatorname{Hom}_{\mathbf Z}(\Lambda_i,\mathbf Z)\). Since \(\Lambda_i\) is finite free over \(\mathbf Z\), tensoring with \(\mathcal O\) identifies its integral coefficient lattice with \(\operatorname{Hom}_{\mathbf Z}(\Lambda_i,\mathcal O)\). Thus \(u e_{i,\tau}\) is integral exactly when \[
u\tau(\Lambda_i)\subseteq\mathcal O.
\] Because \(\mathcal O\) is a ring, this is equivalent to \(u\tau(\mathcal O\Lambda_i)\subseteq\mathcal O\). The maximal-order lattice \(\Lambda_i^+\) is an invertible fractional ideal, so this condition is exactly \(u\in\tau(\Lambda_i^+)^{-1}\). This proves the first formula of (T18.3), with no discriminant denominator or saturation assumption.

The map corresponding to \(\alpha\in I_i\) acts on the selected eigenline by \[
h_\alpha^*e_{i,\tau}=\tau(\alpha)e_{i,0,\tau}.
\] The additive sum of the pullback lattices is therefore \[
\begin{aligned}
\tau(I_i)\tau(\Lambda_i^+)^{-1}e_{i,0,\tau}
&=\tau(\Lambda_i^-(\Lambda_i^+)^{-1})
  \tau(I_{i,0})^{-1}e_{i,0,\tau}\\
&=\tau(\mathfrak r_i)\mathcal L_{i,0,\tau}.
\end{aligned}
\] Taking the index in a rank-one projective \(\mathcal O\)-module gives \(N(\mathfrak r_i)\). The usual ideal-index formula gives \[
N(\mathfrak r_i)=[\Lambda_i^+:\Lambda_i^-]=H_iJ_i.
\] These identities hold for fractional ideals as well as integral ones.

For the principal-polarization assertion, use the alternating Riemann form to define the dual lattice \(M^\#\). Principal polarization says \(\Lambda_i^\#=\Lambda_i\). Rosati conjugation gives \(\langle ax,y\rangle=\langle x,\bar a y\rangle\) for every \(a\in E\), including rational endomorphisms. Since \(\bar{\mathcal O}=\mathcal O\), \[
(\mathcal O\Lambda_i)^\#=(\Lambda_i^\#)^-=\Lambda_i^-.
\] Duality reverses lattice indices, so \[
H_i=[\Lambda_i^+:\Lambda_i]
=[\Lambda_i^\#:(\Lambda_i^+)^\#]
=[\Lambda_i:\Lambda_i^-]=J_i.
\] This proves all of (T18.3).

For the period statement, put \[
k_0=\int_Y f_{1,0}^*e_{1,0,1}\wedge f_{2,0}^*e_{2,0,1}
                  \wedge f_{3,0}^*e_{3,0,c}\wedge f_{4,0}^*e_{4,0,c}
\in E^\times.
\tag{T18.5}
\] The maps have integral pullback on first cohomology and the fundamental class is integral, so evaluation of four integral \(\mathcal O\)-coefficient classes belongs to \(\mathcal O\). Multilinearity gives the exact reference ideal \[
\mathfrak P_0
=k_0 I_{1,0}^{-1}I_{2,0}^{-1}
          \overline{I_{3,0}}^{-1}\overline{I_{4,0}}^{-1}
\subseteq\mathcal O.
\tag{T18.6}
\] It is nonzero by the initial period hypothesis.

Each varying pulled-back eigenvector lies in the transport module in (T18.3). Conversely every element of that transport module is a finite sum of actual pullbacks. Expanding a fourfold product of such sums shows that the additive period group is exactly the image of the four transport modules under the fixed multilinear functional. Substituting (T18.3) proves the ideal identity (T18.4). This is equality for the additive span actually specified, not an enlargement introduced without justification: multiplication of one allowed target eigenvector by \(a\in\mathcal O\) already gives another allowed vector. Even with fixed target eigenvectors, varying \(h_1\) by precomposition with the actual reference endomorphism \([a]\) multiplies its period by \(a\); the maximal-order reference action is used here. The ideal-index formula now proves the final equality. QED.

\textbf{Fixed compatible generators versus integral target eigenvectors.}

If the four target rational eigenvectors are fixed to be the universal \(e_{i,\tau_i}\), without imposing integrality, the additive span of their transported periods is instead \[
k_0 I_1I_2\overline{I_3}\,\overline{I_4}.
\tag{T18.7}
\] Passing from (T18.7) to the genuinely integral eigenline statement inserts the four dual hull ideals. Thus neither a choice of complex unit forms nor a compatible rational basis can be substituted for the integral target eigenvectors without recording that change of lattice.

In T6.16 all source and target eigenline coefficients can be kept in this same fixed Galois CM field \(E\): the maps selected after CM extraction are rational Hodge maps and, after denominator clearing, integral morphisms; the eigenvectors are the \(E\)-valued embedding functionals above. The auxiliary theta spectral projector may require a larger field for its own description, but it does not enlarge the field of the final four-factor Betti coefficient \(k_0\). No rationality of an analytic theta normalization is asserted.

\textbf{What this does and does not force on integer normalization.}

Fix \(0\ne\kappa\in\mathfrak P_0\), and put \(d=[E:\mathbf Q]\) and \(c_\kappa=[\mathfrak P_0:\mathcal O\kappa]\). If an integral linear combination of the allowed period values equals \(N\kappa\) for a positive integer \(N\), then \[
\prod_iJ_iH_i\ \mid\ N^d c_\kappa,
\qquad
N\ge\left(\frac{\prod_iJ_iH_i}{c_\kappa}\right)^{1/d}.
\tag{T18.8}
\] Indeed \(\mathcal O(N\kappa)\subseteq\mathfrak P_A\subseteq\mathfrak P_0\), and multiplication by the rational integer \(N\) has index \(N^d\). This is a membership obstruction for the stated coefficient module. It is not a lower bound for arbitrary target cycles or for every possible trace construction with varying correspondences.

Even taking a trace to \(\mathbf Q\) can lose information from an ideal-index obstruction. For example, in \(\mathbf Z[i]\), all ideals \((2+i)^r\) have trace ideal \(2\mathbf Z\), although their indices are \(5^r\). To check this, write \((2+i)^r=a+bi\). Any common prime divisor of \(a,b\) must divide \(a^2+b^2=5^r\); divisibility by \(5=(2+i)(2-i)\) is impossible because the two Gaussian prime ideals are distinct. Hence \(\gcd(a,b)=1\), and the traces of \((a+bi)(u+vi)\) generate \(2\mathbf Z\). Therefore one cannot infer an unbounded rational trace denominator from (T18.4) alone. A trace argument must retain its full target tensor or cycle lattice.

\textbf{Relation with B23.}

B23 concerns middle-degree \textbf{pushforwards} on the rank-two generic \(K\)-Weil determinant line and includes its relative target determinant ideal. Its common content is the conjugate ideal \(\bar{\mathfrak j}\), of norm \(J\). The present result concerns four independent \textbf{pullbacks} into specified full-CM eigenlines and has the core/hull ideals \(\mathfrak r_i\). The auxiliary product in T6.16 has dimension sixteen and its tested switch tensor has degree four, so it is not itself B23's middle Weil plane. Passing from that switch to a specified generic Weil target must retain the actual correspondence and its integral lattice map. The two formulas explain different arithmetic costs and must not be identified without that additional map.

The ideal-index input is Milne, \emph{Algebraic Number Theory}, v3.08, Proposition 4.2, printed pp.69--70, \url{https://www.jmilne.org/math/CourseNotes/ANT.pdf}; its statement and complete short proof have been read. The Hom-ideal and Rosati duality identities are the checked B19/T6.17 identities, restated and proved in the needed form above. The entire saved proof was read by the author and independently by root and the bounds agent; all three reviews passed before promotion. The checks verified the actual integral coefficient lattice, the core/hull index and principal-polarization duality, the multilinear additive-span equality, and the scope of the integer and trace obstructions.

\subsection{14. Exact trace presentation and actual principal CM families}\label{exact-trace-presentation-and-actual-principal-cm-families}

\subsubsection{DEG-L.T6.19 --- Exact scalar and tensor trace presentation indices}\label{DEG-L.T6.19}

Let \(E\) be a Galois CM field of degree \(d=2g\), let \(\mathcal O\) be its ma\(\xi\)mal order, and let \(c\) be complex conjugation. Retain the four fixed CM types, ma\(\xi\)mal-order references, fixed surface, and nonzero four-factor Betti period of T6.18. Write \(\mathfrak P_A\) for its exact integral period ideal. The following statements hold.

\subparagraph{1. Scalar trace has an exact minimum-valuation formula}\label{scalar-trace-has-an-exact-minimum-valuation-formula}

Let \(\mathfrak a\) be any nonzero fractional \(\mathcal O\)-ideal. Write \(\mathfrak D\) for the different of \(E/\mathbf Q\) and write \[
 a_{\mathfrak p}=v_{\mathfrak p}(\mathfrak a),\qquad
 d_{\mathfrak p}=v_{\mathfrak p}(\mathfrak D),\qquad
 e_{\mathfrak p}=e(\mathfrak p/p).
\] The nonzero fractional ideal of \(\mathbf Q\) generated by the scalar traces is \[
 \operatorname{Tr}_{E/\mathbf Q}(\mathfrak a)
 =\prod_p p^{\,\min_{\mathfrak p\mid p}
       \lfloor(a_{\mathfrak p}+d_{\mathfrak p})/e_{\mathfrak p}\rfloor}\mathbf Z.
 \tag{T19.1}
\] In particular, when \(\mathfrak a\) is integral and its trace ideal is \(m\mathbf Z\), the least positive common rational denominator needed to obtain the scalar 1 from finite rational linear combinations of traces of elements of a is exactly m. This assertion concerns a scalar trace, not the full rational tensor obtained by tracing a coefficient-valued cycle.

\subparagraph{2. Principal CM lattices impose conjugation symmetry}\label{principal-cm-lattices-impose-conjugation-symmetry}

For a principally polarized CM target with rational E-action and Rosati conjugation, let \[
 \mathfrak r=\Lambda^-(\Lambda^+)^{-1}
\] be the core/hull ideal of T6.18. Then \[
 \overline{\mathfrak r}=\mathfrak r.
 \tag{T19.2}
\] Thus the Gaussian-prime ideal in T6.18's abstract trace example is not itself a possible principal-target core/hull ideal. This restriction does not force trace content to grow with its norm when \(g\ge2\).

\subparagraph{3. Two actual principal-polarization families}\label{two-actual-principal-polarization-families}

Fix, for each of the four CM types, any principally polarized lattice \(\Lambda_{i,*}\) with Rosati conjugation. Such a lattice e\(\xi\)sts in each of the relevant rational CM isogeny classes. Let \(\mathfrak P_*\) be the period ideal of T6.18 for these four targets. Choose a rational prime p that splits completely in E and avoids a fixed finite set: the discriminant, all places where the four base lattices are not ma\(\xi\)mal-order-stable, and the numerator and denominator ideals occurring in the fixed data. Additional fixed nonzero coefficients may also be required to be units at p.

For every integer \(r\ge0\) there are actual principally polarized target lattices \(\Lambda_{i,r}\), with the same CM types, the same rational E-actions, and the same polarization volumes, with either of the following properties.

\textbf{Uniform family.} Modify all \(g\) conjugate pairs at p in all four factors. Then \[
 \mathfrak r_{i,r}=p^r\mathfrak r_{i,*},\qquad
 J_{i,r}=p^{gr}J_{i,*},\qquad
 \mathfrak P_r=p^{4r}\mathfrak P_*.
 \tag{T19.3}
\] Consequently \[
 \operatorname{Tr}(\mathfrak P_r)
   =p^{4r}\operatorname{Tr}(\mathfrak P_*).
 \tag{T19.4}
\] The least denominator for scalar trace normalization therefore grows by exactly \(p^{4r}\).

\textbf{Concentrated family, \(g\ge2\).} Modify only one conjugate pair in the first factor and leave the other factors fixed. If \(\mathfrak q=\mathfrak p\,c(\mathfrak p)\) denotes the product of the corresponding two primes above p, then \[
 \mathfrak r_{1,r}=\mathfrak q^r\mathfrak r_{1,*},\qquad
 J_{1,r}=p^rJ_{1,*},\qquad
 \mathfrak P_r=\mathfrak q^r\mathfrak P_*.
 \tag{T19.5}
\] Although its additive index grows by \(p^{2r}\), \[
 \operatorname{Tr}(\mathfrak P_r)
    =\operatorname{Tr}(\mathfrak P_*).
 \tag{T19.6}
\] These are realized CM abelian varieties, rather than arbitrary ideal examples.

\subparagraph{4. Full four-factor trace retains the coefficient ideal}\label{full-four-factor-trace-retains-the-coefficient-ideal}

Fix output CM factors \(B_i\) of the four reference types, and fix nonzero integral \(\mathcal O\)-coefficient eigenvectors \(\beta_i\) whose types are opposite to those of the corresponding input eigenvectors \(\alpha_i\). Thus every scalar conjugate of \(\alpha_i\otimes\beta_i\) is of type (1,1). Put \[
 t=\beta_1\otimes\beta_2\otimes\beta_3\otimes\beta_4.
\] The tensor lies in the degree-four Kunneth summand of \(B=\prod_iB_i\). Its \(d\) scalar Galois conjugates lie in distinct eigenspaces and are linearly independent over \(E\). Let \(W\) be their rational descent and define \[
 \mathcal T_t:E\longrightarrow W,\qquad
 a\longmapsto\sum_{\sigma\in\operatorname{Gal}(E/\mathbf Q)}
                      \sigma(a)\sigma(t).
 \tag{T19.7}
\] Then \(\mathcal T_t\) is a \(\mathbf Q\)-linear isomorphism.

For each integral input eigenvector and each actual integral map admitted by T6.18, choose \(\mathcal O\)-coefficient divisor kernels with pure classes \(\alpha_i\otimes\beta_i\) and apply their product to the surface pushforward. The additive group of rational cohomology classes obtained by coefficientwise trace and integral linear combinations is exactly \[
 \mathcal T_t(\mathfrak P_A).
 \tag{T19.8}
\] In particular, integral \(\mathcal O\) multipliers inside the traces do not enlarge this group. For fixed \(0\ne\kappa\in\mathfrak P_*\) and \(w=\mathcal T_t(\kappa)\), the least positive common denominator of a rational linear combination of these traced classes that equals w is exactly the positive generator of \[
 (\kappa^{-1}\mathfrak P_A)\cap\mathbf Z.
 \tag{T19.9}
\] Equivalently its p-valuation is \[
 \max\left(0,\max_{\mathfrak p\mid p}
     \left\lceil\frac{v_{\mathfrak p}(\mathfrak P_A)
                           -v_{\mathfrak p}(\kappa)}{e_{\mathfrak p}}\right\rceil\right).
 \tag{T19.10}
\] This is a ma\(\xi\)mum over primes above p, in contrast to the minimum in the scalar trace formula.

Take \(p\) additionally prime to \(\kappa\). In the uniform family (T19.3) the exact denominator is \(p^{4r}\). In the concentrated family (T19.5) it is \(p^r\), even though the scalar trace ideal is constant. Dividing w once by its fixed integral cohomology content gives a fixed primitive rational Hodge class; excluding the primes dividing that fixed content preserves these exact p-valuations. Thus a primitive scalar trace does not supply a bounded-denominator presentation of the primitive four-factor tensor in the specified unscaled output group. This does not preclude Chow divisibility and cancellation after normalization.

\paragraph{Proof}\label{rallis_t6_proofs-proof}

\subparagraph{Scalar traces and dual ideals}\label{scalar-traces-and-dual-ideals}

Define the inverse different by \[
 \mathfrak D^{-1}=\{x\in E:\operatorname{Tr}(x\mathcal O)\subseteq\mathbf Z\}.
\] For an \(\mathcal O\)-ideal \(\mathfrak a\) and a rational prime \(p\), localization gives \[
 \operatorname{Tr}(\mathfrak a)\subseteq p^m\mathbf Z_{(p)}
 \quad\Longleftrightarrow\quad
 \mathfrak a\mathcal O_{(p)}\subseteq p^m\mathfrak D^{-1}\mathcal O_{(p)}.
 \tag{T19.11}
\] The reverse implication follows by taking the product with 1. For the forward implication, every \(x\in\mathfrak a\) and every \(b\in\mathcal O\) satisfy \(xb\in\mathfrak a\), so the trace condition holds for \(\operatorname{Tr}(xb)\), which is exactly membership in the trace dual. The same argument after localization handles fractional ideals and every integer m. Comparing ideal valuations in (T19.11) gives \[
 a_{\mathfrak p}\ge me_{\mathfrak p}-d_{\mathfrak p}
 \quad\hbox{for every }\mathfrak p\mid p.
\] The largest permitted \(m\) is the exponent in (T19.1). This proves the formula.

The trace image of an integral ideal is a nonzero subgroup \(m\mathbf Z\). Its generator is attained as a finite integral combination of traces; indeed the trace image itself is already an additive subgroup. If a common denominator \(D\) produces \(1\), clearing it gives \(D\in m\mathbf Z\), so \(m\mid D\). Conversely dividing a trace equal to \(m\) by \(m\) produces \(1\). This proves the asserted optimal scalar denominator.

\subparagraph{Conjugation symmetry of the core/hull ideal}\label{conjugation-symmetry-of-the-corehull-ideal}

Write the rational Riemann form as \[
 \psi(x,y)=\operatorname{Tr}_{E/\mathbf Q}(\xi x\bar y),
 \qquad \bar\xi=-\xi.
 \tag{T19.12}
\] Such a form follows directly from Rosati conjugation: \(\psi(x,y)=\ell(x\bar y)\), the nondegenerate trace pairing writes \(\ell(z)=\operatorname{Tr}(\xi z)\), and alternatingness makes \(\bar\xi=-\xi\). For any fractional \(\mathcal O\)-ideal \(I\), its \(\psi\)-dual is \[
 I^\#=\xi^{-1}\mathfrak D^{-1}\bar I^{-1}.
\] The sign convention for the alternating form can insert a minus sign, which does not alter this ideal. Principal polarization and T6.18 give \(\Lambda^-=(\Lambda^+)^\#\), whence \[
 \mathfrak r
 =\xi^{-1}\mathfrak D^{-1}(\Lambda^+\overline{\Lambda^+})^{-1}.
 \tag{T19.13}
\] The different is conjugation-stable, the last factor is manifestly so, and conjugating \(\xi^{-1}\) changes its sign. This proves (T19.2).

\subparagraph{Constructing the actual lattices}\label{constructing-the-actual-lattices}

First, there is no e\(\xi\)stence assumption on a principal ma\(\xi\)mal-order polarization. Start with a rational Riemann form of the CM type on E tensor R. One can choose \(\xi\) in the purely imaginary part of E with the required signs at the finitely many archimedean embeddings: write \(\xi=\delta a\) with a in \(E^+\), and use density of \(E^+\) in its product of real embeddings to prescribe the signs. A rational symplectic basis for this nondegenerate alternating Q-form spans a lattice \(\Lambda_*\) on which \(\psi\) is integral and unimodular. The positivity and complex-structure compatibility of \(\psi\) are unchanged. Thus the Riemann criterion gives a principally polarized complex abelian variety of the specified CM type, with rational E-action and Rosati conjugation. It need not have integral ma\(\xi\)mal-order action, which is not needed.

There are only finitely many rational primes at which \(\Lambda_*\) differs locally from a fractional ma\(\xi\)mal-order ideal or at which its fixed rational Riemann coefficients are not units. Choose p split completely and outside all these exceptional sets for all four factors. Infinitely many completely split primes e\(\xi\)st; only this qualitative classical fact is used, with no effective or GRH input.

At p the algebra E tensor Q\_p is a product of 2g copies of Q\_p, and O tensor \(\mathbf Z_p\) is the corresponding product of \(\mathbf Z_p\). Complex conjugation pairs the factors. Choose unit generators e\_j,f\_j of each pair so that the local base lattice and Riemann form are \[
 \Lambda_{*,p}=\bigoplus_{j=1}^g(\mathbf Z_p e_j\oplus\mathbf Z_p f_j),
 \qquad \psi(e_j,f_k)=\delta_{jk},
\] with every other pairing between generators zero. The unit rescalings do not change the local ma\(\xi\)mal-order lattice.

Replace one selected pair by \[
 L_r=\mathbf Z_p e\ +\ \mathbf Z_p(f+p^{-r}e).
 \tag{T19.14}
\] The displayed basis has symplectic pairing 1, so \(L_r\) is self-dual. Its ma\(\xi\)mal-order hull and core are, respectively, \[
 L_r^+=p^{-r}\mathbf Z_p e\oplus\mathbf Z_p f,
 \qquad
 L_r^-=\mathbf Z_p e\oplus p^r\mathbf Z_p f.
 \tag{T19.15}
\] For the core, the e-a\(\xi\)s intersection is \(\mathbf Z_p e\), while a vector b f belongs to \(L_r\) exactly when \(p^{-r}b\) lies in \(\mathbf Z_p\). These a\(\xi\)s intersections give the largest submodule for the product ring \(\mathbf Z_p\times\mathbf Z_p\). The hull follows by projecting the two generators onto the two axes. Therefore \[
 [L_r:L_r^-]=[L_r^+:L_r]=p^r,
 \qquad L_r^-(L_r^+)^{-1}=(p^r,p^r).
 \tag{T19.16}
\]

Keep the base lattice at every other finite prime. These local prescriptions determine a global full lattice in E. Concretely, for sufficiently large a, the new local lattice lies between \(p^a\Lambda_{*,p}\) and \(p^{-a}\Lambda_{*,p}\); its image in the finite quotient \(p^{-a}\Lambda_{*,p}/p^a\Lambda_{*,p}\) corresponds to a subgroup of \(p^{-a}\Lambda_*/p^a\Lambda_*\). Its inverse image is the required global lattice. Its completions are exactly the prescribed ones. Duality commutes with completion, so it remains self-dual everywhere and hence globally. The fixed Riemann form remains positive for the fixed CM type, giving actual principally polarized CM abelian varieties. This also proves all required isogeny and fixed-volume assertions.

Applying (T19.16) in every pair gives a local core/hull ideal \(p^r\mathcal O_p\) and core index \(p^{gr}\); outside p the ideal and index are unchanged. Applying it in a single pair gives \(\mathfrak q^r\) and core index \(p^r\). Core and hull commute with localization because their defining lattice conditions involve a finite integral basis of O. T6.18 now gives (T19.3) and (T19.5). In the uniform family the period ideal is multiplied by the rational integer \(p^{4r}\), proving (T19.4). In the concentrated family the p-adic trace is the sum of the 2g coordinates. At least one unmodified coordinate remains when g is at least two, so its trace image is \(\mathbf Z_p\), just as at r=0. At every other prime the ideals are unchanged. This proves (T19.6).

\subparagraph{Full tensor trace and divisor-kernel outputs}\label{full-tensor-trace-and-divisor-kernel-outputs}

Choose \(\beta_i\) in the opposite embedding line of the fixed output factor of the same CM type. Every scalar conjugate of \(\alpha_i\otimes\beta_i\) is then a (1,1) class. The first factor alone distinguishes all Galois conjugates of t: applying \(\sigma\) changes its embedding label by \(\sigma\), and no nonidentity automorphism fixes that label. Thus the d tensors \(\sigma(t)\) are E-linearly independent. Their span is Galois-stable. An invariant vector in it has coefficients \(\sigma(a)\) for one a in E, because invariance identifies every coefficient with the corresponding conjugate of the identity coefficient. This proves surjectivity of (T19.7) onto its rational descent W. Injectivity follows because a zero sum of the independent conjugates forces every \(\sigma(a)\) to vanish. The switch relation makes every conjugate tensor of type (2,2), so W consists of rational Hodge classes.

For completeness, the integral divisor-kernel step uses exactly the all-conjugates hypothesis of T6.7. The pure class \(\alpha_i\otimes\beta_i\) belongs to integral cohomology tensored with O. Every scalar conjugate has type (1,1), so its rational coefficient components lie in the rational Neron--Severi space. The integral Neron--Severi lattice is saturated in integral \(H^2\), by the integral (1,1) theorem and the absence of torsion for an abelian variety. Intersecting with the integral coefficient lattice therefore puts the class in \(\operatorname{NS}\otimes\mathcal O\). Choose an O-linear combination of divisor classes representing it. This introduces no varying denominator. Its pure mixed Kunneth component permits the usual birigidified representative; this is the same divisor-kernel framework as T6.5 and T6.7.

Let K be the product of these four kernels, and z the actual surface pushforward. The projection formula gives \(\operatorname{cl}(K_*z)=kt\), where k is the selected integral coefficient period. Taking the trace in the coefficient ring gives an integral cycle whose class is \(\mathcal T_t(k)\). Here scalar Galois acts on the coefficient field, not on the complex algebraic cycles themselves. By T6.18 the additive period group is exactly \(\mathfrak P_A\), so expanding finite sums proves both inclusions in (T19.8). Multipliers in O preserve \(\mathfrak P_A\) and hence do not enlarge the image.

If a rational combination with common denominator D represents \(w=\mathcal T_t(\kappa)\), clearing denominators gives \(Dw\in\mathcal T_t(\mathfrak P_A)\). Injectivity is equivalent to \(D\kappa\in\mathfrak P_A\). Conversely that membership and (T19.8) produce such an integral combination after multiplying w by D. Hence the least denominator is (T19.9). Ideal valuations give (T19.10), including the ma\(\xi\)mum with zero because D must be an integer.

Choose p avoiding the fixed nonzero coefficient \(\kappa\) and the fixed base ideal. In the uniform family the condition is \(D\kappa\in p^{4r}\mathfrak P_*\), whose least positive solution is D=\(p^{4r}\). In the concentrated family the two selected p-adic valuations require \(v_p(D)\) at least r, while D=\(p^r\) satisfies every ideal inclusion; the least denominator is \(p^r\). The fixed primitive normalization of w changes these valuations by at most a fixed constant, and by none when p is chosen away from its fixed content. This proves all assertions. QED.

\paragraph{What this resolves and the remaining scope}\label{what-this-resolves-and-the-remaining-scope}

The metric-normalized maps of T6.17 e\(\xi\)st along both families: the surface, CM fields and types, fixed theta presentations, and output tensor remain fixed. Their bounded analytic energies and bounded Schwartz seminorms coe\(\xi\)st with unbounded denominators of presentations in the additive group of \textbf{unscaled} traced outputs. The exact local formula measures the saturation needed to pass from that group to a fixed rational tensor. It does not measure the intrinsic denominator of the resulting Chow class.

Indeed, there is exact cancellation at fixed output. Choose fixed source pure birigidified divisor kernels \(D_{i,0}\). If \(h_i^*\alpha_i=\lambda_i\alpha_{i,0}\), then injectivity of the first Chern class on the birigidified Picard group, after extension of coefficients to E, gives \[
 (h_i\times1)^*D_i=\lambda_iD_{i,0}
 \quad\hbox{in }\operatorname{CH}^1\otimes E.
\] This is an equality of the entire coefficient-valued class, not just one complex embedding. Projection and multiplication give raw output \(C_A=(\prod_i\lambda_i)C_0\), while its selected period is \(k_A=(\prod_i\lambda_i)k_0\). Therefore \[
 \frac{C_A}{k_A}=\frac{C_0}{k_0}
 \quad\hbox{in }\operatorname{CH}^2(B)\otimes E.
\] Coefficientwise trace after this cancellation is fixed, so it has a fixed intrinsic rational Chow denominator. Candidate T6.20 proves this carefully and treats simultaneous variation of the output. The growing index in (T19.8) is fully compatible with this cancellation: an unscaled output may be divisible in the larger Chow group even though its division is unavailable inside the chosen unscaled subgroup.

The exact formulas identify a positive arithmetic mechanism as well: scalar trace presentation is possible with precisely the trace-ideal generator, and full-tensor presentation is possible with precisely the integer ideal-membership generator. They do not bound the sizes of integer coefficients in a chosen Bezout presentation or the degrees of the signed cycles. Allowing weights in a fixed fractional coefficient lattice changes only a fixed set of prime denominators; it cannot change (T19.10) at the chosen good prime. Allowing variable weights whose denominators contain p is precisely the division being measured.

The full tensor trace must not be replaced by scalar trace. Its injectivity preserves the coefficient ideal before saturation. The concentrated principal family has unchanged scalar trace ideal but unbounded tensor-presentation index. Any claim about an intrinsic moving-output denominator requires the actual Chow transport and divisibility calculation, rather than this index alone.

\paragraph{Primary sources used and exact check scope}\label{primary-sources-used-and-exact-check-scope}

\begin{enumerate}
\def\labelenumi{\arabic{enumi}.}
\tightlist
\item
  Keith Conrad, \emph{The Different Ideal}, Definition 3.2, printed p.4; Theorem 3.9 and complete proof, and Theorem 3.10 and complete proof, printed p.8; Definition 4.1, printed p.9. Author source: \url{https://kconrad.math.uconn.edu/blurbs/gradnumthy/different.pdf}. These give the trace-dual conventions. Formula (T19.1) and its local proof are supplied above, not asserted to be a quotation of a numbered result there.
\item
  J.S. Milne, \emph{Abelian Varieties}, v2.00, Chapter I, Theorems 2.8--2.9, printed p.14, and the polarization construction and definition, printed p.15. Author source: \url{https://www.jmilne.org/math/CourseNotes/AV.pdf}. The Riemann criterion and algebraicity of the complex-torus maps are imported; the principal CM lattice and deformation constructions above are proved directly.
\item
  J.S. Milne, \emph{Algebraic Number Theory}, v3.08, Theorem 8.31 and Corollary 8.32, printed p.147. Author source: \url{https://www.jmilne.org/math/CourseNotes/ANT.pdf}. Only the e\(\xi\)stence of a completely split prime outside a finite set is imported. This is qualitative Chebotarev; no quantitative estimate or GRH audit is involved.
\item
  T6.7 and T6.18 supply the already checked integral divisor-kernel descent and exact period ideal. T6.17 supplies the independently checked bounded metric normalization in the fixed-source setting. Their proofs are not silently strengthened here.
\end{enumerate}

\paragraph{Check status}\label{rallis_t6_proofs-check-status}

Author, root, and bounds have read the complete initial arithmetic proof and found no gap in the formulas. The scope has been revised to distinguish presentation denominators from intrinsic Chow denominators; the fixed-output cancellation was independently checked with bounds. The revised complete proof and its scope/cancellation addition passed final root and bounds review before promotion. No local compilation is claimed.

\subsection{15. Chow cancellation removes input-period denominators}\label{chow-cancellation-removes-input-period-denominators}

\subsubsection{DEG-L.T6.20 --- Exact Chow cancellation and simultaneous input/output transport}\label{DEG-L.T6.20}

Work over the complex numbers and fix the Galois CM coefficient field \(E\). All eigenvectors below are \(E\)-coefficient Betti vectors, not polarization-unit vectors obtained by arbitrary complex rescaling. For four factors let \[
 h_i:A_{i,0}\longrightarrow A_i,\qquad
 g_i:B_{i,0}\longrightarrow B_i
\] be actual \(E\)-linear isogenies between fixed-type CM abelian varieties. Write \(h=\prod_i h_i\), \(g=\prod_i g_i\), and let \(z_0=f_{0*}[Y]\) for fixed actual maps from a smooth projective surface to \(A_0=\prod_iA_{i,0}\). Choose nonzero eigenvectors with \[
 h_i^*\alpha_i=\lambda_i\alpha_{i,0},\qquad
 g_i^*\beta_i=\mu_i\beta_{i,0},\qquad
 \lambda_i,\mu_i\in E^\times.
 \tag{T20.1}
\] Assume every scalar conjugate of each \(\alpha_i\otimes\beta_i\) has Hodge type \((1,1)\), and similarly for the reference vectors. Choose the uniquely determined birigidified divisor-correspondence classes \[
 D_i\in\operatorname{CH}^1(A_i\times B_i)\otimes E,
 \qquad
 D_{i,0}\in\operatorname{CH}^1(A_{i,0}\times B_{i,0})\otimes E
\] with these pure mixed Betti classes. Let \(K=\prod_iD_i\), \(K_0=\prod_iD_{i,0}\), with the usual pullbacks to the full products. Assume \[
 k_0=\int_Y f_{1,0}^*\alpha_{1,0}\wedge\cdots\wedge
                         f_{4,0}^*\alpha_{4,0}\ne0,
 \qquad C_0=K_{0*}z_0/k_0.
\] Set \[
 \Lambda=\prod_i\lambda_i,\qquad M=\prod_i\mu_i,
 \qquad k_h=\Lambda k_0,\qquad
 C_{A,B}=K_*h_*z_0/k_h.
\] All surface pushforwards retain their actual multiplicities.

\subparagraph{A. Exact cancellation and simultaneous transport}\label{a.-exact-cancellation-and-simultaneous-transport}

There are identities in Chow groups with \(E\) coefficients \[
 (h_i\times g_i)^*D_i=\lambda_i\mu_iD_{i,0},
 \tag{T20.2}
\] \[
 g^*C_{A,B}=M C_0,
 \qquad
 C_{A,B}=\frac{M}{\deg g}\,g_*C_0.
 \tag{T20.3}
\] Every input factor \(\lambda_i\) cancels. In particular, at fixed output \(B=B_0\), with fixed output eigenvectors and \(g=1\), \[
 C_{A,B_0}=C_0.
 \tag{T20.4}
\] This extends T6.5 from clearing integers to arbitrary actual \(E\)-linear input isogenies, provided that the source surface maps and output are fixed and the kernels are the pure birigidified classes stated here. Coefficientwise traces of the normalized cycle are then fixed rational Chow classes and have fixed intrinsic rational denominators. A growing index of the unscaled outputs, as in T6.19, does not contradict (T20.4).

\subparagraph{B. Genuine dual outputs and the direction of the dual isogeny}\label{b.-genuine-dual-outputs-and-the-direction-of-the-dual-isogeny}

Suppose \(B_i=A_i^\vee\), \(B_{i,0}=A_{i,0}^\vee\), and \(D_i\), \(D_{i,0}\) are the selected \(E\)-eigencomponents of the normalized Poincare bundles. Normalize \(\beta_i\), \(\beta_{i,0}\) by these Poincare components, so their pure classes are exactly \(\alpha_i\otimes\beta_i\) and \(\alpha_{i,0}\otimes\beta_{i,0}\). Then the actual dual isogeny points in the reverse direction: \[
 q=\prod_i h_i^\vee:B=\prod_iA_i^\vee
                      \longrightarrow B_0=\prod_iA_{i,0}^\vee.
\] The exact formulas are \[
 (h_i^\vee)^*\beta_{i,0}=\lambda_i\beta_i,
 \qquad
 K_*h_*z_0=q^*(K_{0*}z_0),
 \tag{T20.5}
\] \[
 C_{A,B}=\Lambda^{-1}q^*C_0,
 \qquad
 q_*C_{A,B}=\frac{\deg q}{\Lambda}\,C_0.
 \tag{T20.6}
\] Thus the moving-output problem is an actual isogeny pullback followed by a specified coefficient division. It is not a residual input-period ambiguity.

\subparagraph{C. Uniform denominator for factorwise integral output tensors}\label{c.-uniform-denominator-for-factorwise-integral-output-tensors}

There is a useful construction that does not require varying the input factors at all. Fix integral \(\mathcal O_E\)-coefficient input eigenvectors \(\alpha_{i,0}\), fixed actual \(f_{i,0}\), and their nonzero period \(k_0\in\mathcal O_E\). Let the output CM factors \(B_i\) vary arbitrarily within the required fixed CM types. For every choice of integral \(\mathcal O_E\)-coefficient output eigenvectors \(\beta_i\) of the opposite types, there exist integral \(\mathcal O_E\)-coefficient birigidified divisor kernels on \(A_{i,0}\times B_i\) with class \(\alpha_{i,0}\otimes\beta_i\).

Put \[
 Z_B=K_{B*}z_0\in\operatorname{CH}^2(B)\otimes\mathcal O_E,
 \qquad C_B=Z_B/k_0.
\] For every fixed \(a\in\mathcal O_E\), the rational Chow class \[
 \Gamma_B(a)=\operatorname{Tr}_{E/\mathbf Q}(a C_B)
\] has an integral multiple with the same denominator for all output lattices: \[
 N_0\Gamma_B(a)\in\operatorname{im}
 \bigl(\operatorname{CH}^2(B)\longrightarrow
                 \operatorname{CH}^2(B)\otimes\mathbf Q\bigr),
 \qquad
 N_0=|N_{E/\mathbf Q}(k_0)|.
 \tag{T20.7}
\] A fixed finite set of fractional coefficient weights only adds one fixed clearing integer. This proves a genuine uniform Chow denominator for these \textbf{factorwise integral eigenvector tensors}. It does not by itself make their rational Galois-orbit lattice saturated in integral cohomology or give a degree bound for the signed cycles.

\subparagraph{D. Exact denominator of the displayed dual-pullback representative}\label{d.-exact-denominator-of-the-displayed-dual-pullback-representative}

Choose a finite representative \[
 C_0=\sum_{\ell=1}^u c_\ell[Z_\ell],\qquad c_\ell\in E,
 \tag{T20.8}
\] with distinct integral subvarieties \(Z_\ell\subset B_0\). In the dual situation (B), a representative of \(\Gamma_B(a)\) has coefficient \[
 \operatorname{Tr}_{E/\mathbf Q}(a c_\ell/\Lambda)
 \tag{T20.9}
\] on every irreducible component of \(q^{-1}(Z_\ell)\). Therefore the exact common denominator of this displayed rational cycle is the least common multiple of the denominators of (T20.9).

Let \(\mathfrak C=\sum_\ell\mathcal O_Ec_\ell\) and let \(\mathfrak D\) be the different of \(E/\mathbf Q\). The least positive integer clearing this representative simultaneously for every \(a\in\mathcal O_E\) is the positive generator of \[
 (\Lambda\mathfrak D^{-1}\mathfrak C^{-1})\cap\mathbf Z.
 \tag{T20.10}
\] Its exponent at a rational prime \(p\) is \[
 \max\left(0,\max_{\mathfrak p\mid p}
 \left\lceil\frac{v_{\mathfrak p}(\Lambda)
                      -v_{\mathfrak p}(\mathfrak D)
                      -v_{\mathfrak p}(\mathfrak C)}{e_{\mathfrak p}}\right\rceil\right).
 \tag{T20.11}
\] These are exact representative denominators and hence upper bounds for the intrinsic Chow denominator. Rational equivalence or divisibility in the Chow group may reduce the latter; no converse lower bound is asserted.

\paragraph{Proof}\label{proof-1}

\subparagraph{Birigidification makes the entire divisor class faithful}\label{birigidification-makes-the-entire-divisor-class-faithful}

For complex abelian varieties \(U,V\), write \(\operatorname{Pic}^{\mathrm b}(U\times V)\) for divisor-correspondence classes trivialized along the two zero axes. Its first Chern class is injective. Indeed, a class in its kernel lies in \(\operatorname{Pic}^0(U\times V)\), and \[
 \operatorname{Pic}^0(U\times V)
 =\operatorname{pr}_U^*\operatorname{Pic}^0(U)
  \oplus\operatorname{pr}_V^*\operatorname{Pic}^0(V).
\] Restriction to the two axes recovers these two summands, so birigidification makes both trivial. Equivalently, the divisor-correspondence group is \(\operatorname{Hom}(U,V^\vee)\), and the induced map on first cohomology detects a homomorphism. This is the correspondence/biextension identification of SGA7 I, Expose VII, Examples 2.9.5--2.9.6. Tensoring the injective first Chern class map with the characteristic-zero field \(E\) preserves injectivity.

Pullback by a homomorphism in either argument preserves birigidification. The first Chern class of the left side of (T20.2) is, in the full \(E\)-coefficient Betti space, \[
 h_i^*\alpha_i\otimes g_i^*\beta_i
   =\lambda_i\mu_i(\alpha_{i,0}\otimes\beta_{i,0}).
\] This is not merely an equality after evaluating one complex embedding. Applying the just-proved injectivity gives (T20.2) in the entire birigidified Picard group with \(E\) coefficients, and hence in \(\operatorname{CH}^1\otimes E\).

\subparagraph{Projection and flat base change}\label{projection-and-flat-base-change}

Let \(p_A,p_B\) denote the two projections from \(A\times B\). By definition, \[
 K_*h_*z_0=(p_B)_*\bigl(p_A^*h_*z_0\cdot K\bigr).
\] Flat base change along \(g\) and then along the product square for \(h\), followed by the projection formula, gives \[
 \begin{aligned}
 g^*(K_*h_*z_0)
 &= (K|_{A\times B_0})_*h_*z_0\\
 &= ((h\times g)^*K)_*z_0\\
 &=\Lambda M\,K_{0*}z_0.
 \end{aligned}
 \tag{T20.12}
\] Here the first kernel is the pullback by \(1\times g\), not a literal restriction to a subvariety. All isogenies in characteristic zero are finite etale, hence flat, so these are ordinary Chow pullbacks and proper pushforwards. The product of the four identities (T20.2) proves the last equality. The period scales by \(\Lambda\), proving the first equation of (T20.3). Applying \(g_*\) and using \(g_*g^*=(\deg g)\operatorname{id}\) proves the second. This also proves (T20.4).

The fixed rational Chow denominator in (T20.4) is concrete. Choose a finite cycle representative of \(\operatorname{Tr}(a C_0)\) and clear its finitely many rational coefficients once. The same representative on the same output works for every input isogeny. This is stronger than a bound on the selected period and is not contradicted by the nonsaturation of an unscaled subgroup of Chow classes.

\subparagraph{Duality and the reversed map}\label{duality-and-the-reversed-map}

For a homomorphism \(h_i:A_{i,0}\to A_i\), the normalized Poincare bundles satisfy \[
 (h_i\times1)^*\mathcal P_{A_i}
   \simeq (1\times h_i^\vee)^*\mathcal P_{A_{i,0}}.
 \tag{T20.13}
\] This is the defining universal property of the dual homomorphism. The selected \(E\)-eigencomponents agree because \(h_i\) is \(E\)-linear. In Betti cohomology, the selected equality reads \[
 \lambda_i\alpha_{i,0}\otimes\beta_i
   =\alpha_{i,0}\otimes(h_i^\vee)^*\beta_{i,0},
\] so it gives the first identity of (T20.5). Birigidified Chern-class injectivity also gives (T20.13) for the selected divisor components themselves. These components are legitimate \(E\)-coefficient divisor classes even when their projection from the integral Poincare class has denominators. This part does not assert their integrality.

Taking the fourfold product of (T20.13), pulling it against \(z_0\), and applying flat base change on the output gives the second identity of (T20.5). Divide by \(k_h=\Lambda k_0\) to obtain the first identity of (T20.6). Its pushforward follows from \(q_*q^*=(\deg q)\operatorname{id}\). In particular the scalar in (T20.6) is \(\Lambda\), not an assumed field norm or an arbitrarily chosen complex metric period.

\subparagraph{Integral factorwise tensors require only one fixed period denominator}\label{integral-factorwise-tensors-require-only-one-fixed-period-denominator}

For part (C), each pure mixed class \(\alpha_{i,0}\otimes\beta_i\) belongs to integral cohomology tensored with \(\mathcal O_E\), and all its scalar conjugates are of type \((1,1)\). The integral descent argument of T6.7 and T6.18 puts it in the saturated mixed Neron--Severi lattice tensored with \(\mathcal O_E\). Passing to the birigidified representative removes axis restrictions and introduces no coefficient denominator. Products and proper pushforward preserve integral \(\mathcal O_E\)-coefficient cycles, proving that \(Z_B\) has the stated coefficients.

For a nonzero algebraic integer \(k_0\), the element \(N_0/k_0\) is an algebraic integer: up to the sign of its rational norm it is the product of the other scalar conjugates. Consequently \[
 N_0\Gamma_B(a)
   =\operatorname{Tr}_{E/\mathbf Q}
       \bigl(a(N_0/k_0)Z_B\bigr)
\] is an integral cycle class, proving (T20.7). This argument is uniform over the target lattices and does not use a bound for an isogeny degree. It proves uniform division for the explicitly stated factorwise integral tensors; primitivity in a larger rational orbit lattice is a separate question.

\subparagraph{The explicit representative and the trace dual}\label{the-explicit-representative-and-the-trace-dual}

A finite etale pullback of an integral subvariety is generically reduced; each irreducible component has multiplicity one in the flat pullback cycle. Every component maps onto the original subvariety. Thus the inverse images of distinct \(Z_\ell\) have no common irreducible components. Apply (T20.6) to the representative (T20.8) and then trace the coefficients. This gives exactly (T20.9), with no cancellation among supports of different \(Z_\ell\).

An integer \(D>0\) clears these coefficients for all \(a\in\mathcal O_E\) if and only if \[
 D c_\ell/\Lambda\in\mathfrak D^{-1}
 \quad\hbox{for every }\ell.
\] Equivalently, \(D\mathfrak C\subseteq\Lambda\mathfrak D^{-1}\). Since \(\mathfrak C\) is an invertible fractional maximal-order ideal, this is precisely \(D\in\Lambda\mathfrak D^{-1}\mathfrak C^{-1}\). Taking its intersection with the integers proves (T20.10); comparing valuations gives (T20.11). The denominator of a particular representative always clears its Chow class. The proof supplies no reason that a rationally equivalent representative cannot have a smaller denominator. This completes the proof. QED.

\paragraph{Comparison with the fixed-K Weil determinant plane}\label{comparison-with-the-fixed-k-weil-determinant-plane}

In B23 the coefficient space is a rank-one module over an imaginary quadratic field \(K\), hence a two-dimensional rational Weil plane. For an \(E\)-linear isogeny with scalar \(\alpha\), the pullback scalar on that determinant line is \[
 s=N_{E/K}(\alpha).
\] If \(\nu\) is the homology covolume ratio, its degree is \(\delta=\nu^{-1}s\bar s\). A normalization compatible with the chosen target determinant generator therefore has the same general form \[
 C_A=\frac{s}{\delta}\,h_*C_0
     =\frac{\nu}{\bar s}\,h_*C_0.
 \tag{T20.14}
\] This is the formal comparison of scalar conventions; it does not identify B23's determinant plane with the four-factor switch orbit. B23's common integral pushforward content \(\bar{\mathfrak j}\) records the image before saturation. Multiplication by \(s/\delta\) can leave that integral image, and a claim that the result is integral in Chow would require an additional divisibility argument. Conversely, the image index alone does not prove that the normalized class has an unbounded intrinsic denominator.

For the four separate CM factors, the scalar \(M\) in (T20.3) is a product of four selected embedding eigenvalues; it is generally not \(N_{E/K}(\alpha)\). Its full rational Galois orbit has dimension \([E:\mathbf Q]\), rather than two. Any comparison with a generic \(K\)-Weil target must retain the actual correspondence and its integral lattice map.

The positive conclusion is precise: all varying input isogenies and their period normalization cancel in Chow, and direct kernels with fixed integral input vectors give a uniform denominator for every factorwise integral output tensor. The unresolved moving-output primitive question is whether the integral orbit classes needed for the intended Weil target can be represented by these tensors with uniformly controlled saturation, or whether their normalized isogeny pullbacks have additional Chow divisibility. Neither a metric Rallis norm nor scalar trace contraction settles that question.

\paragraph{Primary sources and exact locators}\label{primary-sources-and-exact-locators}

\begin{itemize}
\tightlist
\item
  Grothendieck, SGA7 I, Expose VII, Example 2.9.5, printed p.175 (PDF p.180; VII43), and Remark 2.9.6, printed p.176 (PDF p.181; VII44), identify birigidified line bundles, biextensions and divisor correspondences \(\operatorname{Hom}(U,V^\vee)\). Original scan: \url{https://library.slmath.org/nonmsri/sga/sga/pdf/sga7-1.pdf}. This agent independently read both original page images in this phase.
\item
  Milne, \emph{Abelian Varieties}, v2.00, Chapter I, Poincare universal property and Remark 8.7, printed p.37, and the opening paragraph of Section 9, printed p.41. The latter explicitly states (T20.13) and defines the dual homomorphism. \url{https://www.jmilne.org/math/CourseNotes/AV.pdf}.
\item
  Stacks Project, Lemma 42.15.1 (tag 02RG) and Lemma 42.15.2 (tag 02RH), in Section 42.15 (tag 02RF), state flat base change and finite locally free push-pull. \url{https://stacks.math.columbia.edu/tag/02RF}. Exact statements and complete short proofs read.
\item
  Stacks Project, Lemma 43.22.1 (tag 0B10), supplies the flat proper projection formula. \url{https://stacks.math.columbia.edu/tag/0B10}. Statement and full displayed proof read; on Chow groups the usual moving/intersection formalism gives the formula used above.
\item
  Conrad, \emph{The Different Ideal}, Definition 3.2 p.4, Theorems 3.9--3.10 p.8 and Definition 4.1 p.9, supplies the trace-dual convention in (T20.10). \url{https://kconrad.math.uconn.edu/blurbs/gradnumthy/different.pdf}. The denominator calculation itself is proved here.
\end{itemize}

\paragraph{Check status}\label{check-status-1}

The entire saved proof was read by the author, root and bounds; all reviews passed before promotion. Source checks are recorded separately. No local compilation is claimed.

\subsection{16. The full integral orbit lattice and uniform division with one maximal-order factor}\label{the-full-integral-orbit-lattice-and-uniform-division-with-one-maximal-order-factor}

\subsubsection{DEG-L.T6.21 --- Exact four-factor saturation and uniform division with an anchor}\label{DEG-L.T6.21}

Let \(E\) be a Galois CM field, \(\mathcal O=\mathcal O_E\), \(d=[E:\mathbf Q]=2g\), and let \(c\) denote complex conjugation. Let \[
 B_i=\mathbf C^g/\Phi_i(\Lambda_i),\qquad i=1,2,3,4,
\] where each \(\Lambda_i\subset E\) is a full \(\mathbf Z\)-lattice. The four CM types satisfy the switch relation, so the diagonal orbit specified below consists of Hodge classes. The first homology has rank one over \(E\); simplicity and primitivity of the types are unnecessary.

Inside the degree-four Kunneth summand \(\bigotimes_i H^1(B_i,\mathbf Q)\), identify the rational diagonal orbit \(W\) with \(E\) by the injective map \[
 \Theta(\kappa)(x_1,x_2,x_3,x_4)
   =\operatorname{Tr}_{E/\mathbf Q}
       (\kappa x_1x_2\bar x_3\bar x_4).
 \tag{T21.1}
\] Exchanging the barred and unbarred labels gives the equivalent opposite switch convention. Put \[
 \begin{aligned}
 M&=\left\langle x_1x_2\bar x_3\bar x_4:
                         x_i\in\Lambda_i\right\rangle_{\mathbf Z},\\
 M^+&=\mathcal O M
   =\Lambda_1^+\Lambda_2^+\overline{\Lambda_3^+}\,
                                  \overline{\Lambda_4^+},\\
 L&=(M^+)^{-1},\qquad
 M^\vee=\{\kappa\in E:\operatorname{Tr}(\kappa M)\subseteq\mathbf Z\}.
 \end{aligned}
 \tag{T21.2}
\] Here \(M\) is a full additive lattice, \(M^+\) is an invertible fractional maximal-order ideal, and \(L\) is its ideal inverse.

\subparagraph{A. The two integral lattices and their exact index}\label{a.-the-two-integral-lattices-and-their-exact-index}

The actual integral intersection lattice of the rational orbit is \[
 W\cap H^4(B_1\times\cdots\times B_4,\mathbf Z)
     =\Theta(M^\vee).
 \tag{T21.3}
\] The additive span of coefficient traces of products of the four \textbf{individually integral \(\mathcal O\)-coefficient eigenvectors} is exactly \[
 \Theta(L).
 \tag{T21.4}
\] Consequently \[
 [M^\vee:L]
   =|\operatorname{Disc}(E)|\,[M^+:M].
 \tag{T21.5}
\] This is an equality of actual integral intersection lattices, including at ramified primes; it does not assume integral eigenprojectors on the varying factors.

\subparagraph{B. Exact saturation exponent and a uniform anchor case}\label{b.-exact-saturation-exponent-and-a-uniform-anchor-case}

Let \(n(M)\) be the exponent of the finite group \(M^\vee/L\), equivalently the least positive integer with \(n(M)M^\vee\subseteq L\). If \(\mathfrak D\) is the different of \(E/\mathbf Q\), then \[
 n(M)\mathbf Z
 =\{n\in\mathbf Z:n\mathfrak D^{-1}M^+\subseteq M\}.
 \tag{T21.6}
\] Let \(e(M)\) be the exponent of \(M^+/M\), and let \(d_E>0\) generate \(\mathfrak D\cap\mathbf Z\). Then \[
 e(M)\mid n(M)\mid d_E e(M).
 \tag{T21.7}
\] In particular, at every prime unramified in \(E\), the saturation exponent equals the exponent of the corresponding part of the product-lattice hull quotient.

If \textbf{at least one} \(\Lambda_i\) is stable under \(\mathcal O\), then \[
 M=M^+,
 \qquad [M^\vee:L]=|\operatorname{Disc}(E)|,
 \qquad n(M)=d_E.
 \tag{T21.8}
\] Thus fixing one maximal-order factor eliminates every target-dependent saturation cost, even if the other three lattices have arbitrarily large core indices.

\subparagraph{C. Uniform integral division of primitive orbit classes with one anchor}\label{c.-uniform-integral-division-of-primitive-orbit-classes-with-one-anchor}

Fix the input surface maps and integral input eigenvectors of T6.20(C), with nonzero period \(k_0\in\mathcal O\), and put \(N_0=|N_{E/\mathbf Q}(k_0)|\). Every integral class \(w\) in the orbit (T21.3), primitive or otherwise, has \[
 N_0 n(M)\,w
\] as the cohomology class of an integral algebraic cycle constructed from the fixed surface and the divisor kernels of T6.20(C). If one output factor has maximal-order action, the same integer \[
 N_0d_E
 \tag{T21.9}
\] works for every primitive integral orbit class and every choice of the other three output lattices. In particular it is at most \(N_0|\operatorname{Disc}(E)|\). This supplies a genuine uniform division bound for this stated family of integral Hodge classes; it does not bound the degrees or the sizes of integer coefficients in a chosen cycle representative.

\subparagraph{D. All four mixed principal lattices can have growing saturation}\label{d.-all-four-mixed-principal-lattices-can-have-growing-saturation}

Take the actual principal-polarization families of T6.19, at a completely split prime \(p\) outside the fixed bad set. In one conjugate pair use, for all four output homology lattices, \[
 L_r=\mathbf Z_p(1,0)+\mathbf Z_p(p^{-r},1),
\] with complex conjugation swapping the two coordinates. Apply the bars in (T21.2) to the third and fourth factors. The resulting product lattice and hull in this pair are \[
 M_{r,p}
   =p^{-r}\mathbf Z_p^2+
           \mathbf Z_p p^{-2r}(1,1),
 \qquad
 M_{r,p}^+=p^{-2r}\mathbf Z_p^2.
 \tag{T21.10}
\] Their quotient is cyclic of order \(p^r\). Equivalently, \[
 \begin{aligned}
 M_{r,p}^\vee
  &=\{(x,y)\in p^r\mathbf Z_p^2:x+y\in p^{2r}\mathbf Z_p\},\\
 L_{r,p}^{\mathrm{fact}}&=p^{2r}\mathbf Z_p^2,\\
 M_{r,p}^\vee/L_{r,p}^{\mathrm{fact}}&\simeq\mathbf Z/p^r\mathbf Z.
 \end{aligned}
 \tag{T21.11}
\] The class of \(p^r(1,-1)\) generates the last quotient. Modifying all \(g\) conjugate pairs gives a direct sum of \(g\) such quotients. Thus \[
 v_p([M_r^+:M_r])=gr,
 \qquad v_p(n(M_r))=r.
 \tag{T21.12}
\] All other prime contributions are fixed. There exist globally primitive integral Hodge classes whose image in this saturation quotient has order \(p^r\). Their expression in the factorwise integral trace lattice requires that \(p\)-denominator. This is a sharp statement about the factorized construction and its saturation, not a lower bound for the intrinsic algebraic-cycle denominator over all possible constructions.

\paragraph{Proof}\label{proof-2}

\subparagraph{The orbit as an integral multilinear form}\label{the-orbit-as-an-integral-multilinear-form}

Each \(H^1(B_i,\mathbf Q)\) is \(\operatorname{Hom}_{\mathbf Q}(E,\mathbf Q)\). Its integral lattice is \(\operatorname{Hom}_{\mathbf Z}(\Lambda_i,\mathbf Z)\). Since the four lattices are free of finite rank, their tensor product identifies with the integral multilinear forms on \(\Lambda_1\times\cdots\times\Lambda_4\). Integral Kunneth makes this tensor product a direct summand of \(H^4(\prod_i B_i,\mathbf Z)\), with no factorial or torsion correction.

Formula (T21.1) is injective: if it vanishes identically, set the last three rational arguments equal to 1 and use nondegeneracy of the field trace pairing on the first argument. After extending coefficients to \(E\), it is the sum of the Galois conjugates of the tensor with embedding labels \((1,1,c,c)\). These conjugates lie in distinct eigenspaces, so its image is exactly the specified rational diagonal orbit. The switch relation says that every one of these conjugates has type \((2,2)\).

Integrality of the multilinear form is precisely the condition \(\operatorname{Tr}(\kappa M)\subseteq\mathbf Z\), proving (T21.3). The lattice \(M\) is finitely generated by products of four chosen integral bases. It has full rational rank: fixing nonzero elements in three factors puts a nonzero scalar multiple of the full first lattice inside it. Expanding the maximal-order hulls shows directly that \(\mathcal O M\) is the displayed product ideal in (T21.2).

\subparagraph{Individually integral eigenvectors and the index}\label{individually-integral-eigenvectors-and-the-index}

By T6.18, the actual integral coefficient line for the identity embedding in the first factor is \((\Lambda_1^+)^{-1}e_1\); the second is analogous. The conjugate lines in the last two factors have coefficient ideals \(\overline{\Lambda_3^+}^{-1}\) and \(\overline{\Lambda_4^+}^{-1}\). The additive span of products of their coefficients is exactly the product of these four ideals, namely \(L\). Coefficientwise field trace of their tensor products is \(\Theta\) applied to those products. This proves (T21.4), as equality of the actual additive spans. No division by \([E:\mathbf Q]\) is used in this trace convention.

Trace duality reverses indices of full additive lattices, giving \[
 [M^\vee:(M^+)^\vee]=[M^+:M].
\] For the fractional maximal-order ideal \(M^+\), trace duality gives \[
 (M^+)^\vee=\mathfrak D^{-1}(M^+)^{-1}
                    =\mathfrak D^{-1}L.
\] Therefore \[
 [(M^+)^\vee:L]=N_{E/\mathbf Q}(\mathfrak D)
                     =|\operatorname{Disc}(E)|.
\] Multiplying the two indices proves (T21.5). This argument incorporates the different exactly, so it remains valid at ramified primes.

\subparagraph{The saturation exponent and an integral maximal-order factor}\label{the-saturation-exponent-and-an-integral-maximal-order-factor}

For a positive integer \(n\), dualizing the inclusion \(nM^\vee\subseteq L\) gives \[
 L^\vee\subseteq n^{-1}M.
\] The fractional ideal duality formula says \(L^\vee=\mathfrak D^{-1}M^+\). This proves (T21.6). Since \(\mathfrak D^{-1}\supseteq\mathcal O\), any such \(n\) also satisfies \(nM^+\subseteq M\); hence \(e(M)\mid n(M)\). Conversely, \[
 d_E e(M)\mathfrak D^{-1}M^+
  \subseteq e(M)M^+\subseteq M,
\] so \(n(M)\mid d_E e(M)\). This proves (T21.7).

If, for example, \(\Lambda_1\) is \(\mathcal O\)-stable, multiplication of a generating product in \(M\) by any \(a\in\mathcal O\) can be absorbed into its first factor. Thus \(M\) itself is \(\mathcal O\)-stable and equals \(M^+\). The same argument works in a barred factor because conjugation preserves \(\mathcal O\). Since this ideal is invertible, condition (T21.6) then reduces to \[
 n\mathfrak D^{-1}\subseteq\mathcal O,
 \quad\hbox{equivalently }n\in\mathfrak D.
\] Its least positive integer is exactly \(d_E\), proving (T21.8). The bound \(d_E\le|\operatorname{Disc}(E)|\), and indeed divisibility by it, follows because the finite group \(\mathcal O/\mathfrak D\) has order \(|\operatorname{Disc}(E)|\), so its order kills it.

For part (C), write \(w=\Theta(\kappa)\) with \(\kappa\in M^\vee\). By definition, \(n(M)\kappa\in L\). Formula (T21.4) writes \(n(M)w\) as a finite integral sum of coefficient traces of factorwise integral eigenvector products. T6.20(C) supplies, for each term, a rational algebraic cycle whose class is that term and whose denominator divides the same \(N_0\). Summing these cycles proves that \(N_0n(M)w\) is represented by an integral algebraic cycle. In the anchor case \(n(M)=d_E\), proving (T21.9). No bounds on the finite sum's coefficients or degrees are inferred.

\subparagraph{The explicit principal family and its saturation quotient}\label{the-explicit-principal-family-and-its-saturation-quotient}

Choose the split prime outside the finitely many primes at which any base lattice differs from \(\mathcal O\), or any fixed Riemann form is not unimodular on that lattice. Use the common idempotent coordinates of \(E\otimes\mathbf Q_p\); the trace is the sum of these coordinates and conjugation swaps each conjugate pair. On a given pair the four Riemann forms have matrices \(u_i\left(\begin{smallmatrix}0&1\\-1&0\end{smallmatrix}\right)\), with possibly different units \(u_i\in\mathbf Z_p^\times\). The determinant-one shear defining \(L_r\) preserves each such form, so the same displayed \(L_r\) is self-dual for all four pairings. Thus no separate unit rescaling of the four coordinate systems is needed. The local-to-global construction of T6.19 realizes exactly these four modifications by actual principally polarized CM varieties.

In one pair let \(e=(1,0)\), \(f=(0,1)\), and \(q=p^{-r}\). The first two factors have generators \(e\) and \((q,1)\), and the conjugated last two have generators \(f\) and \((1,q)\). Multiply componentwise. A product containing both an \(e\) and an \(f\) is zero. The remaining products generate \(qe\), \(qf\), and \(q^2(e+f)\); they also generate \(e,f\), which already belong to \(q\mathbf Z_p^2\). This proves (T21.10).

After multiplication by \(p^{2r}\), the hull becomes \(\mathbf Z_p^2\) and the product lattice becomes \[
 p^r\mathbf Z_p^2+\mathbf Z_p(1,1).
\] The quotient is identified with \(\mathbf Z/p^r\mathbf Z\) by subtracting the two coordinates. This proves the hull quotient assertion. Taking its trace dual before scaling gives the first formula of (T21.11): pairing with \(p^{-r}e,p^{-r}f\) forces \(x,y\in p^r\mathbf Z_p\), and pairing with \(p^{-2r}(e+f)\) forces \(x+y\in p^{2r}\mathbf Z_p\). The inverse hull ideal is \(p^{2r}\mathbf Z_p^2\). Its quotient in the trace dual is generated by \(p^r(1,-1)\) and has order \(p^r\).

Different coordinate pairs contribute independent direct summands. The base lattice at \(p\) is the full product maximal-order lattice, so vectors supported on a single pair are available in every factor. Thus the local product module is the direct sum of the computed modules. Its hull index is \(p^{gr}\) and its quotient exponent is \(p^r\), proving (T21.12). Outside \(p\), all input lattices in this product calculation remain fixed.

Finally, the global finite group \(M^\vee/L\) has this prescribed \(p\)-primary part. Lift an element whose \(p\)-primary order is \(p^r\) to the integral global orbit lattice. The lift is not divisible by \(p\): otherwise its image would be \(p\) times an element in a group whose \(p\)-primary exponent is only \(p^r\), so its \(p\)-primary order would be at most \(p^{r-1}\). Divide the lift by its full integral content. That content is prime to \(p\), so the order \(p^r\) in the saturation quotient is preserved. The resulting class is primitive in the orbit intersection lattice, and hence in ambient integral cohomology because intersection with a rational subspace is saturated. This gives a globally primitive integral Hodge class with the claimed necessary presentation denominator. It does not prove that every algebraic cycle representing this class has that denominator. QED.

\paragraph{Relation to the main normalization question}\label{relation-to-the-main-normalization-question}

T6.20 removes every variable input-isogeny period denominator and gives a fixed denominator on the factorized integral trace lattice. T6.21 identifies the remaining integral-lattice comparison exactly. It closes the primitive-division gate when one output factor has maximal-order action, and it isolates a computable product-lattice exponent when none does.

This is a four-factor full-CM orbit. It is not B23's middle \(K\)-Weil determinant plane, whose integral pushforward content is a relative norm ideal. The exact comparison with such a target still requires its actual correspondence. Nor does a growing factorized-lattice saturation exponent prove an intrinsic integral Hodge obstruction: an additional Chow divisibility relation or a different cycle construction can enlarge the available integral cycle lattice.

\paragraph{Primary sources and checks}\label{primary-sources-and-checks}

The new identities are proved above from integral Kunneth, the field trace pairing and lattice duality. The trace-dual conventions and fractional-ideal duality are Keith Conrad, \emph{The Different Ideal}, Definition 3.2 p.4, Theorem 3.9 p.8, Definition 4.1 p.9; the norm of the different equals the absolute discriminant by Theorem 4.8 p.11, whose full trace-Gram-determinant proof was independently read. Author URL: \url{https://kconrad.math.uconn.edu/blurbs/gradnumthy/different.pdf}.

Actual principal lattice realization is the complete local-to-global construction in candidate T6.19; no additional moduli existence assertion is imported. The uniform cycle denominator uses precisely part (C) of candidate T6.20. Both predecessor candidates retain their stated coefficient and Chow-divisibility scope.

\paragraph{Check status}\label{check-status-2}

The entire revised saved proof was read by the author, root and bounds; all reviews passed before promotion. The common-idempotent and global primitivity refinements were included in the final reads. No computation in a proof assistant is claimed.

\subsection{17. Mixed kernels and every proper block of auxiliary maximal-order factors}\label{mixed-kernels-and-every-proper-block-of-auxiliary-maximal-order-factors}

\subsubsection{DEG-L.T6.22 --- Mixed divisor kernels and proper-block correspondences retain the saturation cost}\label{DEG-L.T6.22}

Let \(E\) be the Galois CM field of T6.21, of degree \(d=2g\), and let \(c\) be complex conjugation. Retain an \textbf{indecomposable} four-factor switch: the selected eigenlines have CM sign functions \(s_1,s_2,s_3,s_4\) with \[
 s_1+s_2+s_3+s_4=0,
 \qquad
 \sum_{i\in I}s_i\ne0
 \quad(\varnothing\ne I\subsetneq\{1,2,3,4\}).
 \tag{T22.1}
\] These are equalities and inequalities of functions on the scalar Galois group, not merely at one complex place. The actual four types in T6.16 satisfy this hypothesis.

Choose a completely split prime \(p\) as in T6.21. On each conjugate pair of coordinates in \(E\otimes\mathbf Q_p\), all four output homology lattices are \[
 L_r=\mathbf Z_p e+\mathbf Z_p(p^{-r}e+f),
 \qquad r\ge1,
 \tag{T22.2}
\] in the common idempotent coordinates, and conjugation interchanges \(e,f\). These are the actual principally polarized CM varieties constructed in T6.19 and T6.21. Put \(B_r=\prod_{i=1}^4B_{i,r}\). Let \(W_r\) be the rational diagonal orbit with labels \((1,1,c,c)\), identified with \(E\) by T6.21's trace map \(\Theta\). Write \(\pi_W\) for the rational projection which first selects the degree-\((1,1,1,1)\) Kunneth summand and then selects this diagonal orbit. This is the specified eigencomponent projection; no integrality or bounded denominator of a correspondence representing \(\pi_W\) is assumed.

At \(p\), the factorwise trace lattice is \[
 L^{\mathrm{fact}}_{r,p}=p^{2r}(\mathcal O_E\otimes\mathbf Z_p),
 \tag{T22.3}
\] and the quotient of the integral orbit lattice by it is \((\mathbf Z/p^r\mathbf Z)^g\), as proved in T6.21.

\subparagraph{A. Arbitrary integral mixed divisor kernels still have the same local content}\label{a.-arbitrary-integral-mixed-divisor-kernels-still-have-the-same-local-content}

Let \(U\) be a fixed product of CM abelian varieties whose rational first homologies have rank one over \(E\). Choose \(p\) also outside the fixed bad set for these source lattices. Thus \[
 H^1(U,\mathbf Z_p)
\] has a basis of pure CM Hodge eigenlines for the separate source \(E\)-actions. This basis assertion is part of the hypothesis; an unspecified higher-rank \(E\)-module with mixed Hodge multiplicities would not suffice.

Let \(z\in\operatorname{CH}_2(U)\) be any integral algebraic cycle. The cycle and its degree need not be fixed. Let \(K\in\operatorname{CH}^4(U\times B_r)\) belong to the integral subring generated by divisor classes. Thus \(K\) is a finite integral sum of products of four integral divisor classes. The divisors may have arbitrary mixed components, need not be \(E\)-linear, and may connect any source factor to any output factor. Source-source and output-output components are also permitted.

Then \[
 \pi_W\operatorname{cl}(K_*z)
       \in\Theta(L^{\mathrm{fact}}_{r,p})
       \subset W_r\otimes\mathbf Q_p.
 \tag{T22.4}
\] In particular, allowing impure eigencomponents, imprimitive CM types, or cancellations between off-orbit terms does not enlarge this local image.

Let \(w_r\) be one of the primitive integral orbit classes of T6.21 whose saturation-quotient image has order \(p^r\). If a rational linear combination of the classes in (T22.4), with common scalar denominator \(D\), equals \(w_r\), then \[
 v_p(D)\ge r.
 \tag{T22.5}
\] More generally, if it equals \(Nw_r\) for an integer \(N>0\), then \[
 v_p(D)+v_p(N)\ge r.
 \tag{T22.6}
\] The same conclusion holds for any finite sum of such constructions, even with unbounded source-cycle degrees. Multiplying or dividing by fixed rational normalization constants changes the inequalities only by their fixed \(p\)-valuations.

The assertion also holds for \(z\in\operatorname{CH}_m(U)\) of any dimension and an integral divisor polynomial \(K\in\operatorname{CH}^{m+2}(U\times B_r)\), so that its output still has codimension two. Source-only divisor factors merely enter the same integral contraction. If \(m<2\), the selected output is zero.

\subparagraph{B. No proper block of maximal-order factors can be returned with bounded denominator}\label{b.-no-proper-block-of-maximal-order-factors-can-be-returned-with-bounded-denominator}

Let \(I\subset\{1,2,3,4\}\) be nonempty and proper, and put \(k=|I|\). For \(i\in I\), let \(C_i\) be a CM abelian variety with rank-one rational \(E\)-homology and integral maximal-order action at \(p\). These varieties may vary. Put \[
 H_I=\bigotimes_{i\in I}H^1(C_i,\mathbf Q),\qquad
 V_I=\bigotimes_{i\in I}H^1(B_{i,r},\mathbf Q),\qquad
 V_{I^c}=\bigotimes_{j\notin I}H^1(B_{j,r},\mathbf Q).
 \tag{T22.7}
\] Let \(u\in H_I\otimes V_{I^c}\) be a rational Hodge class of type \((2,2)\) that is integral at \(p\). No assumption that \(u\) comes from the divisor-kernel construction is needed.

Suppose the rational Hodge homomorphism \(F:H_I\to V_I\) satisfies \[
 F(H_{I,\mathbf Z_p})\subseteq p^{-a}V_{I,\mathbf Z_p},
 \qquad a\ge0,
 \tag{T22.8}
\] where the lattices are the tensor products of the corresponding integral first cohomologies. Then \[
 \pi_W(F\otimes1)u
       \in p^{-a}\Theta(L^{\mathrm{fact}}_{r,p}).
 \tag{T22.9}
\] Consequently this procedure cannot produce \(w_r\) when \(r>a\). If \(u\) itself has an additional denominator of valuation \(b\), the necessary inequality is \(r\le a+b\).

This permits an arbitrary \textbf{joint} rational Hodge map on a block of one, two or three factors. In particular it applies to a degree-preserving algebraic correspondence on that block whose coefficient denominator is bounded at \(p\): take its degree-\(k\) action and the indicated Kunneth component. Integral Kunneth projections preserve the cohomology lattice, even though no integral Chow projector is being asserted. The first-factor isogeny pullbacks and pushforwards of the original one-factor case are included. Thus the positive uniform division result with maximal-order factors does not automatically transport back to all four mixed factors by bounded-denominator correspondences acting on a proper block.

\paragraph{Proof}\label{proof-3}

\subparagraph{1. The local dual lattice and pure Hodge rows}\label{the-local-dual-lattice-and-pure-hodge-rows}

Write \(e^\vee,f^\vee\) for the coordinate functionals on a single conjugate pair. Direct evaluation on the two generators in (T22.2) gives \[
 L_r^*=
 \{a e^\vee+b f^\vee:
       a\in\mathbf Z_p,\quad p^{-r}a+b\in\mathbf Z_p\}.
 \tag{T22.10}
\] The intersections with the two pure coordinate lines are therefore \[
 L_r^*\cap\mathbf Q_p e^\vee=p^r\mathbf Z_p e^\vee,
 \qquad
 L_r^*\cap\mathbf Q_p f^\vee=\mathbf Z_p f^\vee.
 \tag{T22.11}
\] Their direct sum is the dual of the maximal-order hull. The larger lattice (T22.10) allows cancellation between the two conjugate coordinates, whereas (T22.11) records precisely the individual integrality conditions.

Let \(D\) be an integral divisor on \(U\times B_{i,r}\), and take its mixed \(H^1(U)\otimes H^1(B_{i,r})\) component. Integral Kunneth makes this an integral tensor. Expand it in the pure integral source eigenbasis. Each source row is then an integral functional on the target lattice.

For a fixed source eigenline, that row cannot have nonzero entries in both of a target's conjugate eigenlines. Indeed, after extending scalars to \(E\) and then to \(\mathbf C\), all these tensor components are independent eigenspaces. The source eigenline has one pure type, either \((1,0)\) or \((0,1)\), while the two conjugate target lines have opposite types. At most one of their tensor products with that source line has type \((1,1)\). Since a divisor class is of type \((1,1)\), the other component vanishes. The same assertion is preserved under scalar conjugation. The coefficients are algebraic, so a vanished complex component is also zero in the split \(p\)-adic eigenbasis.

Thus, on each target conjugate pair, the row is supported on at most one coordinate. Substituting \(b=0\) or \(a=0\) in (T22.10) proves that every row lies in the direct sum (T22.11), pair by pair. In particular every individual target eigencomponent of the mixed divisor is integral separately, at \(p\).

This argument does not claim that a fixed source projector selects exactly one target eigenline. That claim could fail for imprimitive types. A row can contain several target eigenlines; the point is that it contains no conjugate pair, and the particular target lattice (T22.2) couples only conjugate pairs.

\subparagraph{2. Four source-output divisors}\label{four-source-output-divisors}

First suppose that each of four divisor factors has one source \(H^1\) leg and one output \(H^1\) leg, with the four output legs in the four distinct factors. By the preceding argument, every target eigencomponent of each factor lies in the appropriate pure integral line (T22.11).

Expand their product and contract the four source legs against the integral cycle class of \(z\). All contraction coefficients are \(p\)-adic integers: the source eigenbasis is integral, cup product is integral, and evaluation on an integral cycle is integral. Reordering the exterior factors only introduces signs. Thus each surviving output eigenmonomial has its coefficient in the product of its four individual pure-line coefficient lattices.

For a monomial in the selected diagonal orbit, a given conjugate pair has two \(e^\vee\) entries and two \(f^\vee\) entries, because its labels are \((1,1,c,c)\). Formula (T22.11) therefore forces its coefficient into \[
 (p^r)^2\mathbf Z_p=p^{2r}\mathbf Z_p.
\] This holds in all \(g\) pairs and for the conjugate monomials. The trace identification \(\Theta\) has no division by \(d\); its coefficients in these monomials are exactly the split coordinates of the field coefficient. Hence the selected projection belongs to \(\Theta(L^{\mathrm{fact}}_{r,p})\).

The same conclusion holds if a mixed divisor initially couples one output factor to several source factors: expand the finite sum of source rows and apply the preceding argument term by term. No \(E\)-linearity of the divisor correspondence was used.

\subparagraph{3. General divisor polynomials and off-orbit cancellation}\label{general-divisor-polynomials-and-off-orbit-cancellation}

The degree-four output Kunneth block under consideration has degree one in each \(B_{i,r}\). A divisor component with two legs in a single output factor cannot contribute to this block. A component with one leg in each of two distinct output factors can contribute only if its two selected eigenlabels form a rational \((1,1)\) subrelation.

More explicitly, a rational divisor component in the two selected lines is nonzero only if each scalar conjugate of their product has type \((1,1)\). In sign-function notation, this is exactly \(s_i+s_j=0\) as functions. Hypothesis (T22.1) excludes it. Every output-output divisor component therefore has zero contribution to the selected orbit.

There are four divisors, giving eight cohomological legs in total. To contribute after contraction with the surface cycle, four legs must lie on the source and four on the output. If there are \(a\) source-source components, \(b\) source-output components and \(c\) output-output components, then \[
 a+b+c=4,\qquad 2a+b=4,\qquad b+2c=4,
\] so \(a=c\). Once output-output components have been excluded, source-source components are excluded as well. The only possible contribution is therefore the four source-output case already treated. For the extension with source dimension \(m\) and \(m+2\) divisors, the output-leg condition is still \(b+2c=4\). The same indecomposability excludes \(c>0\), leaving \(b=4\) and \(a=m-2\). The extra source-source factors have integral cohomology and simply multiply the source cup product before evaluation on \(z\). The same target coefficient bound follows; when \(m<2\), no such nonnegative leg count is possible.

Integral Kunneth decompositions and finite sums preserve the asserted containment, which proves (T22.4) for any integral degree-four divisor polynomial. The calculation is performed before applying the selected eigenprojection. It neither assumes an integral projector on the varying output nor pays an unspecified target-dependent projector denominator. Cancellation of off-orbit components in a sum cannot change the local coefficient containment of the orbit components themselves.

If a rational combination with common denominator \(D\) equals \(Nw_r\), clearing its scalar denominators gives \[
 DNw_r\in\Theta(L^{\mathrm{fact}}_{r,p}).
\] The class of \(w_r\) in the saturation quotient has exact order \(p^r\), so \(p^r\mid DN\) in \(\mathbf Z_p\). This proves (T22.5) and (T22.6).

\subparagraph{4. Returning any proper block of maximal-order factors}\label{returning-any-proper-block-of-maximal-order-factors}

Multiply \(F\) by \(p^a\); it suffices to prove (T22.9) for \(a=0\). Since every \(C_i\) has maximal-order action at the split prime, \(H_{I,\mathbf Z_p}\) has a basis \(\{v_\tau\}\) of pure tensor eigenlines. Expand \[
 u=\sum_\tau v_\tau\otimes u_\tau .
\] Each \(u_\tau\) is integral in \(V_{I^c}\). Fix one selected diagonal monomial in the target and write \[
 S_I=\sum_{i\in I}s_i.
\] Only source rows with a nonzero coefficient in that selected block can contribute. For such a row, the complete Hodge sign function of \(v_\tau\) equals \(S_I\), because every scalar conjugate of the rational Hodge map \(F\) preserves its pure type. Since both source and target have weight \(k\), equality of sign sums is exactly equality of Hodge types.

Restrict the target block homology to the conjugate coordinate pair containing the selected label in each factor \(i\in I\). The product lattices are direct sums by these pairs, so the restriction retains integrality. Changing the selected labels to their conjugates in a nonempty subset \(J\subseteq I\) changes the sign function by \[
 -2\sum_{j\in J}s_j.
 \tag{T22.12}
\] Every such \(J\) is a proper subset of all four factors. Indecomposability therefore says that the selected vertex is the \textbf{unique vertex with its own sign function} in this block's conjugate-coordinate cube. No assertion about arbitrary pairs of other vertices is needed. The restricted vector \(F(v_\tau)\) consequently has only the selected coordinate monomial.

An integral multilinear form on a product of the lattices (T22.2) supported on one coordinate monomial has coefficient in the product of the individual pure coefficient lattices (T22.11). For necessity, insert \(p^{-r}e+f\) in every factor where the selected coordinate is \(e^\vee\), and a vector on which \(f^\vee\) is a unit in every selected \(f^\vee\) factor. If there are \(m\) selected \(e^\vee\) coordinates, the value is the coefficient times \(p^{-mr}\). Integrality forces the coefficient into \(p^{mr}\mathbf Z_p\). Sufficiency follows by multiplying the individual bounds. Thus the selected coefficient of \(F(v_\tau)\) has the product dual-hull weight for the factors in \(I\).

Now restrict \(u_\tau\) to the selected conjugate pairs in the complementary factors. The original complementary vertex has sign function \(-S_I\), because \(u\) has type \((2,2)\). Flipping a nonempty subset \(J\subseteq I^c\) again changes this function by (T22.12); every such subset is proper because \(I\ne\varnothing\). Thus the selected complementary vertex is also the only allowed vertex. Integrality of \(u_\tau\) forces its coefficient into the product of the complementary individual pure-line lattices by the same direct evaluation.

Multiplying the two coefficients gives the full fourfold pure-line bound. Summing over \(\tau\) preserves it. For the selected diagonal orbit this is exactly \(p^{2r}\) in every split coordinate, proving (T22.9) for integral \(F\), and hence for general \(a\). A denominator of \(u\) is handled by first multiplying \(u\) by its clearing integer.

An actual integral correspondence induces an integral cohomological map because its class, cup product and proper pushforward are integral. It is a Hodge map in the appropriate degree. Taking the indicated Kunneth components preserves these assertions, so the geometric applications follow with the correspondence's coefficient denominator as a permissible bound for \(a\). QED.

\paragraph{Exact scope of this additional reduction}\label{exact-scope-of-this-additional-reduction}

The theorem rules out two specific escapes from the computed saturation defect on the actual indecomposable principal family:

\begin{enumerate}
\def\labelenumi{\arabic{enumi}.}
\tightlist
\item
  replacing pure divisor kernels by arbitrary integral divisor kernels, while using the specified diagonal eigenprojection;
\item
  first constructing integral Hodge classes with a proper block of maximal-order output factors and then returning that block by a joint Hodge correspondence of bounded coefficient denominator.
\end{enumerate}

It does not prove that \(w_r\) lacks an integral algebraic-cycle representative. It does not control arbitrary correspondences acting jointly on all four output factors. For that full block, the selected vertex and its antipode are both balanced, so the unique-vertex argument fails at exactly that point. It also does not replace the missing actual integral map from the four-factor full-CM orbit to a generic \(K\)-Weil determinant plane. Such a map may have its own lattice content and is not supplied by a formal scalar relative trace.

An arbitrary non-diagonal rational projector that mixes other Hodge orbits into this one is not covered merely by calling it a projector. Its actual algebraic correspondence and its integral lattice map would have to be specified. The canonical diagonal eigenprojection in this theorem has been specified explicitly, and its integrality has not been assumed.

\paragraph{Inputs and check status}\label{inputs-and-check-status}

The proof uses the explicitly checked T6.16 indecomposable sign relation, the actual principal lattice realization of T6.19, and the exact lattice and primitive quotient classes of T6.21. The scalar-conjugate CM type convention is the supplied 032 paper, Section 2.2, printed p.7, equations (2.2)--(2.3); Lemma 2.4 and its complete proof, printed p.8, give the corresponding complete-type criterion for degree-one Hodge maps. Those exact passages were reread for this candidate. The remaining arguments are integral Kunneth, pure Hodge types of CM eigenlines, the \((1,1)\) property of divisor classes, and the displayed direct evaluations. No additional analytic or automorphic input is used.

The author has read the complete saved proof, including the general proper-block extension and the arbitrary-source-dimension divisor-kernel extension. The original one-factor version passed root and bounds review. The strengthened version passed author and root whole-proof review before promotion; the bounds strengthened review was pending at that time. No computation in a proof assistant is claimed.

\section{Bounds, Weil projectors, and the integral map lattice}

% Include-ready complete bounds proofs through B19; generated from lvalue_bounds.md. Requires amsmath, amssymb, hyperref and a URL/path package.
\providecommand{\tightlist}{\setlength{\itemsep}{0pt}\setlength{\parskip}{0pt}}
\section{DEG\_L research: L-value bounds, fixed Weil projector, bounded witness theorem}\label{deg_l-research-l-value-bounds-fixed-weil-projector-bounded-witness-theorem}

Date: 21 November 2026. Actor: GPT-6 Astra Pro, subagent \texttt{/root/lvalue\_bounds}; effort setting not exposed. Run: DEG\_L\_20261121. This is a working research contribution to the root paper, not a claim to solve T6 or T7. Every assertion below is either explicitly conditional, an elementary derived finding with proof, or an identified cited background result. All degrees are geometric intersection degrees, without factorial normalization.

\subsection{Executive finding}\label{executive-finding}

The strongest positive reduction obtained here is \textbf{B8}: in a polarized family of abelian varieties of Weil type of dimension \(2n\), a uniformly bounded-degree effective \(n\)-cycle with \emph{any nonzero Weil component} supplies a bounded-division representative of every fixed flat rational Weil class. The multiplier and both signed degrees are \(O(D^2)\), with explicit constants depending only on the fixed family, polarization type, field action, and target class. No varying full CM field or varying Hecke projector is required. The proof gives a fixed polynomial correspondence in the generic imaginary-quadratic action, a comass estimate on invariant harmonic forms, and an integral two-dimensional Cramer calculation. This removes the denominator and isolation gates once bounded geometric witnesses exist; the existence of those witnesses in the open families remains an actual hypothesis.

The analytic findings are complementary. The genuine \(U(1)\to U(3)\) Rallis formula uses \(L_E(3/2,\chi)\), not \(L_E(1,\chi)\); its global factor has elementary conductor-independent upper/lower bounds for fixed field degree. It does not compute the degree of the unnormalized 032 image surface. Edge-of-strip bounds under the admitted Alpöge input and conditional Masullo input apply only after identifying the correct normalized \(L\)-function, special point, local vectors, measure, and CM extraction coefficients.

\subsection{Source findings and exact scope}\label{source-findings-and-exact-scope}

\textbf{S032.} OpenAI family 032, \emph{The rational Hodge conjecture for CM abelian varieties}, §4.4 Lemma 4.4 (printed pp.20--21) fixes finite Schwartz data and permits all sufficiently deep principal levels; §7 proof of Theorem 2.5, printed p.46, explicitly states that pulling to a deeper principal cover multiplies the mixed integral by the cover degree. The theorem then selects a nonzero term and rational Hodge maps by span and finite expansion; it supplies no estimate for coefficients of those maps. These are statements used and rechecked in this run, not independent reproofs of 032.

\textbf{S-Rallis1.} Jaeho Haan, \emph{The Bessel Period of U(3) and U(2) involving a non-tempered representation}, arXiv:1403.5061v7, §3.3, Theorem 3.4, printed p.8, gives, in that paper's Tamagawa and splitting conventions, \[
\langle\theta(\bar f_1,\varphi_1),\theta(\bar f_2,\varphi_2)\rangle
 = \frac{L_E(3/2,\operatorname{BC}(\tau)\otimes\gamma^3)}{L_F(3,\epsilon_{E/F})}
 \prod_v Z_v^{\#}(1,f_{1,v},f_{2,v},\Phi_{1,v}).
\] Here \(\tau\) is an automorphic character of \(U(1)\) and \(\gamma\) is the chosen unitary splitting character; §3.1 (3.1)--(3.2) defines the normalization of local zeta integrals. §3.6.1, printed p.9, explicitly records absolute convergence and nonvanishing of the numerator at \(3/2\). This is a formula for integrated theta lifts with source character, and cannot be substituted for an arbitrary theta-kernel evaluation or wedge in 032 without comparison.

URL: https://arxiv.org/pdf/1403.5061 .

\textbf{S-Rallis2.} Gan--Qiu--Takeda, \emph{The Regularized Siegel--Weil Formula (The Second Term Identity) and the Rallis Inner Product Formula}, arXiv:1207.4709v3, Theorem 11.4, printed pp.54--55, under its explicit range and cuspidality/first-occurrence hypotheses, gives the unitary point \((m-n+1)/2\), with twisting character and normalized doubling integrals. Its range \(n<m\leq2n\) excludes \(U(1)\to U(3)\). The \(U(2)\to U(3)\) point is \(1\), but the standard degree-two \(L\)-function is not thereby a Hecke-character \(L\)-function. A source spectral decomposition or endoscopic factorization must be supplied. The theorem is used as background and its range checked, not reproved.

URL: https://arxiv.org/pdf/1207.4709 .

\textbf{S-Alpöge.} The user's Dirichlet GRH input is admitted. Exact supplied Lean declarations: \texttt{ALPOGE\_Main.lean}, \path|LFunction_eq_zero_re_of_isPrimitive : PrimitiveDirichletLZerosOnLine|, and \path|LFunction_ne_zero_of_half_lt_re : DirichletLNonvanishingRightOfHalf|; definitions in \texttt{ALPOGE\_Statements.lean}. The supplied \emph{Remark}, step 20, printed p.4, announces GRH; steps 9, 12, 18 concern lift norms, Rallis, and indices respectively. The full imported archive is not supplied, and this run did not compile it. No claimed new divisibility theorem is inferred from those steps or from the mapping file.

\textbf{S-Masullo.} The full draft is not supplied. The prompt attributes Hecke GRH to §1.2. The supplied \texttt{MASULLO\_STATEMENTS\_INDEX.md} is an index of openings, not theorem statements or proofs. It records Chapter 12, Theorem 12.5.1.1, PDF p.519, and Chapter 13, Theorem 13.3.1.1, PDF p.522, as trace formulas, and records an unfinished-looking opening labelled 13.3.1.2 (at the text-layer level). None is silently promoted to a precise theorem of Hecke GRH in this report. All use of Hecke GRH is conditional on the claim specified by the user, not certified from the index.

\subsection{B1. Elementary uniform bounds in the absolutely convergent half-plane}\label{b1.-elementary-uniform-bounds-in-the-absolutely-convergent-half-plane}

\textbf{Statement.} Let \(K\) have degree \(d\), and let \(\chi\) be a unitary Hecke character. For real \(\sigma>1\), using the finite Euler product with ramified factors equal to one, \[
 \left(\frac{\zeta(2\sigma)}{\zeta(\sigma)}\right)^d
 \leq |L_K(\sigma,\chi)|\leq \zeta(\sigma)^d.
\] More sharply, \(\zeta_K(2\sigma)/\zeta_K(\sigma)\leq|L_K(\sigma,\chi)|\leq\zeta_K(\sigma)\). No RH or GRH is needed.

\textbf{Proof.} For every unramified prime ideal \(\mathfrak p\), unitarity gives \(|\chi(\mathfrak p)|=1\). Thus \[
 (1+N\mathfrak p^{-\sigma})^{-1}\leq
 |1-\chi(\mathfrak p)N\mathfrak p^{-\sigma}|^{-1}
 \leq(1-N\mathfrak p^{-\sigma})^{-1}.
\] The same inequalities hold with the middle term replaced by one at ramified primes. Absolute convergence permits multiplication. The left product is \(\zeta_K(2\sigma)/\zeta_K(\sigma)\). Above each rational prime there are at most \(d\) prime ideals, each with norm at least that prime. Comparing each factor gives the stated bounds by rational zeta functions. All factors are positive and the products converge. QED.

\textbf{Application to the finite Euler factors.} In S-Rallis1 with \([F:\mathbb Q]=d\) and \([E:\mathbb Q]=2d\), the absolute value of the ratio of the finite Euler products lies between \[
 \frac{\zeta(3)^d}{\zeta(3/2)^{2d}}
 \quad\hbox{and}\quad
 \zeta(3/2)^{2d}\left(\frac{\zeta(3)}{\zeta(6)}\right)^d.
\] This follows by applying B1 separately to numerator and denominator. It controls no local \(Z_v^\#\) and no change of measure, input scale, or cycle normalization. When the \(L\)-symbols include their archimedean local factors, those factors must also be included; the actual cohomological situation is evaluated next.

\subsubsection{B1a. Actual infinity type and the location of the rank-one value}\label{b1a.-actual-infinity-type-and-the-location-of-the-rank-one-value}

For an actual \(H^{1,0}\) oscillator summand, choose the target splitting character \(\gamma=\mu\) from the supplied van Hoften--Sempliner §3.1.3, printed p.9. It has infinity type \(\gamma_v(z)=(z/|z|)^{\varepsilon_v}\) with \(\varepsilon_v\in\{1,-1\}\). Retwist the source as in T6.11 to Haan's compatible pair \((\gamma^3,\gamma)\). Then the source character \(\tau\) is trivial at every archimedean place. Consequently \[
 \chi=\operatorname{BC}(\tau)\gamma^3,
 \qquad \chi_v(z)=(z/|z|)^{3\varepsilon_v},
 \qquad L_v(s,\chi_v)=\Gamma_{\mathbb C}(s+3/2).
\tag{B1a.1}
\] Here \(\Gamma_{\mathbb C}(s)=2(2\pi)^{-s}\Gamma(s)\) and \(\Gamma_{\mathbb R}(s)=\pi^{-s/2}\Gamma(s/2)\).

\textbf{Proof.} The center of \(U(3)\) acts trivially on its symmetric domain and on its cotangent representation; at the other real places the differential-form coefficient has the trivial compact representation. Thus every target archimedean center acts trivially on the actual cohomological vector and on its generated representation. Haan, §2.1, Remark 2.1, printed p.4, states that for a one-dimensional member the images of that member and of the center of the other member coincide in the metaplectic group under his power-compatible splittings. Equivalently their scalar actions in a nonzero local theta quotient match, possibly after inversion according to the theta convention. It follows that the source \(U(1)\) character is trivial. Its base change \(z\mapsto\tau_v(z/\bar z)\) is therefore trivial as well. The displayed infinity type follows. The standard complex local factor can also be read directly from the Gaussian Mellin integral: after canceling the angular exponent by the appropriate polynomial, its radial integral is a nonzero fixed multiple of \(\int_0^\infty r^{2s+|3\varepsilon_v|-1}e^{-2\pi r^2}\,dr\), giving \(\Gamma_{\mathbb C}(s+3/2)\). QED.

Since \(\epsilon_{E/F,v}=\operatorname{sgn}\) at every real place of \(F\), its local factor at \(3\) is \(\Gamma_{\mathbb R}(4)\). Thus the all-place product of archimedean ratios in the rank-one formula is exactly \[
 \left(\frac{\Gamma_{\mathbb C}(3)}{\Gamma_{\mathbb R}(4)}\right)^d
 =(2\pi)^{-d}.
\tag{B1a.2}
\] With these explicit local \(L\) conventions, multiply both finite-product bounds above by \((2\pi)^{-d}\) when using all-place \(L\)-functions. This is a fixed factor, independent of the finite conductor. A separately completed functional equation may put a further \(q^{s/2}\) into a symbol \(\Lambda(s)\); that completion is not the all-place local product appearing in the Rallis normalization and must not be inserted silently.

The unitary functional-equation center is \(1/2\), while the Rallis point is \(3/2\), in the absolutely convergent half-plane. If \(\gamma=\psi\|\cdot\|_E^{1/2}\) with the algebraic character \(\psi\) supplied in the proof of van Hoften--Sempliner Lemma 3.1.5, pp.9--10, put \(\Psi=\operatorname{BC}(\tau)\psi^3\). Then \[
 L(s,\chi)=L(s+3/2,\Psi).
\] The arithmetic normalization of weight three has center \(2\), and the Rallis value becomes \(L(3,\Psi)\). Its gamma factors and those of the reflected conjugate are \(\Gamma_{\mathbb C}(s)\) and \(\Gamma_{\mathbb C}(4-s)\), respectively, so the critical integers are \(1,2,3\). Thus \(3\) is the rightmost critical integer, not the central one. This arithmetic description is a normalization statement; the unconditional bound uses the unitary finite Euler product at \(3/2\).

\subsection{B2. Exact edge bounds and what each GRH input buys}\label{b2.-exact-edge-bounds-and-what-each-grh-input-buys}

Lamzouri--Li--Soundararajan, \emph{Conditional bounds for the least quadratic non-residue and related problems}, arXiv:1309.3595v3, Theorem 1.5 (printed pp.3--4), gives for primitive Dirichlet characters of conductor \(q\geq10^{10}\) under GRH, with \[
 A(q)=\log\log q-\log2+\tfrac12+1/\log\log q,
\] \[
 |L(1,\chi)|\leq 2e^\gamma A(q),\qquad
 |L(1,\chi)|^{-1}\leq\frac{12e^\gamma}{\pi^2}
 \left(A(q)+\frac{14\log\log q}{\log q}\right).
\] The constant in the second inequality is \(12e^\gamma/\pi^2\), not \(12e^{2\gamma}/\pi^2\). An older hosted text extraction returned the latter; the latest arXiv version and Lumley's Theorem 1.1 provide the correct former value. Small conductors require their own finite verification or a separate bound. Under the user's admitted Alpöge theorem, this consequence no longer has a GRH condition. URL: https://arxiv.org/pdf/1309.3595.

For primitive nontrivial unitary Hecke characters whose completed \(L\)-function is entire, write the completed function in degree \(d=[K:\mathbb Q]\) over \(\mathbb Q\), with conductor \(q(\chi)=|D_K|N\mathfrak f_\chi\) and archimedean parameters \(\kappa_j\). Allysa Lumley, \emph{Explicit bounds for L-functions on the edge of the critical strip}, arXiv:1804.09850v1, §1.1 and Theorems 2.1--2.2 (printed pp.3--4), applies to \[
 C(\chi)=\frac{q(\chi)}{\pi^d}\prod_{j=1}^d\left|\frac{1+\kappa_j}{2}\right|.
\] For fixed degree, its bounds have orders \((\log\log C)^{d}\) and \((\log\log C)^{-d}\); its explicit formulas apply when \(\log C\geq23d\). Unit modulus of unramified Hecke parameters supplies the required Ramanujan hypothesis; the remaining GRH input is conditional on Masullo for general \(K\). Pole cases (including the trivial character) are excluded. Field degree and archimedean type cannot be hidden while allowing them to vary. URL: https://arxiv.org/pdf/1804.09850.

These are known results applied, not new theorems. In particular, this report does not assert an edge bound for a nonabelian \(U(2)\) standard \(L\)-function from Hecke GRH. Nor does GRH give a positive lower bound at the center: the functional equation can force a central zero when a self-dual root number is \(-1\).

\subsection{B3. Level, vector scale, and arithmetic conductor are distinct}\label{b3.-level-vector-scale-and-arithmetic-conductor-are-distinct}

\textbf{Statement.} If \(\pi:Y\to X\) is a finite étale cover of degree \(d\) and \(\eta\) is a top-degree form on \(X\), then \(\int_Y\pi^*\eta=d\int_X\eta\). Consequently a nonzero pulled-back theta form has unnormalized Petersson norm squared multiplied by \(d\). If its underlying automorphic representation, infinity type, and finite vector are fixed, the representation's arithmetic and analytic conductors stay fixed.

\textbf{Proof.} Integration over evenly covered coordinate sets and a partition of unity gives the integral identity. Apply it to the positive top form defining the norm. Conductors are invariants of the fixed representation and archimedean parameters, not of which smaller compact open subgroup happens to fix its vector. QED.

This is also explicitly recorded for the mixed period in 032 §7 p.46. It is an exact level-growth statement for a chosen presentation, not a lower bound for the intrinsic minimum complexity \(c(s)\) in Route 2. Similarly scaling a Schwartz input by \(u\) scales a theta lift by \(u\) and its norm squared by \(|u|^2\), while leaving its representation's \(L\)-function unchanged. The local integrals and/or input pairing carry that factor.

\subsection{B4. Polarized isogenies: rational degree can stay fixed while denominators grow}\label{b4.-polarized-isogenies-rational-degree-can-stay-fixed-while-denominators-grow}

\textbf{Statement.} Let \(\psi:A\to B\) be an isogeny of \(g\)-dimensional polarized abelian varieties. Suppose \(\psi^*c_1(L_B)=c\,c_1(L_A)\) and \(L_A,L_B\) have the same polarization degree. For a codimension-\(p\) cycle \(Z\), \[
 \deg\psi=c^g,\qquad
 \deg_{L_A}(\psi^*Z)=c^p\deg_{L_B}Z.
\]

\textbf{Proof.} Integrating the \(g\)th power of the polarization identity gives the first equality. By projection formula, \[
 \int_A\psi^*[Z]c_1(L_A)^{g-p}
 =c^{-(g-p)}\deg\psi\int_B[Z]c_1(L_B)^{g-p},
\] which is the second. QED.

The rational cycle \(c^{-p}\psi^*Z\) has unchanged degree but its displayed coefficient has denominator \(c^p\) when \(c\) is an integer. For multiplication by \(m\) on a fixed variety, \(c=m^2\). This is not a proof that all representatives have large denominators: it is a precise diagnosis of the natural pullback representative.

\subsection{B5. A rigorous norm-to-degree bridge for surface maps}\label{b5.-a-rigorous-norm-to-degree-bridge-for-surface-maps}

\textbf{Statement.} Let \(Y\) be a smooth projective surface, \(\omega\) a Kähler class, \(V=\int_Y\omega^2\), and \(f:Y\to A\) a morphism with two-dimensional image \(Z\) and generic degree \(\delta\). Choose a translation-invariant representative \(\omega_A\in c_1(L)\) of a polarization, and invariant holomorphic forms \(\alpha_1,\ldots,\alpha_g\) such that \(\omega_A=i\sum_j\alpha_j\wedge\bar\alpha_j\). Put \[
 E_f=\int_Y f^*\omega_A\wedge\omega
 =\sum_j\|f^*\alpha_j\|_\omega^2,\qquad
 \|u\|_\omega^2=i\int_Yu\wedge\bar u\wedge\omega.
\] Then \[
 \deg_L Z=\frac{\int_Y(f^*\omega_A)^2}{\delta}
 \leq\frac{E_f^2}{\delta V}.
\]

\textbf{Proof.} Proper pushforward gives \(f_*[Y]=\delta[Z]\) and projection formula gives the equality. With \(\beta=f^*\omega_A\), decompose \(\beta=(E_f/V)\omega+\beta_0\), where \(\int\beta_0\omega=0\). The Hodge index theorem on the surface gives \(\int\beta_0^2\leq0\), so \(\int\beta^2\leq E_f^2/V\). The displayed energy identity follows by linearity from the specified polarization representative. QED.

This is a direct application of the classical surface Hodge index theorem in a form adapted to the present problem, not a claim to reprove that theorem. Under a cover of degree \(d\), \(E_f\) and \(V\) both scale by \(d\). Thus the bound \(E_f^2/V\) for the degree of the pushforward class scales by \(d\). For the reduced image, the generic mapping degree \(\delta\) also scales by \(d\), and the displayed bound remains unchanged; this matches B3 exactly.

For the sufficiently deep principal ball-quotient surfaces used in 032, the actual pullbacks \(f^*\alpha_j\) do have finite expansions \(\sum_a b_{ja}\theta_a\) in genuine integrated theta lifts from characters of one-dimensional unitary groups. The existence of these expansions is supplied by Bergeron--Millson--Mœglin, \emph{The Hodge conjecture and arithmetic quotients of complex balls}, Acta Math.216 (2016), Theorem 9.1, pp.73--75, with the range in §9.1, p.72: for \(m=3\) and \((a,b)=(0,1)\), its condition is \(4<6\). The supplied van Hoften--Sempliner paper makes the all-\(H^{1,0}\) spanning conclusion explicit in the proof of Proposition 4.1.8, pp.15--16; Proposition 3.1.7, p.11, identifies the oscillator constituents. The published BMM proof and the supplied Proposition 4.1.8 proof were read. This is imported classification/spanning, not a new proof of those theorems.

For completeness, only finitely many direct-limit theta vectors are needed for each fixed form. Put them at a common finer level and average over the desired compact open subgroup. The original form is invariant, and equivariance moves the finite average to the Schwartz inputs, so the expansion exists at that level. Equality in \(H^{1,0}\) is equality of holomorphic forms by the injectivity into Betti cohomology for a compact Kähler variety. If necessary one first deepens the principal level so the arithmetic \(U\) and \(SU\) groups agree: their determinant values are norm-one algebraic units, hence roots of unity, and sufficiently deep congruence kills the finite set of nontrivial roots. T6.9 gives the full argument and its scope.

With norm squares \(Q_a\) computed by a valid rank-one Rallis formula, and with the scalar/geometric metric and measure conversions included in \(Q_a\), the triangle inequality gives the concrete bound \[
 \deg_L Z\leq\frac{1}{\delta V}
 \left\{\sum_j\left(\sum_a|b_{ja}|\sqrt{Q_a}\right)^2\right\}^2.
\] The existence gate is therefore removed. What remains is a quantitative bound on the coefficients \(b_{ja}\) relative to the integral Hom lattice, together with CM periods, test-vector factors and metric/measure normalization. Orthonormalizing the theta basis does not control those integral coefficients. B10 below reduces them to an integral energy-lattice covolume estimate, and B11 shows exactly why that estimate cannot be inferred from finite generation or orthonormalization alone.

\subsection{B6. Fixed generic-field projector onto Weil classes}\label{b6.-fixed-generic-field-projector-onto-weil-classes}

\textbf{Setup.} Let \(A\) be a polarized abelian variety of dimension \(2n\), carrying an action of an order \(\mathcal O\subset K\) in an imaginary quadratic field. Assume the Rosati involution restricts to conjugation on \(K\) and the two eigenspaces in \(H^{1,0}\) both have dimension \(n\) (Weil type). Let \[
 W=\bigwedge_K^{2n}H^1(A,\mathbb Q)\subset H^{2n}(A,\mathbb Q).
\] Here the inclusion means the rational rank-two subspace whose complexification is the sum of the two extreme \(K\)-eigenlines. Choose a nonzero \(a\in\mathcal O\) such that \(a/\bar a\) is not a root of unity, and let \(T=[a]^*\) on cohomology. Such an \(a\) exists: only finitely many arguments arise from roots of unity in \(K\), whereas the arguments of nonreal elements \(M+\tau\) for a fixed nonreal \(\tau\in\mathcal O\) vary with integers \(M\) and tend to zero.

\textbf{Statement.} There is a polynomial \(P(X)\in\mathbb Z[X]\) and a positive integer \(R\), depending only on \((K,n,a)\), such that \(e=P(T)/R\) is the projector onto \(W\) on \(H^{2n}\). It is induced by a fixed rational algebraic correspondence and is self-adjoint for the Poincaré pairing on that degree. Its denominator is independent of the CM point or Hecke level. One possible choice is \[
 \begin{aligned}
 \lambda_r&=a^r\bar a^{2n-r}\qquad(0\leq r\leq2n),\\
 f_W(X)&=(X-a^{2n})(X-\bar a^{2n}),\\
 f_R(X)&=\prod_{r=1}^{2n-1}(X-\lambda_r),\\
 R&=\operatorname{Res}(f_W,f_R)>0,
 \end{aligned}
\] with \(P=vf_R\) for integral polynomials satisfying \(u f_W+v f_R=R\).

\textbf{Proof.} Decomposing \(H^1_\mathbb C=V_a\oplus V_{\bar a}\), each summand has dimension \(2n\), and \[
 H^{2n}_\mathbb C=\bigoplus_{r=0}^{2n}\bigwedge^rV_a\otimes\bigwedge^{2n-r}V_{\bar a}.
\] The corresponding eigenvalues are \(\lambda_r\). They are distinct since \(a/\bar a\) has infinite order. Each polynomial displayed is fixed by conjugation and has algebraic-integer coefficients, hence lies in \(\mathbb Z[X]\); the two factors are coprime. The Sylvester matrix for these monic polynomials has determinant their resultant. Multiplying that matrix by its integral adjugate and extracting the column corresponding to the constant polynomial gives integral \(u,v\) with \(u f_W+v f_R=R\) (a simultaneous sign change adjusts the determinant convention). Thus \(P/R\) takes value one at the two extreme eigenvalues and zero at every other eigenvalue. Since \(T\) is diagonalizable, it is the asserted projector.

Every power \(T^j\) is induced by pullback along the algebraic isogeny \([a^j]\), so the polynomial defines an algebraic correspondence. Its cohomological projector property is what is asserted: \textbf{no Chow-group identity \(e^2=e\) is claimed} merely from this calculation. For a full homological-motive statement, compose it with the algebraic degree projector \(\pi_{2n}=\prod_{j\ne2n}([2]^*-2^j)/(2^{2n}-2^j)\), where \(0\leq j\leq4n\); this acts as the identity on \(H^{2n}\) and as zero in all other degrees. Thus it defines an idempotent on \(h^{2n}\) in homological motives, and on numerical cycle classes, without claiming a Chow idempotence relation.

To prove self-adjointness on this degree, the Poincaré pairing of two \(T\)-eigenvectors can be nonzero only if \(\lambda_r\lambda_s=\deg[a]=N(a)^{2n}\): this follows by applying \([a]^*\) to their product in top degree. Infinite order of \(a/\bar a\) makes this condition exactly \(r+s=2n\). Consequently the two extreme eigenlines pair only with each other and are orthogonal to all the intervening eigenspaces. Projection onto their sum is Poincaré-self-adjoint. Finally, \(R=f_R(a^{2n})\overline{f_R(a^{2n})}>0\). QED.

A useful explicit expression is \[
 R=N(a)^{n(2n-1)}\prod_{k=1}^{2n-1}|a^k-\bar a^k|^2.
\] It follows by factoring \(f_R(a^{2n})=a^{n(2n-1)}\prod_{k=1}^{2n-1}(a^k-\bar a^k)\) and taking the field norm.

If \(P=\sum_jb_jX^j\), put \[
 B_P=\sum_j|b_j|N(a)^{jn}.
\] For effective codimension-\(n\) cycles \(Z\) of degree \(D\), \(P([a]^*)Z\) is an integral signed cycle whose positive and negative parts both have degree at most \(B_PD\). Indeed \([a]^*c_1(L)=N(a)c_1(L)\) by the Rosati condition, so B4 gives degree multiplier \(N(a)^n\) for pullback on codimension \(n\); use it term by term. This works fibrewise in any family with the fixed action and polarization data.

\subsection{B7. Uniform Weil-component bound from effective degree}\label{b7.-uniform-weil-component-bound-from-effective-degree}

\textbf{Statement.} In B6 put \(\omega=c_1(L)\), \(V_A=\int_A\omega^{2n}/(2n)!=\chi(L)\), and \[
 q(u,v)=(-1)^n\int_Au\wedge v\quad(u,v\in W_\mathbb R).
\] Then \(q\) is positive definite. For every effective \(n\)-cycle \(Z\) of degree \(D=\int_Z\omega^n\), \[
 \|e[Z]\|_q\leq\frac{D}{n!\sqrt{V_A}}.
\] The constant depends only on polarization type and \(n\), hence is uniform on the whole Weil family, without a compactness restriction on its base.

\textbf{Proof.} The Rosati identity makes \(\omega\) a mixed-eigenvalue class in \(V_a\otimes V_{\bar a}\): its \([a]^*\) eigenvalue is \(N(a)\) on degree two, different from \(a^2\) and \(\bar a^2\). Wedge product by \(\omega\) kills each extreme degree-\(2n\) eigenline because it would require \(2n+1\) vectors in one \(2n\)-dimensional eigenspace. Thus \(W\) is primitive. Weil type makes both lines of Hodge type \((n,n)\). The Hodge--Riemann bilinear relations give positivity of \(q\); the Hodge star on these primitive middle-degree real classes is \(*=(-1)^n\).

Give \(A\) its flat translation-invariant metric associated with \(\omega\). Every harmonic form is translation invariant, so for its harmonic representative \(\eta\in W_\mathbb R\), the pointwise Euclidean norm is constant and equals \(\|\eta\|_q/\sqrt{V_A}\). Its comass is at most its Euclidean norm. The complex analytic cycle \(Z\) has volume \(D/n!\) by Wirtinger's identity, with multiplicities; its singular locus has measure zero. Therefore \[
 \left|\int_Z\eta\right|\leq\frac{D}{n!\sqrt{V_A}}\|\eta\|_q.
\] By B6, \(e\) is self-adjoint for Poincaré pairing, so for \(w=e[Z]\), \[
 \|w\|_q^2=|q(w,w)|=\left|\int_A[Z]\wedge w\right|=\left|\int_Zw\right|.
\] If \(w=0\) the result holds; otherwise apply the preceding inequality with \(\eta=w\) and divide by \(\|w\|_q\). QED.

\subsection{B8. Effective witness implies bounded division, with quadratic bounds}\label{b8.-effective-witness-implies-bounded-division-with-quadratic-bounds}

\textbf{Statement.} Keep a connected polarized Weil-type family as in B6, with a globally defined \(\mathcal O\)-action (after finite base change if necessary), fixed \(n,K,a,L\), and fixed rational flat target \(t\in W\). Put \[
 \Lambda=W\cap H^{2n}(A,\mathbb Z),\qquad \delta_\Lambda=\det(q(e_i,e_j))>0
\] for any integral basis of this rank-two saturated lattice. Choose \(b\geq1\) with \(\tau=bt\in\Lambda\), and write \(h=N(a)^n\). The constants \(V_A\), \(\delta_\Lambda\), and \(\|\tau\|_q\) are locally constant along the family and so fixed on a connected base.

Suppose a fibre contains an effective \(n\)-cycle \(Z\) of degree \(D\) with \(e[Z]\ne0\). Then there are an integer \(N\geq1\) and effective \(Z^+,Z^-\) with \[
 Nt=[Z^+]-[Z^-],
\] where, with \(M=RD/(n!\sqrt{V_A})\), \[
 N\leq\frac{b h M^2}{\sqrt{\delta_\Lambda}},\qquad
 \deg Z^\pm\leq\frac{2h\|\tau\|_q M B_P D}{\sqrt{\delta_\Lambda}}.
\] Thus both are \(O(D^2)\), uniformly in the fibre. In particular a uniform bound for degrees of nonzero Weil witnesses on a Zariski-dense set gives bounded division on that set. The proper parameter-space argument of Route 2 then applies, once supplied separately with its stated hypotheses.

\textbf{Proof.} Let \(V_Z=P([a]^*)Z\) as an integral signed cycle and \(v=[V_Z]=R e[Z]\). This is a nonzero element of \(\Lambda\), with \(\|v\|_q\leq M\) by B7. The characteristic polynomial of \(T\) on \(W\) is \(f_W\). Since \(a/\bar a\) is not a root of unity, \(a^{2n}\) is nonreal, so \(f_W\) is irreducible over \(\mathbb Q\). Hence \(v,Tv\) are a rational basis of \(W\): otherwise \(v\) would be a rational eigenvector of \(T\) with a rational eigenvalue.

Use any oriented integral basis of \(\Lambda\). Its determinant \[
 \Delta=\det(v,Tv)\in\mathbb Z\setminus\{0\}
\] therefore satisfies \(\Delta\tau=m v+\ell Tv\), where the integral Cramer numerators are \(m=\det(\tau,Tv)\) and \(\ell=\det(v,\tau)\). On \(W\), \(T\) is a similitude of \(q\) of multiplier \(\deg[a]=N(a)^{2n}=h^2\), so \(\|Tv\|_q=h\|v\|_q\). The area of a parallelogram in this metric is \(\sqrt{\delta_\Lambda}\) times its lattice determinant. The elementary area bound consequently gives \[
 |\Delta|\leq hM^2/\sqrt{\delta_\Lambda},\quad
 |m|\leq\|\tau\|_q hM/\sqrt{\delta_\Lambda},\quad
 |\ell|\leq\|\tau\|_q M/\sqrt{\delta_\Lambda}.
\] Let \(N=b|\Delta|\), and let \(C=\operatorname{sgn}(\Delta)(mV_Z+\ell[a]^*V_Z)\). Then \([C]=Nt\). By the final part of B6 the positive and negative degrees of \(V_Z\) are at most \(B_PD\), and those of \([a]^*V_Z\) are at most \(hB_PD\). Splitting the signed cycle \(C\) into positive and negative parts thus bounds each by \[
 (|m|+h|\ell|)B_PD\leq 2h\|\tau\|_q M B_PD/\sqrt{\delta_\Lambda}.
\] This proves the claim. The intersection lattice and its pairing vary as an integral local system, so its discriminant and the norm of a fixed flat target are constant; \(V_A\) is fixed by polarization type. No choice of a global basis is needed to make the constants uniform. QED.

\textbf{Stronger exact index identity.} If \(\mu=a^{2n}\) under the fixed complex embedding, then \[
 |\Delta|=\frac{|\operatorname{Im}\mu|}{\sqrt{\delta_\Lambda}}\,\|v\|_q^2.
\] Indeed \(T\) is multiplication by the nonreal scalar \(\mu\) on the real plane \(W\), its characteristic polynomial has trace \(2\operatorname{Re}\mu\), and it is a positive-orientation \(q\)-similitude. In a \(q\)-orthonormal basis it is the rotation-dilation matrix for \(\mu\) or \(\bar\mu\); the area of \((v,Tv)\) is \(|\operatorname{Im}\mu|\|v\|^2\). Dividing by the lattice covolume proves the formula. This makes the last arithmetic index precise, without claiming it is automatically an automorphic \(L\)-value.

\textbf{The quadratic size is sharp for extraction from one lattice vector.} In the positive lattice \(\Lambda=\mathbb Z[i]\), let \(v=m+i\) and target \(\tau=1\). If \(N\tau\) belongs to the principal ideal \(v\mathbb Z[i]\), then \[
 \frac{N}{m+i}=\frac{Nm}{m^2+1}-i\frac{N}{m^2+1}\in\mathbb Z[i],
\] so \(m^2+1\mid N\). Since \(\|v\|^2=m^2+1\), extraction using \(v\) and any fixed integral multiplication operator can require an integer of quadratic size in \(\|v\|\). This is an arithmetic sharpness statement for that extraction method. It does not establish growth of intrinsic cycle complexity, since other geometric cycles might represent the target more efficiently.

\textbf{What remains.} The theorem does not create \(Z\) in the non-CM fibres, and 032 gives neither a uniform degree bound for such \(Z\) on a dense orbit nor a deformation of one actual full cycle class. It proves that a separate unbounded full-CM-field projector is unnecessary in the Weil-family route. Once geometric witness degrees are bounded, the two-dimensional integral extraction is automatically bounded.

\subsection{B9. A precise failure of archimedean norm control for arbitrary projectors}\label{b9.-a-precise-failure-of-archimedean-norm-control-for-arbitrary-projectors}

On the lattice \(\mathbb Z^2\) with Euclidean pairing, the orthogonal projector to \(\mathbb Q(m,1)\) is \[
 E_m=\frac1{m^2+1}\begin{pmatrix}m^2&m\\m&1\end{pmatrix}.
\] It has operator norm one, while its exact common denominator is \(m^2+1\), because its lower-right entry is \(1/(m^2+1)\). The calculation follows by the rank-one projector formula \(vv^t/(v^tv)\). This example already lives in \(\operatorname{End}(\mathbf1\oplus\mathbf1)\) in homological or numerical motives: all matrices are algebraic correspondences between two copies of a point. It refutes any abstract inference from bounded archimedean projector norm to bounded integral denominator. It does not refute B6, which fixes the generic-field projector, or the existence of other bounded representatives in a geometric family.

Similarly \(x_m=(1,1/m)\) has norm in \([1,\sqrt2]\) while its pairing with \((0,1)\) is \(1/m\). Lower bounds for two individual norms therefore give no lower bound for the mixed period needed in 032. The positive continuous comparison used by 032 is a separate input, whose quantitative approximation costs must be controlled.

\subsection{B10. An integral Hom covolume bound replaces individual coefficient bounds}\label{b10.-an-integral-hom-covolume-bound-replaces-individual-coefficient-bounds}

\textbf{Setup.} Let \(Y\) be a smooth projective variety of dimension \(d\ge1\), with an integral ample class \(H\), and fix a base point for its Albanese map \(\alpha:Y\to B=\operatorname{Alb}(Y)\). Let \((A_i,L_i)\) be polarized abelian varieties, for \(1\le i\le k\). Put \[
 \Lambda_i=\operatorname{Hom}(B,A_i),\qquad
 E_i(u)=\int_Y (u\alpha)^*c_1(L_i)H^{d-1},\qquad
 V=\int_YH^d.
\] Here maps are homomorphisms; translations do not affect the forms, energies or degrees. Assume \(r_i=\operatorname{rank}\Lambda_i>0\). Let \(q_i\) be the real quadratic metric with square norm \(E_i\), let \(C_i\) be the Euclidean covolume of \(\Lambda_i\) for \(q_i\), and let \(\mu_i>0\) be any known lower bound for \(E_i(u)\) for nonzero integral \(u\). One may always take \(\mu_i=1\). Set \[
 v_r=\frac{\pi^{r/2}}{\Gamma(1+r/2)},\qquad
 R_i=\frac{2^{r_i}}{v_{r_i}}\frac{C_i}{\mu_i^{(r_i-1)/2}}.
\]

\textbf{Statement.} The lattice \(\Lambda_i\) contains \(r_i\) independent homomorphisms \(u_{ij}\), spanning \(\Lambda_i\otimes\mathbb Q\), such that \(E_i(u_{ij})\le R_i^2\). Suppose a specified real-multilinear period \(P\) on \(\prod_i(\Lambda_i\otimes\mathbb R)\) is nonzero. Some tuple of these homomorphisms has \(P(u_{1j_1},\ldots,u_{kj_k})\ne0\). If the period is the integral of a wedge of \(2d\) pulled-back one-forms, its nonvanishing implies that the product map \[
 F=(u_{1j_1}\alpha,\ldots,u_{kj_k}\alpha):Y\longrightarrow\prod_i A_i
\] has a \(d\)-dimensional image \(Z\). For the product polarization \(L=\bigotimes_i\operatorname{pr}_i^*L_i\), and the actual generic degree \(\delta_F\) of \(F\) onto \(Z\), \[
 \deg_L Z\le\frac{(\sum_iR_i^2)^d}{\delta_F V^{d-1}}.
\tag{B10.1}
\] For 032, \(d=2\) and \(k=4\). More generally, if a target polarization \(L'\) satisfies \(B c_1(L)-c_1(L')\) nef for a known \(B>0\), the right side is multiplied by \(B^d\).

\textbf{Proof.} Homomorphisms inject into homomorphisms of their integral first homology lattices, so \(\Lambda_i\) is a finitely generated free abelian group. Pullback on invariant one-forms depends linearly on the homomorphism. Writing a polarization as a positive sum of \(i\eta\wedge\bar\eta\) shows that \(E_i\) is a positive semidefinite real quadratic form. There is no nonzero integral vector in its kernel: if all pulled-back one-forms vanish on \(Y\), the map \(u\alpha\) has zero derivative and is constant; because \(\alpha(Y)\) generates \(B\) and both maps preserve the chosen origins, \(u=0\). The intersection expression for \(E_i\) is integral, so polarization gives a rational Gram matrix in every integral basis. A nonzero real kernel of a rational matrix contains a nonzero rational vector, and clearing its denominator would give a nonzero integral kernel vector. Thus the real kernel is zero and \(q_i\) is positive definite. Its positive values on nonzero lattice vectors are at least one.

Let \(\lambda_1\le\cdots\le\lambda_{r_i}\) be the successive minima of the Euclidean unit ball for this metric. Minkowski's second theorem gives \[
 v_{r_i}\prod_j\lambda_j\le 2^{r_i}C_i.
\] Since every \(\lambda_j\ge\sqrt{\mu_i}\), one obtains \(\lambda_{r_i}\le R_i\). The minima are attained, giving the required independent maps. They need not be an integral basis, and none is asserted: a rational spanning set is enough. If \(P\) vanished on every tuple from these spanning sets, multilinearity would make it vanish identically, a contradiction. A map with image dimension smaller than \(d\) pulls every smooth real \(2d\)-form back to zero, so the nonzero wedge period gives the image-dimension assertion.

For the degree bound, write \(\beta=F^*c_1(L)\), a nef class. Its energy is \(E=\beta H^{d-1}=\sum_iE_i(u_{ij_i})\le\sum_iR_i^2\). The mixed-intersection inequality \[
 \beta^d\le\frac{(\beta H^{d-1})^d}{(H^d)^{d-1}}
\tag{B10.2}
\] follows directly from surface Hodge index: for rational ample \(\beta,H\), take smooth complete-intersection surfaces of suitable very ample multiples of \(j-1\) copies of \(\beta\) and \(d-j-1\) copies of \(H\), for \(1\le j\le d-1\). Hodge index on each surface gives \(a_j^2\ge a_{j-1}a_{j+1}\) for \(a_j=\beta^jH^{d-j}\). These inequalities imply \(a_d\le a_1^d/a_0^{d-1}\). Approximating nef \(\beta\) by \(\beta+\varepsilon H\) proves the required nef case by continuity. For \(d=1\) it is equality, and for \(d=2\) it is B5. Projection formula gives \(\beta^d=\delta_F\deg_L Z\), proving (B10.1). Finally nefness of \(Bc_1(L)-c_1(L')\) gives \(E_{L'}\le B E_L\) after pullback and intersection with \(H^{d-1}\), so the same argument proves the last assertion. QED.

\textbf{Exact imported lattice source.} Martin Henk, \emph{Successive minima and lattice points}, arXiv:math/0204158v1, Theorem 1.3, printed p.2, and its proof in §3, printed pp.5--7, \url{https://arxiv.org/pdf/math/0204158}. The statement and complete proof were read; only the displayed upper-product inequality is imported here.

\textbf{What determinant is required.} On the real Hom space, let \(w_1,\ldots,w_r\) be a basis of normalized analytic vectors, with real Gram matrix \(G\). If the columns of \(B\) express an integral basis of \(\Lambda\) in these coordinates, then exactly \[
 C_\Lambda^2=\det(B^{\mathsf t}GB)=|\det B|^2\det G.
\tag{B10.3}
\] Replacing the analytic basis by an orthonormal one replaces \(B\) by \(G^{1/2}B\) and preserves this covolume. Rallis can contribute to \(G\) after the CM-period and measure comparisons are fixed; it does not identify \(B\). If \(\Gamma=\sum\mathbb Zw_j\) is commensurable with \(\Lambda\), then \[
 \frac{C_\Lambda}{C_\Gamma}
 =\frac{[\Gamma:\Gamma\cap\Lambda]}{[\Lambda:\Gamma\cap\Lambda]}.
\] Thus a proved integral analytic sublattice \(\Gamma\subset\Lambda\) with bounded Gram determinant already gives \(C_\Lambda\le C_\Gamma\); no saturation-index upper bound is then needed. If instead \(\Lambda\subset\Gamma\), its index contributes in the numerator and must be bounded. A uniform integer \(M\) with \(M\Gamma\subset\Lambda\) gives the sufficient bound \(C_\Lambda\le M^r\sqrt{\det G}\). This identifies the exact direction of the integral comparison that a theta-index theorem would have to supply.

\textbf{Variation and normalization.} The ranks \(r_i\) may grow when genuinely new cohomology appears at deeper source levels. Neither a rank-independent constant nor a uniform covolume is supplied by finite generation. For the fixed sublattice of old maps pulled through a degree-\(e\) cover of the source, pull back \(H\) as well. Then \(E_i\mapsto eE_i\), \(C_i\mapsto e^{r_i/2}C_i\), and a known \(\mu_i\) becomes \(e\mu_i\). Consequently \(R_i\mapsto\sqrt e R_i\), \(V\mapsto eV\), and \(\delta_F\mapsto e\delta_F\); the reduced-image bound (B10.1) is unchanged. The actual generic degree in that bound must be retained. Replacing it by the trivial lower bound one can lose an essential cancellation. This proposition gives a sufficient determinant criterion for producing a small witness from prescribed source lattices. It is not a necessary condition for intrinsic bounded witness degree.

\subsection{B11. A fixed-order CM Hecke example separates the Hom regulator from witness degree}\label{b11.-a-fixed-order-cm-hecke-example-separates-the-hom-regulator-from-witness-degree}

Let \[
 E=\mathbb C/(\mathbb Z+i\mathbb Z),\qquad
 E_m=\mathbb C/(\mathbb Z+mi\mathbb Z),\quad m\ge1,
\] with their degree-one polarizations. Then \[
 \operatorname{Hom}(E,E_m)=m\mathbb Z[i],\qquad
 \deg(m(a+bi))=m(a^2+b^2).
\tag{B11.1}
\] In particular the least positive degree grows with \(m\), despite fixed rank and fixed CM field.

\textbf{Proof.} A homomorphism is multiplication by a complex number \(\gamma\) with \(\gamma\mathbb Z[i]\subset\mathbb Z+mi\mathbb Z\). Write \(\gamma=x+mi y\) with \(x,y\in\mathbb Z\). The further requirement \(i\gamma\in\mathbb Z+mi\mathbb Z\) forces \(m\mid x\). Thus \(\gamma=m(a+bi)\), and the converse is immediate. The degree is the index of \(\gamma\mathbb Z[i]\) in the target lattice, namely \(|\gamma|^2/m=m(a^2+b^2)\). These elliptic curves are algebraic and have CM; multiplication by \(m\) gives a degree-\(m\) isogeny, so they lie in one polarized Hecke orbit. QED.

Here is a version with a fixed integral quadratic action and actual Weil witnesses. Put \(A_m=E_m^4\) with product principal polarization \(L_m\), and let \(\mathbb Z[i]\) act by sending \(i\) to \[
 J=\operatorname{diag}\left(
 \begin{pmatrix}0&-1\\1&0\end{pmatrix},
 \begin{pmatrix}0&-1\\1&0\end{pmatrix}\right).
\] This is an integral action for every \(m\), has Rosati involution \(J^\dagger=-J\), and has multiplicities \((2,2)\) on \(H^{1,0}\), so \(A_m\) is of Weil type for the same order. The diagonal isogeny \(E^4\to E_m^4\) given by multiplication by \(m\) commutes with this action and pulls the target polarization back to \(m\) times the source polarization. Thus the example stays in the corresponding polarized quadratic-action Hecke orbit.

Let \(Y=E^2\), \(H\) its product degree-one polarization, and consider \[
 f_m(x,y)=(mx,0,my,0),\qquad
 Z_m=E_m\times\{0\}\times E_m\times\{0\}\subset A_m.
\] The image has \(\deg_{L_m}Z_m=2\). Its Weil projection is nonzero: the extreme \(i\)-eigenline contains the nonzero \((2,2)\)-form \[
 (dz_1+i\,dz_2)\wedge(dz_3+i\,dz_4)
 \wedge(d\bar z_1+i\,d\bar z_2)
 \wedge(d\bar z_3+i\,d\bar z_4),
\] whose restriction to \(Z_m\) is \(dz_1\wedge dz_3\wedge d\bar z_1\wedge d\bar z_3\) and has nonzero integral. Pairing with this Weil form detects the nonzero projection, by the Poincaré self-adjoint projector in B6.

The full fixed-source Hom lattice is \[
 \operatorname{Hom}(E^2,A_m)=M_{4\times2}(m\mathbb Z[i]),
\] of rank \(16\). In the integral basis furnished by the real and imaginary matrix units, its energy metric is \(m I_{16}\): the polarization form on \(E_m\) is \(\frac{i}{2m}dz\wedge d\bar z\), so multiplication by \(m\gamma\) contributes \(m|\gamma|^2\) to the energy. Consequently its least nonzero energy is \(m\), its Gram determinant is \(m^{16}\), and its metric covolume is \(m^8\). Orthonormalizing the analytic coordinates transfers the factor \(\sqrt m\) to the integral coefficients and changes none of these intrinsic metric quantities.

For the selected two-dimensional map, however, \[
 f_m^*c_1(L_m)=mH,\quad
 E_{f_m}=2m,\quad V=H^2=2,\quad \delta_{f_m}=m^2,
\] and therefore \[
 \frac{E_{f_m}^2}{\delta_{f_m}V}=2=\deg_{L_m}Z_m.
\] All assertions follow by the lattice-index calculation, product intersection, and the displayed invariant form. This supplies a precise counterexample to any claim that finite generation or an orthonormal analytic basis uniformly bounds integral map energy on a CM Hecke orbit. At the same time it supplies bounded nonzero Weil witnesses. It therefore does \textbf{not} refute intrinsic bounded division; it shows that a raw Hom-regulator bound is sufficient but unnecessary, and that generic-degree cancellation is mathematically essential.

\subsection{B12. The exact geometric quantity is an infimum of normalized energies}\label{b12.-the-exact-geometric-quantity-is-an-infimum-of-normalized-energies}

For a polarized Weil-type abelian \(2n\)-fold \((A,L)\), let \(D_W(A,L)\) be the least degree of an effective integral \(n\)-cycle with nonzero Weil projection, or \(+\infty\) if none exists. Then \[
 D_W(A,L)=\inf_{(Y,f,H)}
 \frac{\bigl(f^*c_1(L)\cdot H^{n-1}\bigr)^n}
 {\delta_f (H^n)^{n-1}},
\tag{B12.1}
\] where \(Y\) runs over smooth projective \(n\)-folds, \(f:Y\to A\) over morphisms generically finite onto an image \(Z\) with nonzero Weil projection, \(\delta_f\) is the generic mapping degree, and \(H\) runs over integral ample classes on \(Y\). The infimum of the empty set is understood to be \(+\infty\).

\textbf{Proof.} The mixed-intersection inequality B10.2 and projection formula show that every expression on the right is at least \(\deg_L Z\ge D_W(A,L)\). If a witness exists, at least one irreducible component has nonzero Weil projection and no larger degree. Positive degrees are integers, so a least-degree irreducible witness \(Z\) exists. Take a projective resolution \(f:Y\to Z\subset A\), with \(\delta_f=1\), and put \(\beta=f^*c_1(L)\). This is integral, big and nef, and \(\beta^n=\deg_L Z=D_W(A,L)\). Fix an integral ample \(H_0\). The classes \(H_j=j\beta+H_0\) are integral ample. As \(j\to\infty\), \[
 \beta H_j^{n-1}=j^{n-1}D_W+O(j^{n-2}),\qquad
 H_j^n=j^nD_W+O(j^{n-1}).
\] For \(n=1\) the expression is already equality; for \(n\ge2\) the quotient in B12.1 tends to \(D_W\). This proves the reverse inequality and hence equality. QED.

This is an exact variational characterization, not an effective construction. It explains what a successful automorphic estimate would have to control: normalized energies for some witnesses, allowing their actual mapping degrees and source volumes. Requiring every theta coefficient or the entire Hom regulator to remain bounded is stronger than necessary. B10 provides one workable sufficient criterion with explicit constants. B11 shows why it is not an equivalence on a prescribed source. The infimum in B12 allows resolutions and polarizations of the unknown witnesses, so it does not close the open existence or uniformity problem by itself.

\subsection{B13. A fixed general-type surface bounds mapping degree to a simple target}\label{b13.-a-fixed-general-type-surface-bounds-mapping-degree-to-a-simple-target}

\textbf{Statement.} Let \(Y\) be a fixed smooth projective surface of general type, and let \(f:Y\to A\) have surface image \(Z\) in a simple abelian variety of dimension strictly greater than two. Then \[
 \delta_f\operatorname{vol}(K_{\widetilde Z})
 \le \operatorname{vol}(K_Y),\qquad
 \delta_f\le \operatorname{vol}(K_Y),
\tag{B13.1}
\] where \(\widetilde Z\) is any smooth projective model of \(Z\). In particular, if \(Y\) is a compact ball quotient, then \(\delta_f\le K_Y^2\). No smoothness assumption on \(Z\) and no unramifiedness assumption on \(f\) is required.

\textbf{Proof.} A proper positive-dimensional subvariety of a simple abelian variety is of general type by Ueno's theorem. Its Kodaira dimension is defined using a nonsingular model, so this applies even when \(Z\) is singular. The induced dominant rational map \(Y\dashrightarrow\widetilde Z\) has degree \(\delta_f\). The canonical-volume inequality for generically finite rational maps of varieties of general type gives the first inequality in (B13.1). One way to see the mechanism is to resolve the map, use the effective ramification divisor to compare the canonical class with the pullback of the target canonical class, and then use birational invariance, monotonicity, and finite-map scaling of volume. A smooth minimal model of a surface of general type has nef and big canonical divisor, whose square is a positive integer; its canonical volume is that square. Hence \(\operatorname{vol}(K_{\widetilde Z})\ge1\), proving the second inequality. For a compact ball quotient, \(K_Y\) is ample and its volume is \(K_Y^2\). QED.

\textbf{Exact imported sources.} Kenji Ueno, \emph{On Kodaira Dimensions of Certain Algebraic Varieties}, Proc. Japan Acad.47 (1971), pp.157--159, §3, Corollary after Proposition 2, printed p.158; the birational definition appears in §2, p.157. The note announces the theorem and says details will appear elsewhere; the whole note and the exact corollary were read, without claiming to have checked an absent full proof. Primary original: \url{https://www.jstage.jst.go.jp/article/pjab1945/47/2/47_2_157/_pdf}. The word ``proper'' is retained explicitly here, since the ambient abelian variety itself is not of general type. Lucio Guerra and Gian Pietro Pirola, \emph{On the finiteness theorem for rational maps on a variety of general type}, arXiv:0807.2980v2, §3.2, Proposition 3.2(2), printed pp.10--11, gives the canonical-volume inequality; its complete short proof was read. Primary: \url{https://arxiv.org/pdf/0807.2980}. Their later paper \emph{Parametrization of rational maps on a variety of general type, and the finiteness theorem}, arXiv:1203.2246v1, §1(b), printed pp.2--3, explicitly records the surface lower bound \(\epsilon_2=1\) and the degree-volume inequality. Primary: \url{https://arxiv.org/pdf/1203.2246}.

This is a strict refinement for a fixed ball-quotient source and a simple target: the unbounded generic-degree cancellation in B11 cannot occur there. It is not an energy estimate or a bound for all image degrees. For an étale source cover of degree \(e\), the canonical volume itself scales by \(e\), consistently with B3 and B10. The actual 032 target is generally a product, so simplicity is an additional hypothesis and must not be silently imposed. More generally the first inequality holds whenever the image has a general-type smooth model, without assuming the ambient target is simple.

\subsection{B14. Exact local bounds, and a fixed-level unit-input counterexample}\label{b14.-exact-local-bounds-and-a-fixed-level-unit-input-counterexample}

The local integrals in B1 cannot be bounded merely by fixing the unitary representation conductor, the target compact level, and the \(L^2\) norm of the Schwartz input. The following calculation makes the obstruction explicit and also gives a sharp bound with a stated support hypothesis.

\textbf{Statement at a split finite place.} Let \(F_v\) be nonarchimedean, with ring of integers \(\mathcal O_v\), uniformizer \(\varpi\), residue cardinality \(q\), additive Haar measure \(\operatorname{vol}(\mathcal O_v^3)=1\) and multiplicative Haar measure \(\operatorname{vol}(\mathcal O_v^\times)=1\). For an unramified compatible splitting, the rank-one member \(U(1)(F_v)=F_v^\times\) acts in \(\mathcal S(F_v^3)\) by \[
 \omega(a)=\rho(a)D(a),\qquad
 D(a)\phi(x)=|a|^{3/2}\phi(ax),
\] where \(\rho\) is a fixed unramified unitary character. Replacing \(a\) by \(a^{-1}\) in the model changes only the harmless phase convention below. Fix an unramified unitary source character \(\tau\). At the \(U(1)\to U(3)\) Rallis point, the diagonal local zeta integral, with a unit source vector, is \[
 Z_v(1;\phi)=\int_{F_v^\times}
       \langle\omega(a)\phi,\phi\rangle\tau(a)\,d^\times a.
\tag{B14.1}
\] For every integer \(M\ge1\) there is a \(GL_3(\mathcal O_v)\)-invariant Schwartz function \(\phi_M\) of \(L^2\) norm one such that \[
 Z_v(1;\phi_M)=M.
\tag{B14.2}
\] Thus \(Z_v^\#(1;\phi_M)\) is a fixed nonzero local Euler-factor ratio times \(M\).

\textbf{Proof.} The normalized linear model is Haan §2.3, equation (2.1), printed p.5; restricting the target \(GL_3\) action to its scalar center gives the displayed dilation, up to its fixed unitary character. The compatible center identification is Remark 2.1, p.4. Equation (3.3), printed p.8, expresses the Siegel--Weil section as the oscillator matrix coefficient times \(\Delta^{s-s_m}\). Here the source dimension is one, the target dimension is three and \(s_m=1\), so the extra factor is one. With the local source matrix coefficient this is exactly (B14.1). If a convention conjugates that coefficient, replace \(\tau\) by its inverse in the phase choice.

Put \(S_j=\varpi^j\mathcal O_v^3\setminus\varpi^{j+1}\mathcal O_v^3\) and \(e_j=\operatorname{vol}(S_j)^{-1/2}\mathbf1_{S_j}\) for \(j\in\mathbb Z\). The \(e_j\) are orthonormal, and \(D(\varpi)e_j=e_{j-1}\). Choose \(c=(\rho(\varpi)\tau(\varpi))^{-1}\), of modulus one, and set \[
 \phi_M=M^{-1/2}\sum_{j=0}^{M-1}c^j e_j.
\] It has norm one. Its support is contained in \(\mathcal O_v^3\) and it vanishes on \(\varpi^M\mathcal O_v^3\), so it is a Schwartz function. Every shell is \(GL_3(\mathcal O_v)\)-stable, and the splitting character is unramified; hence the full function is invariant under that same compact subgroup for every \(M\).

With the inner product linear in the first variable, \[
 \langle D(\varpi^k)\phi_M,\phi_M\rangle
 =c^k\left(1-\frac{|k|}{M}\right)
 \quad (|k|<M),
\] and the coefficient is zero for \(|k|\ge M\). The phases \(\rho(\varpi)^k\tau(\varpi)^kc^k\) cancel. Units contribute no phase and preserve every \(e_j\). Decomposing \(F_v^\times\) into \(\varpi^k\mathcal O_v^\times\), one obtains \[
 Z_v(1;\phi_M)=\sum_{|k|<M}\left(1-\frac{|k|}{M}\right)=M.
\] The sum is finite, so no continuation or interchange of divergent integrals is involved. The normalization in Haan Theorem 3.4 is multiplication by the fixed nonzero number \(L_{F_v}(3,\epsilon)/L_{E_v}(3/2,\operatorname{BC}(\tau)\gamma^3)\). At a split unramified place, writing \(\lambda\) for the first of the two inverse unitary Satake parameters, this number is \[
 \frac{|1-\lambda q^{-3/2}|^2}{1-q^{-3}}>0.
\] Thus the normalized local factors grow linearly too. QED.

\textbf{Sharp positive bound with radial support specified.} For \(\phi=\sum_{j\in J}a_j e_j\) supported on a finite set \(J\) of shells, absorb the same phases into \(a_j\). Summing its matrix coefficients gives \[
 Z_v(1;\phi)=\left|\sum_{j\in J}a_jc^{-j}\right|^2
 \le |J|\sum_{j\in J}|a_j|^2
 =|J|\|\phi\|_2^2.
\tag{B14.3}
\] This is Cauchy--Schwarz, and the \(M\)-shell example attains equality. Thus a support-size hypothesis gives an exact local bound, while the common group-theoretic level does not supply that support bound. The additive local constancy scale of \(\phi_M\) does become finer with \(M\); it has never been asserted fixed.

There is also a uniform bound using an explicit radial seminorm. For every \(\varepsilon>0\) and every integer center \(j_0\), weighted Cauchy--Schwarz gives \[
 Z_v(1;\phi)\le
 \bigl(2\zeta(1+\varepsilon)-1\bigr)
 \sum_{j\in\mathbb Z}(1+|j-j_0|)^{1+\varepsilon}|a_j|^2.
\tag{B14.5}
\] Indeed the sum of the inverse weights is \(1+2\sum_{k\ge1}(1+k)^{-1-\varepsilon}=2\zeta(1+\varepsilon)-1\). One may take the infimum over \(j_0\). Every radial Schwartz function has coefficients of this kind: only finitely many negative \(j\) occur, and local constancy near zero makes the positive tail decay as \(q^{-3j/2}\). Hence its coefficients are absolutely summable, the matrix-coefficient integral is absolutely convergent, and the exact radial formula and (B14.5) continue to hold for that tail. The estimate therefore converts the local-factor question into concrete support or weighted-seminorm control in this split radial model.

A pure dilation replaces the coefficient sequence by an index shift and a unit phase. It preserves both \(Z_v(1;\phi)\) and \(\|\phi\|_2\), and preserves the infimum of the seminorms over \(j_0\). Growth in (B14.2) comes from a coherent sum of many orthogonal shell vectors, rather than from translating one shell.

\textbf{Compact places.} If the local source torus \(T_v\) is compact, let its Haar volume be \(v_T\). The same diagonal section identity gives instead \[
 Z_v(1;\phi)=v_T\|P_{\tau^{-1}}\phi\|_2^2
 \le v_T\|\phi\|_2^2,
\tag{B14.4}
\] where \(P_{\tau^{-1}}=v_T^{-1}\int_{T_v}\tau(t)\omega(t)\,dt\) is the orthogonal character projector. This follows from compact-group averaging and unitarity. Multiplication by the explicitly retained local \(L\)-ratio gives the corresponding bound for \(Z_v^\#\). There is no positive lower bound for arbitrary inputs, since that projector may kill an input.

These are statements about the chosen theta inputs and local norm formulas. They do not imply that the intrinsic least degree of a geometric witness, or the intrinsic bounded-division complexity, grows. In particular the counterexample does not invalidate B8 or DIV1. It identifies precisely why an \(L\)-value estimate, even combined with unit input norms and fixed target level, is insufficient without additional control of the actual finite Schwartz data or an integral energy-lattice comparison.

\subsubsection{B14a. The split-place estimate for arbitrary angular inputs}\label{b14a.-the-split-place-estimate-for-arbitrary-angular-inputs}

The positive support and weighted bounds extend to every Schwartz input and to ramified unitary characters in the same split Schrödinger model. Let \(\mathcal H_0=L^2(S_0)\) and let \[
 (I_j b)(x)=q^{3j/2}b(\varpi^{-j}x)\mathbf1_{S_j}(x).
\] Then \(I_j:\mathcal H_0\to L^2(S_j)\) is unitary. Write \(\phi=\sum_j I_jb_j\), so \(\|\phi\|_2^2=\sum_j\|b_j\|_2^2\). Put \(\eta=\rho\tau\), \(u=\eta(\varpi)\), and let \(R(v)b(x)=b(vx)\) for \(v\in\mathcal O_v^\times\). The compact-unit operator \[
 P_\eta=\int_{\mathcal O_v^\times}\eta(v)R(v)\,d^\times v
\] is the orthogonal projector to the \(\eta^{-1}\) character. Then exactly \[
 Z_v(1;\phi)=
 \left\|\sum_{j\in\mathbb Z}u^jP_\eta b_j\right\|_{\mathcal H_0}^2.
\tag{B14a.1}
\] Consequently, for every \(\varepsilon>0\) and \(j_0\in\mathbb Z\), \[
 \begin{aligned}
 Z_v(1;\phi)
 &\le \bigl(2\zeta(1+\varepsilon)-1\bigr)
 \sum_j(1+|j-j_0|)^{1+\varepsilon}\|P_\eta b_j\|_2^2\\
 &\le \bigl(2\zeta(1+\varepsilon)-1\bigr)
 \int_{F_v^3}(1+|\nu(x)-j_0|)^{1+\varepsilon}|\phi(x)|^2\,dx,
 \end{aligned}
\tag{B14a.2}
\] where \(\nu(x)=\min_i v(x_i)\) for \(x\ne0\); the value assigned at zero is immaterial. If \(\phi\) is supported on at most \(M\) shells, then \(Z_v(1;\phi)\le M\|\phi\|_2^2\).

\textbf{Proof.} Normalized dilation satisfies \(D(\varpi^k)I_jb=I_{j-k}b\), while a unit acts by \(I_jR(v)b\). Integrating the unit part of (B14.1) replaces each inner product by \(\langle P_\eta b_j,b_{j-k}\rangle=\langle P_\eta b_j,P_\eta b_{j-k}\rangle\). Summing the valuation part and setting \(l=j-k\) gives \[
 \sum_{j,l}u^{j-l}\langle P_\eta b_j,P_\eta b_l\rangle
 =\left\|\sum_j u^jP_\eta b_j\right\|^2.
\] All interchanges are legitimate: compact support makes \(b_j=0\) for \(j\ll0\), and local constancy at zero gives \(b_j=q^{-3j/2}\phi(0)\mathbf1_{S_0}\) for \(j\gg0\). Hence \(\sum_j\|b_j\|<\infty\), and the double sum of absolute matrix-coefficient bounds is finite. Hilbert-space Cauchy--Schwarz with the weights from B14.5 proves the first estimate; the projection is a contraction and \(I_j\) is unitary, proving the integral expression and the second estimate. Ordinary finite-sum Cauchy--Schwarz proves the support bound. QED.

The unramified radial hypothesis was needed only to obtain the particularly simple fixed-\(GL_3(\mathcal O_v)\) counterexample in B14. It is not needed for this local bound. For a ramified character, the sufficiently deep constant tail may be annihilated by \(P_\eta\), which is already accounted for by the exact identity. Multiplying by the retained local Euler ratio gives the corresponding explicit estimate for the normalized zeta integral.

\subsection{B15. A linear bound in polarized isogeny degree on a Weil Hecke orbit}\label{b15.-a-linear-bound-in-polarized-isogeny-degree-on-a-weil-hecke-orbit}

\textbf{Setup.} Keep the connected polarized Weil-type family and fixed nonzero rational flat class \(t\) of B8. Fix one fibre \(A_0\) and effective integral cycles \(Z_0^+,Z_0^-\) such that \[
 N_0t_0=[Z_0^+]-[Z_0^-],\qquad
 \deg_{L_0}Z_0^\pm\le D_0,
\] with \(N_0\ge1\). Suppose a fibre \(A_x\) admits an \(\mathcal O\)-linear polarized isogeny \(\psi:A_x\to A_0\) with \[
 \psi^*c_1(L_0)=c\,c_1(L_x),\qquad c\in\mathbb Q_{>0},
\] and the same polarization type. The isogeny is not assumed to carry \(t_0\) to a rational multiple of \(t_x\).

\textbf{Statement.} Let \(b,\tau=bt,h,R,B_P,V_A,\delta_\Lambda\) be the fixed constants from B6--B8 and set \[
 \begin{aligned}
 C_N&=\frac{b h R^2}{(n!)^2 V_A\sqrt{\delta_\Lambda}},\\
 C_D&=\frac{2h\|\tau\|_q R B_P}{n!\sqrt{V_A}\sqrt{\delta_\Lambda}}.
 \end{aligned}
\] Then \(t_x\) has an integral signed-cycle representative \(N_xt_x=[Z_x^+]-[Z_x^-]\) satisfying \[
 \begin{aligned}
 1\le N_x&\le C_N D_0^2\deg\psi,\\
 \deg_{L_x} Z_x^\pm&\le C_D D_0^2\deg\psi.
 \end{aligned}
\tag{B15.1}
\] In particular its intrinsic bounded-division complexity is at most \(D_0^2\max(C_N,C_D)\deg\psi\). One may minimize this upper bound over all isogenies with the stated properties.

\textbf{Proof.} Polarization integration gives \(\deg\psi=c^{2n}\), and the codimension-\(n\) pullback calculation B4 gives \[
 \deg_{L_x}\psi^*Z_0^\pm
 =c^n\deg_{L_0}Z_0^\pm
 \le c^nD_0.
\] Flat pullback under an isogeny preserves effectiveness and integral coefficients. Because \(\psi\) is \(\mathcal O\)-linear, its pullback on \(H^1\) commutes with the two \(K\) eigenspaces, and hence induces an isomorphism \(W_0\to W_x\) on their extreme exterior powers. Equivalently it commutes with the fixed polynomial projector: \(e_x\psi^*=\psi^*e_0\). Therefore \[
 e_x[\psi^*Z_0^+]-e_x[\psi^*Z_0^-]
 =N_0\psi^*t_0\ne0.
\] At least one of these effective pullbacks is a nonzero Weil witness of degree \(D\le c^nD_0\). Apply B8 to that witness and to the prescribed target \(t_x\), using its two-dimensional extraction rather than assuming alignment of \(\psi^*t_0\) with \(t_x\). Substituting \(D^2\le c^{2n}D_0^2=D_0^2\deg\psi\) into the exact B8 constants gives (B15.1).

Uniformity follows from the family hypotheses, rather than from the isogeny: the order action makes \(W\) and its saturated integral lattice a local system, \(q\) is flat, \(\delta_\Lambda\) and \(\|bt\|_q\) are constant on the connected base, and \(V_A\) is fixed by polarization type. The fixed \(a\), \(P\), \(R\), \(h\) and \(B_P\) also do not vary. This proves the stated uniform constant. QED.

\textbf{Scalar special case.} If \(\psi^*t_0=r t_x\) for a rational scalar \(r\), then flatness gives \(\|t_0\|_q=\|t_x\|_q\), while Poincaré pairing under pullback gives \(\|\psi^*t_0\|_q^2=\deg\psi\,\|t_0\|_q^2\). Consequently \(r=\pm c^n\). Since \(c\) is rational and \(c^{2n}=\deg\psi\) is an integer, \(c\) is a positive integer. Pulling back the original signed representative, with the two signs interchanged when necessary, yields directly \[
 N_x=N_0c^n,\qquad
 \deg Z_x^\pm\le c^nD_0,
\] so the complexity is at most \(\max(N_0,D_0)\sqrt{\deg\psi}\). This is the scalar refinement of DIV4 in the fixed flat Weil-family normalization. The linear bound (B15.1) remains valid when the isogeny rotates the Weil plane and no such rational scalar exists.

This gives an unconditional quantitative upper bound once the fixed starting representative is supplied, for example at a fibre covered by an admitted Hodge-conjecture theorem. It does not assert a bound independent of isogeny degree and does not supply a lower bound on intrinsic complexity. Neither Dirichlet GRH nor Hecke GRH is used.

\subsection{B16. Exact Chow idempotence after the canonical degree projector}\label{b16.-exact-chow-idempotence-after-the-canonical-degree-projector}

The cohomological caution in B6 is essential for that elementary proof. For abelian varieties, an additional known motivic theorem does supply an exact Chow upgrade, with a fixed denominator. It is not obtained by inferring a Chow identity from Betti eigenvalues.

\textbf{Statement.} Work with rational Chow motives of an abelian variety \(A\) as in B6, or with relative Chow motives of an abelian scheme over a smooth connected quasi-projective complex algebraic base where the \(\mathcal O\) action is defined. Let \(\Pi_{2n}\) be the canonical Deninger--Murre degree projector. Then \[
 \mathcal E_W=\Pi_{2n}\circ\frac{P([a]^*)}{R}
\tag{B16.1}
\] is an idempotent correspondence modulo rational equivalence, whose Betti realization is the projection onto \(W\) in degree \(2n\) and zero in all other degrees. With \(M=[2]^*\), one may take \[
 \Pi_{2n}=\prod_{\substack{0\le j\le4n\\j\ne2n}}
              \frac{M-2^j}{2^{2n}-2^j}.
\tag{B16.2}
\] Thus a denominator for the displayed correspondence divides \[
 R D_n,\qquad
 D_n=\prod_{\substack{0\le j\le4n\\j\ne2n}}|2^{2n}-2^j|,
\] a fixed integer independent of the fibre and its extra CM endomorphisms.

\textbf{Imported motivic input.} Ben Moonen, \emph{On the Chow motive of an abelian scheme with non-trivial endomorphisms}, arXiv:1110.4264v2, §2.1 and Proposition 2.5--Corollary 2.6, PDF pp.7--8, states the exact canonical degree decomposition, its functoriality, the algebra homomorphism on \(h^1\), and the fact that \[
 r^{(i)}:\operatorname{End}^0(A)^{\mathrm{op}}
       \longrightarrow\operatorname{End}(h^i(A))
\] is a multiplicative homogeneous polynomial map of degree \(i\). The complete short proof of Corollary 2.6 was read. It uses the canonical exterior-power description, citing Künnemann, \emph{On the Chow motive of an abelian scheme} (1994), Theorem (3.3.1); Proposition 2.5 cites Kings (1998), Proposition 2.2.1. These motivic inputs are imported, not newly proved here. Primary source: \url{https://arxiv.org/pdf/1110.4264}.

\textbf{Proof.} The exact degree decomposition says that \(M\) acts by \(2^j\) on \(h^j(A)\) as a Chow-motive endomorphism. Consequently its interpolation polynomial in (B16.2) is exactly \(\Pi_{2n}\), not just its realization. Functoriality makes \(\Pi_{2n}\) commute with every \([a]^*\).

It remains to prove the polynomial relation for \(T=[a]^*\) on \(h^{2n}(A)\). Restrict the imported polynomial map \(r^{(2n)}\) to \(K\subset\operatorname{End}^0(A)\) and extend the coefficient field from \(\mathbb Q\) to \(K\). Under \(K\otimes_{\mathbb Q}K\simeq K\times K\), this becomes a multiplicative homogeneous polynomial map \[
 A(x,y)=\sum_{r=0}^{2n}E_r x^r y^{2n-r}
\] with coefficients \(E_r\) in the endomorphism algebra of \(h^{2n}(A)\) with \(K\) coefficients. Comparing coefficients in \(A(x,y)A(u,v)=A(xu,yv)\) gives \[
 E_rE_s=0\ (r\ne s),\qquad E_r^2=E_r.
\] The identity \(A(1,1)=1\) gives \(\sum_r E_r=1\). Evaluating at \((a,\bar a)\) yields \[
 T=\sum_{r=0}^{2n}a^r\bar a^{2n-r}E_r,
 \qquad \prod_{r=0}^{2n}(T-\lambda_r)=0
\] as exact identities of Chow-motive endomorphisms. Since scalar extension of a rational vector space is faithful, the rational polynomial relation holds already over \(\mathbb Q\).

The integral Bezout polynomial from B6 takes value \(R\) on the two extreme eigenvalues and zero on all others. Therefore \(P(T)/R=E_0+E_{2n}\) on \(h^{2n}(A)\) and is exactly idempotent there. Composing with \(\Pi_{2n}\) makes it zero on the other degrees, proving (B16.1) on the full Chow motive. Both factors are rational polynomials in algebraic pullback correspondences, and clearing their displayed denominators gives the asserted \(R D_n\) bound. Its realization agrees with B6 by the same extreme-eigenspace calculation. QED.

This does not assert that the cycle-class map is injective on the resulting Chow summand, or that every rational Weil class has an algebraic representative. It provides a canonical-degree Chow projector once the established abelian-motive decomposition is used. B8 and B15 already worked with the smaller cohomological polynomial \(P\) and do not require this additional motivic input.

\subsection{B17. Controlled local representatives and the remaining arithmetic comparison}\label{b17.-controlled-local-representatives-and-the-remaining-arithmetic-comparison}

This proposition refines B14--B14a. Arbitrarily many coherent radial shells can enlarge a local norm, but this presentation feature can be removed for a \textbf{specified existing} theta vector. The operation preserves that vector; it does not make its prescribed output small. A separate construction gives bounded local test vectors without prescribing their arithmetic realization.

\subsubsection{B17.1. An exact one-shell section}\label{b17.1.-an-exact-one-shell-section}

Use B14a's split model, measures and notation. In particular, \(\eta=\rho\tau\), \(u=\eta(\varpi)\), \(\mathcal H_0=L^2(S_0)\), and \[
 A_\eta\phi=\sum_j u^jP_\eta b_j,\qquad
 Z_\eta(\phi,\psi)=\langle A_\eta\phi,A_\eta\psi\rangle.
\] The second identity follows from the same absolutely convergent double sum as B14a, now with two inputs. Let \(\mathcal H_\eta^\infty\) be the locally constant functions on \(S_0\) transforming under scalar units by \(\eta^{-1}\). Define \[
 s_\eta(h)=I_0h,\qquad C_\eta=s_\eta A_\eta.
\] Then \(C_\eta\phi\) is a Schwartz function supported on the single shell \(S_0\), and \[
 \begin{aligned}
 A_\eta s_\eta&=1,& C_\eta^2&=C_\eta,\\
 A_\eta(\phi-C_\eta\phi)&=0,&
 Z_\eta(\phi-C_\eta\phi,\phi-C_\eta\phi)&=0,\\
 \|C_\eta\phi\|_2^2
 &=Z_\eta(C_\eta\phi,C_\eta\phi)
 =Z_\eta(\phi,\phi).
 \end{aligned}
\tag{B17.1}
\] For every \(\varepsilon>0\) its squared weighted shell norm from B14a, centered at zero, equals this same quantity. Moreover \(C_\eta\) commutes with \(GL_3(\mathcal O)\) in the linear target model, including a fixed smooth unitary determinant twist. It therefore preserves every compact subgroup level and compact type within this group.

\textbf{Proof.} Compact support gives \(b_j=0\) for \(j\ll0\). For some integer \(J\), local constancy at the origin gives \[
 b_j=q^{-3j/2}\phi(0)\mathbf1_{S_0}\quad(j\ge J).
\] Consequently \[
 A_\eta\phi
 =\sum_{j<J}u^jP_\eta b_j
 +\phi(0)P_\eta\mathbf1_{S_0}
       \frac{(u q^{-3/2})^J}{1-uq^{-3/2}}.
\tag{B17.2}
\] The first sum is finite, and the second is a locally constant function. Thus \(A_\eta\phi\in\mathcal H_\eta^\infty\). Extension by zero from the compact open shell is Schwartz. Its only nonzero shell coefficient is \(h\) at index zero, so \(A_\eta s_\eta h=P_\eta h=h\). This proves the section and idempotence identities. The bilinear form identity proves the radical and norm identities. For \(k\in GL_3(\mathcal O)\) the action \(V(k)\phi(x)=\kappa(\det k)\phi(k^t x)\) preserves each shell and acts identically on every shell coordinate. It commutes with scalar units and hence with \(P_\eta\), \(A_\eta\) and \(s_\eta\). The weighted norm assertion follows because only index zero remains. QED.

\textbf{Preservation of the actual global theta vector.} Assume the actual square-integrable rank-one theta lift and fixed splittings are in Haan's Theorem 3.4, printed p.8. For a pure tensor \(\Phi=\phi_v\otimes\Phi^v\), change one split finite factor to \(C_{\eta_v}\phi_v\). The difference has zero local raw zeta integral by (B17.1), and hence zero normalized integral: the normalizing Euler factors at \(3/2\) and \(3\) are finite and nonzero. Haan's \textbf{global diagonal norm identity} therefore gives \[
 \|\theta(\bar f,(\phi_v-C_{\eta_v}\phi_v)\otimes\Phi^v)\|_{\mathrm{Pet}}^2=0.
\] Positive definiteness and smoothness imply that this global automorphic theta vector is identically zero. Thus the replacement preserves the existing vector, not merely its norm. Apply this argument term by term for finite sums of pure tensors and sequentially at any finite set of places. It also applies componentwise to the finite-dimensional differential-form coefficient spaces in T6.11. In particular, if the original form realizes an integral algebraic homomorphism, its actual Betti periods and integral-Hom normalization are unchanged. This proof does not infer the kernel of a local theta quotient from a local norm alone.

Only finitely many places are changed; the distinguished unramified vectors are retained elsewhere. An infinite replacement by one-shell vectors need not be an adelic Schwartz function. If \(h=A_\eta\phi\ne0\), \(I_0h\) preserves the prescribed output and has norm \(\|h\|\); \(I_0(h/\|h\|)\) has unit output norm but changes the output by the scalar \(\|h\|^{-1}\). This distinction is essential for an arithmetic normalization.

\subsubsection{B17.2. Least-norm questions at split and compact source places}\label{b17.2.-least-norm-questions-at-split-and-compact-source-places}

For fixed nonzero \(h\in\mathcal H_\eta^\infty\), there is no least-\(L^2\)-norm Schwartz representative of its split output. Indeed \[
 \phi_{M,h}=\frac1M\sum_{j=0}^{M-1}u^{-j}I_jh
\] satisfies \[
 A_\eta\phi_{M,h}=h,\qquad
 \|\phi_{M,h}\|_2^2=\frac{\|h\|^2}{M}.
\tag{B17.3}
\] These are Schwartz functions with the same compact target level as \(h\). The identities follow from \(A_\eta I_jh=u^jh\) and orthogonality of different shells. Their squared norms tend to zero; zero cannot have output \(h\). Thus a minimal-input proposal requires an additional support, weighted, or arithmetic-lattice constraint.

At a compact source torus \(T\) with Haar volume \(v_T\), by contrast, \[
 P_{\tau^{-1}}=v_T^{-1}\int_T\tau(t)\omega(t)\,dt,
 \qquad Z_T(\phi,\phi)=v_T\|P_{\tau^{-1}}\phi\|^2.
\tag{B17.4}
\] The average is an orthogonal projector on smooth vectors. It commutes with the target action, and replacing \(\phi\) by \(P_{\tau^{-1}}\phi\) preserves the global theta vector by the preceding global norm argument. For a fixed projected output \(h\), any representative is \(h+(1-P)\phi\), an orthogonal sum, so \(h\) is the unique least-norm representative. This proves the assertion including its normalization.

\subsubsection{B17.3. Uniform analytic test vectors and the standard adelic tail}\label{b17.3.-uniform-analytic-test-vectors-and-the-standard-adelic-tail}

For any smooth unitary \(\eta\) the explicit function \[
 \phi_\eta(x)=\eta(x_1)^{-1}
       \mathbf1_{\mathcal O^\times\times\mathcal O\times\mathcal O}(x)
\tag{B17.5}
\] is Schwartz, has cyclotomic-integral values, and obeys \[
 A_\eta\phi_\eta=\phi_\eta|_{S_0},\qquad
 \|\phi_\eta\|_2^2=Z_\eta(\phi_\eta,\phi_\eta)=1-q^{-1}.
\tag{B17.6}
\] Thus its raw local zeta factor is between \(1/2\) and \(1\) independently of conductor. Normalization by \((1-q^{-1})^{-1/2}\) gives unit squared norm, unit raw zeta factor and unit squared weighted shell norm.

\textbf{Proof and level accounting.} The unit restriction of a smooth character has finite image, hence roots of unity. Scalar units act on (B17.5) by \(\eta^{-1}\); its support is one shell and has additive volume \(1-q^{-1}\). This proves (B17.6). For \(m\ge1\), \(k\in1+\varpi^mM_3(\mathcal O)\) preserves the chart and changes its first coordinate from \(x_1\) to \(x_1(1+\varpi^m z)\). Therefore \(m\ge\max(1,a(\eta))\) suffices for invariance in the untwisted model; with determinant twist \(\kappa\) it suffices that \(m\ge\max(1,a(\eta),a(\kappa))\). Conversely, for ramified \(\eta\) in the untwisted model, a nonzero \(A_\eta\)-image of a principal-level-\(m\) vector forces \(\eta\) to be trivial on \(1+\varpi^m\mathcal O\), by the scalar matrices in that subgroup. At level zero the same argument uses all units. Thus the conductor enters the level condition even though it does not enlarge this norm. These are Schwartz test vectors; no representation-theoretic newform identification is needed. QED.

For an unramified character the standard lattice vector \(\phi_0=\mathbf1_{\mathcal O^3}\) has \[
 Z_\eta(\phi_0,\phi_0)=\frac{1-q^{-3}}{|1-uq^{-3/2}|^2},\qquad
 \|\phi_0\|_{\mathrm{wt},2}^2
 =\frac{1+q^{-3}}{(1-q^{-3})^2}\le\frac{72}{49},
\tag{B17.7}
\] where the squared weight on shell \(j\) is \((1+|j|)^2\). Indeed \(b_j=q^{-3j/2}\mathbf1_{S_0}\) for \(j\ge0\), and zero otherwise. Summing the geometric series proves the first formula; differentiating it twice, or multiplying absolutely convergent series, gives \(\sum_{j\ge0}(j+1)^2t^j=(1+t)/(1-t)^3\), proving the second. Haan's standard unramified normalization gives normalized local factor one.

There is also a useful uniform restricted-product bound. Over a fixed number field \(F\), use the unit-normalized chart vector at any finite set of split places and the standard vector at all other split places. The product of their squared weighted norms is at most \[
 \prod_{v<\infty}\frac{1+q_v^{-3}}{(1-q_v^{-3})^2}
 =\frac{\zeta_F(3)^3}{\zeta_F(6)}.
\tag{B17.8}
\] Every factor on the right is greater than one. The Euler product converges absolutely, since its logarithm is \(O(q_v^{-3})\), and the equality follows from \(1+t=(1-t^2)/(1-t)\). Thus neither the number nor conductor of the selected split places forces growth in this purely analytic choice of local weighted norms. The statement prescribes no global integral-Hom class or Weil period.

\subsubsection{B17.4. Arithmetic coefficients retained by the section}\label{b17.4.-arithmetic-coefficients-retained-by-the-section}

The section preserving an existing vector need not preserve an integral lattice of Schwartz representatives. If \(U\subset\mathcal O^\times\) kills the unit action on a given vector and the character, normalized averaging is \[
 P_\eta b=\frac1{[\mathcal O^\times:U]}
       \sum_{v\in\mathcal O^\times/U}\eta(v)R(v)b.
\tag{B17.9}
\] This denominator occurs, for example, on \(b=\mathbf1_{U\times\mathcal O\times\mathcal O}\), whose projection is \([\mathcal O^\times:U]^{-1}\phi_\eta\). Shell coordinates also contain \(q^{-3j/2}\). Formula (B17.2) displays the exact geometric-tail denominator \(1-uq^{-3/2}\); the tail vanishes if \(\eta\) is ramified on units. For unramified \(\eta\) and \(u=1\), \[
 C_\eta(\mathbf1_{\mathcal O^3})
       =\frac{\mathbf1_{S_0}}{1-q^{-3/2}}.
\tag{B17.10}
\] When \(q=p^2\), the coefficient is the reduced rational number \(p^3/(p^3-1)\), with bounded real size and denominator \(p^3-1\). This is a local coefficient calculation, not a lower bound on geometric cycle complexity: the existing global vector is unchanged. The analogous compact-torus projector has an averaging denominator \([T:U]\) on any finite quotient supporting the vector.

\textbf{Main-gate consequence.} B17 removes arbitrary radial support complexity from an already specified rank-one theta vector without changing its level or integral realization. It also supplies analytic test inputs of uniformly bounded weighted norm. Neither assertion bounds the norm of the actual theta output realizing a primitive integral algebraic homomorphism. The remaining comparison is between that integral Hom lattice and the positive theta-output metric. B18 below removes a further small-period concern in the fixed generic rank-two Weil plane only after an integral middle-dimensional witness is realized. It does not supply the preceding lattice comparison.

The complete local derivation and a separate finite-ring model Fourier-period example are in \texttt{local\_newvector\_followup.md}, LN1--LN7. The model example is not identified with the geometric Weil period. Haan's imported exact formula is Theorem 3.4, with §2.3 equation (2.1), §3.1 equation (3.1), and §3.2 equation (3.3) fixing the local model and normalization. Source: \url{https://arxiv.org/pdf/1403.5061}.

\subsection{B18. A nonzero integral Weil witness has no arbitrarily small normalized Weil period}\label{b18.-a-nonzero-integral-weil-witness-has-no-arbitrarily-small-normalized-weil-period}

\textbf{Setup.} Keep the fixed generic rank-two Weil space \(W\subset H^{2n}(A,\mathbb Q)\), integral resultant projector \(Re=P([a]^*)\), saturated lattice \(\Lambda=W\cap H^{2n}(A,\mathbb Z)\) and positive integral form \[
 q(v,w)=(-1)^n\int_Av\smile w
\] from B6--B8. Choose \(\Omega\in W_\mathbb C\) in one of the two extreme \(K\)-eigenlines, normalized by \[
 q(\Omega,\Omega)=0,\qquad q(\Omega,\overline\Omega)=1.
\tag{B18.1}
\] The first identity is forced by the eigencharacters; the second can always be imposed by a positive real scaling. Either eigenline and any remaining phase give the same absolute values below. Let \(Z\) be an effective integral middle-dimensional cycle with \(w=e[Z]\ne0\).

\textbf{Statement.} There is the exact formula and uniform lower bound \[
 \left|\int_Z\Omega\right|
 =\frac{\|e[Z]\|_q}{\sqrt2}
 \ge\frac1R.
\tag{B18.2}
\] More precisely, \[
 R^2\left|\int_Z\Omega\right|^2
      =\frac12q(P[Z],P[Z])\in\mathbb Z_{>0}.
\tag{B18.3}
\] If \(D=\deg_L Z\) and \(V_A=\chi(L)\), B7 also gives \[
 \frac1R
 \le\left|\int_Z\Omega\right|
 \le\frac{D}{\sqrt2\,n!\sqrt{V_A}}.
\tag{B18.4}
\] For a proper generically finite map \(f:Y^n\to Z\subset A\) of generic degree \(\delta\), the corresponding source period satisfies \[
 \left|\int_Yf^*\Omega\right|
 =\delta\left|\int_Z\Omega\right|
 \ge\frac{\delta}{R}.
\tag{B18.5}
\] The constants are uniform in the connected family of B8 because \(R\), the generic Weil lattice and its integral pairing are fixed local-system data; no bound on the extra CM endomorphism field is required.

\textbf{Proof.} The extreme eigenlines are conjugate and paired nondegenerately by \(q\), while each is isotropic. They form all of \(W_\mathbb C\). Thus every real \(w\in W_\mathbb R\) has a unique expansion \(w=c\Omega+\bar c\overline\Omega\). Equation (B18.1) gives \[
 q(w,w)=2|c|^2,\qquad q(w,\Omega)=\bar c.
\] The projector \(e\) is Poincare-selfadjoint by B6 and fixes \(\Omega\). Consequently \[
 \int_Z\Omega
 =\int_A[Z]\smile\Omega
 =\int_Ae[Z]\smile\Omega
 =(-1)^nq(w,\Omega).
\] Taking absolute values proves the exact equality in (B18.2). The class \(v=Rw=P[Z]\) is a nonzero element of \(\Lambda\). Its square under \(q\) is a positive even integer. Positivity is Hodge--Riemann on \(W\). For evenness, the underlying real torus of \(A\) has integral cohomology ring \(\Lambda^*H^1(A,\mathbb Z)\), by the integral product computation for circles. Expand the positive even-degree class \(v\) in its ordered exterior-monomial basis. Each monomial has square zero, and two distinct monomials commute in this even degree, so every cross term in \(v^2\) occurs twice. Thus \(v^2\) is divisible by two in integral cohomology and \(q(v,v)\in2\mathbb Z_{>0}\). This proves (B18.3) and the lower bound. The B7 comass estimate gives the upper bound. Finally \(f_*[Y]=\delta[Z]\) as integral fundamental cycles, so the projection formula proves (B18.5). QED.

\textbf{Scope of the normalization.} This concerns the fixed generic two-dimensional Weil plane and a cycle already realized with integral coefficients. It does not concern a single full-CM embedding coefficient in the 032 construction, an arbitrary coordinate functional on a higher-dimensional CM space, or a complex-normalized theta form before arithmetic realization. Once such an integral witness exists, the scalar forcing this nonzero period to have absolute value one has absolute value at most \(R\) (at most \(R/\delta\) for (B18.5)). This eliminates a separate arbitrarily-small-period scaling factor at that stage. Complex period normalization need not preserve rational or integral cycle coefficients, and does not bound the original degree or integral-Hom energy. The positive bound in B8 still requires a degree bound for the effective witness.

\subsection{B19. An exact CM Hom covolume formula and a conductor bound for witnesses}\label{b19.-an-exact-cm-hom-covolume-formula-and-a-conductor-bound-for-witnesses}

This result controls an actual integral Hom lattice in a fixed simple CM isogeny class. It gives an arithmetic alternative to bounding the coefficients of a theta expansion. The relevant integer is the index of the largest lattice stable under the maximal CM order; it can vary along the Hecke orbit.

\subsubsection{B19.1. The exact metric covolume}\label{b19.1.-the-exact-metric-covolume}

Let \(E\) be a CM field of degree \(2g\), \(\Phi\) a fixed primitive CM type, \(\mathcal O=\mathcal O_E\), and \(D_E\) the absolute discriminant. Fix \[
 A_0=\mathbb C^g/\Phi(I_0),
\] where \(I_0\) is an invertible fractional \(\mathcal O\)-ideal, with ample polarization \(L_0\). Let \[
 A=\mathbb C^g/\Phi(\Lambda)
\] be an abelian variety in the same \(E\)-linear isogeny class, where \(\Lambda\subset E\) is a full integral lattice. Fix an ample polarization \(L\) on \(A\). Assume, as holds for polarizations of simple CM abelian varieties with full endomorphism field \(E\), that both Rosati involutions restrict to complex conjugation on \(E\). Put \[
 V_0=\chi(L_0),\qquad V=\chi(L),\qquad
 \Lambda^-=\{x\in\Lambda:\mathcal O x\subset\Lambda\},\qquad
 J=[\Lambda:\Lambda^-].
\tag{B19.1}
\] The lattice \(\Lambda^-\) is the largest \(\mathcal O\)-stable sublattice of \(\Lambda\). It has finite index because \(\Lambda\) is commensurable with \(\mathcal O\).

Let \(\omega_0,\omega\) be the invariant Kähler representatives of the two polarizations. For an \(E\)-linear homomorphism \(h:A_0\to A\), define the squared Hilbert--Schmidt norm \[
 Q_A(h)=\operatorname{tr}_{\omega_0}(h^*\omega).
\tag{B19.2}
\] The real rank of \(\operatorname{Hom}_E(A_0,A)\) is \(2g\). Its covolume in the positive metric \(Q_A\) is exactly \[
 \operatorname{covol}_{Q_A}\operatorname{Hom}_E(A_0,A)
 =2^{-g}\sqrt{|D_E|}\,\frac{V}{V_0}\,J.
\tag{B19.3}
\]

\textbf{Proof.} In the chosen CM coordinates, every \(E\)-linear rational homomorphism is multiplication by an element \(\alpha\in E\). It is an actual homomorphism precisely when \(\alpha I_0\subset\Lambda\). Thus its integral Hom lattice is the fractional \(\mathcal O\)-ideal \[
 I=(\Lambda:I_0)=\{\alpha\in E:\alpha I_0\subset\Lambda\}.
\] Invertibility of \(I_0\) gives \(I I_0=\Lambda^-\). Indeed \(I I_0\) is \(\mathcal O\)-stable and contained in \(\Lambda\), hence in \(\Lambda^-\). Conversely \(\Lambda^- I_0^{-1}\subset I\), and multiplying by \(I_0\) gives the reverse inclusion.

Use ordinary Euclidean measure \(\prod_{\sigma\in\Phi}dx_\sigma\,dy_\sigma\) on \(\mathbb C^g\). Write \(c(M)\) for the Euclidean covolume of a lattice \(M\subset E\). For fractional \(\mathcal O\)-ideals, multiplicativity of ideal norm gives \[
 c(I I_0)=\frac{c(I)c(I_0)}{c(\mathcal O)}.
\] This is also immediate by writing fractional ideals as products of prime ideals and using their lattice indices. Since \(c(\Lambda^-)=Jc(\Lambda)\), it follows that \[
 c(I)=J\,c(\Lambda)\frac{c(\mathcal O)}{c(I_0)}.
\tag{B19.4}
\] The standard ideal-covolume formula is Milne, \emph{Algebraic Number Theory}, version 3.08, Proposition 4.26, printed pp.79--80; the following determinant calculation verifies its convention here. The discriminant determinant in all \(2g\) embeddings has absolute value \(\sqrt{|D_E|}\). Passing each conjugate pair of embedding rows to its real and imaginary parts changes the absolute determinant by \(1/2\). Therefore \[
 c(\mathcal O)=2^{-g}\sqrt{|D_E|}.
\tag{B19.5}
\]

Rosati conjugation makes distinct CM eigenlines orthogonal. Hence \[
 \omega_0=\frac i2\sum_{\sigma\in\Phi}t_{0,\sigma}\,dz_\sigma\wedge d\bar z_\sigma,
 \qquad
 \omega=\frac i2\sum_{\sigma\in\Phi}t_\sigma\,dz_\sigma\wedge d\bar z_\sigma
\] for positive real coefficients. Explicitly, the adjoint relation for every \(a\in E\) forces an off-diagonal Hermitian coefficient between distinct embeddings to vanish. Integration over a fundamental domain gives \[
 V_0=c(I_0)\prod_\sigma t_{0,\sigma},\qquad
 V=c(\Lambda)\prod_\sigma t_\sigma.
\tag{B19.6}
\] Multiplication by \(\alpha\) has \[
 Q_A(\alpha)=\sum_\sigma\frac{t_\sigma}{t_{0,\sigma}}|\sigma(\alpha)|^2.
\] Changing from Euclidean metric to this real metric multiplies covolume by \(\prod_\sigma(t_\sigma/t_{0,\sigma})\): each complex coordinate contributes two real coordinates. Combining (B19.4)--(B19.6) proves (B19.3). In particular the period-lattice covolume cancels exactly. QED.

The analytic CM model used here follows from the rank-one \(E\)-module \(H_1(A,\mathbb Q)\) and the chosen CM type; an \(E\)-linear isogeny identifies the rational modules of all targets. Changing the identification by an element of \(E^\times\) scales \(\Lambda\) and \(\Lambda^-\) together, so \(J\) is unchanged. No comparison with an unspecified theta lattice has entered the calculation.

\subsubsection{B19.2. A small actual isogeny}\label{b19.2.-a-small-actual-isogeny}

Write \(\mu_A=\min_{0\ne h\in\operatorname{Hom}_E(A_0,A)}Q_A(h)\). The least nonzero squared Hom norm has the sharp order \[
 \begin{aligned}
 g\left(\frac{V}{V_0}J\right)^{1/g}
 &\le\mu_A\le K_E\left(\frac{V}{V_0}J\right)^{1/g},\\
 K_E&=\frac2\pi\bigl(g!\sqrt{|D_E|}\bigr)^{1/g}.
 \end{aligned}
\tag{B19.7}
\] A homomorphism \(h\) attaining the minimum is an actual isogeny and simultaneously satisfies \[
 J\le\deg h\le \frac{K_E^g}{g^g}\,J.
\tag{B19.8}
\] The lower bound \(\deg h\ge J\) holds for every nonzero integral \(E\)-linear homomorphism.

\textbf{Proof.} Any nonzero \(\alpha\in I\) acts invertibly on rational homology, so the resulting homomorphism is an isogeny, with \[
 \deg h=[\Lambda:\alpha I_0]\ge[\Lambda:\Lambda^-]=J,
\] because the \(\mathcal O\)-stable lattice \(\alpha I_0\) is contained in \(\Lambda^-\). Let \(\lambda_\sigma=(t_\sigma/t_{0,\sigma})|\sigma(\alpha)|^2\) be the positive eigenvalues of \(h^*\omega\) relative to \(\omega_0\). Then \[
 \sum_\sigma\lambda_\sigma=Q_A(h),\qquad
 \prod_\sigma\lambda_\sigma=\deg h\,\frac{V}{V_0}.
\] Arithmetic--geometric mean proves the lower bound in (B19.7).

Apply Minkowski's convex body theorem in real dimension \(r=2g\) to the closed ball in the metric \(Q_A\). The precise imported statement is J.S. Milne, \emph{Algebraic Number Theory}, version 3.08, Theorem 4.19, printed p.76: a compact symmetric convex set of volume at least \(2^r\) times the lattice covolume contains a nonzero lattice point. Since \(v_{2g}=\pi^g/g!\), (B19.3) gives \[
 \min_{h\ne0}Q_A(h)
 \le\left(\frac{2^{2g}\operatorname{covol}_{Q_A} I}{v_{2g}}\right)^{1/g}
 =K_E\left(\frac{V}{V_0}J\right)^{1/g}.
\] The minimum exists by positive definiteness and discreteness. For a minimizing \(h\), arithmetic--geometric mean in the other direction gives \[
 \deg h\le\frac{V_0}{V}\left(\frac{Q_A(h)}g\right)^g
 \le \frac{K_E^g}{g^g}J.
\] Together with the already proved lower bound, this proves all assertions. QED.

The cited statements and complete short proofs of Milne's Theorem 4.19 and Proposition 4.26 were read; these classical results are imported. Source: \url{https://www.jmilne.org/math/CourseNotes/ANT.pdf}.

\subsubsection{B19.3. Effective Weil witnesses and bounded division}\label{b19.3.-effective-weil-witnesses-and-bounded-division}

Now assume \(g=2n\), that the fixed generic imaginary quadratic field \(K\) of B6--B8 lies in \(E\), and that the CM type restricts to Weil signature \((n,n)\) on \(K\). Assume the order action and polarizations have the compatibility required there. Fix an effective integral cycle \(Z_0\subset A_0\) of dimension \(n\) with nonzero generic Weil projection and degree at most \(D_0\). Suppose the targets \(A\) belong to the fixed polarized Weil family of B8, with rational flat target class \(t\).

Then there is an effective integral nonzero Weil witness \(Z\) on each such target satisfying \[
 \deg_L Z
 \le D_0 K_E^n\left(\frac{V}{V_0}J\right)^{1/2}.
\tag{B19.9}
\] Using the explicit B8 constants \(C_N,C_D\) defined in B15, there is a signed integral representative \(N_A t_A=[Z_A^+]-[Z_A^-]\) with \[
 \begin{aligned}
 1\le N_A&\le C_N D_0^2 K_E^{2n}\frac{V}{V_0}J,\\
 \deg_L Z_A^\pm&\le C_D D_0^2 K_E^{2n}\frac{V}{V_0}J.
 \end{aligned}
\tag{B19.10}
\] These constants are independent of the endomorphism order of \(A\).

\textbf{Proof.} Choose the isogeny \(h\) from (B19.7) and take the effective integral pushforward cycle \(Z=h_*Z_0\), with its actual multiplicities. Positivity of Hermitian forms gives \[
 h^*\omega\le Q_A(h)\omega_0.
\] On each complex \(n\)-plane, the determinant is monotone for positive Hermitian forms; hence integration on the smooth loci of the components of \(Z_0\) gives \[
 \deg_L(h_*Z_0)
 =\int_{Z_0}(h^*\omega)^n
 \le Q_A(h)^n\deg_{L_0}Z_0.
\tag{B19.11}
\] Here no smoothness of the cycle or bound on the restriction's generic degree is needed. Since \(n/g=1/2\), (B19.7) proves (B19.9).

The isogeny is \(K\)-linear. Thus \(h^*\) commutes with the fixed \(K\) action, identifies the two generic Weil planes, and commutes with their polynomial projectors. Its pushforward is \(\deg h\,(h^*)^{-1}\) in cohomology, so it also commutes with those projectors and is an isomorphism on the Weil plane. Therefore \(e_A[h_*Z_0]=h_*e_0[Z_0]\ne0\). The effective cycle is a valid B8 witness. Substituting the square of (B19.9) into B8 gives (B19.10). The constants of B8 are flat-family constants as explained in B15; \(E\), \(V_0\) and \(D_0\) are fixed starting data. QED.

A fixed signed starting representative suffices: at least one of its two effective parts has nonzero Weil projection, and that part may be chosen for \(Z_0\). If the starting CM variety is first replaced by a fixed isogenous one with maximal-order action, transfer a fixed nonzero witness through that one fixed isogeny. Its resulting degree is then part of \(D_0\).

\subsubsection{B19.4. Conductor and polarized-isogeny comparisons}\label{b19.4.-conductor-and-polarized-isogeny-comparisons}

Let \[
 \mathcal R_A=\{a\in E:a\Lambda\subset\Lambda\},\qquad
 \mathfrak f_A=\{b\in\mathcal O:b\mathcal O\subset\mathcal R_A\}
\] be the CM endomorphism order and its conductor in \(\mathcal O\). Then \[
 J\le N_{E/\mathbb Q}(\mathfrak f_A).
\tag{B19.12}
\] Thus (B19.10) gives an unconditional upper bound linear in this conductor norm for the fixed CM isogeny class.

\textbf{Proof.} Put \(\Lambda^+=\mathcal O\Lambda\), a fractional \(\mathcal O\)-ideal. For \(b\in\mathfrak f_A\), \(a\in\mathcal O\) and \(x\in\Lambda\), the element \(ba\) belongs to \(\mathcal R_A\), so \(bax\in\Lambda\). Hence \(\mathfrak f_A\Lambda^+\subset\Lambda\), and this sublattice is \(\mathcal O\)-stable. It is therefore contained in \(\Lambda^-\). Consequently \[
 J\le[\Lambda:\mathfrak f_A\Lambda^+]
 \le[\Lambda^+:\mathfrak f_A\Lambda^+]
 =N_{E/\mathbb Q}(\mathfrak f_A),
\] as asserted. QED.

There is a useful comparison with B15 when both polarizations are principal. Suppose \(\psi:A\to A_0\) is an \(E\)-linear polarized isogeny with \(\psi^*\omega_0=c\omega\) for a positive rational \(c\). Then \[
 J\le\deg\psi.
\tag{B19.13}
\] Indeed identify \(\psi\) with multiplication by \(\alpha\) and put \(M=\alpha\Lambda\subset I_0\). Let \(M^\#\) denote the dual lattice for the integral alternating form of \(L_0\). Principal polarization gives \(I_0^\#=I_0\) and \(M^\#=c^{-1}M\). Write \(M^+=\mathcal O M\) and \(M^-\) for its largest \(\mathcal O\)-stable sublattice. The Rosati relation gives \[
 (M^+)^\#=(M^\#)^-=c^{-1}M^-.
\] Dualizing again and comparing finite indices yields \[
 J=[M:M^-]=[(M^-)^\#:M^\#]
   =[M^+:M]\le[I_0:M]=\deg\psi.
\] This proves (B19.13), and also identifies the core and hull indices exactly in this principal polarized situation. No corresponding comparison is asserted here without those polarization hypotheses.

\subsubsection{B19.5. Period-normalized analytic energy on a fixed source}\label{b19.5.-period-normalized-analytic-energy-on-a-fixed-source}

Suppose in addition that \(Y\) is a fixed smooth connected projective variety of dimension \(n\), with a fixed Kähler form \(\omega_Y\), and \(f_0:Y\to A_0\) is a morphism whose image has dimension \(n\). Let \(Z_0=f_{0*}[Y]\), with its actual generic multiplicity, and assume \(e_0[Z_0]\ne0\). Put \(g=2n\) as in B19.3, and choose \(h\) attaining the minimum in (B19.7). Define \[
 E_h=\int_Y(hf_0)^*\omega\,\omega_Y^{n-1}.
\] For each \(\sigma\in\Phi\) let \[
 a_\sigma=\int_Yf_0^*\left(\frac i2t_{0,\sigma}\,dz_\sigma\wedge d\bar z_\sigma\right)
                  \omega_Y^{n-1},\qquad
 a_{\min}=\min_\sigma a_\sigma,\quad a_{\max}=\max_\sigma a_\sigma.
\] These are fixed positive constants, and \[
 a_{\min}Q_A(h)\le E_h\le a_{\max}Q_A(h).
\tag{B19.14}
\] In particular the smallest energy of a nonzero integral map of this fixed form is bounded above and below by fixed positive multiples of \(J^{1/g}\).

Choose q-unit generic Weil forms \(\Omega_0,\Omega_A\) in corresponding extreme \(K\)-eigenlines as in B18, and put \[
 p_0=\left|\int_Yf_0^*\Omega_0\right|>0,\qquad
 p_h=\left|\int_Y(hf_0)^*\Omega_A\right|.
\] Then \[
 \begin{aligned}
 p_h&=\sqrt{\deg h}\,p_0\ge\sqrt J\,p_0,\\
 \frac{E_h}{p_h^{1/n}}
 &\le a_{\max}K_E\left(\frac{V}{V_0}\right)^{1/g}p_0^{-1/n}.
 \end{aligned}
\tag{B19.15}
\] Thus the period-normalized analytic energy is uniformly bounded in this fixed CM isogeny class and fixed-source construction.

\textbf{Proof.} Since \(A_0\) is simple and \(f_0\) is nonconstant, the induced Albanese homomorphism surjects onto \(A_0\). Pullback of invariant holomorphic one-forms is therefore injective. Every \(f_0^*dz_\sigma\) is nonzero, and its positive Kähler energy is strictly positive. This proves \(a_{\min}>0\). Diagonalization already used in B19.1 gives the exact identity \[
 E_h=\sum_\sigma a_\sigma\lambda_\sigma,
\] which proves (B19.14). Combining it with (B19.7) proves the asserted two-sided order for the minimum over all nonzero integral \(h\).

On middle cohomology, integration under the isogeny gives \[
 q_0(h^*u,h^*v)=\deg h\,q_A(u,v).
\] The \(K\)-linear isogeny preserves the chosen extreme eigenline, so \(h^*\Omega_A=\lambda\Omega_0\) for some \(\lambda\in\mathbb C^\times\). The q-unit normalization forces \(|\lambda|^2=\deg h\), proving the period identity. The lower bound uses (B19.8). Finally (B19.14) and (B19.7), together with \(g=2n\), give \[
 \frac{E_h}{p_h^{1/n}}
 \le\frac{a_{\max}K_E((V/V_0)J)^{1/g}}
          {p_0^{1/n}J^{1/(2n)}}
 =a_{\max}K_E(V/V_0)^{1/g}p_0^{-1/n}.
\] This proves (B19.15). QED.

For precision about the normalization, choose a polarization-unit invariant holomorphic coframe on \(A\). The sum of the fixed-source squared energies of its pullbacks is \(E_h\). Multiplying every pulled-back one-form by the positive scalar \(p_h^{-1/(2n)}\) makes the absolute value of the selected \(2n\)-fold generic Weil period equal to one and makes the summed squared energy equal to \(E_h/p_h^{1/n}\). The generic Weil form is a homogeneous wedge of \(2n\) one-forms and their conjugates; positive scalar scaling therefore has exactly this exponent.

When \(n=2\) and \(Y\) is a fixed compact unitary surface with the full rank-one theta spanning property of T6.9, this also gives uniformly controlled \textbf{analytic theta presentations}. Any one auxiliary finite cover required for that spanning statement is fixed at the outset, together with the pulled-back \(f_0\) and its multiplicity. Fix once and for all a finite basis \(\theta_1,\ldots,\theta_r\) of genuine theta one-forms on this source, with all its Schwartz data, source characters and levels. Define the positive energy Gram matrix \(G\) by \(\|\sum_k c_k\theta_k\|^2=c^*Gc\), and let \(\lambda_{\min}(G)>0\) be its least eigenvalue. Expansion of the normalized coordinate pullbacks in this basis gives \[
 \sum_{j,k}|c_{j,k}|^2
 \le \lambda_{\min}(G)^{-1}\frac{E_h}{p_h^{1/n}}.
\tag{B19.16}
\] Indeed the squared energy of each coefficient vector is \(c_j^*Gc_j\), and summing gives the numerator. A fixed linear section from this finite coefficient space to the finite collection of Schwartz presentations is bounded for every fixed Schwartz seminorm, including the weighted shell seminorms of B14a at the finitely many nonstandard places. Standard unramified inputs are kept outside that fixed set. Thus bounded period-normalized presentations exist in this explicitly fixed-source setting.

These normalized forms generally no longer arise from integral algebraic homomorphisms: the common real scaling changes their Betti lattice. The theorem supplies a positive answer to the analytic normalization question for this fixed reference construction, while retaining the arithmetic cost \(J\). It does not assert a uniform integral Hom norm, a uniform effective degree, a dense bounded-\(J\) locus, or a theorem for an arbitrary prescribed full-CM embedding coefficient in 032.

\textbf{Consequence and limit.} The Hom metric has an exact arithmetic size parameter; its determinant is not controlled solely by the finite Hecke \(L\)-value at \(3/2\). The integral Hom covolume is unbounded precisely when \(J\) is unbounded in this fixed-field, fixed-volume model, and the least nonzero squared Hom norm has the sharp order \(J^{1/g}\). Nonetheless a small actual isogeny and a Weil witness with degree \(O(\sqrt J)\) exist, giving complexity \(O(J)\). The elliptic lattices \(\mathbb Z+mi\mathbb Z\) of B11 already have \(J=m\) and minimum nonzero Hom squared norm \(m\), showing why a conductor-independent integral-Hom norm bound is unavailable in that example. Its image-cycle degree stays bounded, so this is still not a lower bound on intrinsic Weil complexity. No uniformity in \(E\) or its discriminant is asserted, and no GRH is used.

\subsection{Computation attached to B6}\label{computation-attached-to-b6}

\texttt{weil\_projector\_certificate.py} constructs the projector for \(K=\mathbb Q(\omega)\), \(\omega^2+\omega+1=0\), \(a=3+\omega\) (trace 5, norm 7), for \(n=2,3,4\), and checks exact polynomial idempotence modulo \(f_W f_R\). It also checks the norm-product resultant formula. This is a computation checking those examples, not a formal proof of B6 for all \(n\). For \(n=2\), for example, \[
 f_W=X^2-23X+2401,\quad f_R=X^3-126X^2+6174X-117649,
\] \[
 R=25\,729\,836\,300.
\] Reducing rational polynomial coefficients gives a smaller possible denominator \(476\,478\,450\) in this example. Minimality here means the least common denominator of this unique degree-\(\leq4\) interpolation polynomial, not a claim of a minimal denominator among all algebraic projector correspondences.

\subsection{Review of approach paths}\label{review-of-approach-paths}

\begin{enumerate}
\def\labelenumi{\arabic{enumi}.}
\tightlist
\item
  Tried direct GRH bound at the asserted edge. Refined by checking the actual Rallis dimensions: genuine torus-to-three-dimensional lift is at \(3/2\), and B1 is stronger and unconditional for the global factor.
\item
  Tried transfer from theta norm to degree. Direct identity is not justified; B5 supplies a valid inequality conditional on controlling all actual polarization coordinate pullbacks.
\item
  Tried conductor control along principal levels. B3 gives exact presentation growth and shows conductor is not level; no inference to intrinsic \(c(s)\) is made.
\item
  Tried use of bounded arbitrary projector norms. B9 refutes it. Replaced it by the fixed polynomial projector B6 in the generic \(K\)-action.
\item
  Tried treating the resulting rank-two coefficient index as a separate unbounded gate. B7 and B8 remove it once an effective degree bound is given, with explicit quadratic bounds. The residual gate is uniformly bounded effective witnesses with nonzero Weil component, or deformation of a full class that supplies them.
\end{enumerate}

\subsection{Suggested exact Grothendieck placement, to be independently source-checked by root}\label{suggested-exact-grothendieck-placement-to-be-independently-source-checked-by-root}

B6 concerns algebraic correspondences and primitive projectors in homological/numerical motives; B7 uses the Hodge standard positivity and the primitive Lefschetz decomposition; B8 concerns integral lattices underlying those motives together with proper cycle parameter spaces. The attached sketch points to cote 16 scans 44--50 for the standard-conjecture correspondence framework. No original archive scans are supplied here; this report does not claim to have read or corrected their transcription. Root must cite exact EGA/SGA/FGA and archive locators only after actual verification.

\subsection{Additional checked classical sources for B5 and B7}\label{additional-checked-classical-sources-for-b5-and-b7}

\begin{itemize}
\tightlist
\item
  Eduardo Cattani, \emph{Introduction to Kähler manifolds}, ICTP June 2010 preliminary notes, §5.3, Proposition 5.15 and Theorem 5.17, printed p.38: the primitive Hodge-star formula and Hodge--Riemann positivity. Example 5.19 and equation (5.23), printed pp.39--40, specialize to the surface Hodge inequality. The displayed statements and the adjoining positivity proof were read. Source: \url{https://people.maths.ox.ac.uk/joyce/Nairobi2019/CattaniKahlerManifolds.pdf}. These are classical imported results; this report does not claim an independent proof of the full Hodge--Riemann theorem.
\item
  Jean-Pierre Demailly, \emph{Complex Analytic and Differential Geometry}, version 21 June 2012, Chapter III, §1.D, equation (1.24) and Wirtinger's inequality (1.25), printed pp.136--137: analytic-cycle volume and its linear-algebra proof. Source: \url{https://people.math.harvard.edu/~demarco/Math274/Demailly_ComplexAnalyticDiffGeom.pdf}. Demailly uses \(\omega_D=i\sum h_{jk}dz_j\wedge d\bar z_k\) and the real metric \(\operatorname{Re}h\), giving \(\omega_D^p/(2^p p!)\). For the same real metric our convention is \(\omega=\omega_D/2\), so this is precisely \(\omega^p/p!\) as used in B7. The formula and proof were read and the factor of two checked.
\end{itemize}

\section{The limit of transporting a fixed witness}

Run \texttt{DEG\_L\_20261121}, author \texttt{/root}, GPT-6 Astra Pro,
effort not exposed. This result concerns the specified transport
construction. It supplies no lower bound for the minimum degree among
all cycles on a target.

\subsection{B20. Finite stabilizers prevent cancellation of the CM index
for a fixed
witness}\label{fixed-witness-transport-b20.-finite-stabilizers-prevent-cancellation-of-the-cm-index-for-a-fixed-witness}

Keep the fixed primitive CM type, maximal-order reference \(A_0\),
target \(A\), index \(J=[\Lambda:\Lambda^-]\), and compatible imaginary
quadratic Weil action of B19. Let \(\dim A_0=2n\) and \(V=\chi(L)\) be
fixed. Let \(Z_0\subset A_0\) be an integral closed subvariety of
dimension \(n\) with \[
 w_0=e_0[Z_0]\ne0,
 \qquad s_0=|\operatorname{Stab}_{A_0}(Z_0)|.
\] Here \(e_0\) is the generic Weil projector of B6 and its norm is the
positive form \(q_0=(-1)^n\int_{A_0}(-)\smile(-)\). Since the primitive
CM type makes \(A_0\) simple and \(Z_0\) is proper and
positive-dimensional, \(s_0\) is finite. For every nonzero \(E\)-linear
homomorphism \(h:A_0\to A\), put \(Z_h=h(Z_0)_{\mathrm{red}}\). Then \[
 \deg(h|_{Z_0})
   =|\ker h\cap\operatorname{Stab}_{A_0}(Z_0)|\le s_0,
\tag{B20a}
\] \[
 \|e_A h_*[Z_0]\|_{q_A}
       =\sqrt{\deg h}\,\|w_0\|_{q_0},
\tag{B20b}
\] and consequently \[
 \begin{aligned}
 \deg_L(h_*Z_0)&\ge n!\sqrt V\,\|w_0\|_{q_0}\sqrt{\deg h},\\
 \deg_L Z_h&\ge
   \frac{n!\sqrt V\,\|w_0\|_{q_0}}{s_0}\sqrt J.
 \end{aligned}
\tag{B20c}
\] Together with B19.9 this pins the least degree among these reduced
images between two positive fixed multiples of \(\sqrt J\). In
particular, transporting this one witness, or a fixed finite list of
such witnesses, cannot give uniformly bounded degrees on a sequence of
targets with \(J\to\infty\).

\textbf{Proof.} The finiteness of the stabilizer is the simple-variety
specialization of Ullmo--Yafaev, \emph{O-minimal flows on abelian
varieties}, arXiv:1610.01488v1, Proposition4.1(1), printed p8: the
identity component of the stabilizer is an abelian subvariety, and a
positive-dimensional one in a simple abelian variety would force the
proper subvariety to equal the ambient variety. The statement and its
proof were read at
\url{https://discovery.ucl.ac.uk/id/eprint/1524974/1/Yafaev_1610.01488v1.pdf}.

Write \(K_h=\ker h(\mathbf C)\). For every \(a\in K_h\) not stabilizing
\(Z_0\), the intersection \(Z_0\cap(Z_0-a)\) is a proper closed subset
of the integral variety \(Z_0\). Remove the finite union of these
subsets. For any point \(x\) of the remaining nonempty open set, the
points of \(Z_0\) mapping to \(h(x)\) are exactly \(x+a\) with
\(a\in K_h\cap\operatorname{Stab}(Z_0)\). They are distinct. Since an
isogeny over \(\mathbf C\) is finite etale, the generic fibre is reduced
and this count equals the generic degree. This proves (B20a) without a
smoothness hypothesis on \(Z_0\).

Let \(m=\deg h\). The \(K\)-linear isogeny commutes with the polynomial
Weil projectors, and \(h_*=(m)(h^*)^{-1}\) on rational cohomology. For
any two middle-degree classes \(u,v\) on \(A_0\), the projection formula
gives \[
 \begin{aligned}
 m\int_A h_*u\smile h_*v
 &=\int_{A_0}h^*h_*u\smile h^*h_*v\\
 &=m^2\int_{A_0}u\smile v.
 \end{aligned}
\] With \(u=v=w_0\), this gives the exact positive norm identity (B20b).
Apply B7 to the effective integral pushforward cycle \(h_*Z_0\) to
obtain the first inequality of (B20c). Since
\(h_*Z_0=\deg(h|_{Z_0})Z_h\), equation (B20a) bounds the division needed
to pass to its reduced image.

Finally represent \(h\) by multiplication by \(\alpha\in E^\times\). The
image lattice \(\alpha I_0\) is \(\mathcal O_E\)-stable and lies in
\(\Lambda\), hence lies in \(\Lambda^-\). Thus \[
 \deg h=[\Lambda:\alpha I_0]
     =[\Lambda:\Lambda^-]\,[\Lambda^-:\alpha I_0]\ge J.
\] This proves the second inequality. For the small isogeny chosen in
B19, the reduced image has degree at most its integral pushforward
degree, which is bounded by the right side of B19.9. The two constants
depend only on the fixed reference witness and polarization volumes. For
finitely many reference witnesses, take the minimum of their positive
lower constants and the maximum of their upper constants. QED.

\textbf{Scope and next gate.} The irreducibility assumption ensures the
exact generic-fibre count. Any effective witness has an irreducible
component with nonzero Weil projection; one may fix that component at
the start. The simplicity or finite-stabilizer condition is essential:
B11 gives an explicit product example with growing mapping multiplicity
and constant reduced-image degree. The result also does not exclude
target-dependent cycles, rational normalizations followed by different
integral extraction, or deformation of a new object. It locates the
failure of the fixed finite-reference transport route precisely, even
after generic degree is retained. No Dirichlet or Hecke GRH input enters
this result.

\textbf{Grothendieck placement.} These are finite correspondences,
proper pushforward of effective cycles, and their middle-degree Poincare
realization; the positive Weil summand is the B6--B8 motive. Cote16,
Coates letter, PDF45--46 (archive footers44--45), is the checked
algebraic-Lefschetz/positivity framework; it supplies no claim of
bounded cycles. FGA, Bourbaki221, Lemme2.4 p255 and Theoreme3.2 p260,
supplies the precise bounded-cycle parameter spaces whose degrees this
construction fails to bound. The generic-fibre count and the explicit
index inequality are the present calculation, not statements attributed
to those archival pages.

\subsection{Check status}\label{fixed-witness-transport-check-status}

Root and the bounds agent have read the complete proof. Both checked the
projection-formula factor, the kernel/stabilizer count, the
degree-to-index inequality, and the precise fixed-menu scope. The
independent checker also read the full cited stabilizer argument. Actual
check events are recorded; no formal proof-assistant verification is
claimed.

\section{Transporting a bounded family of source witnesses}

% Include-ready B21 proof; generated from bounded_source_transport.md.
\providecommand{\tightlist}{\setlength{\itemsep}{0pt}\setlength{\parskip}{0pt}}
\section{DEG\_L / bounded families of source witnesses under CM transport}\label{deg_l-bounded-families-of-source-witnesses-under-cm-transport}

Run DEG\_L\_20261121; author /root/lvalue\_bounds; GPT-6 Astra Pro; effort not exposed. This statement extends B20 from a fixed finite list to every source witness of bounded degree on the fixed simple reference variety. It concerns these transport constructions, not the intrinsic minimum over all cycles on the target.

\subsection{B21. Bounded source degree still forces square-root index growth}\label{b21.-bounded-source-degree-still-forces-square-root-index-growth}

Keep the fixed simple CM reference \((A_0,L_0)\), primitive type \((E,\Phi)\), maximal-order reference lattice, target \(A\), core index \(J=[\Lambda:\Lambda^-]\), and compatible imaginary quadratic Weil action of B19. Write \(\dim A_0=2n\), let \(V=\chi(L)\) be the fixed target polarization volume, and let \(R\) be the fixed denominator from the integral generic Weil projector \(Re_0=P([a]^*)\). Fix a positive integer \(D_0\).

There is a finite integer \[
 S_{A_0,L_0,n}(D_0)\ge1
\] such that every integral closed subvariety \(Z_0\subset A_0\) of dimension \(n\) and degree at most \(D_0\) satisfies \[
 |\operatorname{Stab}_{A_0}(Z_0)|\le S_{A_0,L_0,n}(D_0).
\tag{B21.1}
\] No numerical or algorithmic formula for this integer is asserted without presentations of the relevant Hilbert schemes and finite modules.

For every such \(Z_0\) with \(e_0[Z_0]\ne0\) and every nonzero integral \(E\)-linear homomorphism \(h:A_0\to A\), the reduced image \(Z_h=h(Z_0)_{\mathrm{red}}\) has nonzero generic Weil projection and satisfies \[
 \deg_L Z_h
 \ge
 \frac{\sqrt2\,n!\sqrt V}
      {R\,S_{A_0,L_0,n}(D_0)}\sqrt J.
\tag{B21.2}
\] If the bounded source family contains at least one such witness, let \(\mathcal D_A(D_0)\) be the infimum of \(\deg_L h(Z_0)_{\mathrm{red}}\) over all source witnesses and nonzero integral \(E\)-linear homomorphisms just specified. There is the two-sided bound \[
 \begin{aligned}
 \frac{\sqrt2\,n!\sqrt V}
      {R\,S_{A_0,L_0,n}(D_0)}\sqrt J
 &\le \mathcal D_A(D_0)
 \le D_0K_E^n\sqrt{\frac{V}{V_0}}\sqrt J,\\
 K_E&=\frac2\pi\bigl((2n)!\sqrt{|D_E|}\bigr)^{1/(2n)}.
 \end{aligned}
\tag{B21.3}
\] Here \(h\) ranges over integral \(E\)-linear homomorphisms to the specified target \(A\). The constants are independent of how the source witness is selected within the entire degree-bounded family.

\subsubsection{Proof of the uniform stabilizer bound}\label{proof-of-the-uniform-stabilizer-bound}

Choose a fixed positive integer \(k\) such that \(L_0^{\otimes k}\) is very ample. The given source subvarieties then have degree at most \(k^nD_0\) in one fixed projective embedding. Grothendieck's boundedness result for reduced pure-dimensional subschemes gives only finitely many Hilbert polynomials for these subvarieties. This is FGA, Bourbaki exposé 221, the bounded-family discussion and Lemme 2.4, printed p.255, as used in DIV1. Theorem 3.2, printed p.260, supplies the corresponding projective Hilbert schemes. Let \(\mathcal H\) be the finite disjoint union of those Hilbert schemes for subschemes of \(A_0\).

Every point of \(\mathcal H\) represents a nonempty subscheme of dimension \(n<2n\). Boundary points may be nonreduced, nonpure, or have embedded components. Nothing in the following argument assumes otherwise.

Translation by \(A_0\) acts algebraically on \(\mathcal H\). To verify that it preserves each chosen Hilbert polynomial, fix a subscheme and consider its translated family over \(A_0\). This family is flat, since it is obtained from its product with \(A_0\) by an automorphism of \(A_0\times A_0\). Its Hilbert polynomial is locally constant, and \(A_0\) is connected. The polynomial therefore equals that of the original subscheme at every translation. The universal translated family gives the algebraic action by the Hilbert functor.

Consider the universal stabilizer \[
 \mathcal S=\{(a,z)\in A_0\times\mathcal H:a\cdot z=z\}.
\tag{B21.4}
\] It is a closed subscheme: it is the inverse image of the diagonal of the separated scheme \(\mathcal H\) under \((a,z)\mapsto(a\cdot z,z)\). Hence \(\mathcal S\to\mathcal H\) is proper, since \(A_0\) is proper.

Every geometric fibre is zero-dimensional. Indeed, the identity component of the stabilizer of a subscheme is a connected algebraic subgroup of \(A_0\). If it had positive dimension, simplicity would make it all of \(A_0\). The support of the nonempty subscheme would then be invariant under every translation, hence equal to \(A_0\), contradicting its dimension \(n<2n\). This argument applies to every Hilbert boundary point as well. The stabilizer fibres are therefore finite.

A proper morphism with finite fibres is finite: Stacks Project, Lemma 37.44.1, tag 02LS, equivalence of (1) and (2). Consequently \[
 \mathcal M=(\mathcal S\to\mathcal H)_*\mathcal O_{\mathcal S}
\] is a coherent finite \(\mathcal O_{\mathcal H}\)-module. Cover the quasi-compact scheme \(\mathcal H\) by finitely many affine opens \(U_i\), and choose a finite number \(N_i\) of module generators for \(\mathcal M|_{U_i}\). At a point \(z\in U_i\) the scheme length of the stabilizer fibre is \[
 \dim_{\kappa(z)}(\mathcal M\otimes\kappa(z))\le N_i.
\] The number of its geometric points is no larger than this length. Therefore one may take \[
 S_{A_0,L_0,n}(D_0)=\max(1,N_1,\ldots,N_s).
\] This proves (B21.1). The finite-generator construction establishes existence of the constant; effective computation would require explicit charts and presentations and is not claimed.

\subsubsection{Proof of the transport estimates}\label{proof-of-the-transport-estimates}

For any integral nonzero generic Weil witness, the strengthened B18 gives \[
 q_0(P[Z_0],P[Z_0])\in2\mathbb Z_{>0},
 \qquad
 \|e_0[Z_0]\|_{q_0}\ge\frac{\sqrt2}{R}.
\tag{B21.5}
\] This lower bound uses the fixed integral Weil lattice and applies separately to every member of the bounded family. It does not require a chosen flat class on that Hilbert family.

B20 gives the exact generic-degree and norm identities \[
 \begin{aligned}
 \delta_h&=\deg(h|_{Z_0})
   =|\ker h\cap\operatorname{Stab}_{A_0}(Z_0)|
   \le S_{A_0,L_0,n}(D_0),\\
 \|e_A h_*[Z_0]\|_{q_A}
   &=\sqrt{\deg h}\,\|e_0[Z_0]\|_{q_0}.
 \end{aligned}
\tag{B21.6}
\] The first remains valid for singular integral \(Z_0\) by the generic-fibre argument in B20. Since \(h_*Z_0=\delta_h Z_h\), the nonzero Weil component of the pushforward implies that of the reduced image. The B7 comass estimate, (B21.5), and (B21.6) give \[
 \begin{aligned}
 \delta_h\deg_L Z_h
 &=\deg_L h_*Z_0\\
 &\ge n!\sqrt V\,\sqrt{\deg h}\,\|e_0[Z_0]\|_{q_0}\\
 &\ge \frac{\sqrt2\,n!\sqrt V}{R}\sqrt{\deg h}.
 \end{aligned}
\] Every nonzero \(E\)-linear integral homomorphism has \(\deg h\ge J\) by B19.8. Division by the uniform upper bound for \(\delta_h\) proves (B21.2).

For the upper bound, select any one source witness of degree at most \(D_0\) and use the small homomorphism from B19.7. The reduced-image degree is no larger than its effective pushforward degree, and B19.9 gives the right side of (B21.3). The lower bound holds for every choice in that infimum. This proves (B21.3). QED.

\subsubsection{Scope of the improvement}\label{scope-of-the-improvement}

The chosen source witness may now depend arbitrarily on the target, as long as it is an integral subvariety on the fixed \(A_0\) and its source degree stays at most \(D_0\). It may move through infinitely many Hilbert points; a finite reference list is unnecessary. An effective integral cycle of bounded degree with nonzero Weil projection has an irreducible component with nonzero Weil projection and no larger degree, so the same obstruction applies to selecting such a component before transport. The reduced support of the image of the whole effective cycle contains that component's image and has at least its degree.

Thus a uniformly bounded-degree source family cannot yield uniformly bounded-degree effective transported witnesses when \(J\to\infty\), even after generic mapping multiplicities are removed. A successful uniform construction must permit unbounded source complexity, move beyond this fixed-reference isogeny transport, or supply a separate target-dependent cycle or deformation mechanism. This conclusion supplies no lower bound for the intrinsic degree minimum over all target cycles, and no lower bound for intrinsic bounded-division complexity. Dirichlet GRH and Hecke GRH do not enter the proof.

\subsubsection{Exact imported inputs and check scope}\label{exact-imported-inputs-and-check-scope}

\begin{itemize}
\tightlist
\item
  Grothendieck, FGA, Bourbaki exposé 221, bounded-family discussion and Lemme 2.4 p.255; Théorème 3.2 p.260: bounded reduced pure-dimensional degree and projective Hilbert schemes. These are the already verified inputs of DIV1. The universal-stabilizer construction above is the present deduction.
\item
  Stacks Project, Lemma 37.44.1, tag 02LS, \url{https://stacks.math.columbia.edu/tag/02LS}: proper with finite fibres is finite. The exact statement and entire short proof were read.
\item
  Ullmo--Yafaev, \emph{o-minimal flows on abelian varieties}, arXiv:1610.01488v1, Proposition 4.1(1), p.8: the abelian-subvariety stabilizer input for reduced irreducible varieties, already checked in B20. The support argument above extends the needed finiteness to the arbitrary Hilbert boundary subschemes.
\item
  B6--B8, strengthened B18, B19, and B20: the exact fixed projector, positive integral norm, comass estimate, CM lattice index and generic-fibre identities are retained with their stated hypotheses.
\end{itemize}

The complete proof has been read, including the arbitrary nonreduced Hilbert boundary, the passage from finite fibres to a finite morphism, the uniform module-generator bound, and both constants in (B21.3). No formal proof-assistant verification is claimed.

\section{An explicit source-complexity obstruction}

\providecommand{\tightlist}{\setlength{\itemsep}{0pt}\setlength{\parskip}{0pt}}
\section{DEG\_L / explicit stabilizer and source-degree bounds}\label{deg_l-explicit-stabilizer-and-source-degree-bounds}

Run DEG\_L\_20261121; author /root/lvalue\_bounds; GPT-6 Astra Pro; effort not exposed. This refines the existence constant in B21. The geometric input is the classical generalized Bézout inequality, imported with its exact source; the argument below tracks its constants in the present transport problem.

\subsection{B22. Explicit stabilizer bound and a necessary source-degree growth rate}\label{b22.-explicit-stabilizer-bound-and-a-necessary-source-degree-growth-rate}

Let \((A_0,L_0)\) be a polarized complex abelian variety, and fix an integer \(k\ge1\) such that \(L_0^{\otimes k}\) is very ample. Let \(Z\subset A_0\) be an integral subvariety of dimension \(n\ge1\) with finite translation stabilizer. Set \[
 D=\deg_{L_0}Z=\int_Zc_1(L_0)^n.
\] Then \[
 |\operatorname{Stab}_{A_0}(Z)|
 \le (k^nD)^{n+1}
 =k^{n(n+1)}D^{n+1}.
\tag{B22.1}
\] In particular, one may use \(k=3\) for every ample line bundle on an abelian variety. If \(L_0\) is already very ample, one may use \(k=1\). Simplicity of \(A_0\) is sufficient for the finite-stabilizer hypothesis whenever \(0<n<\dim A_0\), but is not needed for (B22.1) itself.

Now retain the fixed simple CM reference and compatible generic Weil data of B19--B21, with \(\dim A_0=2n\), target polarization volume \(V=\chi(L)\), fixed integral-projector denominator \(R\), and target maximal-order core index \(J\). For every integral source witness \(Z\subset A_0\) with \(e_0[Z]\ne0\) and every nonzero integral \(E\)-linear homomorphism \(h:A_0\to A\), write \[
 D'=\deg_Lh(Z)_{\mathrm{red}}.
\] There is the explicit inequality \[
 D'D^{n+1}
 \ge
 \frac{\sqrt2\,n!\sqrt V}{R\,k^{n(n+1)}}\sqrt J.
\tag{B22.2}
\] Consequently a transport construction satisfying \(D'\le B\), with \(B>0\) fixed, must use \[
 D\ge k^{-n}
 \left(\frac{\sqrt2\,n!\sqrt V}{RB}\right)^{1/(n+1)}
 J^{1/(2(n+1))}.
\tag{B22.3}
\] For surfaces in abelian fourfolds, the universal choice \(k=3\) gives \[
 D\ge\frac19
 \left(\frac{2\sqrt2\sqrt V}{RB}\right)^{1/3}J^{1/6}.
\tag{B22.4}
\] No optimality of the exponent in (B22.1) or (B22.3) is asserted. These are necessary bounds for the specified transport mechanism, not lower bounds for the intrinsic least degree of a target Weil witness.

\subsubsection{Proof of (B22.1)}\label{proof-of-b22.1}

Write \(H=\operatorname{Stab}_{A_0}(Z)\). As sets of complex points, \[
 H=\bigcap_{z\in Z}(Z-z).
\tag{B22.5}
\] Indeed, membership on the right says that \(a+Z\subseteq Z\). Both sides of this inclusion are irreducible closed subvarieties of the same dimension, so the inclusion is equality. Conversely, a stabilizing translation belongs to every displayed translate.

We select at most \(n+1\) of the translates in (B22.5) whose intersection is finite. Start with any \(z_0\in Z\) and \(X_0=Z-z_0\). Suppose that \(X_j\) has a positive-dimensional irreducible component \(C\). Since \(H\) is finite, choose \(c_C\in C\setminus H\). By (B22.5), there is a point of \(Z\) whose translate does not contain \(c_C\). Equivalently, \[
 B_C=\{z\in Z:c_C\in Z-z\}=Z\cap(Z-c_C)
\] is a proper closed subset of \(Z\). There are only finitely many positive-dimensional components \(C\) of \(X_j\). Irreducibility of \(Z\) allows a point \(z_{j+1}\) outside their finitely many proper closed sets \(B_C\). The translate \(Z-z_{j+1}\) then contains none of those components in its entirety. Thus \[
 X_{j+1}=X_j\cap(Z-z_{j+1})
\] has strictly smaller dimension on every previously positive-dimensional component. Starting at dimension \(n\), after at most \(n\) further intersections all components have dimension zero. The resulting nonempty finite intersection contains \(H\). If it becomes finite sooner, repeat any one of the already chosen translates to have exactly \(n+1\) factors.

Translation on the connected complex torus \(A_0\) acts trivially on integral cohomology, hence preserves numerical degree. In the fixed embedding defined by \(L_0^{\otimes k}\), all of these \(n+1\) translates have degree \(k^nD\).

We now apply Fulton, \emph{Intersection Theory}, second edition, Example 8.4.7, printed p.148: for subvarieties of a fixed projective scheme with finite intersection, the number of intersection points is at most the product of their projective degrees. Its input, Example 8.4.6 on the same page, bounds the sum of degrees of all irreducible components of an arbitrary intersection and imposes no proper-intersection hypothesis. Therefore \[
 |H|\le\#\bigcap_{j=0}^{n}(Z-z_j)
 \le\prod_{j=0}^{n}(k^nD)=(k^nD)^{n+1}.
\] This proves (B22.1). The classical fact that \(L_0^{\otimes3}\) is very ample is stated, for example, in Milne, \emph{Abelian Varieties}, version 2.00, Chapter I, section 6, printed p.30, in the paragraph immediately preceding Example 6.5. We import this fact; no new proof of very ampleness or of Bézout is intended here. QED.

\subsubsection{Proof of (B22.2)--(B22.4)}\label{proof-of-b22.2b22.4}

The generic degree of \(h\) on \(Z\) is, by the exact B20 finite-kernel count, \[
 \delta_h=|\ker h\cap\operatorname{Stab}_{A_0}(Z)|
 \le k^{n(n+1)}D^{n+1}.
\tag{B22.6}
\] The strengthened B18 even-square integrality gives \[
 \|e_0[Z]\|_{q_0}\ge\frac{\sqrt2}{R}.
\] Poincaré pairing under the finite isogeny and the comass bound B7 give \[
 \begin{aligned}
 \delta_hD'
 &=\deg_Lh_*Z\\
 &\ge n!\sqrt V\,\|e_Ah_*[Z]\|_{q_A}\\
 &=n!\sqrt V\sqrt{\deg h}\,\|e_0[Z]\|_{q_0}\\
 &\ge\frac{\sqrt2\,n!\sqrt V}{R}\sqrt J.
 \end{aligned}
\tag{B22.7}
\] The last inequality uses \(\deg h\ge J\), the exact integral Hom lattice bound B19. No choice of a small homomorphism is needed for this lower bound. Combining (B22.6) and (B22.7) proves (B22.2); taking the \((n+1)\)st root proves (B22.3). Substitution of \(n=2\) and \(k=3\) proves (B22.4). QED.

\subsubsection{Consequences and limits}\label{consequences-and-limits}

One may replace the abstract constant of B21 by the explicit value \[
 S_{A_0,L_0,n}(D_0)=k^{n(n+1)}D_0^{n+1}.
\] Thus every target-dependent choice of source witnesses of degrees at most \(D_0\) satisfies the fully explicit lower bound obtained by substituting this expression into B21.2. The Hilbert-scheme proof B21 remains useful for arbitrary boundary subschemes and for other bounded families; the present sharper constant applies to the reduced integral witnesses needed in the transport construction.

The result does not say that growing the source degree at the indicated rate is sufficient. It also does not exclude new target cycles, or a construction using a changing reference family. In particular, an obstruction to these reference-isogeny transports cannot be promoted to a lower bound for intrinsic bounded-division complexity. Dirichlet GRH and Hecke GRH do not enter the argument.

\subsubsection{Primary sources and actual reading scope}\label{primary-sources-and-actual-reading-scope}

\begin{itemize}
\tightlist
\item
  William Fulton, \emph{Intersection Theory}, second edition, Springer, 1998, Examples 8.4.5--8.4.7, printed pp.147--148. The original page images were visually read, including the full ruled-join reduction in Example 8.4.5, the complete short proof of the generalized Bézout inequality in Example 8.4.6, and the finite-intersection specialization in Example 8.4.7. The source scan is \url{https://djvu.online/file/Pl4avuZzLcI6I}; the exact original page images are \url{https://djvu.online/jpg2/P/l/4/Pl4avuZzLcI6I/161.webp} and \url{https://djvu.online/jpg2/P/l/4/Pl4avuZzLcI6I/162.webp}. Local checked images are \texttt{work/external/Fulton\_page161.webp} and \texttt{work/external/Fulton\_page162.webp}.
\item
  J.S. Milne, \emph{Abelian Varieties}, version 2.00, 16 March 2008, Chapter I, section 6, printed p.30, the paragraph immediately preceding Example 6.5: the cube of any ample line bundle on an abelian variety is very ample. Official source \url{https://www.jmilne.org/math/CourseNotes/AV.pdf}, PDF index 35. The exact assertion was read; this source states it without proof, and its upstream proof is imported.
\item
  Marc Hindry, \emph{Autour d'une conjecture de Serge Lang}, Invent. Math. 94 (1988), Lemmas 5--8, printed pp.579--581, gives the related degree-and-equation formulation for stabilizers. The original author-hosted scan \url{https://webusers.imj-prg.fr/~marc.hindry/Conj-Lang.pdf} was read on those pages, including the full proof of Lemma 8. Its extra equation-degree hypotheses are unnecessary in the direct numerical-degree application of Fulton above; no unverified conversion from defining-equation degree is used.
\item
  B7, strengthened B18, B19.8, B20 and B21 supply the already checked polarization normalization, integral Weil lower bound, integral Hom index, and generic-degree identity. The whole argument above, including the finite intersection selection and exponents, has been read and checked. No proof-assistant verification is claimed.
\end{itemize}

\subsection{B22a. Two translates and an intrinsic quadratic stabilizer bound}\label{b22a.-two-translates-and-an-intrinsic-quadratic-stabilizer-bound}

Assume that \(A_0\) is simple of complex dimension \(2n\), and let \(Z\subset A_0\) be an integral subvariety of dimension \(n\ge1\). Write \[
 D=\deg_{L_0}Z,\qquad V_0=\chi(L_0),\qquad H=\operatorname{Stab}_{A_0}(Z).
\] Then the difference map \(\Delta:Z\times Z\to A_0\), \((z_1,z_2)\mapsto z_1-z_2\), is dominant and generically finite. Its degree is the positive self-intersection of the middle-dimensional class. More precisely, \[
 \begin{aligned}
 d_\Delta&=\int_{A_0}[Z]^2>0,\\
 d_\Delta/|H|&\in2\mathbf Z_{>0},\\
 |H|&\le\frac12\int_{A_0}[Z]^2
 \le\frac{D^2}{4(n!)^2V_0}.
 \end{aligned}
\tag{B22a.1}
\] The last inequality uses the polarization metric directly; no projective-embedding power is needed.

In the fixed simple CM transport setting of B19--B23, for a source witness with nonzero generic Weil projection and target reduced-image degree \(D'\), this improves (B22.2) to \[
 D'D^2\ge
 \frac{4\sqrt2\,(n!)^3V_0\sqrt V}{R}\sqrt J.
\tag{B22a.2}
\] Thus \(D'\le B\) requires \[
 D\ge
 \left(\frac{4\sqrt2\,(n!)^3V_0\sqrt V}{RB}\right)^{1/2}J^{1/4}.
\tag{B22a.3}
\] For surfaces in abelian fourfolds the constant inside the square root is \(32\sqrt2\,V_0\sqrt V/(RB)\). These bounds retain exactly the reference-transport scope of B22.

\subsubsection{Proof of dominance and the self-intersection formula}\label{proof-of-dominance-and-the-self-intersection-formula}

We use the characteristic-zero sum-dimension theorem in its standard geometric-nondegeneracy formulation. Debarre, \emph{On the geometry of abelian varieties}, Chapter 1, Definition 2.1 and Examples 2.2(2), printed p.8, states that every irreducible subvariety of a simple complex abelian variety is geometrically nondegenerate. Corollary 2.5(a), printed pp.9--10, then gives \[
 \dim(V+W)=\min(\dim V+\dim W,\dim A_0)
\] when \(V\) is geometrically nondegenerate. Smoothness of the subvarieties is not assumed. Apply this to \(V=Z\) and \(W=-Z\). It gives \(Z-Z=A_0\) and hence dominance of \(\Delta\). The dimensions of its domain and target are both \(2n\), so it is generically finite. These classical results are imported, not reproved here.

We spell out the generic-fibre issue for possibly singular \(Z\). The singular boundary \[
 (Z_{\mathrm{sing}}\times Z)\cup(Z\times Z_{\mathrm{sing}})
\] has dimension at most \(2n-1\). Its image under the proper map \(\Delta\) is a proper closed subset of \(A_0\). Remove it. Also restrict to a nonempty open set where the proper generically finite map is finite. In characteristic zero, generic smoothness, followed by removing the image of the remaining ramification locus of this finite map, gives a further nonempty open set over which it is étale. Fibres there consist of exactly \(d_\Delta\) reduced points, all lying in \(Z_{\mathrm{reg}}\times Z_{\mathrm{reg}}\).

For such a point \(u\in A_0\), the first projection identifies \(\Delta^{-1}(u)\) with \(Z\cap(Z+u)\). At a corresponding pair \((z_1,z_2)\) the kernel of the differential of \(\Delta\) is the intersection of the translated tangent spaces \(T_{z_1}Z\cap T_{z_2}Z\). Étaleness makes this zero; since both tangent spaces have dimension \(n\), the intersection is transverse in the ambient \(2n\)-fold. Every local intersection multiplicity is therefore one. Translation acts trivially on cohomology, so \[
 d_\Delta=\#\bigl(Z\cap(Z+u)\bigr)=\int_{A_0}[Z]^2>0.
\tag{B22a.4}
\]

This also gives the proposed two-translate argument explicitly. Choose \((z_1,z_2)\in\Delta^{-1}(u)\). The finite intersection \((Z-z_1)\cap(Z-z_2)\) is isomorphic to the same fibre and contains \(H\). Equivalently, \(H\) acts freely by simultaneous translation on every such difference fibre. In particular \(|H|\) divides \(d_\Delta\). Fulton Example 8.4.7 would already give the quadratic bound \(|H|\le k^{2n}D^2\); the next two steps improve its constant.

\subsubsection{Proof of the factor two and the metric bound}\label{proof-of-the-factor-two-and-the-metric-bound}

Let \(q:A_0\to A_0/H\) be the quotient isogeny, of degree \(|H|\), and set \(Z'=q(Z)_{\mathrm{red}}\). The full inverse image of \(Z'\) is \(Z\): set-theoretically it is the union of the translates \(Z+h\), all equal to \(Z\). The isogeny is étale, and the inverse image of a reduced complex subvariety under an étale map is reduced. Thus this is also the scheme-theoretic equality needed for flat pullback of fundamental cycles. Consequently \[
 [Z]=q^*[Z'],\qquad
 \int_{A_0}[Z]^2=|H|\int_{A_0/H}[Z']^2.
\tag{B22a.5}
\] The rightmost self-intersection is a positive integer by (B22a.4). It is even: integral cohomology of an abelian variety is an exterior algebra on its degree-one lattice; a positive even-degree class has zero squares on individual exterior monomials, while its distinct cross terms occur in pairs. This is the general even-square argument already checked in B18. Hence the quotient in (B22a.5) belongs to \(2\mathbf Z_{>0}\), proving \(|H|\le d_\Delta/2\).

Choose the invariant Kähler form \(\omega_0\) representing \(c_1(L_0)\) and its flat metric. The total metric volume is \[
 \operatorname{Vol}(A_0)=\int_{A_0}\frac{\omega_0^{2n}}{(2n)!}=V_0.
\] Let \(\alpha\) be the invariant harmonic representative of \([Z]\), and write \(\alpha_+=(\alpha+*\alpha)/2\) and \(\alpha_-=(\alpha-*\alpha)/2\). On real middle-degree forms in real dimension \(4n\), the Hodge star is an isometric involution. If \(v\) is a unit decomposable real middle-degree covector, then \(\langle v,*v\rangle=0\), because \(v\wedge v=0\). Hence its self-dual projection has norm \(1/\sqrt2\). For every self-dual form \(\eta\) this proves the pointwise estimate \[
 \operatorname{comass}(\eta)\le |\eta|/\sqrt2.
\] All harmonic forms for this flat torus metric are invariant, so the pointwise norm of \(\alpha_+\) is \(\|\alpha_+\|_{L^2}/\sqrt{V_0}\). Poincaré duality, self-duality, orthogonality of the two star eigenspaces, and Wirtinger's formula give \[
 \begin{aligned}
 \|\alpha_+\|_{L^2}^2
 &=\int_{A_0}\alpha\wedge\alpha_+
 =\int_Z\alpha_+\\
 &\le\operatorname{comass}(\alpha_+)\operatorname{Vol}(Z)
 \le\frac{\|\alpha_+\|_{L^2}}{\sqrt{2V_0}}\frac{D}{n!}.
 \end{aligned}
\tag{B22a.6}
\] For singular \(Z\), the integral is over its smooth locus, or equivalently its integration current; the singular locus has zero volume and Wirtinger's formula still gives \(D/n!\). The self-dual part is nonzero because the self-intersection is positive. Cancelling its norm yields \[
 \int_{A_0}[Z]^2
 =\|\alpha_+\|_{L^2}^2-\|\alpha_-\|_{L^2}^2
 \le\|\alpha_+\|_{L^2}^2
 \le\frac{D^2}{2(n!)^2V_0}.
\] Together with the preceding factor two this proves (B22a.1).

Finally, the exact generic degree \(\delta_h\) of the reference isogeny on \(Z\) is no larger than \(|H|\). Substitute (B22a.1) into the already proved inequality (B22.7). This gives (B22a.2), and taking a square root gives (B22a.3). QED.

\subsubsection{Improved integer constant for the bounded-family and ideal-content results}\label{improved-integer-constant-for-the-bounded-family-and-ideal-content-results}

For a positive integer source-degree bound \(D_0\), one may now use the explicit positive integer \[
 S_2(D_0)=\max\left\{1,
 \left\lfloor\frac{D_0^2}{4(n!)^2V_0}\right\rfloor\right\}
\tag{B22a.7}
\] in B21 and in B23.17--B23.18. The maximum with one covers the empty small-degree family and ensures that \(\operatorname{lcm}(1,\ldots,S_2(D_0))\) is defined in the same convention. If there is any such source subvariety, (B22a.1) itself ensures that the displayed real upper bound is at least one. The proof of B23 uses only an integer upper bound on every generic mapping multiplicity, so its exact divisibility and fixed-order comparison remain valid with this improved choice.

\subsubsection{Exact source locators and check scope for B22a}\label{exact-source-locators-and-check-scope-for-b22a}

\begin{itemize}
\tightlist
\item
  Olivier Debarre, \emph{On the geometry of abelian varieties}, author-hosted lecture notes, \url{https://perso.imj-prg.fr/wp-content/uploads/debarre-pub/Taiwan.pdf}, Chapter 1: Lemmas 1.1--1.2 p.8; Definition 2.1 and Examples 2.2(2) p.8; Theorem 2.3, Lemma 2.4, and Corollary 2.5(a) pp.9--10. The exact statements and all the adjoining short proofs were read. All varieties in this chapter are irreducible, reduced, and complex; no smoothness assumption is added to Corollary 2.5.
\item
  The same simple-ambient nondegeneracy assertion appears in Debarre, \emph{Fulton--Hansen and Barth--Lefschetz theorems for subvarieties of abelian varieties}, J. reine angew. Math. 467 (1995), section 1, Examples (1), pp.188--189; original author-hosted scan \url{https://perso.imj-prg.fr/wp-content/uploads/debarre-pub/30.pdf}.
\item
  A direct current statement is Olivier Debarre and Ben Moonen, \emph{Bézout's theorem for abelian varieties}, Theorem 1.1 p.1, \url{https://www.math.ru.nl/~bmoonen/Papers/BezoutAV.pdf}; section 2.1 p.3 reviews the characteristic-zero argument. These statements and that complete short review were read. The stronger positive-characteristic and semismallness results are unnecessary here and are not independently audited.
\item
  The middle-class estimate uses the invariant-metric/comass argument of B7, sharpened by the explicit self-dual projection calculation above. The full proof was checked, including singular-boundary removal, reduced generic fibres, the stabilizer quotient, even self-intersection, the extra factor two in the self-dual comass bound, and the improved integer constant. It supplies no assertion for a merely finite-stabilizer subvariety of a nonsimple ambient variety without the stated geometric nondegeneracy input.
\end{itemize}

\section{Integral content of all reference transports}

\providecommand{\tightlist}{\setlength{\itemsep}{0pt}\setlength{\parskip}{0pt}}
\section{DEG\_L / integral ideal content of all CM transports}\label{deg_l-integral-ideal-content-of-all-cm-transports}

Run DEG\_L\_20261121; author /root/lvalue\_bounds; GPT-6 Astra Pro; effort not exposed. This tests whether integral combinations of many reference isogenies can remove the index appearing in B19. The answer below concerns pushforwards with their actual multiplicities. It does not apply to replacing each pushforward by its reduced image.

\subsection{B23. Relative determinant content survives every integral combination}\label{b23.-relative-determinant-content-survives-every-integral-combination}

Let \(K\) be an imaginary quadratic field, let \(E\) be a CM field containing \(K\), and put \[
 [E:K]=g=2n.
\] Fix a CM type of \(E\) whose restriction to \(K\) has signature \((n,n)\). Suppose that the rational homology has rank one over \(E\). Let \[
 A_0=\mathbf C^g/\Phi(I_0),\qquad
 A=\mathbf C^g/\Phi(\Lambda),
\] where \(I_0\) is an invertible fractional \(\mathcal O_E\)-ideal and \(\Lambda\) is a full lattice stable under \(\mathcal O_K\). The target is not assumed to have integral \(\mathcal O_E\) action. Simplicity and primitivity of the CM type are unnecessary for the lattice assertions below.

Let \(W_K(A)\subset H^g(A,\mathbf Q)\) be the usual two-dimensional rational Weil summand, the sum of the two extreme \(K\)-eigenspaces after extending scalars. Define its actual integral intersection lattice by \[
 \mathscr W_A=W_K(A)\cap H^g(A,\mathbf Z),
 \qquad
 \mathscr W_0=W_K(A_0)\cap H^g(A_0,\mathbf Z).
\tag{B23.1}
\] On this rational Weil plane use the determinant \(K\)-module structure: \(\mathrm{GL}_K(H^1)\) acts by its determinant. Thus the geometric endomorphism \([a]^*\) acts by the determinant scalar \(a^g\), not by \(a\). We identify the rational cohomology spaces using the trace pairing \(\operatorname{Tr}_{E/\mathbf Q}(xy)\). This fixes the conjugation convention in the formulas below.

Put \[
 \begin{aligned}
 \Lambda^-&=\{x\in\Lambda:\mathcal O_E x\subset\Lambda\},
 &J&=[\Lambda:\Lambda^-],\\
 I&=(\Lambda:I_0)=\operatorname{Hom}_E(A_0,A),
 &\mathfrak d&=\frac{\det_{\mathcal O_K}\Lambda}
                       {\det_{\mathcal O_K}I_0},\\
 \mathfrak j&=\frac{\det_{\mathcal O_K}\Lambda^-}
                       {\det_{\mathcal O_K}\Lambda}.
 \end{aligned}
\tag{B23.2}
\] Determinant modules are rank-one projective modules; their ratios are fractional ideals of \(K\). In particular, \(\mathfrak j\) is integral and \[
 N_{K/\mathbf Q}(\mathfrak j)=J.
\]

The following formulas are exact, including at primes ramified in \(K\): \[
 \begin{aligned}
 \mathscr W_A&=\mathfrak d^{-1}\mathscr W_0,\\
 N_{E/K}(I)&=\mathfrak j\mathfrak d,\\
 \sum_{0\ne h\in\operatorname{Hom}_E(A_0,A)}h_*\mathscr W_0
 &=\overline{\mathfrak j}\,\mathscr W_A.
 \end{aligned}
\tag{B23.3}
\] Here the sum means the additive subgroup generated by all the indicated integral images; it is already an \(\mathcal O_K\)-submodule. Thus its index in the target Weil lattice is exactly \(J\).

For a fixed nonzero \(w_0\in\mathscr W_0\), set \[
 \Gamma_{w_0}=\sum_{h}\mathbf Z\,h_*w_0,
 \qquad
 c_0=[\mathscr W_0:\mathcal O_Kw_0].
\] Choose once and for all \(a\in\mathcal O_K\) such that \(a^g\notin\mathbf R\), and put \[
 C=[\mathcal O_K:\mathbf Z[a^g]].
\tag{B23.4}
\] Then \[
 \begin{aligned}
 \mathcal O_K\Gamma_{w_0}
 &=\overline{\mathfrak j}\,\mathfrak d^{-1}\mathcal O_Kw_0,\\
 C\,\mathcal O_K\Gamma_{w_0}
 &\subseteq\Gamma_{w_0}
 \subseteq\mathcal O_K\Gamma_{w_0},\\
 Jc_0&\le[\mathscr W_A:\Gamma_{w_0}]\le C^2Jc_0.
 \end{aligned}
\tag{B23.5}
\] The number \(C\) is fixed independently of the target lattice. Hence even the additive span of every transport of one fixed nonzero integral Weil class retains index of order \(J\).

Such an \(a\) exists: choose a nonreal \(\theta\in\mathcal O_K\) and take \(a=m+\theta\) with \(m\in\mathbf Z\). The imaginary part of \((m+\theta)^g\) is a nonzero real polynomial in \(m\), with leading coefficient \(g\operatorname{Im}\theta\), so only finitely many integers are excluded.

\subsubsection{Consequence for algebraic cycles and division integers}\label{consequence-for-algebraic-cycles-and-division-integers}

Assume the fixed integral generic Weil projector of B6 is available: \[
 P_0=Re_0\quad\text{on }H^g(A_0),\qquad
 P_A=Re_A\quad\text{on }H^g(A),
\] with the same positive integer \(R\). Let \(\xi\) be any integral linear combination of the cohomology classes of multiplicity-preserving pushes \(h_*Z\), where \(h:A_0\to A\) is integral and \(E\)-linear and \(Z\) is any integral codimension-\(n\) cycle on \(A_0\). The source cycles may vary without any degree bound. Then \[
 P_A\xi\in\overline{\mathfrak j}\,\mathscr W_A.
\tag{B23.6}
\] In particular, if \(t\in\mathscr W_A\) is nonzero and \(\xi=Nt\) for a positive integer \(N\), put \[
 c_t=[\mathscr W_A:\mathcal O_Kt].
\] There is the necessary divisibility and lower bound \[
 J\mid R^2N^2c_t,
 \qquad
 N\ge\frac{\sqrt J}{R\sqrt{c_t}}.
\tag{B23.7}
\] In a fixed generic Weil family with integral \(\mathcal O_K\) action, after fixing the integral local system and a nonzero flat integral section \(t\), the integer \(c_t\) is constant. Thus no integral combination of these reference transports can represent a fixed such class using uniformly bounded division integers as \(J\to\infty\).

This is an obstruction for the indicated construction. It is not an intrinsic lower bound over all algebraic cycles on \(A\).

\subsection{Proof}\label{transport_ideal_content_proofs-proof}

\subsubsection{1. The actual integral Weil lattice is a determinant-lattice functor}\label{the-actual-integral-weil-lattice-is-a-determinant-lattice-functor}

Let \(M\) be a full \(\mathcal O_K\)-lattice of rank \(g\) in a \(K\)-vector space \(V\). Let \(W(V)\) be its two-dimensional rational extreme exterior summand in \(\bigwedge^g_{\mathbf Q}V\). Define \[
 L(M)=W(V)\cap\bigwedge^g_{\mathbf Z}M.
\] Fix a rational prime \(p\). The ring \(\mathcal O_K\otimes\mathbf Z_p\) is a product of discrete valuation rings, and \(M\otimes\mathbf Z_p\) is free of rank \(g\) over this ring. Under a chosen local basis the two-dimensional rational summand is the determinant representation of \(\mathrm{GL}_g(K\otimes\mathbf Q_p)\).

The integral intersection in the standard local lattice is stable under the full ring \(\mathcal O_K\otimes\mathbf Z_p\). Indeed, for any \(b\) in that ring the \(\mathcal O_K\otimes\mathbf Z_p\)-linear endomorphism with matrix \(\operatorname{diag}(b,1,\ldots,1)\) preserves the standard lattice. Its ordinary \(g\)th exterior power preserves the integral exterior lattice and acts on the extreme summand by the determinant scalar \(b\). This argument works for nonunits and for the split product ring as well. The standard intersection is therefore a rank-one projective module over the maximal order.

If \(T\) is a \(K\otimes\mathbf Q_p\)-linear change of lattice, its ordinary exterior power takes the full integral exterior lattice isomorphically onto that of the new lattice and preserves the rational extreme summand. Consequently \[
 L(TM_p)=\det_K(T)L(M_p).
\tag{B23.8}
\] The possible scalar ideal distinguishing the standard intersection from the abstract determinant module is the same before and after the change of lattice, so it cancels in (B23.8). No unramifiedness assertion or division by the discriminant of \(K\) is being used.

These local identities glue to the corresponding determinant-ideal transformation law over \(\mathcal O_K\). For \(H^1\) one takes the integral \(\mathbf Z\)-dual of homology. Under the fixed trace pairing, taking this dual reverses the relative determinant ideal: locally a homology change matrix \(T\) changes its dual lattice by \((T^{-1})^{\mathsf t}\). The trace-dual factor of \(\mathcal O_K\) is common to source and target and again cancels. Since integral cohomology of an abelian variety is the exterior algebra on \(H^1(A,\mathbf Z)\), we obtain \[
 \mathscr W_A=\mathfrak d^{-1}\mathscr W_0.
\] This also proves that these actual integral intersection lattices are stable under the determinant \(\mathcal O_K\) action, a fact stronger than stability merely under the geometric scalars \(a^g\).

\subsubsection{2. Relative ideal norms and the action of a pushforward}\label{relative-ideal-norms-and-the-action-of-a-pushforward}

The B19 Hom ideal identity gives \(II_0=\Lambda^-\). Taking determinant modules over \(\mathcal O_K\) yields \[
 N_{E/K}(I)=\frac{\det_{\mathcal O_K}(II_0)}
                    {\det_{\mathcal O_K}I_0}
 =\mathfrak j\mathfrak d.
\tag{B23.9}
\] Here the determinant norm formula can be checked prime by prime: an \(\mathcal O_E\) prime \(\mathfrak P\) above \(\mathfrak p\) changes the determinant valuation over \(\mathcal O_K\) by its residue degree \(f(\mathfrak P/\mathfrak p)\). This is exactly the standard definition of the relative ideal norm. Similarly, local elementary divisors show that the numerical norm of \(\mathfrak j\) is the additive index \([\Lambda:\Lambda^-]=J\).

Let \[
 \nu=N_{K/\mathbf Q}(\mathfrak d)>0.
\] It is the ratio of the real covolumes of \(\Lambda\) and \(I_0\). For \(0\ne\alpha\in I\), let \(h_\alpha\) be the corresponding isogeny and write \(s_\alpha=N_{E/K}(\alpha)\). On rational Weil cohomology, \[
 h_\alpha^*=s_\alpha,\qquad
 \deg h_\alpha=\nu^{-1}s_\alpha\overline{s_\alpha}.
\] The identity \(h_{\alpha*}h_\alpha^*=(\deg h_\alpha)\mathrm{id}\) therefore gives \[
 h_{\alpha*}=\nu^{-1}\overline{s_\alpha}.
\tag{B23.10}
\] These are middle-degree cohomological pushforwards with their actual multiplicities.

The \(\mathcal O_K\)-ideal generated by \(N_{E/K}(\alpha)\) for \(\alpha\in I\) is exactly \(N_{E/K}(I)\). One inclusion follows from \((\alpha)\subseteq I\). For the reverse inclusion, fix a prime \(\mathfrak p\) of \(K\). By the Chinese remainder theorem for the invertible ideal \(I\), choose an element \(\alpha\in I\) attaining its minimal ideal valuation at every prime \(\mathfrak P\) of \(E\) above \(\mathfrak p\). Its norm attains precisely the valuation of \(N_{E/K}(I)\) at \(\mathfrak p\). This proves equality after localizing at every prime, hence equality of the ideals. Fractional ideals reduce to integral ideals by multiplication by a fixed nonzero denominator.

Each image \(h_*\mathscr W_0\) is an \(\mathcal O_K\)-module. Using (B23.9)--(B23.10), their additive sum is \[
 \begin{aligned}
 \sum_hh_*\mathscr W_0
 &=\nu^{-1}\overline{N_{E/K}(I)}\,\mathscr W_0\\
 &=\nu^{-1}\overline{\mathfrak j}\,\overline{\mathfrak d}\,\mathscr W_0\\
 &=\overline{\mathfrak j}\,\mathfrak d^{-1}\mathscr W_0
 =\overline{\mathfrak j}\,\mathscr W_A.
 \end{aligned}
\tag{B23.11}
\] The third equality uses \(\mathfrak d\overline{\mathfrak d}=(\nu)\) in the imaginary quadratic field. It proves all of (B23.3).

There is also a geometric factorization explaining the common content. Every \(h\) factors through \[
 A^-=\mathbf C^g/\Phi(\Lambda^-)
 \xrightarrow{q}A,
 \qquad\deg q=J.
\] The first part of the proof gives \(\mathscr W_{A^-}=\mathfrak j^{-1}\mathscr W_A\). In the common rational cohomology coordinates, \(q^*\) is the identity and \(q_*=J\). Therefore \[
 q_*\mathscr W_{A^-}=J\mathfrak j^{-1}\mathscr W_A
 =\overline{\mathfrak j}\,\mathscr W_A.
\tag{B23.12}
\] Thus the common sublattice arises before any choice of short isogeny, norm estimate, or theta input.

\subsubsection{3. A fixed source vector and its additive span}\label{a-fixed-source-vector-and-its-additive-span}

The same calculation with the principal rank-one submodule \(\mathcal O_Kw_0\) gives the first formula in (B23.5), of index \(Jc_0\) in \(\mathscr W_A\).

For every \(\alpha\in I\) one also has \(a\alpha\in I\), and (B23.10) gives \[
 h_{a\alpha*}w_0=\overline a^{\,g}h_{\alpha*}w_0.
\] Consequently \(\Gamma_{w_0}\) is a module over \(\mathbf Z[\overline a^{\,g}]\). The latter order has index \(C\) in \(\mathcal O_K\), so \[
 C\mathcal O_K\subseteq\mathbf Z[\overline a^{\,g}].
\] It follows that \(C\mathcal O_K\Gamma_{w_0}\subseteq\Gamma_{w_0}\). The other inclusion is tautological. Since all these lattices have rank two over \(\mathbf Z\), multiplication by \(C\) has index \(C^2\). This proves the two-sided index bound in (B23.5).

\subsubsection{4. Projection of integral cycles and the division bound}\label{projection-of-integral-cycles-and-the-division-bound}

The homomorphisms \(h\) are \(K\)-linear. They commute on cohomology with the fixed polynomial projector operators, including pushforwards: \[
 P_Ah_*=h_*P_0.
\] For each integral source cycle \(Z\), the class \(P_0[Z]\) belongs to \(\mathscr W_0\). Formula (B23.3) therefore proves (B23.6) for every integral linear combination, regardless of the degrees or number of source cycles used.

If \(\xi=Nt\) lies in the Weil plane, then \(P_A\xi=RNt\). Its membership in \(\overline{\mathfrak j}\mathscr W_A\) implies the lattice inclusion \[
 \mathcal O_K(RNt)\subseteq\overline{\mathfrak j}\mathscr W_A.
\] Taking additive indices in the rank-two lattice \(\mathscr W_A\) gives \[
 [\mathscr W_A:\mathcal O_K(RNt)]
 =R^2N^2c_t
 =J[\overline{\mathfrak j}\mathscr W_A:\mathcal O_K(RNt)].
\] This proves both assertions of (B23.7). The determinant \(K\) action, integral Weil intersection lattice, and a fixed integral section \(t\) are flat in the indicated generic Weil family. Their finite index \(c_t\) is therefore locally constant and hence constant on a connected base after fixing the local system. QED.

\subsection{More precise ideal divisibility}\label{more-precise-ideal-divisibility}

The necessary projected-class condition is more precise than its numerical norm. Define \[
 \mathfrak b_t=\{b\in\mathcal O_K:bt\in\overline{\mathfrak j}\mathscr W_A\},
 \qquad
 \mathfrak b_t\cap\mathbf Z=m_t\mathbf Z,
\tag{B23.13}
\] where \(m_t\) is positive. Then \(RNt\in\overline{\mathfrak j}\mathscr W_A\) is equivalent to \[
 \frac{m_t}{\gcd(m_t,R)}\mid N.
\tag{B23.14}
\] This is an exact test for membership in the common projected lattice. It is not a sufficiency assertion for algebraicity, the full cohomology class, or a prescribed effective cycle.

For example, if \(t\) generates \(\mathscr W_A\) over \(\mathcal O_K\), then \(\mathfrak b_t=\overline{\mathfrak j}\). If \(\mathfrak j=(m)\), the projected membership condition requires \(m/\gcd(m,R)\mid N\), and \(J=m^2\). If instead \(\mathfrak j=\mathfrak p^r\) for one of the primes above a split rational prime \(p\), then \(J=p^r\) and \(m_t=p^r\). For \(p\nmid R\), the condition requires \(p^r\mid N\). These are conditional computations for the displayed determinant ideals, not claims that every such ideal occurs in a prescribed polarized moduli family.

\subsection{A fixed nonmaximal quadratic order}\label{a-fixed-nonmaximal-quadratic-order}

The exact formulas above assume integral \(\mathcal O_K\) action. If the fixed generic action is only by an order \(\mathcal O\subset\mathcal O_K\), choose a fixed positive integer \(f\) with \(f\mathcal O_K\subset\mathcal O\). Set \[
 \Lambda^+=\mathcal O_K\Lambda,
 \qquad A^+=\mathbf C^g/\Phi(\Lambda^+),
 \qquad u:A\to A^+.
\] Then \(f\Lambda^+\subset\Lambda\), so \(q=\deg u=[\Lambda^+:\Lambda]\le f^{2g}\). Let \((\Lambda^+)^-\) denote its maximal \(\mathcal O_E\) core and put \(J^+=[\Lambda^+:(\Lambda^+)^-]\). There are inclusions \[
 f(\Lambda^+)^-\subseteq\Lambda^-\subseteq(\Lambda^+)^-.
\] If \(s=[(\Lambda^+)^-:\Lambda^-]\), then \(s\le f^{2g}\) and \[
 qJ=sJ^+,
 \qquad J^+\ge J/f^{2g}.
\tag{B23.15}
\]

Apply the proved maximal-order theorem to every composite \(uh\). If an integral combination of the original pushes has class \(Nt\), its image under \(u_*\) has class \(N t^+\), where \(t^+=u_*t\ne0\) is integral. On a family with fixed integral \(\mathcal O\) action and fixed flat \(t\), the saturation by \(\mathcal O_K\), the finite isogeny \(u\), the induced Weil lattices, and \(t^+\) are all determined by the same fixed rational local system. In particular \[
 c_{t^+}=[\mathscr W_{A^+}:\mathcal O_Kt^+]
\] is fixed. The projector commutes with \(u\) and keeps denominator \(R\). Thus \[
 N\ge\frac{\sqrt{J^+}}{R\sqrt{c_{t^+}}}
 \ge\frac{\sqrt J}{R f^g\sqrt{c_{t^+}}}.
\tag{B23.16}
\] This supplies a uniform comparison for a fixed quadratic order. It does not silently replace a varying quadratic order by its maximal order at no cost.

\subsection{Reduced images from an arbitrary bounded source family}\label{reduced-images-from-an-arbitrary-bounded-source-family}

Assume now that \(A_0\) is simple, as in B21--B22. Fix positive integers \(D_0\) and \(k\), with \(L_0^{\otimes k}\) very ample, and put \[
 \begin{aligned}
 S&=k^{n(n+1)}D_0^{n+1}\in\mathbf Z_{>0},\\
 L_S&=\operatorname{lcm}(1,2,\ldots,S).
 \end{aligned}
\] The sharper intrinsic result B22a permits the smaller choice \[
 S=\max\left\{1,\left\lfloor
 \frac{D_0^2}{4(n!)^2\chi(L_0)}\right\rfloor\right\}.
\] Every assertion below holds with either choice, or with any other positive integer upper bound for the stabilizer orders in this source family. Let \(\xi\) be any integral linear combination of reduced-image classes \[
 [h(Z)_{\mathrm{red}}],
 \qquad
 Z\subset A_0\text{ integral of dimension }n,
 \qquad \deg_{L_0}Z\le D_0,
\] with arbitrary integral \(E\)-linear isogenies \(h:A_0\to A\). There is no bound on the number of these summands or their integer coefficients, and the source subvarieties may vary throughout the entire bounded-degree family. If \(\xi=Nt\) for a nonzero integral target Weil class \(t\), then in the maximal-order setting \[
 J\mid R^2L_S^2N^2c_t.
\tag{B23.17}
\] Consequently the two order conventions give respectively \[
 \begin{aligned}
 N&\ge\frac{\sqrt J}{R L_S\sqrt{c_t}}
 &&\text{for integral }\mathcal O_K\text{ action},\\
 N&\ge\frac{\sqrt J}{R L_S f^g\sqrt{c_{t^+}}}
 &&\text{for the fixed order }\mathcal O.
 \end{aligned}
\tag{B23.18}
\]

To prove this, B20 and B22 give, for each summand, \[
 h_*Z=\delta_hh(Z)_{\mathrm{red}},\qquad
 \delta_h=|\ker h\cap\operatorname{Stab}(Z)|\le S.
\] Thus \(\delta_h\) divides \(L_S\). Multiplying the entire proposed relation by \(L_S\) expresses \(L_S\xi=L_SN t\) as an integral linear combination of multiplicity-preserving source pushes. Applying (B23.6)--(B23.7) proves (B23.17) and the first bound in (B23.18). For the fixed nonmaximal order, apply the same argument after pushing this multiplied relation through \(u:A\to A^+\). It gives \(J^+\mid R^2L_S^2N^2c_{t^+}\), and (B23.15) proves the second bound. No new multiplicity denominator is introduced by this final pushforward.

The sharper condition is \(m_t\mid RL_SN\) in the notation (B23.13). In particular, for a rational prime \(p>S\) one has \(v_p(L_S)=0\), so division by generic mapping multiplicities removes none of the \(p\)-primary condition \(v_p(N)\ge\max(0,v_p(m_t)-v_p(R))\). Integral greatest-common-divisor extraction across arbitrarily many reduced images therefore still cannot produce a fixed flat target class with bounded \(N\) along unbounded \(J\), provided the source degrees remain bounded. All constants refer to the fixed source polarization, quadratic order, integral local system and class. This assertion still imposes no intrinsic lower bound on cycles obtained by a different construction.

\subsection{Scope and the remaining route}\label{scope-and-the-remaining-route}

Taking greatest common divisors across all \(E\)-linear isogenies removes individual principal relative norm factors only down to the relative norm ideal of the full Hom lattice. Its remaining target content is the conjugate determinant ideal \(\overline{\mathfrak j}\), of norm \(J\). The same obstruction applies to integral sums described as traces if every summand is one of the indicated multiplicity-preserving reference transports. Operations leaving this class of summands require a separate argument.

Replacing \(h_*Z=\delta_h h(Z)_{\mathrm{red}}\) by the reduced image divides its class by \(\delta_h\). Such a division can leave the sublattice in (B23.6); this is precisely why B20--B22 separately control stabilizers and generic multiplicities. The present result imposes no new obstruction on arbitrary target cycles or new deformation constructions. It is also distinct from a single full-CM tensor coefficient in the product construction of 032: applying it there requires identifying an appropriate generic rank-two \(K\)-Weil summand and its integral lattice.

\subsubsection{Exact classical inputs and reading scope}\label{exact-classical-inputs-and-reading-scope}

\begin{itemize}
\tightlist
\item
  J.S. Milne, \emph{Algebraic Number Theory}, version 3.08, 19 July 2020, section 4, printed pp.68--70, definition of the relative ideal norm and Propositions 4.1--4.2. Official source \url{https://www.jmilne.org/math/CourseNotes/ANT.pdf}. The definition, compatibility with principal norms, ideal-index formula, and the complete adjoining short proofs were read. These ideal-theoretic facts are imported; the determinant-lattice and transport applications above are proved here.
\item
  Integral cohomology of a complex abelian variety as the exterior algebra on its integral first cohomology, and finite-isogeny pull-push, are the same inputs already used and checked in B7, B18 and B20. The local maximal-order lattice argument above explicitly handles ramified and split primes; no unrecorded saturation denominator is assumed.
\item
  B6 supplies the common integral projector; B19 supplies \(II_0=\Lambda^-\) and the identification of integral \(E\)-linear Hom with the ideal \(I\). Simplicity is not used in these algebraic identities, as also checked in T6.17.
\end{itemize}

This proof has been read in full, including the determinant action versus geometric scalar distinction, the actual integral intersection lattice, the relative norm ideal generated by all element norms, and the fixed-order comparison. No proof-assistant verification is claimed.

\section{The saturated stabilizer quotient and residual isogeny}

\providecommand{\tightlist}{\setlength{\itemsep}{0pt}\setlength{\parskip}{0pt}}
\section{B25: quotient saturation and the precise stabilizer escape}\label{b25-quotient-saturation-and-the-precise-stabilizer-escape}

Status: author and root independently read the complete proof and passed it on 2026-11-21. Run DEG\_L\_20261121; author /root/lvalue\_bounds; GPT-6 Astra Pro; effort not exposed. This addresses unbounded source witnesses and the effect of passing from isogeny pushforwards to reduced images. The exact maximal-order scope and the distinction from intrinsic cycle complexity remain explicit.

\subsection{B25. A bounded reduced image forces a target-sized intermediate CM core}\label{b25.-a-bounded-reduced-image-forces-a-target-sized-intermediate-cm-core}

Retain the middle-dimensional CM/Weil setting of B23: \(K\subset E\) is imaginary quadratic inside a CM field, \([E:K]=g=2n\), the fixed CM type has \(K\)-signature \((n,n)\), and the rational homology has rank one over \(E\). Let \[
 A_0=\mathbf C^g/\Phi(I_0),\qquad
 A=\mathbf C^g/\Phi(\Lambda_A),
\] where \(I_0\) is an invertible fractional \(\mathcal O_E\)-ideal and \(\Lambda_A\) is \(\mathcal O_K\)-stable. Let \(L\) be a polarization on \(A\) compatible with conjugation on \(K\), and write \(V=\chi(L)\). Use the same integral Weil projector \(P=Re\) from B6 on all the \(\mathcal O_K\)-abelian varieties in question.

Let \(h:A_0\to A\) be any nonzero integral \(E\)-linear homomorphism, and let \(Z\subset A_0\) be an integral subvariety of dimension \(n\) with \(e_0[Z]\ne0\). There is no bound on its degree or stabilizer. Put \[
 H=\ker h\cap\operatorname{Stab}_{A_0}(Z),
 \qquad
 G=\mathcal O_K H\subseteq\ker h.
\tag{B25.1}
\] The second expression means the finite subgroup generated by the images of \(H\) under the integral \(\mathcal O_K\) action on \(A_0\). It is contained in \(\ker h\) because the target also has integral \(\mathcal O_K\) action. Form the quotient and factorization \[
 A_0\xrightarrow{q}B=A_0/G\xrightarrow{g}A,
 \qquad h=gq,
 \qquad Z_B=q(Z)_{\mathrm{red}},\quad Z_A=h(Z)_{\mathrm{red}}.
\tag{B25.2}
\] Then \(B\) has integral \(\mathcal O_K\) action, the same rational \(E\) action and CM type, and \[
 \begin{aligned}
 \deg(q|_Z)&=|H|,&g_*Z_B&=Z_A,\\
 \deg g&=[\ker h:G],&e_B[Z_B]&\ne0.
 \end{aligned}
\tag{B25.3}
\] In particular, the residual isogeny \(g\) has generic degree one on this intermediate image cycle.

For \(X=A,B\), define its maximal-order core and determinant ideal by \[
 \Lambda_X^-=\{x\in\Lambda_X:\mathcal O_E x\subset\Lambda_X\},
 \quad
 \mathfrak j_X=\frac{\det_{\mathcal O_K}\Lambda_X^-}
                     {\det_{\mathcal O_K}\Lambda_X},
 \quad J_X=N_{K/\mathbf Q}(\mathfrak j_X).
\] Here \(\Lambda_B\supseteq I_0\) is the lattice defining the quotient \(B\) in the common CM coordinates. If \(g\) is multiplication by \(\alpha\in E^\times\), set \[
 \mathfrak k_g=
 \frac{\det_{\mathcal O_K}(\alpha\Lambda_B)}
      {\det_{\mathcal O_K}\Lambda_A}\subseteq\mathcal O_K,
 \qquad N_{K/\mathbf Q}(\mathfrak k_g)=\deg g.
\tag{B25.4}
\] There are the exact lattice image and ideal containment \[
 \begin{aligned}
 g_*\mathscr W_B&=\overline{\mathfrak k_g}\,\mathscr W_A,\\
 \mathfrak k_g\mathfrak j_B&\subseteq\mathfrak j_A,\\
 P_A[Z_A]&\in\overline{\mathfrak k_g}\,\mathscr W_A
 \subseteq
 \overline{\bigl(\mathfrak j_A\mathfrak j_B^{-1}
                         \cap\mathcal O_K\bigr)}\,\mathscr W_A.
 \end{aligned}
\tag{B25.5}
\] Consequently \[
 (\deg g)J_B\ge J_A.
\tag{B25.6}
\]

Let \(D'=\deg_L Z_A\). Then \[
 \begin{aligned}
 D'^2&\ge \frac{2(n!)^2V}{R^2}\deg g,\\
 [\ker h:G]=\deg g
 &\le\frac{R^2D'^2}{2(n!)^2V},\\
 J_B&\ge\frac{2(n!)^2V}{R^2D'^2}J_A.
 \end{aligned}
\tag{B25.7}
\] Thus a bounded-degree reduced image must factor through an intermediate source whose CM core index is at least a fixed positive multiple of the target index, with a residual isogeny of bounded degree. This statement does not claim that such an intermediate witness can be constructed with bounded degree.

\subsubsection{The maximal-order stabilizer case}\label{the-maximal-order-stabilizer-case}

If \(G\) is \(\mathcal O_E\)-stable, then \(B\) retains integral \(\mathcal O_E\) action and \(J_B=1\). This includes the case when \(H\) itself is \(\mathcal O_E\)-stable. In this situation \[
 D'\ge\frac{\sqrt2\,n!\sqrt V}{R}\sqrt{J_A}.
\tag{B25.8}
\] Moreover \(P_A[Z_A]\in\overline{\mathfrak j_A}\mathscr W_A\). Hence all integral combinations of such reduced images satisfy the original B23 determinant-content obstruction, even if the source degrees, stabilizer orders, and number of transports are unbounded.

\subsection{Proof}\label{quotient_escape_proofs-proof}

\subsubsection{1. The saturated stabilizer quotient and generic degree one}\label{the-saturated-stabilizer-quotient-and-generic-degree-one}

The group \(H\) is finite because it is contained in the kernel of an isogeny. The integral \(\mathcal O_K\) action preserves \(\ker h\), so \(G\) is a finite subgroup of that kernel and the quotient \(q\) and residual isogeny \(g\) in (B25.2) exist. The group \(G\) is \(\mathcal O_K\)-stable, so \(B\) inherits this integral action. In CM coordinates the quotient replaces \(I_0\) by a commensurable full \(\mathcal O_K\)-stable lattice \(\Lambda_B\); the rational \(E\) action and the CM type remain the same.

For a finite isogeny \(u\) and an integral subvariety \(T\), the generic degree of \(u|_T\) is \(|\ker u\cap\operatorname{Stab}(T)|\), by the finite-kernel generic-fibre argument of B20. It applies even if the full stabilizer of \(T\) is positive-dimensional. Since \[
 H\subseteq G\subseteq\ker h,
 \qquad G\cap\operatorname{Stab}(Z)=H,
\] both \(h|_Z\) and \(q|_Z\) have generic degree \(|H|\). Degrees of the finite maps between the integral image varieties multiply, so \(g|_{Z_B}\) has generic degree one. Thus \(g_*Z_B=Z_A\) with no remaining generic multiplicity. This proves the degree assertions in (B25.3).

The two isogenies are \(K\)-linear, so their cohomological pushforwards commute with the Weil projectors. We have \(q_*[Z]=|H|[Z_B]\); pushforward is invertible on rational cohomology. Therefore \(e_0[Z]\ne0\) implies \(e_B[Z_B]\ne0\), and then \(e_A[Z_A]\ne0\) as well.

\subsubsection{2. Exact determinant content of the residual isogeny}\label{exact-determinant-content-of-the-residual-isogeny}

The integral Weil-lattice covariance proved in B23 applies to any two \(\mathcal O_K\)-stable homology lattices, whether or not either is \(\mathcal O_E\)-stable. Identify \(B\) with the abelian variety whose homology lattice is \(\alpha\Lambda_B\). The residual isogeny is then induced by the inclusion \[
 \alpha\Lambda_B\subseteq\Lambda_A.
\] The associated determinant ideal is \(\mathfrak k_g\). The inverse-determinant formula for integral Weil cohomology and pull-push by this inclusion give \[
 g_*\mathscr W_B
 =N(\mathfrak k_g)\mathfrak k_g^{-1}\mathscr W_A
 =\overline{\mathfrak k_g}\,\mathscr W_A.
\] This is the first assertion of (B25.5), including at split and ramified primes of \(K\).

Because multiplication by \(\alpha\) is \(E\)-linear, \[
 (\alpha\Lambda_B)^-=\alpha\Lambda_B^-.
\] This is an \(\mathcal O_E\)-stable sublattice of \(\Lambda_A\), hence it is contained in \(\Lambda_A^-\). Taking determinants of this lattice inclusion yields \[
 \mathfrak k_g\mathfrak j_B
 =\frac{\det(\alpha\Lambda_B^-)}{\det\Lambda_A}
 \subseteq\frac{\det\Lambda_A^-}{\det\Lambda_A}
 =\mathfrak j_A.
\] Taking numerical norms proves (B25.6). Combining this containment with the integral ideal condition \(\mathfrak k_g\subseteq\mathcal O_K\) gives \[
 \mathfrak k_g\subseteq
 \mathfrak j_A\mathfrak j_B^{-1}\cap\mathcal O_K.
\] Finally, the equality of image cycles and projector functoriality give \[
 P_A[Z_A]=g_*P_B[Z_B],\qquad P_B[Z_B]\in\mathscr W_B.
\] This proves every assertion of (B25.5).

\subsubsection{3. The degree of the residual isogeny}\label{the-degree-of-the-residual-isogeny}

Give \(B\) the pullback polarization \(g^*L\). Its Rosati involution restricts to conjugation on \(K\), and its Weil plane is primitive and positive for the signed middle Poincaré form used in B7. Let \[
 v_B=P_B[Z_B]\in\mathscr W_B\setminus\{0\},
 \qquad v_A=P_A[Z_A]=g_*v_B.
\] The integral even-square argument B18 gives \(\|v_B\|_{q_B}^2\ge2\). The exact finite-isogeny Poincaré scaling gives \[
 \|v_A\|_{q_A}^2=(\deg g)\|v_B\|_{q_B}^2\ge2\deg g.
\] On the other hand, the existing B7 comass estimate on the target says \[
 \|v_A\|_{q_A}=R\|e_A[Z_A]\|_{q_A}
 \le\frac{RD'}{n!\sqrt V}.
\] Combining these inequalities proves the first two lines of (B25.7), and then (B25.6) proves the third. This argument does not require the source \(B\) to have fixed polarization volume or fixed degree, and it does not modify the previously checked B7 constant.

If \(G\) is \(\mathcal O_E\)-stable, its inverse-image lattice \(\Lambda_B\) is an \(\mathcal O_E\)-ideal because \(I_0\) already is. Thus \(\Lambda_B^- =\Lambda_B\), \(\mathfrak j_B=\mathcal O_K\), and \(J_B=1\). Formula (B25.5) then gives the original target determinant content, and (B25.7) gives (B25.8). The same content inclusion is closed under arbitrary integral sums. QED.

\subsection{Prime-by-prime criterion for several reduced images}\label{prime-by-prime-criterion-for-several-reduced-images}

The exact containment in (B25.5) identifies which prime powers can be removed by a stabilizer quotient. At a prime \(\mathfrak p\) of \(K\), \[
 v_{\mathfrak p}(\mathfrak k_g)
 \ge\max\{v_{\mathfrak p}(\mathfrak j_A)
              -v_{\mathfrak p}(\mathfrak j_B),0\}.
\tag{B25.9}
\] The target projected class has the conjugate prime content. A large numerical stabilizer order alone therefore does not specify which target prime factors are removed; the determinant ideal of the intermediate core does.

For a finite integral combination of reduced-image classes with intermediate quotients \(B_1,\ldots,B_r\), set \[
 \mathfrak b=\bigcap_{i=1}^r\mathfrak j_{B_i},
 \qquad
 \mathfrak c=\overline{\bigl(\mathfrak j_A\mathfrak b^{-1}
                                      \cap\mathcal O_K\bigr)}.
\tag{B25.10}
\] Then the projected integral combination belongs to \(\mathfrak c\mathscr W_A\). Indeed the sum of the containing ideals \(\mathfrak j_A\mathfrak j_{B_i}^{-1}\cap\mathcal O_K\) from (B25.5) has valuation \[
 \min_i\max\{v_{\mathfrak p}(\mathfrak j_A)
                    -v_{\mathfrak p}(\mathfrak j_{B_i}),0\}
 =\max\{v_{\mathfrak p}(\mathfrak j_A)
                    -\max_i v_{\mathfrak p}(\mathfrak j_{B_i}),0\},
\] which is exactly the valuation of \(\mathfrak j_A\mathfrak b^{-1}\cap\mathcal O_K\). If the full combination represents \(Nt\) for a nonzero integral target Weil class, then \[
 N(\mathfrak c)\mid R^2N^2[\mathscr W_A:\mathcal O_Kt].
\tag{B25.11}
\] For an algebraic combination of bounded total absolute degree, each surviving integral component has bounded degree. Every component with nonzero Weil projection that is obtained by the specified transport therefore obeys (B25.7) individually. Without such a component-degree bound, different intermediate quotients may absorb different target primes; (B25.9)--(B25.11) are the appropriate exact necessary conditions. They are not existence assertions for cycles with prescribed prime cancellation.

\subsection{Relation to the earlier source-degree bounds}\label{relation-to-the-earlier-source-degree-bounds}

The subgroup construction also clarifies the exponent in B22a. Choose an integral basis \(\mathcal O_K=\mathbf Z+\mathbf Z\theta\). Then \[
 G=H+\theta H,\qquad |G|\le|H|^2.
\tag{B25.12}
\] The equality follows because this sum is already stable under \(\theta\), using the quadratic integral equation of \(\theta\). The cardinality inequality follows from the elementary formula for the order of a sum of finite subgroups. Together with \(\deg h\ge J_A\) and (B25.7), this recovers the need for \(|H|\) at least a constant times \(\sqrt{J_A}/D'\). In the simple source case, B22a bounds \(|H|\) by a constant times the square of the source degree, yielding its \(J_A^{1/4}\) requirement.

If \(H\) is already \(\mathcal O_K\)-stable, then \(G=H\) and the stronger condition \(|H|\ge 2(n!)^2VJ_A/(R^2D'^2)\) holds. Combining it with B22a gives, for a simple polarized source of volume \(V_0\) and source degree \(D\), \[
 DD'\ge\frac{2\sqrt2\,(n!)^2\sqrt{V_0V}}{R}\sqrt{J_A}.
\tag{B25.13}
\] Thus bounded target degree then requires source degree of order at least \(J_A^{1/2}\). If the saturated subgroup is \(\mathcal O_E\)-stable, the stronger no-escape statement (B25.8) applies instead.

These deductions concern the specific isogeny construction. The intermediate variety \(B\) may depend on the target and on the source cycle. Saying that it must have target-sized CM core does not prove that all target-dependent constructions fail; it locates exactly where that arithmetic complexity has to reappear. It also avoids a circular construction: merely taking a full inverse image of an already known target witness produces such an intermediate witness, but does not construct the target witness in the first place.

\subsection{Inputs and check status}\label{inputs-and-check-status}

The proof uses the exact integral determinant-lattice functor of B23, the finite-kernel generic-degree identity B20, the existing target comass estimate B7, the integral even-square lower bound B18, and B22a only in the last optional source-degree corollary. The ideal operations and numerical indices use Milne, \emph{Algebraic Number Theory}, version 3.08, Propositions 4.1--4.2, printed pp.69--70, \url{https://www.jmilne.org/math/CourseNotes/ANT.pdf}, already checked in B23. No new unverified analytic or automorphic theorem enters this quotient argument.

The maximal quadratic order acts integrally on the target in this exact statement. A target carrying only a fixed nonmaximal order can first be passed through the bounded-degree maximal-order saturation of B23, but translating its polarization degree requires retaining that extra bounded isogeny. That extension is not silently included in the displayed constants above.

\section{A sharp nonsimple stabilizer-quotient example}

\providecommand{\tightlist}{\setlength{\itemsep}{0pt}\setlength{\parskip}{0pt}}
\section{B25a: an exact fourfold realization of the quotient escape}\label{b25a-an-exact-fourfold-realization-of-the-quotient-escape}

Status: author and root independently read the complete proof and passed it on 2026-11-21. Author: /root/lvalue\_bounds, GPT-6 Astra Pro; effort not exposed. Root proposed the fourfold refinement of the author's surface example. This example has a full rank-one CM field and a fixed integral quadratic action, but its CM type is imprimitive and its abelian varieties are not simple.

\subsection{Statement}\label{statement}

Let \[
 F=\mathbf Q(\zeta_3),\qquad K=\mathbf Q(i),\qquad
 E=\mathbf Q(\zeta_{24}),\qquad t=\zeta_{24},\quad \omega=\zeta_3.
\] Use the CM type of \(E\) induced from the fixed complex embedding of \(F\) sending \(\omega\) to \(e^{2\pi i/3}\). Put \(\mathcal O_m=\mathbf Z+m\mathcal O_F\) and identify the rational homology with the one-dimensional \(E\)-space using the \(F\)-basis \[
 (e_1,e_2,e_3,e_4)=(1,i,t,it).
\] Then \[
 A_0=\mathbf C^4/\mathcal O_F^4,\qquad
 A_m=\mathbf C^4/\mathcal O_m^4\quad(m\ge1)
\] are abelian fourfolds with product principal polarizations \(L_0,L_m\), rational \(E\) action with Rosati complex conjugation, and integral \(\mathcal O_K\) action of signature \((2,2)\). Their source homology lattice is \(I_0=\mathcal O_E\) and their target lattice is \[
 \Lambda_m=\mathcal O_m e_1\oplus\mathcal O_m e_2
              \oplus\mathcal O_m e_3\oplus\mathcal O_m e_4.
\] Its maximal-order core and hull, numerical core index, and relative \(K\) determinant ideal are \[
 \Lambda_m^-=m\mathcal O_E,\qquad \Lambda_m^+=\mathcal O_E,
 \qquad J_m=m^4,
 \qquad \mathfrak j_m=(m^2)\mathcal O_K.
\tag{B25a.1}
\]

Multiplication by \(m\) gives an \(E\)-linear polarized isogeny \[
 h_m:A_0\longrightarrow A_m,\qquad
 \deg h_m=m^4,
 \qquad h_m^*c_1(L_m)=m c_1(L_0).
\tag{B25a.2}
\] Let \(Z_0\) be the coordinate plane with coordinates 2 and 4 equal to zero, and let \(Z_m=h_m(Z_0)_{\rm red}\). These are integral surfaces with nonzero \(K\)-Weil projection and \[
 \deg_{L_0}Z_0=\deg_{L_m}Z_m=2,
 \qquad (h_m)_*Z_0=m^2 Z_m.
\tag{B25a.3}
\] For the intermediate quotient of B25, \[
 H_m=\ker h_m\cap\operatorname{Stab}(Z_0),\qquad
 G_m=\mathcal O_K H_m,
\] one has \[
 |H_m|=m^2,\qquad
 G_m=\ker h_m,\qquad |G_m|=m^4,
 \qquad B_m=A_0/G_m\simeq A_m.
\tag{B25a.4}
\] Thus the residual isogeny is an isomorphism and the intermediate core index equals the entire target index. For \(m>1\), \(G_m\) is not \(\mathcal O_E\)-stable. The scalar multiplicity \(m^2\) in (B25a.3) removes exactly the determinant-ideal content \((m^2)\) of the multiplicity-preserving pushforward. This supplies a geometric realization of the escape allowed by B25, without asserting any such escape on simple CM fourfolds.

\subsection{Proof}\label{quotient_escape_example_proofs-proof}

\subsubsection{1. The maximal order and the induced CM coordinates}\label{the-maximal-order-and-the-induced-cm-coordinates}

The cyclotomic integral-basis theorem gives \(\mathcal O_E=\mathbf Z[t]\) and \(\mathcal O_F=\mathbf Z[\omega]\). Here \(i=t^6\), \(\omega=t^8\), and \[
 t^2=-i\omega.
\] Consequently \[
 \mathbf Z[t]=\mathcal O_F[i]+t\mathcal O_F[i].
\] The two summands are independent over \(F(i)\) since \([E:F(i)]=2\). The summands \(\mathcal O_F+i\mathcal O_F\) are independent over \(F\) since \([F(i):F]=2\). Hence \[
 \mathcal O_E=\mathcal O_F\oplus i\mathcal O_F
               \oplus t\mathcal O_F\oplus it\mathcal O_F
\tag{B25a.5}
\] with the displayed basis. The degree assertions follow either from the cyclotomic degrees \(\varphi(24)=8\), \(\varphi(12)=4\), and \(\varphi(3)=2\), or from the distinct imaginary quadratic fields \(F,K\) and the quadratic relation for \(t\).

In the CM type induced from \(F\), the natural representation of \(E\) on \(\mathbf C^4\) is the regular \(F\)-linear representation in this basis, after the fixed embedding of \(F\) into \(\mathbf C\). It has precisely the four embeddings of \(E\) extending that embedding of \(F\). Thus it is the stated CM type, and the rational homology is rank one over \(E\).

Multiplication by \(i\) has the integral matrix \[
 J=\operatorname{diag}\left(
 \begin{pmatrix}0&-1\\1&0\end{pmatrix},
 \begin{pmatrix}0&-1\\1&0\end{pmatrix}\right).
\tag{B25a.6}
\] Its analytic eigenvalues are \(i,-i\), each twice. Multiplication by \(t\) maps \[
 e_1\mapsto e_3,\quad e_2\mapsto e_4,\quad
 e_3\mapsto-\omega e_2,\quad e_4\mapsto\omega e_1.
\tag{B25a.7}
\] This matrix is monomial and unitary for the product Hermitian form because \(|\omega|=1\). Therefore its Rosati adjoint is multiplication by \(t^{-1}=\bar t\). Since \(t\) generates \(E\), the full rational \(E\) action has Rosati involution equal to CM conjugation. The same assertion holds on every \(A_m\): the product polarization metric changes by a common positive scalar in the four coordinates, so the adjoint matrices are unchanged. The matrix \(J\) preserves \(\mathcal O_m^4\) for every \(m\), giving the fixed integral maximal quadratic-order action.

Each factor \(\mathbf C/\mathcal O_m\) is an elliptic curve with its degree-one polarization; the product gives the asserted principally polarized fourfold. Every multiplication matrix for \(E\) preserves a suitable integer multiple of the homology lattice, so its action remains rational and algebraic. No claim of integral \(\mathcal O_E\) action on \(A_m\) is made when \(m>1\).

\subsubsection{2. Core, hull, and determinant ideal}\label{core-hull-and-determinant-ideal}

The largest \(\mathcal O_F\)-stable sublattice in \(\mathcal O_m\) is \(m\mathcal O_F\). Indeed write \(x=a+mb\omega\) with \(a,b\in\mathbf Z\). Since \(\omega^2=-1-\omega\), \[
 \omega x=-mb+(a-mb)\omega.
\] The condition \(\omega x\in\mathcal O_m\) forces \(m\mid a\), and conversely every element of \(m\mathcal O_F\) remains in \(\mathcal O_m\) under \(\mathcal O_F\). As \(\mathcal O_E\) contains \(\mathcal O_F\), the maximal-order core of \(\Lambda_m\) is contained in \(m\mathcal O_F^4=m\mathcal O_E\). The latter is itself an \(\mathcal O_E\)-ideal contained in \(\Lambda_m\), so equality follows. The hull is \(\mathcal O_E\) since \(\Lambda_m\subset\mathcal O_E\) contains 1. Each coordinate contributes index \([\mathcal O_m:m\mathcal O_F]=m\), proving \(J_m=m^4\).

For the relative determinant, use the \(\mathcal O_K\)-basis \((1,\omega,t,\omega t)\) of \(\mathcal O_E\). It gives \[
 \Lambda_m=\mathcal O_K\oplus m\omega\mathcal O_K
                    \oplus t\mathcal O_K\oplus m\omega t\mathcal O_K.
\] Thus \(\det_K\Lambda_m=(m^2)\det_K\mathcal O_E\) while \(\det_K(m\mathcal O_E)=(m^4)\det_K\mathcal O_E\). Their quotient is \((m^2)\mathcal O_K\), as in (B25a.1), whose norm is \(m^4\).

\subsubsection{3. Isogeny, witness, and quotient}\label{isogeny-witness-and-quotient}

Multiplication by \(m\) maps \(\mathcal O_E\) into \(\Lambda_m\). Its degree is \([\Lambda_m:m\mathcal O_E]=m^4\). On each elliptic factor the target fundamental area is \(m\) times that of the source. A degree-one polarization has invariant two-form inversely proportional to this area, and multiplication by \(m\) multiplies that two-form by \(m^2\). Hence pullback multiplies the polarization by \(m\), proving (B25a.2).

The restriction to the selected two coordinates is the product of two degree-\(m\) elliptic isogenies. Its generic degree is \(m^2\), and its image is exactly the corresponding coordinate plane in \(A_m\). The product polarization on two elliptic factors has square 2, proving (B25a.3).

One of the two Weil lines, namely \(\bigwedge^4 H^1(A_m,\mathbf C)_{J^*=i}\), contains \[
 \eta=(dz_1+i\,dz_2)\wedge(dz_3+i\,dz_4)
       \wedge(d\bar z_1+i\,d\bar z_2)
       \wedge(d\bar z_3+i\,d\bar z_4).
\] Its restriction to \(Z_m\) is \(dz_1\wedge dz_3\wedge d\bar z_1\wedge d\bar z_3\), with nonzero integral. The self-adjoint Weil projector therefore detects a nonzero Weil component of \([Z_m]\). The same argument applies to \(Z_0\).

The kernel in a single elliptic coordinate is \[
 m^{-1}\mathcal O_m/\mathcal O_F
       =(m^{-1}\mathbf Z+\mathcal O_F)/\mathcal O_F\simeq\mathbf Z/m\mathbf Z.
\] The stabilizer of the coordinate plane \(Z_0\) is \(Z_0\) itself. Its kernel intersection \(H_m\) consists of the first and third kernel coordinates and has order \(m^2\). Applying \(J\) moves these to the second and fourth coordinates, so \[
 H_m\cap iH_m=0,\qquad H_m+iH_m=\ker h_m.
\] This proves (B25a.4) and the residual-isomorphism assertion. For \(m>1\), the class of \(m^{-1}e_3\) belongs to \(G_m\), but multiplication by \(t\) sends it to \(-m^{-1}\omega e_2\). Its image under \(h_m\) is \(-\omega e_2\), which is not in \(\Lambda_m\); hence it is not in \(G_m\). Thus \(G_m\) is not \(\mathcal O_E\)-stable.

Finally \(\mathfrak j_m=(m^2)\) and \(h_{m*}Z_0=m^2Z_m\) prove the exact cancellation of that scalar ideal content. The displayed bounded target witnesses are abelian subvarieties, and their stabilizers have positive dimension. The simple-source hypotheses of B20--B22a therefore do not apply. QED.

\subsection{Primary inputs}\label{primary-inputs}

J.S. Milne, \emph{Algebraic Number Theory}, version 3.08, Theorem 6.4(a)--(b), printed p.99, gives the degree and maximal order of a cyclotomic field. Its full displayed proof on p.99 and the complete supporting Lemma 6.5 proof on p.100 were read in this run; lower prerequisites are imported. Source: \url{https://www.jmilne.org/math/CourseNotes/ANT.pdf}. The identity (B25a.5) is proved explicitly above from that theorem and the powers of \(t\).

J.S. Milne, \emph{Abelian Varieties}, version 2.00, Chapter I, Theorems 2.8--2.9, printed p.14, and polarization definition printed p.15, supply the usual algebraicity of polarized complex tori and of their analytic homomorphisms. These exact statements and the complete brief proof sketch were read in this run; no full Riemann-criterion proof is claimed. Source: \url{https://www.jmilne.org/math/CourseNotes/AV.pdf}.

\section{Sharp Weil calibration in the present normalization}

\providecommand{\tightlist}{\setlength{\itemsep}{0pt}\setlength{\parskip}{0pt}}
\section{B26: exact polarization normalization of the Weil calibration bound}\label{b26-exact-polarization-normalization-of-the-weil-calibration-bound}

Status: author /root and independent reviewer /root/lvalue\_bounds read the entire proof and passed it on 2026-11-21. Run DEG\_L\_20261121; author /root, GPT-6 Astra Pro, effort not exposed. The principal-angle inequality and its equality rigidity were already proved in rStation paper4/v1, Theorem1(b) and Theorem3(a),(c), by gabrielov\_cell, 16 October2026. The finding below is their exact conversion to the signed integral Poincare normalization used in B7, B18 and B25, together with its integral strictness margin. No new calibration theorem is claimed.

\subsection{B26. Sharp Weil-projection mass bound and integral margin}\label{b26.-sharp-weil-projection-mass-bound-and-integral-margin}

Let \((A,L,K)\) be a polarized abelian variety of Weil type, of complex dimension \(2n\), with \(n\ge2\), imaginary quadratic action \(K\) of signature \((n,n)\), and polarization Rosati involution equal to conjugation on \(K\). Put \[
 \theta=c_1(L),\qquad V=\int_A\frac{\theta^{2n}}{(2n)!}=\chi(L).
\] Let \(W\subset H^{2n}(A,\mathbf R)\) be the real generic Weil plane, and let \(e\) be its orthogonal projection, equivalently the extreme-weight projector of B6. On \(W\) put \[
 q(u,u)=(-1)^n\int_Au^2.
\] For every effective integral \(n\)-cycle \(Z\) of degree \(D=\int_Z\theta^n\), one has \[
 \boxed{\quad
 \|e[Z]\|_q\le
 \frac{D}{2^{n-1/2}n!\sqrt V}.
 \quad}
 \tag{B26.1}
\] The constant is optimal across polarized Weil varieties. If \(e[Z]\ne0\), equality implies that every integral component of \(Z\) is a translate of an \(n\)-dimensional abelian subvariety; its tangent plane is a unitary graph for the two \(K\) eigenspaces, with the same maximizing phase. Consequently (B26.1) is strict when \(A\) is simple and \(e[Z]\ne0\).

Assume the fixed projector denominator is \(R>0\), so \(P=Re\) preserves integral cohomology, and put \(v=P[Z]\). If \(A\) is simple and \(v\ne0\), then \[
 R^2D^2-2^{2n-1}(n!)^2Vq(v,v)
 \ \ge\ \gcd\bigl(R^2,2^{2n}(n!)^2V\bigr)>0.
 \tag{B26.2}
\] In particular every nonzero Weil witness satisfies \[
 D\ge\frac{2^n n!\sqrt V}{R},
 \tag{B26.3}
\] strictly on a simple variety. This is a lower bound on a witness, not an existence theorem or an upper degree bound.

\subsection{Proof of the normalization}\label{proof-of-the-normalization}

Use the invariant metric of the positive form representing \(\theta\). Its real volume is \(V\). Polarization compatibility makes the two holomorphic \(K\) eigenspaces orthogonal. Choose local linear coordinates \(z_1,\ldots,z_n,w_1,\ldots,w_n\) so that \[
 \omega=\frac i2\sum_{j=1}^n
 (dz_j\wedge d\bar z_j+dw_j\wedge d\bar w_j).
\] The Weil plane consists of the real forms generated by \[
 \Psi=dz_1\wedge\cdots\wedge dz_n
       \wedge d\bar w_1\wedge\cdots\wedge d\bar w_n,
 \qquad \bar\Psi.
\] These forms are primitive of Hodge type \((n,n)\): wedging with any summand of \(\omega\) repeats a factor. Thus the middle Hodge star acts on the real Weil plane as \((-1)^n\). The extreme-weight projector is orthogonal for this metric, because the normalized \(K\) action is unitary; equivalently the orthogonal eigenspace decomposition gives exactly the projector of B6. Invariant harmonic representatives therefore identify \(q\) with the \(L^2\) metric on \(W\).

The complex one-forms \(dz_j,d\bar w_j\) are mutually orthogonal and have squared norm2. Consequently \(|\Psi|^2=2^{2n}\), and \(\Psi\) is orthogonal to \(\bar\Psi\). If \[
 u=\operatorname{Re}(a\Psi),\qquad a\in\mathbf C,
\] then \[
 |u|^2=2^{2n-1}|a|^2,
 \qquad \|u\|_{L^2}=2^{n-1/2}|a|\sqrt V.
 \tag{B26.4}
\] The principal-angle bound of the cited Theorem1(b), with the differential-form conversion in its Theorem3(c), says that the absolute value of \(\Psi\) on a complex \(n\)-plane, measured against its induced Kähler volume, is at most1. This statement is independent of the particular imaginary quadratic field: its proof uses only the displayed orthogonal splitting. For clarity about its normalization, on a graph \(w=Bz\) the quotient is \[
 \frac{|\Psi|_{\Gamma_B}|}{d\operatorname{vol}_{\Gamma_B}}
 =\frac{2^n|\det B|}{\det(I+B^*B)}\le1.
 \tag{B26.5}
\] This formula also gives zero on a nongraph plane with singular first projection. Equality holds precisely when \(B\) is unitary. Varying its determinant phase shows that for every real Weil form \(u\), the exact comass on positively oriented complex \(n\)-planes is \[
 \sup_T|u(T)|=|a|=
 \frac{\|u\|_q}{2^{n-1/2}\sqrt V}.
 \tag{B26.6}
\] The absolute value of the integral of a smooth form over an effective cycle is bounded by this comass times its mass. Singular loci have measure zero, and Wirtinger's formula gives mass \(D/n!\). Set \(u=e[Z]\). Orthogonality and the middle star identity imply \[
 \left|\int_Zu\right|
 =\left|\int_A[Z]u\right|
 =\left|\int_Au^2\right|=\|u\|_q^2.
\] For \(u\ne0\), (B26.6) now gives \[
 \|u\|_q^2\le
 \frac{\|u\|_q}{2^{n-1/2}\sqrt V}\frac D{n!}.
\] Cancelling the nonzero norm proves (B26.1); the zero case is immediate.

Equality in this integral estimate forces equality pointwise on the smooth locus of each component, with a common phase: the nonnegative continuous defect integrates to zero. The corresponding holomorphic Gauss map takes values in the unitary graph locus. The total-reality argument of the cited Theorem3(a) makes it constant, and the component is a translate of an abelian subvariety. This also follows directly from Debarre's tangent-span Lemma1.1 applied to the smooth locus: its constant tangent plane spans the tangent space of the generated abelian subvariety, of the same dimension. Simplicity rules it out for a positive proper component.

To see optimality, take a CM elliptic curve \(C\) with endomorphism field \(K\), give \(A=C^n\times C^n\) the product polarization, and let \(K\) act by \(a\) on the first half and \(\bar a\) on the second. The diagonal \(C^n\) has unitary graph tangent everywhere, and the relevant Weil form has constant phase. Its projection is nonzero; hence equality holds. This example is nonsimple, as required by the equality analysis. No claim that the projected class itself is integral is needed.

\subsection{Proof of the arithmetic margin and transport consequence}\label{proof-of-the-arithmetic-margin-and-transport-consequence}

The middle class \(v=P[Z]\) is integral. By the torus exterior-algebra parity calculation of B18, its positive signed square is in \(2\mathbf Z_{>0}\). Squaring the strict version of (B26.1), and using \(v=Ru\), gives a strictly positive integer \[
 R^2D^2-2^{2n-1}(n!)^2Vq(v,v).
\] Its first term is divisible by \(R^2\), and its second by \(2^{2n}(n!)^2V\). Their difference is therefore divisible by the positive gcd of these integers, proving (B26.2). Using only \(q(v,v)\ge2\) proves (B26.3), with the strictness stated above. QED.

In the residual-isogeny situation of B25, the already checked identity \(q(P_A[Z_A],P_A[Z_A])\ge2\deg g\) combines with (B26.1) to give \[
 D'^2\ge\frac{2^{2n}(n!)^2V}{R^2}\deg g,
 \qquad
 J_B\ge\frac{2^{2n}(n!)^2V}{R^2D'^2}J_A.
 \tag{B26.7}
\] Both inequalities are strict when the target is simple and the Weil component is nonzero. The original conservative B7/B25 constants remain valid; this is an optional stronger version. It does not change the exponent of the index or construct the intermediate witness.

\subsection{Attribution, framework and limits}\label{attribution-framework-and-limits}

The calibration bound and equality rigidity are imported from rStation paper4/v1, \emph{Calibrated Geometry of Weil Classes\ldots{}}, Theorem1(b), §2, and Theorem3(a),(c), §4, \url{https://rstation.ai/problems/hodge_conjecture/papers/4/v1}. The exact local proofs and differential-form convention were read. Only those parts are used; the other claims in that paper are not audited or imported here. Debarre, \emph{On the geometry of abelian varieties}, Chapter1, Lemma1.1 p8 gives the already checked tangent-span statement. The explicit signed-\(q\), polarization-volume and integer-gcd formulation is proved above.

The statement concerns effective cycle classes in the Betti realization of the middle homological motive, and their numerical degree against a polarization. The framework locators are cote44, PDF scan2/footer1, Theorem4 and scan4/footer3, Theorem5; and cote16, PDF scans45--46/footers44--45 for the separate algebraic-projector question. No calibration theorem is attributed to those manuscript pages. The strictness provides no uniform relative gap as \(D\) grows and does not close the bounded-witness existence gate.

\section{Finiteness of bounded-index polarized CM targets}

Run DEG\_L\_20261121; author /root; GPT-6 Astra Pro; effort not exposed.
B24 was written and whole-proof checked by root and independently by
/root/rallis\_t6 on 2026-11-21. The statement concerns polarization
homomorphisms, or numerical polarization classes, rather than individual
line bundles differing by a member of Picard degree zero.

\subsection{B24. Bounded core index gives finitely many polarized CM
points}\label{bounded-core-moduli-b24.-bounded-core-index-gives-finitely-many-polarized-cm-points}

Fix a CM field \(E\) of degree \(2g\), a CM type \(\Phi\), and a
positive integer \(V\). Consider complex polarized abelian varieties
\((A,\iota,\lambda)\) with rank-one \(E\)-homology and this fixed CM
type, where \(\iota:E\to\operatorname{End}^0(A)\) is the given rational
action, the Rosati involution of \(\lambda\) restricts to complex
conjugation on \(E\), and \(\deg\lambda=V^2\). Write \[
 A=\mathbf C^g/\Phi(\Lambda),\qquad
 \Lambda^- =\{x\in\Lambda:\mathcal O_E x\subseteq\Lambda\},
 \qquad J(A,\iota)=[\Lambda:\Lambda^-].
 \tag{B24.1}
\] These quantities are unchanged by replacing \(\Lambda\) with an
\(E^\times\)-multiple. Let \(h_E\) be the ideal class number of \(E\),
and let \(a_r(m)\) be the finite number of index-\(m\) subgroups of
\(\mathbf Z^r\).

For every positive integer \(B\), the number of \(E\)-linear polarized
isomorphism classes with \(J\le B\) is at most \[
 h_E\left(\sum_{m=1}^B a_{2g}(m)\right)
       a_{2g}(V^2)\,2^g.
 \tag{B24.2}
\] The same finiteness holds after imposing any fixed finite level
structure. In particular every infinite set of distinct polarized CM
points with these fixed field, type, and polarization-volume data has
unbounded \(J\). A bounded-\(J\) subset cannot be dense in a
positive-dimensional polarized moduli space.

Primitivity and simplicity are not needed for this finiteness statement.
The later lower bounds B20--B22 still retain their own simplicity
hypotheses.

\subsubsection{Proof}\label{bounded-core-moduli-proof}

The full \(\mathbf Z\)-lattice \(\Lambda^-\) is an invertible fractional
\(\mathcal O_E\)-ideal: it is a nonzero full lattice stable under
\(\mathcal O_E\), and \(\mathcal O_E\) is Dedekind. Nonzeroness follows,
for example, by choosing a nonzero integer \(c\) with
\(c\mathcal O_E\Lambda\subseteq\Lambda\). Its \(E^\times\)-homothety
class is therefore one of \(h_E\) possibilities. Fix one representative
\(I\) of each such ideal class and multiply \(\Lambda\) by an element of
\(E^\times\) so that its core becomes \(I\). Multiplication is complex
linear for the fixed CM type, commutes with \(E\), and transports the
polarization without changing its degree.

For a fixed representative \(I\), every remaining lattice satisfies \[
 I\subseteq\Lambda\subset E,\qquad [\Lambda:I]=m\le B.
\] There are exactly \(a_{2g}(m)\) overlattices of a free rank-\(2g\)
lattice of index \(m\). One way to see the equality with the subgroup
count is to pass to dual lattices: an overlattice determines and is
determined by an index-\(m\) sublattice of \(I^\vee\). For finiteness
alone, \(m\Lambda\subseteq I\) embeds \(\Lambda/I\) into the finite
group \(m^{-1}I/I\). We may count every overlattice, even though some
have core strictly larger than \(I\); this only increases the upper
bound. Thus at most \(h_E\sum_{m\le B}a_{2g}(m)\) unpolarized
\(E\)-linear isomorphism classes occur.

Fix one such \(A\). A polarization homomorphism \(\lambda:A\to A^\vee\)
of degree \(V^2\) is injective on integral first homology, and its image
is an index-\(V^2\) sublattice of \(H_1(A^\vee,\mathbf Z)\). There are
at most \(a_{2g}(V^2)\) possibilities for this image. Discard those
images not realized by a compatible polarization, and for each remaining
image choose one compatible polarization \(\lambda_0\) with that image.

Let \(\lambda\) be any other compatible polarization with the same image
and put \[
 u=\lambda_0^{-1}\lambda\in\operatorname{End}^0(A).
\] Compatibility means \(\lambda e=\overline e^{\vee}\lambda\) for every
\(e\in E\), and the same identity holds for \(\lambda_0\). Consequently
\(u\) commutes with \(E\). Since \(H_1(A,\mathbf Q)\) has rank one over
\(E\), its \(E\)-linear endomorphism algebra is \(E\) itself; hence
\(u\in E\). Symmetry of the two polarizations makes \(u\) self-adjoint
for the Rosati involution of \(\lambda_0\), so in fact \(u\in F:=E^+\).
Their positivity makes \(u\) totally positive. Concretely, the positive
Hermitian forms on the \(\Phi\)-coordinates have positive ratios equal
to the real embeddings of \(u\).

Equality of the two image lattices gives \(u\Lambda=\Lambda\). If \[
 R_F=\{b\in F:b\Lambda\subseteq\Lambda\},
\] then \(R_F\) is an order in \(F\) and \(u\in R_F^{\times,+}\). For
every \(b\in R_F^\times\), multiplication by \(b\) is an actual
\(E\)-linear automorphism of \(A\). Its pullback sends the polarization
\(\lambda_0u\) to \(\lambda_0ub^2\), because \(b\) is Rosati
self-adjoint and commutes with \(u\). Thus square classes of
\(R_F^\times\) suffice to count polarized isomorphism classes over this
one image lattice.

The group \(R_F^\times\) is a subgroup of \(\mathcal O_F^\times\).
Dirichlet's unit theorem implies that it is a finitely generated abelian
group of rank at most \(g-1\), and its torsion is contained in
\(\{1,-1\}\). Therefore \[
 [R_F^\times:(R_F^\times)^2]\le2^g.
\] The totally positive subgroup has no more cosets modulo these
squares. This proves the polarization count and hence (B24.2). A fixed
finite level structure has only finitely many choices on each polarized
variety. The final assertions follow immediately from finiteness. QED.

\subsubsection{Consequence for the degree and division
gates}\label{bounded-core-moduli-consequence-for-the-degree-and-division-gates}

For the simple reference and fixed generic Weil data in B21--B22, a
bounded-degree source family yields reduced transported witnesses of
degree at least \(c\sqrt J\), with an explicit positive \(c\) after B22.
B24 proves that the alternative of a bounded core index does not
describe an infinite polarized CM Hecke set with these fixed data. Thus
the fixed-reference, bounded-source construction cannot supply a uniform
effective witness bound on such an infinite set, in particular on a
dense positive-dimensional Hecke orbit. This is still a restriction on
the construction, not a lower bound for all cycles on the targets. The
proof does not assert any Hecke-orbit density theorem and does not
require one for this implication.

The field and CM type stay fixed. Allowing them to vary is outside
(B24.2). Picard-degree-zero changes of a line bundle do not change the
polarization homomorphism or any degree used here; they must not be
counted as new polarized moduli points.

\subsubsection{Exact background and framework
placement}\label{bounded-core-moduli-exact-background-and-framework-placement}

\begin{itemize}
\tightlist
\item
  Milne, \emph{Algebraic Number Theory}, version 3.08, Theorems
  4.3--4.4, printed pp.70--71, gives ideal class finiteness; Theorem
  5.1, printed p.85, gives finite generation and rank of units. Original
  author source: \url{https://www.jmilne.org/math/CourseNotes/ANT.pdf}.
  The exact statements and the short proof of Theorem 4.4 were read.
  These classical results are imported, not reproved.
\item
  The CM lattice model and the Rosati metric convention are those
  already checked in B19.1 and T6.17. The image-lattice and square-class
  counting argument above is supplied in full for the form needed here;
  it does not invoke an unlocated finiteness theorem for polarizations.
\item
  Grothendieck framework: polarized abelian schemes and the intersection
  form, cote 44, PDF scan 2 / archive footer 1, Theorem 4; birigidified
  line bundles and polarization homomorphisms in SGA7 I, Exposé VII,
  Example 2.9.5 p.175 and Remark 2.9.6 p.176. These are framework
  placements, not claims that B24 or a CM counting theorem is stated on
  those pages.
\end{itemize}

\subsubsection{Check status}\label{bounded-core-moduli-check-status}

Author and independent reviewer /root/rallis\_t6 have read the complete
proof, including the rank-one centralizer argument, same-image
automorphism, real-unit squares, and the distinction between a line
bundle and a polarization homomorphism. The reviewer independently
reopened the exact Milne class-number and unit statements. Both checks
passed with the explicit Rosati compatibility retained, including in
imprimitive types. No numerical computation or proof-assistant
verification is claimed.

\section{A local Fourier normalization control}

The local shell results LN1--LN6 are incorporated into B14a and B17. The
following separate control is included with its original identifier; it
concerns a model Fourier functional, not a geometric Weil period.

\subsection{LN7. A model period can be small while the integral test
vector has bounded
norm}\label{local-fourier-period-ln7.-a-model-period-can-be-small-while-the-integral-test-vector-has-bounded-norm}

Let \(\eta|_{\mathcal O^\times}\) be primitive of conductor \(m\ge1\),
and let \(\psi:F\to\mathbb C^\times\) be an additive character trivial
on \(\mathcal O\) but nontrivial on \(\varpi^{-1}\mathcal O\). For the
cyclotomic-integral-valued test vector (LN5.1), define the Fourier
functional

\[
 \ell_m(\phi)=\int_{F^3}\phi(x)\psi(\varpi^{-m}x_1)\,dx.
\]

Then

\[
 |\ell_m(\phi_\eta)|=q^{-m/2},
 \qquad \|\phi_\eta\|^2=Z_\eta(\phi_\eta,\phi_\eta)=1-q^{-1}.
 \tag{LN7.1}
\]

\textbf{Proof.} Set \(R=\mathcal O/\varpi^m\mathcal O\) and
\(\Psi(a)=\psi(\varpi^{-m}a)\), a primitive additive character of \(R\).
For \(a\in R\), put

\[
 S(a)=\sum_{x\in R^\times}\eta(x)^{-1}\Psi(ax).
\]

For a unit \(a\), substitution \(y=ax\) gives \(S(a)=\eta(a)S(1)\). If
\(m\ge2\), primitiveness provides \(h\in1+\varpi^{m-1}\mathcal O\) with
\(\eta(h)\ne1\). For nonunit \(a\), the substitution \(x\mapsto hx\)
does not change \(\Psi(ax)\), so \(S(a)=0\). For \(m=1\), the only
nonunit is zero and the nontrivial character sum is zero. Orthogonality
of additive characters on the finite ring gives

\[
 \sum_{a\in R}|S(a)|^2
 =q^m\sum_{x\in R}|\mathbf1_{R^\times}(x)\eta(x)^{-1}|^2
 =q^m|R^\times|.
\]

There are \(|R^\times|\) unit arguments, all with magnitude \(|S(1)|\),
so \(|S(1)|^2=q^m\). Haar integration over unit residue classes gives
\(\ell_m(\phi_\eta)=q^{-m}S(1)\), proving (LN7.1). QED.

This is a model Fourier period, not the selected geometric Weil period.
It proves a narrower point: integral values of the Schwartz input and a
uniform Rallis norm do not imply a uniform lower bound for a
Haar-averaged arithmetic functional. In this example forcing the
displayed period to have absolute value one costs \(q^{m/2}\) in input
norm. No such growth is asserted for the actual construction.

\section{Semiregularity and the replacement-object problem}

% Candidate statements and complete proofs for the DEG_L paper.
% Checked by /root/semireg_t7, GPT-6 Astra Pro; inherited effort not exposed.
% This is an input fragment, not a standalone article.
% Source locators are written explicitly so no guessed bibliography keys are used.

\subsection{The actual surface and its deformation problem}

\begin{proposition}[S-01: Allowed source covers do not preserve an embedding premise]
\label{prop:semireg-cover-domain}
Let $f:Y\to A$ be a surface morphism obtained from the four rational
Hodge maps in 032, Proposition~3.1. If $\pi:Y'\to Y$ is a connected finite
\'etale cover of degree $d>1$, the pulled-back maps still satisfy its
nonzero-period criterion, but $f\pi$ is not a closed immersion.
Moreover $(f\pi)_*[Y']=d f_*[Y]$.
\end{proposition}
\begin{proof}
Pullback preserves rational Hodge maps and multiplies the integral of
a top-degree class by $d$. Every geometric point of $Y$ has $d$ distinct
preimages with the same image under $f\pi$, so this morphism is not
injective on geometric points. The last assertion is the proper
pushforward formula for a finite map. The deeper principal covers
used here are explicitly permitted in 032, Corollary~6.7 and \S7.
\end{proof}

\begin{proposition}[S-02: The exact conormal gate]
\label{prop:semireg-dual-gate}
For a smooth closed immersion $i:Y\hookrightarrow A$ of a surface in
an abelian $g$-fold, $g\geq3$, write $N=N_{Y/A}$ and
$W=H^0(A,\Omega_A^1)$. Bloch semiregularity is equivalent to
surjectivity of
\[
 H^1(A,\mathcal O_A)\otimes\bigwedge^3W
 \longrightarrow H^1(Y,N^\vee\otimes K_Y).
\]
The map is restriction followed by the natural morphism
$\bigwedge^3 i^*\Omega_A^1\to N^\vee\otimes K_Y$.
\end{proposition}
\begin{proof}
Bloch's map is
$H^1(Y,N)\to H^{g-1}(A,\Omega_A^{g-3})$.
Serre duality identifies its dual source with $H^1(A,\Omega_A^3)$
and its dual target with $H^1(Y,N^\vee\otimes K_Y)$.
The conormal exact sequence induces the displayed exterior-power
map, since the quotient $\Omega_Y^1$ has rank two.
This is the defining dual map in Buchweitz--Flenner \S8.1;
Proposition~8.2 identifies it with Bloch's map.
Finally $\Omega_A^3=\mathcal O_A\otimes\bigwedge^3W$,
and finite-dimensional duality gives the equivalence.
\end{proof}

\begin{proposition}[S-03: A nonzero period does not force semiregularity]
\label{prop:semireg-period-counterexample}
There is a smooth surface in a CM abelian threefold with a nonzero
product of two holomorphic and two antiholomorphic one-forms for
which the Bloch semiregularity kernel has dimension at least eight.
\end{proposition}
\begin{proof}
Let $E$ be a CM elliptic curve, $B=E^2$, and
$L=\mathcal O_E(3[0])\boxtimes\mathcal O_E(3[0])$.
A general $C\in|L|$ is smooth and connected. Since $L^2=18$ and
$K_B=\mathcal O_B$, adjunction gives $g(C)=10$ and
$N_{C/B}=K_C$. Put $A=B\times E$ and $Y=C\times E$.
Then $N_{Y/A}=K_C\boxtimes\mathcal O_E$, so K\"unneth gives
$h^1(Y,N)=1+10=11$. The divisor semiregularity target
$H^2(A,\mathcal O_A)$ has dimension three.
The curve $C$ dominates the first elliptic factor: otherwise its
ample divisor class would be a fibre class. Hence a nonzero one-form
from that factor restricts to a nonzero form $\alpha$ on $C$.
For a nonzero one-form $\beta$ on the last $E$, the integral
$\int_Y(\alpha\wedge\beta)\wedge\overline{\alpha\wedge\beta}$
is nonzero, by the product formula and positivity of the two
one-dimensional integrals.
\end{proof}

\begin{proposition}[S-04: Rigid generating sources]
\label{prop:semireg-rigid-source}
Let $Y$ be a smooth projective surface with $H^1(Y,T_Y)=0$ and
$i:Y\hookrightarrow A$ a smooth closed immersion. Let $G\subset A$
be the abelian subvariety generated by a translate of $i(Y)$.
Then the ambient obstruction map satisfies
\[
\ker(\operatorname{ob}_i)=
\operatorname{Lie}(A)\otimes
\ker\bigl(H^1(A,\mathcal O_A)\to H^1(Y,\mathcal O_Y)\bigr).
\]
In particular it is injective when $G=A$. If in addition $i$ is
Bloch semiregular, this is also the infinitesimal Hodge-locus kernel
of the entire class $i_*[Y]$.
\end{proposition}
\begin{proof}
Restriction on $H^1(T_A)$ is $\operatorname{id}_{\operatorname{Lie}(A)}
\otimes i^*$, since $T_A$ is trivial. The normal sequence and
$H^1(T_Y)=0$ make $H^1(i^*T_A)\to H^1(N)$ injective.
This proves the kernel formula. The map $\operatorname{Alb}(Y)\to A$
has image $G$, so $i^*$ on $H^1(\mathcal O)$ is injective precisely
when $G=A$. Buchweitz--Flenner, Corollary~4.12 and Proposition~8.2,
identify
\[
 \sigma_i(\operatorname{ob}_i(\kappa))
 =\kappa\mathbin{\lrcorner}i_*[Y].
\]
If $\sigma_i$ is injective, the kernels of the two sides agree.
\end{proof}

The actual 032 source is a compact smooth complex two-ball quotient.
Calabi--Vesentini, \emph{Sur les vari\'et\'es complexes compactes
localement sym\'etriques}, Th\'eor\`eme~3, p.~315, gives
$H^1(Y,T_Y)=0$. The nonzero four-factor period does not establish
$G=A$. For example, separate injectivity on the two eigenspaces
of a quadratic action is insufficient: in
$A=E^n\times E^n$ with action $(a,\bar a)$, pullback to the diagonal
$E^n$ is injective on each eigenspace, but their images coincide.

\begin{proposition}[S-05: A quantitative obstruction for the actual morphism]
\label{prop:semireg-map-dimension}
Let $Y$ be a compact smooth complex two-ball quotient and
$f:Y\to A$ a morphism with two-dimensional image, where $\dim A=g$.
Let $\tau_f$ be the Buchweitz--Flenner morphism semiregularity map.
Then
\[
 \dim\ker\tau_f\ \geq\ gq(Y)-\binom{2g}{g-2}.
\]
\end{proposition}
\begin{proof}
In characteristic zero $df:T_Y\to f^*T_A=\mathcal O_Y^{\oplus g}$
is generically injective. Its kernel is a torsion-free sheaf of
generic rank zero, hence vanishes. With $N_f=\operatorname{coker}(df)$,
the derived dual of $L_{Y/A}$ is the two-term complex
$[T_Y\to f^*T_A]$ in degrees $0,1$, so
$T^2_{Y/A}(\mathcal O_Y)=H^1(Y,N_f)$.
Calabi--Vesentini and the cokernel sequence give
$H^1(\mathcal O_Y)^{\oplus g}\hookrightarrow H^1(N_f)$.
The target specified in Buchweitz--Flenner, Proposition~7.21,
Lemma~7.22 and the definition before Theorem~7.23, is
\[
 \bigoplus_{r=0}^{g-2}H^{r+2}(A,\Omega_A^r).
\]
Its dimension is
$\sum_{r=0}^{g-2}\binom gr\binom g{r+2}=\binom{2g}{g-2}$
by Vandermonde. Rank--nullity proves the assertion.
\end{proof}

\begin{corollary}[S-06: Failure at sufficiently deep allowed levels]
\label{cor:semireg-level-failure}
Fix a nonzero 032 surface construction $f:Y\to A$. There are
connected finite \'etale covers $\pi_j:Y_j\to Y$, preserving its
four pulled-back Hodge maps and nonzero period, for which
$\tau_{f\pi_j}$ is not injective. Such covers may be chosen among
a cofinal system of principal congruence levels.
\end{corollary}
\begin{proof}
The nonzero period implies $H^1(Y,\mathbf C)\ne0$.
Venkataramana, \emph{Virtual Betti numbers of compact locally
symmetric spaces}, Theorem~1, proves that finite-dimensional
direct-limit cohomology over finite covers equals the cohomology
of the compact dual. Here the compact dual is $\mathbf P^2$,
whose first cohomology is zero. Pullback under finite covers is
injective over $\mathbf Q$ by transfer, so the direct limit is
nonzero and therefore infinite-dimensional. Hence $q(Y_j)$ is
unbounded, and the preceding proposition applies with fixed $A$.
The congruence version is explicitly identified in the final Remark
on p.~3 of that paper. Its proof uses the same rational-conjugation
action, since rational conjugation and finite intersection preserve
congruence subgroups; principal subgroups are cofinal.
Persistence of the period follows from the first proposition.
\end{proof}

This is a statement about the deformation problem of a morphism.
The factorization
$\tau_f=\sigma_{Rf_*\mathcal O_Y}\circ\psi_f$
does not imply that a kernel of $\tau_f$ survives in the perfect
complex obstruction space: $\psi_f$ itself can have a kernel.
Nor does the corollary settle Bloch semiregularity of a particular
smooth reduced image.

\begin{proposition}[S-07: Descent of the cover obstruction space]
\label{prop:semireg-cover-invariants}
For a finite \'etale Galois cover $\pi:Y'\to Y$ with deck group $H$,
$N_{f\pi}=\pi^*N_f$ and
$H^1(Y',N_{f\pi})^H=H^1(Y,N_f)$.
\end{proposition}
\begin{proof}
Flat pullback identifies the tangent cokernel sequences.
Finite pushforward and the projection formula identify
$\pi_*\pi^*N_f=N_f\otimes\pi_*\mathcal O_{Y'}$.
In characteristic zero averaging over $H$ splits off precisely
the invariant summand $N_f$. Cohomology commutes with this direct
summand, giving the assertion.
\end{proof}

\subsection{Derived replacements and their remaining gate}

\begin{proposition}[S-09: An integral Chern-component projector]
\label{prop:semireg-k-projector}
Let $A$ be an abelian $g$-fold and $\xi\in K_0(A)$ with
$\operatorname{ch}_j(\xi)=0$ for $j<p$. For $m\geq2$ put
\[
 M=[m]^*,\qquad
 D=\prod_{j=p+1}^g(m^{2j}-m^{2p}),\qquad
 \eta=\prod_{j=p+1}^g(m^{2j}\operatorname{id}-M)\xi.
\]
Then $\eta\in K_0(A)$ and
$\operatorname{ch}(\eta)=D\operatorname{ch}_p(\xi)$ in rational
Betti cohomology.
\end{proposition}
\begin{proof}
The operator $M$ multiplies $H^{2j}(A,\mathbf Q)$ by $m^{2j}$.
Thus every higher Chern component is annihilated by its corresponding
factor, lower components were already zero, and the component of
index $p$ is multiplied by $D$. The polynomial has integer
coefficients. Integral K-classes are represented by perfect complexes
on the smooth projective variety, using finite locally free
resolutions and shifts to implement negative coefficients.
\end{proof}

Starting from an integral codimension-$p$ cycle of class $N_0t$,
one may use $\xi=\sum a_i[\mathcal O_{Z_i}]$; its leading Chern
character is the fundamental cycle. Alternatively, resolve each
$Z_i$ and use the proper pushforward of its structure sheaf;
Grothendieck--Riemann--Roch gives the same leading component.
Adding line-bundle classes gives the useful character
$Nt+\sum_a b_a e^{a\theta}$, which remains Hodge on the polarized
Route-2 family. This constructs a K-class, not a semiregular object.

\begin{proposition}[S-08: Multiplicity creates trace kernels]
\label{prop:semireg-multiplicity}
If $E$ is perfect and $\operatorname{Ext}^2(E,E)\ne0$, then
$E^{\oplus m}$ is not fully semiregular for any $m\geq2$.
\end{proposition}
\begin{proof}
Its degree-two self-Ext group is
$\operatorname{Ext}^2(E,E)\otimes\operatorname{End}(\mathbf C^m)$.
The Atiyah class is $\operatorname{at}(E)\otimes\operatorname{id}$;
therefore every semiregularity component is the corresponding
component for $E$ tensored with the matrix trace.
A nonzero degree-two Ext class tensored with a nonzero traceless
matrix lies in the kernel.
\end{proof}

\begin{proposition}[S-10: Full Fourier--Mukai transport]
\label{prop:semireg-full-fm}
For the fixed enhanced Fourier--Mukai equivalence in the supplied
032 companion \S5.2, the full trace maps satisfy
\[
 \sigma_{\Phi Q}\circ\Phi_{\rm Ext}
   =\Phi^H_{HH}\circ\sigma_Q
 \quad\hbox{on }\operatorname{Ext}^2_A(Q,Q).
\]
Thus full semiregularity is invariant under this equivalence.
\end{proposition}
\begin{proof}
Use the companion's enhanced adjunctions and notation.
Put
$b_Q(v)=\operatorname{Tr}_{Q,H_A[2]}(A_Q^H[2]\circ v)
\in HH_{-2}(A)$.
Equations (5.20), (5.24) identify its HKR image with $\sigma_Q(v)$.
For $\delta\in T^{-2}(A)$, trace-compatible adjunction gives
\[
 \langle b_Q(v),\delta\rangle
   =\operatorname{Tr}_{A,Q}(\delta_Q[2]\circ v).
\]
Indeed, if $\delta':H_A\to\omega_A[g-2]$ is the adjoint coefficient
map, move $\delta'[2]$ through the partial trace by (5.24);
evaluation then gives the right side. This is exactly (5.37)
with the displayed even shift.
The convolution proof of (5.33), with the same inverse adjunction
data and shifted by two, gives
$(\Psi^{-2}\delta)_{\Phi Q}=\theta_Q[-2]\Phi(\delta_Q)$.
By the trace preservation in (5.31),
\[
 \langle b_{\Phi Q}(\Phi_{\rm Ext}v),\Psi^{-2}\delta\rangle
  =\operatorname{Tr}_{A,Q}(\delta_Q[2]\circ v)
  =\langle b_Q(v),\delta\rangle.
\]
Definition (5.34) and nondegeneracy of the Serre pairing imply
$b_{\Phi Q}(\Phi_{\rm Ext}v)=\Phi_*^{HH}b_Q(v)$.
Apply (5.35) and the fixed HKR maps. The two vertical maps
are isomorphisms, proving the assertion about kernels.
This proves a refinement; Proposition~5.2 itself only asserted
transport on the obstruction-evaluation image.
\end{proof}

\begin{proposition}[S-11: An ample complete intersection need not be semiregular]
\label{prop:semireg-ci-counterexample}
Let $A$ be an abelian fourfold, $L$ very ample, and
$Y=D_1\cap D_2$ a smooth transverse intersection of two members
of $|L|$. Then $h^1(N_{Y/A})=24$, whereas
$h^3(A,\Omega_A^1)=16$. Thus its Bloch map has a kernel of
dimension at least eight.
\end{proposition}
\begin{proof}
The twisted Koszul sequence is
$0\to L^{-1}\to\mathcal O_A^{\oplus2}\to\mathcal I_Y\otimes L\to0$.
Kodaira vanishing and $K_A=\mathcal O_A$ give $H^j(A,L)=0$ for
$j>0$; Serre duality gives $H^j(A,L^{-1})=0$ for $j<4$.
Thus
$H^2(A,\mathcal I_Y\otimes L)=H^2(A,\mathcal O_A)^{\oplus2}$.
The restriction sequence gives
$H^1(Y,L|_Y)=H^2(A,\mathcal I_Y\otimes L)$, of dimension twelve.
Transversality identifies $N_{Y/A}=(L|_Y)^{\oplus2}$.
The Bloch target has dimension
$\binom43\binom41=16$, so rank--nullity finishes the proof.
\end{proof}

\begin{proposition}[S-12: A pure Weil character obstruction for coherent sheaves]
\label{prop:semireg-pure-weil-no-go}
Let $A_0$ be a projective member of Voisin's four-dimensional
marked $\mathbf Q(i)$-Weil torus family, and let $t\ne0$ be a
rational Weil class. No fully Buchweitz--Flenner-semiregular
coherent sheaf $F$ on $A_0$ satisfies
\[
 \operatorname{ch}_1(F)=\operatorname{ch}_3(F)=0,\qquad
 \operatorname{ch}_2(F)=t.
\]
The rank and top-degree character are unrestricted.
\end{proposition}
\begin{proof}
The unpolarized Weil torus parameter space is smooth near $A_0$;
Voisin, \emph{A counterexample to the Hodge conjecture extended
to K\"ahler varieties}, \S3, equations (3.1)--(3.8), describes it.
Lemma~3.1 shows that the rational Weil space remains of Hodge
type $(2,2)$. Hence the components with indices $I=\{1,2,3\}$
in the assertion remain horizontal Hodge classes. In dimension
four these are precisely the components tested by full
semiregularity.
Buchweitz--Flenner Theorem~5.1 and its proof on p.~179 would
therefore extend $F$ as a coherent flat sheaf on a neighborhood.
Voisin Proposition~3.2 and Proposition~2.1 show that very general
fibres in that neighborhood have $\operatorname{NS}=0$ and
$c_2(G)=0$ for every coherent analytic sheaf $G$.
On such a fibre $c_1(F)=c_2(F)=0$, whence
$\operatorname{ch}_2(F)=0$, contrary to the nonzero horizontal
class $t$.
\end{proof}

This last proposition is a coherent-sheaf statement for the
specified fourfold family. Extension to arbitrary perfect complexes
would require an analytic effectivity theorem for their formal
deformations; algebraic/formal Pridham alone is not that theorem.
Allowing a nonzero polarization term in the Chern character
avoids this particular contradiction.

\begin{proposition}[S-13: One exponential polarization term does not repair the coherent gate]
\label{prop:semireg-single-exponential-no-go}
In the fourfold setting of S-12, let $L$ be a polarization whose
Rosati involution restricts to conjugation on $\mathbf Q(i)$, and
write $\theta=c_1(L)$. If $a\in\mathbf Z$, $\lambda\in\mathbf Q^\times$,
and $t\ne0$ is a rational Weil class, no fully semiregular coherent
sheaf $F$ has
\[
 \operatorname{ch}(F)=r e^{a\theta}+\lambda t+c[\mathrm{pt}]
\]
for any $r,c\in\mathbf Q$.
\end{proposition}
\begin{proof}
Write $H^1(A_0,\mathbf C)=V_\sigma\oplus V_{\bar\sigma}$.
The Rosati condition places $\theta$ in
$V_\sigma\otimes V_{\bar\sigma}$, whereas the complexified Weil
space is
$\bigwedge^4V_\sigma\oplus\bigwedge^4V_{\bar\sigma}$.
Both eigenspaces have dimension four. Consequently
$\theta\smile t=0$, and hence $e^{-a\theta}t=t$.
Tensoring by $L^{-a}$ is a Fourier--Mukai autoequivalence.
The functoriality in Perry Remarks~2.3 and~2.5, p.~5, preserves
full semiregularity under this equivalence, as in S-10.
It preserves coherence as well.
For $F'=F\otimes L^{-a}$ the Chern-character formula gives
\[
 \operatorname{ch}(F')=r+\lambda t+c[\mathrm{pt}].
\]
Thus $\operatorname{ch}_1(F')=\operatorname{ch}_3(F')=0$ and
$\operatorname{ch}_2(F')=\lambda t\ne0$, contradicting S-12.
\end{proof}

The integer condition on $a$ ensures that this normalization is an
actual line-bundle twist. The conclusion does not exclude all
polarization polynomials: it excludes the displayed single
exponential, including an arbitrary additional top-degree term.

\begin{proposition}[S-14: An equivariant perfect-complex gate implies bounded division]
\label{prop:semireg-perry-bounded-division}
Let $\pi:A\to S$ be an abelian scheme over a smooth irreducible
complex algebraic variety of finite type, with a relatively very
ample line bundle $H$ and $\theta=c_1(H)$. Let
$t\in H^0(S,R^{2p}\pi_*\mathbf Q(p))$ be a flat section whose values
are Hodge classes, where $0\leq p\leq\dim(A/S)$.
Suppose that for one $s_0\in S(\mathbf C)$ there exist:
\begin{enumerate}
\item a finite subgroup $G\subset A_{s_0}\times A_{s_0}^{\vee}$
with a specified coherent action on $\operatorname{Perf}(A_{s_0})$
by translations and tensoring by line bundles;
\item a weakly $G$-semiregular perfect complex $E_0$ with
$\operatorname{Ext}^{i}(E_0,E_0)=0$ for all $i<0$;
\item an algebraic class $B_0\in H^2(A_{s_0},\mathbf Q(1))$,
a nonzero $\lambda\in\mathbf Q$, and $P\in\mathbf Q[x]$ such that
\[
 e^{B_0}\operatorname{ch}(E_0)=\lambda t_{s_0}+P(\theta_{s_0})
\]
in the full even rational cohomology.
\end{enumerate}
Then there is an integer $C\geq1$ such that, for every
$s\in S(\mathbf C)$, there are an integer $1\leq N_s\leq C$
and effective integral codimension-$p$ cycles $Z_s^+,Z_s^-$,
each of $H$-degree at most $C$, with
\[
 N_s t_s=\operatorname{cl}(Z_s^+)-\operatorname{cl}(Z_s^-).
\]
In particular, there are single integers $N,D\geq1$ for the entire
family such that $Nt_s$ admits such representatives of degrees
at most $D$ on every fibre.
A sufficient form of the second hypothesis is a specified
$G$-equivariant structure on $E_0$ for which the ordinary full
semiregularity map is injective on
$\operatorname{Ext}^2(E_0,E_0)^G$, together with the same vanishing
of negative self-Ext groups.
\end{proposition}
\begin{proof}
The class
$w=\lambda t+P(\theta)$ is a global flat section of the full even
cohomology local system, and every component has the required
Hodge type. It extends
$w_0=e^{B_0}\operatorname{ch}(E_0)$.
Apply Alexander Perry, \emph{The semiregularity theorem for
equivariant noncommutative varieties},
arXiv:2604.00511v2 (24 June 2026), Theorem~1.2, p.~3.
Its second conclusion states that $w_s$ is algebraic for every
$s\in S(\mathbf C)$. Since $P(\theta_s)$ is algebraic and
$\lambda\ne0$, each $t_s$ is algebraic.
DIV1(2), applied to the fixed smooth projective family and $H$,
now supplies a single bound $C$ for the signed integral complexity.
This is exactly the displayed conclusion.
For a common multiplier take $N=\operatorname{lcm}(1,\ldots,C)$
and multiply the cycles on fibre $s$ by $N/N_s$.
Their degrees are at most $D=NC$.
Finally, Perry Definition~2.6 and Lemma~2.7, pp.~5--6, show that
the stated invariant-injectivity criterion implies weak
$G$-semiregularity over $\mathbf C$.
\end{proof}

The gate requires the entire displayed Chern character to remain
Hodge, a coherent categorical action, and vanishing of negative
self-Ext groups. A virtual K-class alone supplies none of these
deformation properties. The theorem supplies algebraicity on
every fibre; no properness assertion for a moduli space of
perfect complexes is used. The bound obtained from DIV1 is
existential. Here $G$ is a subgroup of the target's translation
and degree-zero line-bundle group, not an arbitrary deck group
of the source surface.

\begin{proposition}[S-15: Canonical source-cover equivariance retains a full kernel]
\label{prop:semireg-cover-equivariance}
Let $f:Y\to A$ be a proper morphism from a smooth projective
variety to an abelian variety, and put $E=Rf_*\mathcal O_Y$.
\begin{enumerate}
\item If $\rho:Y'\to Y$ is a finite \'etale Galois cover with
group $K$, the canonical $K$-equivariant lift of
$E'=R(f\rho)_*\mathcal O_{Y'}$ for the trivial action of $K$
on $\operatorname{Perf}(A)$ can be semiregular only if $E$
is fully semiregular.
\item Suppose $\rho$ is obtained by base change from an isogeny
$q:A'\to A$. Set
$G=\ker(q^*:A^\vee\to(A')^\vee)$, acting on
$\operatorname{Perf}(A)$ by the coherently chosen tensor
line bundles in $q_*\mathcal O_{A'}$.
The canonical $G$-equivariant lift of
$q_*q^*E=R(f\rho)_*\mathcal O_{Y'}$ can be semiregular only if
$E$ is fully semiregular.
\end{enumerate}
These assertions concern the canonical lifts supplied by the
covers; they do not exclude every possible new equivariant
structure on the underlying objects.
\end{proposition}
\begin{proof}
For (1), finite pushforward and characteristic-zero averaging give
\[
 (E')^K
 =Rf_*(\rho_*\mathcal O_{Y'})^K
 =Rf_*\mathcal O_Y=E.
\]
Because the categorical action is trivial,
$\operatorname{Perf}(A)^K$ is the finite product of copies of
$\operatorname{Perf}(A)$ indexed by irreducible representations
of $K$. The factor belonging to the trivial representation of
the canonical lift is $E$. Self-Ext groups, Hochschild homology,
and the semiregularity map decompose by these factors.
Semiregularity of the lift therefore implies that of $E$.

For (2), let $R=q_*\mathcal O_{A'}$. The character decomposition
of the finite abelian torsor gives
$R=\bigoplus_{P\in G}P$, with multiplication isomorphisms
$P\otimes Q\cong PQ$. These provide the indicated coherent
tensor action. A $G$-equivariant object is equivalently an
$R$-module object; since $q$ is finite \'etale this identifies
$\operatorname{Perf}(A)^G$ with $\operatorname{Perf}(A')$.
Under this equivalence the canonical lift of $q_*q^*E$
corresponds to $q^*E$.
Flat base change for the square defining $Y'$ identifies the
underlying object with the stated source-cover pushforward.

The pullback map
\[
 q^*:\operatorname{Ext}^2_A(E,E)
 \longrightarrow
 \operatorname{Ext}^2_{A'}(q^*E,q^*E)
\]
is injective: by adjunction its target is
$\operatorname{Ext}^2_A(E,E\otimes R)$, and the unit
$\mathcal O_A\to R$ has the retraction
$(\deg q)^{-1}\operatorname{Tr}_q$.
Semiregularity is natural for the enhanced functor $q^*$,
by Perry Remark~2.3, p.~5. Thus pullback carries
$\ker\sigma_E$ injectively into $\ker\sigma_{q^*E}$.
Semiregularity of the canonical lift would make the latter
kernel zero, proving the assertion.
\end{proof}

For an ordinary source deck transformation $k$ one has
$f\rho k=f\rho$. Consequently an induced translation of the
target satisfying this equivariance must be the identity.
An ordinary deck action therefore supplies internal
automorphisms of the pushed-forward object; it does not by
itself specify the nontrivial target action required in S-14.
Part (2) treats a special case where target tensor symmetries
really do arise, and shows why their canonical lift still
does not remove a pre-existing full semiregularity kernel.

There is also an independent primary-source confirmation of S-10:
Perry Remarks~2.3 and~2.5, p.~5, give functoriality of the
semiregularity morphism on the entire mapping complex and identify
its degree-two part with Buchweitz--Flenner under HKR.
An equivalence induces isomorphisms on both mapping complexes and
Hochschild homology, so the full kernel is invariant. This argument
does not depend on a raw versus Todd-modified HKR convention;
for abelian varieties the Todd class is one.

\begin{proposition}[S-16: A polarization exponential realizes a K-class by a vector bundle]
\label{prop:semireg-kclass-bundle-realization}
Let $X$ be a smooth projective variety, let $L$ be ample, and let
$\xi\in K_0(X)$. There exist integers $m,r\geq1$ and a vector
bundle $V$ such that
\[
 [V]=\xi+r[L^{-m}]\quad\hbox{in }K_0(X).
\]
In particular, if $\operatorname{ch}(\xi)=\lambda t$, then
$\operatorname{ch}(V)=r e^{-m c_1(L)}+\lambda t$.
In the setting of S-13, with $\lambda t\ne0$, every vector bundle
with this character fails full semiregularity.
\end{proposition}
\begin{proof}
On a smooth projective variety a K-class is a difference
$\xi=[U]-[W]$ of vector bundles: take finite locally free
resolutions and group their even and odd terms.
For sufficiently large $m$, the vector bundle $W\otimes L^m$
is generated by finitely many global sections.
The precise ampleness criterion is Stacks Project,
Proposition~28.27.13(7)--(8), tag 01Q3.
Choose a surjection
\[
 0\longrightarrow K\longrightarrow\mathcal O_X^{\oplus r}
   \longrightarrow W\otimes L^m\longrightarrow0 .
\]
One can add a redundant zero section to ensure $r\geq1$.
The kernel $K$ is locally free, since the surjection onto a
vector bundle splits locally. After tensoring by $L^{-m}$,
the K-group relation is
$[K\otimes L^{-m}]=r[L^{-m}]-[W]$.
Thus $V=U\oplus(K\otimes L^{-m})$ has the required class.
The Chern-character formula follows by additivity, and the
last assertion is S-13 with $a=-m$ and $c=0$.
\end{proof}

This gives an actual coherent vector-bundle representative
after one exponential addition. It supplies neither full nor
equivariant semiregularity; in the indicated fourfold case the
full version is ruled out by S-13.

\begin{proposition}[S-17: Tensor--determinant untwisting after a positive-rank deformation]
\label{prop:semireg-tensor-determinant}
Let $\pi:\mathcal X\to X$ be a $\mu_M$-gerbe over a connected
smooth complex scheme, and let $F$ be a weight-one perfect complex
of positive virtual rank $r$. Then
\[
 F^{\otimes^{\mathbf L}r}\otimes\det(F)^{-1}
 \simeq\pi^*K
\]
for an ordinary perfect complex $K$ on $X$.

In the family setting and under all the hypotheses of S-14,
suppose $\dim A_s=2n$, $n\ge2$, that $t$ is a nonzero relative
Weil class of codimension $n$, and that
\[
 w_0=e^{B_0}\operatorname{ch}(E_0)
     =P(\theta_0)+\lambda t_0,
 \qquad P(0)=r>0,\quad\lambda\ne0,\quad\theta_0t_0=0.
\]
Put $p_1=[x]P(x)$. On the local deformation neighborhood supplied
by Perry, the descended complex has character, through
codimension $n$,
\[
 \operatorname{ch}(K_s)
 \equiv e^{-p_1\theta_s}P(\theta_s)^r+r^r\lambda t_s.
 \tag{S17.1}
\]
There are fixed integers $R,m\ge1$ and $a,k\ge0$ from a finite
bundle presentation and generation twists for which the following
numbers give a uniform signed-cycle bound on the entire base.
Set
\[
 Q(x)=e^{kx}\bigl(e^{-p_1x}P(x)^r+R e^{-mx}+a\bigr),
 \qquad\mu=r^r\lambda,
\]
\[
 C_n=[u^n]\exp\!\left(\sum_{j=1}^n
     (-1)^{j-1}(j-1)![x^j]Q(x)u^j\right),
 \qquad\beta=(-1)^{n-1}(n-1)!\mu .
\]
If $d\ge1$ clears the denominators of $C_n$ and $\beta$, and
$v=\int_{A_s}\theta_s^{2n}$, then
\[
 N=|d\beta|>0,\qquad D=dC_nv>0
 \tag{S17.2}
\]
are fixed integers such that $Nt_s$ is the difference of two
effective integral cycles, each of polarization degree at most
$D$, on every fibre. A negative nonzero virtual rank is reduced
to this case by shifting the initial and deformed complexes.
\end{proposition}
\begin{proof}
The determinant functor for perfect complexes exists and commutes
with base change by Knudsen--Mumford, \emph{The projectivity of the
moduli space of stable curves I}, Theorem~2, p.~42, with the
extended determinant definition on pp.~40--42. Apply it on an
\'etale atlas of the gerbe. Functoriality and base-change
compatibility glue the local determinants to a line bundle
$\det(F)$ on $\mathcal X$. Locally a finite locally free
presentation gives
\[
 \det(F)=\bigotimes_i\det(F^i)^{(-1)^i}.
\]
The inertia element $\zeta\in\mu_M$ acts on this line by
$\zeta^{\sum_i(-1)^i\operatorname{rk}F^i}=\zeta^r$.
Consequently $F^{\otimes^{\mathbf L}r}\otimes\det(F)^{-1}$
has weight zero. Weight-zero quasi-coherent complexes on a
$\mu_M$-gerbe are precisely pullbacks from the base: after an
\'etale trivialization this is the weight-zero summand of the
representation category of $\mu_M$, and the identifications
descend. Perfectness is local for this topology, proving the
first assertion.

Use the weight-one deformation $F$ on the gerbe in Perry's proof
of Theorem~1.2, \S6.2, pp.~26--28. At the central fibre its
normalization has character $w_0$ and first component
$w_{0,1}=p_1\theta_0$. Directly, cancellation of the $B$-field
in the tensor--determinant expression gives
\[
 \operatorname{ch}(K_0)
 =\operatorname{ch}(E_0)^r
       e^{-c_1(E_0)}
 =w_0^r e^{-w_{0,1}} .
\]
This also follows from the central root description of the gerbe:
the root-line factors cancel between the $r$ tensor factors and
the determinant. Since $n\ge2$, $w_{0,1}=p_1\theta_0$ contains
no Weil term. Through codimension $n$, the only occurrence of
$t_0$ in $w_0^r$ is
$rP(0)^{r-1}\lambda t_0=r^r\lambda t_0$.
All positive powers of $\theta_0$ annihilate $t_0$.
This proves (S17.1) at the central fibre; flatness of Chern classes
proves it over the connected deformation neighborhood.
Terms in higher codimension, including a possible $t_0^2$ term,
are not used in this formula or in the computation of $c_n$.

Shrink the neighborhood to a smooth affine base $U$.
Its total space $A_U$ is smooth and quasi-projective, and its
relatively very ample bundle $H$ is ample over this affine base.
Stacks Project Lemma~36.37.2, tag 0F8F, represents a perfect
complex by a bounded complex of vector bundles on a scheme
with the resolution property. Thus write
$[K]=[V^+]-[V^-]$. For sufficiently large $m$, choose a finite
surjection $\mathcal O^{\oplus R}\to V^-\otimes H^m$.
Its kernel $J$ is locally free, and
\[
 [V^+\oplus(J\otimes H^{-m})]=[K]+R[H^{-m}].
\]
Adding redundant sections ensures $R\ge1$.
This is the relative version of S-16, using the finite global
generation criterion of Stacks Proposition~28.27.13(7)--(8),
tag 01Q3. Add $a$ trivial summands to make the rank at least $n$,
and twist by a sufficiently large $H^k$ to obtain a globally
generated vector bundle $V$. Its character through codimension
$n$ is $Q(\theta_s)+\mu t_s$.

The identity
\[
 c_u(V)=\exp\!\left(\sum_{j\ge1}
             (-1)^{j-1}(j-1)!\operatorname{ch}_j(V)u^j\right)
\]
gives $c_n(V_s)=C_n\theta_s^n+\beta t_s$.
If $q=\operatorname{rk}V$, the degeneracy cycle of $q-n+1$
general sections is effective of codimension $n$, or empty,
and represents $c_n(V_s)$; see Fulton, \emph{Intersection
Theory}, Theorem~14.4 and Examples~14.4.2--3, pp.~254--256.
Call it $Z_s$. Its Weil component $\beta t_s$ is nonzero,
so the cycle is nonempty. Since $\theta_st_s=0$, its degree is
$C_nv$, which is positive. Let $W_s$ be a complete intersection
of $n$ general members of $|H_s|$; it represents $\theta_s^n$
and has degree $v$. Then
\[
 d\beta t_s=\operatorname{cl}(dZ_s)
                   -\operatorname{cl}(dC_nW_s).
\]
Both cycles are integral and effective of degree $D$.
If $\beta<0$, interchange them. This proves (S17.2) over a
nonempty \'etale image, which is a dense Zariski open subset
of the irreducible base. The fixed-$N$, fixed-$D$ locus of such
signed representatives is Zariski closed by the proper
parameter-space argument in DIV1. It therefore equals the
entire base, preserving these same constants.

Finally, a shift changes the signs of the virtual rank and
Chern character, preserves the self-Ext groups and
semiregularity, and shifts the existing deformation.
This reduces negative rank to positive rank without adding
a new existence hypothesis.
\end{proof}

This construction is used only after the successful twisted
deformation exists. No semiregularity of $K$ or of $V$ is claimed.
The Adams construction DER4 removes the nonzero-rank restriction;
S-17 is a separate, simpler tensor--determinant alternative.

\begin{proposition}[S-18: Canonical induction preserves a full semiregularity kernel]
\label{prop:semireg-induction-kernel}
Let $\mathcal C$ be a smooth proper complex linear enhanced
category with a coherent action of a finite group $G$ by exact
autoequivalences, and assume its equivariant category
$\mathcal C^G$ is smooth proper. Let
\[
 \operatorname{Ind}:\mathcal C\longrightarrow\mathcal C^G
\]
be induction, whose underlying object is
$\bigoplus_{g\in G}gE$, with the canonical permutation
equivariant structure. Induction is injective on every morphism
space. In particular it induces an injection
\[
 \ker\sigma_E\hookrightarrow
       \ker\sigma_{\operatorname{Ind}(E)}.
\]
Thus a nonzero full semiregularity kernel cannot be removed by
canonical induction, including for actions combining
translations and tensoring by degree-zero line bundles on an
abelian variety.
\end{proposition}
\begin{proof}
For $u:E\to E'[j]$, the underlying map of
$\operatorname{Ind}(u)$ is the block diagonal map
$\bigoplus_{g\in G}g(u)$.
If it is zero, compose with the inclusion of the identity
summand of the source and the projection to the identity
summand of the target. The coherence isomorphism for the
identity element identifies this composite with $u$, so $u=0$.
In particular, if $\operatorname{Ind}(u)$ is zero in the
equivariant category, forgetting equivariance gives zero and
the same argument applies. This proves injectivity in every
degree, without requiring the forgetful functor to be full.

Induction is an exact enhanced functor, so naturality of the
full semiregularity morphism gives
\[
 \sigma_{\operatorname{Ind}(E)}(\operatorname{Ind}(u))
       =\operatorname{Ind}_*\sigma_E(u).
\]
This is Perry Remark~2.3, p.~5; Remark~2.5 identifies the
geometric map with Buchweitz--Flenner under HKR.
Consequently a nonzero element of $\ker\sigma_E$ has a
nonzero image lying in the displayed target kernel.
\end{proof}

S-18 rules out semiregularity of the canonical induced lift.
It does not rule out a different equivariant structure on the
same underlying object: weak equivariant semiregularity is
an existence condition over possible lifts, as specified in
Perry Definition~2.6, pp.~5--6.

The qualitative fourfold secant obstruction and its free-translation
variant are prior results of rStation paper~37, version~4,
Theorems~A--C and~$B_G$. The next two applications add
explicit bounds for the supplied Markman constructions.

\begin{proposition}[S-19: The isolated-point genus-four secant seed fails at every parameter]
\label{prop:semireg-markman-genus-four-points}
Consider the first genus-four construction of Markman
Example~8.2.3 in the supplied 31 May 2025 source. Thus $X$ is
a genus-four Jacobian with principal polarization $\theta$,
$d\ge1$, and $G\subset X$ is a cyclic translation group of
order $\ell=d+1$. The $G$-invariant reduced subscheme $Z$ is
the union of $\ell$ translates of $W_2$ in the stated general
position and $2d\ell$ isolated reduced points away from them.
Let $q:X\to Y=X/G$, let $\bar Z=Z/G$, and set
$E=\mathcal I_{\bar Z}$. Then
\[
 \chi(E,E)=8d,\qquad
 \operatorname{ext}^1(E,E)\ge8d+4,\qquad
 \operatorname{ext}^2(E,E)\ge24d+6>28.
 \tag{S19.1}
\]
In particular $E$ is not fully semiregular. The original ideal
$\mathcal I_Z$ is not weakly $G$-semiregular for this translation
action, for any choice of linearization. The same holds for
$\mathcal I_Z(\theta)$ under the action conjugated by tensoring
with $\mathcal O_X(\theta)$.
\end{proposition}
\begin{proof}
The quotient is finite \'etale of degree $\ell$. Its isolated
point set has $2d$ reduced points, and
$q^*\mathcal I_{\bar Z}=\mathcal I_Z$.
Markman's character calculation in Example~8.2.3 gives
\[
 w:=\operatorname{ch}(\mathcal I_Z(\theta))
   =1+\theta-\frac d2\theta^2-\frac d6\theta^3
          +\frac{d^2}{24}\theta^4 .
\]
Since $\int_X\theta^4=24$, direct multiplication yields
$\int_X w^\vee w=8d(d+1)$. Tensoring by the theta line bundle
does not change this Euler pairing. Riemann--Roch and the
\'etale degree formula therefore give $\chi(E,E)=8d$.

We verify the asserted tangent-space lower bound. On a smooth
variety, the endomorphism sheaf of an ideal of codimension at
least two is $\mathcal O$: an endomorphism is multiplication
by a regular function away from that locus, and normality
extends the function. The unit and perfect trace give a split
inclusion
\[
 \mathcal O_Y\longrightarrow R\mathcal Hom(E,E)
       \longrightarrow\mathcal O_Y
\]
whose composite is the virtual rank, namely one. Its complement
has no degree-zero cohomology and no negative cohomology.
Consequently
\[
 \operatorname{Ext}^1(E,E)
 \cong H^1(Y,\mathcal O_Y)
       \oplus H^0(Y,\mathcal Ext^1(E,E)).
\]
At each isolated reduced point the ideal stalk is the maximal
ideal $\mathfrak m$ of a regular local ring $R$ of dimension
four. Apply $\operatorname{Hom}_R(\mathfrak m,-)$ to
$0\to\mathfrak m\to R\to\mathbf C\to0$.
Here $\operatorname{Hom}_R(\mathfrak m,R)=R$ maps to zero in
$\operatorname{Hom}_R(\mathfrak m,\mathbf C)$, and
\[
 \operatorname{Ext}^1_R(\mathfrak m,R)
   =\operatorname{Ext}^2_R(\mathbf C,R)=0.
\]
Thus
$\operatorname{Ext}^1_R(\mathfrak m,\mathfrak m)
\cong(\mathfrak m/\mathfrak m^2)^\vee$ has dimension four.
The isolated supports contribute independent skyscraper
summands. Since $h^1(\mathcal O_Y)=4$, this proves
$\operatorname{ext}^1(E,E)\ge4+8d$.

The sheaf $E$ is simple. Serre duality on the abelian fourfold
gives $\operatorname{ext}^4=\operatorname{ext}^0=1$ and
$\operatorname{ext}^3=\operatorname{ext}^1$; there are no
self-Ext groups outside degrees zero through four. Hence
\[
 \operatorname{ext}^2(E,E)
 =\chi(E,E)-2+2\operatorname{ext}^1(E,E)
 \ge24d+6.
\]
The full semiregularity target on an abelian fourfold has
dimension
\[
 h^2(\mathcal O_Y)+h^3(\Omega_Y^1)+h^4(\Omega_Y^2)
 =6+16+6=28,
\]
proving failure of injectivity.

Equivariant descent identifies
$\operatorname{Perf}(X)^G$ with $\operatorname{Perf}(Y)$
and the canonical lift of $\mathcal I_Z$ with $E$.
Moreover $\operatorname{Aut}(\mathcal I_Z)=\mathbf C^*$,
so any two $G$-linearizations differ by a character of $G$.
Their descents differ by tensoring with the associated
degree-zero torsion line bundle on $Y$. Such a tensor
equivalence preserves the full kernel by S-10.
Thus no linearization is semiregular in the invariant category.
Conjugating this action by the theta tensor equivalence gives
the last assertion; the conjugated action consists of target
translations and tensoring by degree-zero line bundles.
\end{proof}

\begin{proposition}[S-20: The fourfold adjunction test excludes the partial-normalization seed]
\label{prop:semireg-fourfold-partial-normalization}
On an abelian fourfold $Y$, let $E$ be perfect and let
$c_E^2:HT^2(Y)\to H\Omega^{-2}(Y)$ be contraction with
$\operatorname{ch}(E)$, in the conventions of the supplied
032 companion, \S5.1. If $E$ is fully semiregular, then
\[
 \operatorname{ext}^2(E,E)=\operatorname{rank}c_E^2.
 \tag{S20.1}
\]
For the cyclic partial-normalization construction in Markman
Example~8.2.3, let $\bar E$ be the descended complex
\[
 [\mathcal O_Y\longrightarrow
       \bar\nu_*\mathcal O_{\widetilde Z/G}]
\]
in degrees zero and one, where four of the six transverse
meeting points for each pair of translated $W_2$ surfaces
are normalized, and the collection of choices is $G$-invariant.
Then, for every $d\ge1$,
\[
 \operatorname{rank}c_{\bar E}^2=12,
 \qquad
 \operatorname{ext}^2(\bar E,\bar E)\ge8d+6>12.
 \tag{S20.2}
\]
Thus this complex is not fully semiregular, and its pullback
is not weakly $G$-semiregular for the specified free translation
action, with any linearization. In fact
\[
 \dim\ker\sigma_{\bar E}\ge4d-3.
 \tag{S20.3}
\]
\end{proposition}
\begin{proof}
The adjunction criterion is the previously known CY4 test
(rStation paper~37, version~4, Theorem~A). In the conventions
used here, 032 equations (5.20)--(5.24) identify $\sigma_E$
with the Serre transpose of $\operatorname{ob}_E$, and
Proposition~5.1, equation (5.9), gives
$\sigma_E\operatorname{ob}_E=c_E^2$. The Todd correction is
one on an abelian variety. Thus injectivity of $\sigma_E$
forces surjectivity of $\operatorname{ob}_E$ and proves
(S20.1).

The imaginary-quadratic rank calculation is likewise the
one used in that prior criterion. For the character
$w=\alpha+\beta$ of S-19, put $U=H^1(\mathcal O_Y)$,
$A=M\mathbin{\lrcorner}\theta$,
$B=\tfrac12\nu\mathbin{\lrcorner}\theta^2$, and
$D=\eta-dB$, where
$(\eta,M,\nu)\in\wedge^2U\oplus(U\otimes H^0(T_Y))
\oplus\wedge^2H^0(T_Y)$.
The first two output blocks are
$D+A$ and $\theta\wedge(D-dA)$.
Their coefficient determinant is $-(d+1)$, and every output
factors through the two six-dimensional variables $(D,A)$;
hence the rank is twelve. Multiplication of the character
by $e^{b\theta}$ conjugates this calculation by the invertible
shear $\eta'=\eta+bA+b^2B$, $A'=A+2bB$.
This verifies the convention needed for the rational
polarization on the quotient. The complete convention
verification is retained in the working report.

It remains to verify the actual object. Write
$D_0=\bar\nu_*\mathcal O_{\widetilde Z/G}$.
It is a pure sheaf of dimension two: it injects into the
finite pushforward of the structure sheaves of the smooth
normalization components. The map $\mathcal O_Y\to D_0$
has kernel $I=\mathcal I_{\bar Z_{\mathrm{surf}}}$ and
zero-dimensional cokernel $Q$. Thus $\bar E$ has cohomology
$I$ in degree zero and $Q$ in degree one.
Since $I$ is torsion free,
$\operatorname{Hom}(Q,I)=0$. The truncation triangles show
that this implies
$\operatorname{Ext}^{j}(\bar E,\bar E)=0$ for $j<0$.
They also give
$\operatorname{Hom}(\bar E,D_0[-1])=0$, because
$\operatorname{Hom}(Q,D_0)=0$ by purity.
Codimension two of $\operatorname{Supp}D_0$ on the smooth
fourfold gives
$\operatorname{Hom}(D_0,\mathcal O_Y)
=\operatorname{Ext}^1(D_0,\mathcal O_Y)=0$.
Applying $\operatorname{Hom}(-,\mathcal O_Y)$ to
$\bar E\to\mathcal O_Y\to D_0$ therefore yields
$\operatorname{Hom}(\bar E,\mathcal O_Y)=\mathbf C$.
Applying $\operatorname{Hom}(\bar E,-)$ and using the other
vanishing embeds $\operatorname{End}(\bar E)$ into this
one-dimensional space. The identity is nonzero, so
$\bar E$ is simple.

Normalizing a transverse meeting point of two smooth
surfaces adds one skyscraper class to the structure sheaf.
The chosen four points for each of the $\binom\ell2$ pairs
therefore add $4\binom\ell2=2d\ell$ point classes upstairs.
This is exactly the number of isolated points in S-19.
Consequently the partial-normalization complex has the same
K-character as that first ideal seed, as asserted in
Markman Example~8.2.3. The quotient has
$\chi(\bar E,\bar E)=8d$, virtual rank one, and character
$e^{-\theta}w$ in the rational quotient polarization.
The preceding linear calculation gives
$\operatorname{rank}c_{\bar E}^2=12$.
The unit--trace composite is again rank one, so it injects
$H^1(\mathcal O_Y)$ into
$\operatorname{Ext}^1(\bar E,\bar E)$.
The negative-Ext vanishing, Serre duality, and simplicity
now give
\[
 \operatorname{ext}^2(\bar E,\bar E)
 =8d-2+2\operatorname{ext}^1(\bar E,\bar E)
 \ge8d+6>12.
\]
This contradicts (S20.1) if semiregularity is assumed.

For the quantitative assertion, put
$m=\operatorname{ext}^2(\bar E,\bar E)$ and
$a=\operatorname{rank}\sigma_{\bar E}
=\operatorname{rank}\operatorname{ob}_{\bar E}$.
The rank inequality for the composition gives
$12\ge2a-m$. Hence
\[
 \dim\ker\sigma_{\bar E}=m-a
 \ge\frac{m-12}{2}\ge4d-3.
\]
The argument proving simplicity also applies upstairs.
All linearizations therefore differ by characters of $G$,
and their descents by torsion line bundles, as in S-19.
Tensor equivalence preserves the nonzero full kernel,
proving the weak-equivariant assertion for this action.
\end{proof}

For both S-19 and S-20, taking an external product with a
nonzero perfect object preserves the displayed kernel:
the functor sends $u$ to $u\boxtimes\operatorname{id}$,
which is nonzero by K\"unneth, and the full trace is functorial.
The same is true after a derived equivalence.
Thus the canonical product and Orlov-transform route does not
repair these particular factors. This statement does not
exclude a different object or a different finite target action.

\begin{proposition}[S-21: All finite target actions obey the fourfold secant obstruction]
\label{prop:semireg-all-finite-target-actions}
Let $A$ be a complex abelian fourfold, $\mathcal C=
\operatorname{Perf}(A)$, and let a finite subgroup
$G\subset A\times\widehat A$ have a specified coherent action
on $\mathcal C$ by the corresponding translations and
tensor products with degree-zero line bundles. In particular,
no nonidentity group element represents the pair $(0,\mathcal O_A)$.
For a lift $\widetilde E\in\mathcal C^G$ of a perfect object
$E\in\mathcal C$, suppose
$\operatorname{Ext}^{j}_{\mathcal C^G}(\widetilde E,
\widetilde E)=0$ for $j<0$. Suppose further that
\[
 e^B\operatorname{ch}(E)=a\alpha+b\beta,
 \quad
 \alpha=\frac{e^{s\theta}+e^{-s\theta}}2,
 \quad
 \beta=\frac{e^{s\theta}-e^{-s\theta}}{2s},
 \quad s^2=-d,
 \tag{S21.1}
\]
where $d>0$, $\theta$ is a rational ample class,
$B$ is a rational algebraic divisor class, and
$a,b\in\mathbf Q$ are not both zero.
If $\widetilde E$ is semiregular in $\mathcal C^G$, then
\[
 \operatorname{ext}^2_{\mathcal C^G}(\widetilde E,
              \widetilde E)=12,
 \qquad
 \operatorname{ext}^0_{\mathcal C^G}(\widetilde E,
              \widetilde E)
 \ge 2+\frac{\chi_A(E,E)}{2|G|}>2,
 \tag{S21.2}
\]
where
\[
 \chi_A(E,E)=8d(da^2+b^2)\frac{\int_A\theta^4}{24}>0.
 \tag{S21.3}
\]
Consequently no such lift that is simple in the equivariant
category is semiregular. In particular, an underlying simple
perfect object with (S21.1) and no negative self-Ext admits no
weakly semiregular lift for any such finite target action.
The conclusion allows arbitrary mixtures of translations
and degree-zero tensor operations; it does not require a
Fourier--Mukai model on which the whole group acts freely
by translations.
\end{proposition}
\begin{proof}
Write $\phi_g$ for the functor associated to $g=(x,P)$.
The invariant category is smooth and proper in characteristic
zero; see Perry 2026, \S5, opening paragraph, p.~12,
and Bayer--Perry Proposition~3.15, author PDF pp.~20--21.
First we verify that all nonidentity coefficient sectors vanish:
\[
 HH^*(\mathcal C,\phi_g)=0\qquad(g\ne1).
 \tag{S21.4}
\]
The coefficient space is the derived Hom from the diagonal
kernel to the kernel on the graph of translation by $x$,
with its line-bundle coefficient. This is the kernel
description in Perry, \emph{Hochschild cohomology and group
actions}, arXiv:1807.09268v2, Remark~4.1 and the proof of
Corollary~4.8, pp.~9 and~12. If $x\ne0$, the graph is
disjoint from the diagonal, so the derived Hom is zero.
If $x=0$ and $g\ne1$, then $P\ne\mathcal O_A$ and
the sheafified HKR decomposition identifies the coefficient
groups with sums of
$H^q(A,\wedge^pT_A\otimes P)$.
They vanish, because $T_A$ is trivial and a nontrivial
degree-zero line bundle on an abelian variety has no
cohomology. The latter also follows from the transform of
$\mathcal O_A$ in Mukai Theorem~2.2, p.~156, by derived
base change away from the origin.

Translations and degree-zero tensor operations act trivially
on $HH^*(A)$ and $HH_*(A)$: under HKR the former fix invariant
forms and vector fields, and the latter have zero Atiyah
class and Chern character one. Perry's Theorem~4.4 and
Remark~4.5, pp.~10--12, now give
\[
 HH^*(\mathcal C^G)\cong HH^*(\mathcal C).
 \tag{S21.5}
\]
The identity-sector inclusion is compatible with the
Hochschild algebra and with evaluation on equivariant objects.
This conclusion uses the specified coherent action, not
merely an unlinearized list of isomorphic functors.

The invariant category has Serre functor $[4]$.
To check this directly, the ordinary Serre trace pairing on
$\operatorname{Perf}(A)$ is invariant under translations
and degree-zero tensor products: translations preserve the
chosen invariant volume form, and tensoring preserves
perfect trace. Conjugation by the linearizations also
preserves trace. Taking $G$-invariants in the Serre duality
isomorphism remains a perfect duality, since $|G|$ is
invertible. This is exactly the invariant-duality argument
in the proof of Perry Lemma~6.5, equations (6.1)--(6.3),
p.~17; the proof there treats the general finite-group
Serre descent before its stated cyclic specialization.

Perry Definition~6.1 and Remark~6.2, p.~16, identify
Hochschild homology with Serre-coefficient Hochschild
cohomology. Converting his cohomological grading to the
usual homological convention, a category with Serre functor
$[4]$ has $HH_j\cong HH^{4-j}$ as vector spaces.
Thus (S21.5) gives equal finite dimensions for
$HH_j(\mathcal C^G)$ and $HH_j(\mathcal C)$, in particular
for $j=0,-2$.

Let $F:\mathcal C^G\to\mathcal C$ be forgetting and
$I:\mathcal C\to\mathcal C^G$ induction. These functors
are mutually left and right adjoint by Perry Lemma~3.3,
p.~7, and $FI\cong\bigoplus_{g\in G}\phi_g$.
Functoriality and additivity of Hochschild homology imply
\[
 F_*I_*=|G|\operatorname{id}_{HH_*(\mathcal C)}.
\]
One may read this additivity from Shklyarov, arXiv:0710.1937v3,
Proposition~3.2 and Theorem~3.4, pp.~14--16, by applying
the latter to a sum of functor kernels.
Equal dimensions then show that $F_*$ is an isomorphism
in degrees zero and minus two, and that
\[
 I_*F_*=|G|\operatorname{id}_{HH_j(\mathcal C^G)}
       \qquad(j=0,-2).
 \tag{S21.6}
\]
In particular, forgetting cannot hide a semiregularity
kernel in a nonidentity Hochschild sector.

We also need the exact Euler scaling, not an estimate by
averaging dimensions. Hochschild Chern characters commute
with functors, so (S21.6) implies
$\operatorname{ch}_{HH}(IF\widetilde E)
=|G|\operatorname{ch}_{HH}(\widetilde E)$.
The categorical Riemann--Roch formula makes the Euler
pairing linear in these characters: its precise source is
Shklyarov Theorem~3.5, p.~17, and its Hom formulation
equations (3.7) and (3.9), pp.~18--19.
The formula applies here via a proper DG-algebra model:
induction of a generator of $\operatorname{Perf}(A)$ is
a generator of $\mathcal C^G$, by adjunction and the
averaging retraction; its endomorphism algebra is proper.
Using the right adjunction $F\dashv I$, we obtain
\[
 \chi_A(E,E)
 =\chi_{\mathcal C^G}(\widetilde E,IF\widetilde E)
 =|G|\chi_{\mathcal C^G}(\widetilde E,\widetilde E).
 \tag{S21.7}
\]

Write $V_i=\operatorname{Ext}^i_A(E,E)^G$, which is the
degree-$i$ self-Ext in $\mathcal C^G$.
Perry 2026 Remarks~2.3 and~2.5, p.~5, give the commuting
semiregularity diagram for forgetting. Since $F_*$ in
degree minus two is an isomorphism, semiregularity of the
lift is equivalent to injectivity of
$\sigma_E|_{V_2}$.
The ordinary evaluation maps
$\operatorname{ob}_E^i:HH^i(A)\to\operatorname{Ext}^i(E,E)$
have image in $V_i$, because $G$ acts trivially on
$HH^i(A)$ and evaluation is natural.
The fourfold adjunction of S-20 restricts to the perfect
invariant Serre pairing on $V_2$. Consequently, if the
lift is semiregular, evaluation in degree two is onto
$V_2$ and
\[
 \dim V_2=\operatorname{rank}c_E^2.
 \tag{S21.8}
\]
Similarly $\dim V_1\ge\operatorname{rank}c_E^1$.
For clarity, the latter needs only the degree-one trace
identity. For $\gamma:\mathcal O_\Delta\to
\mathcal O_\Delta[1]$, naturality of the same adjunction
unit as in 032 equation (5.23) gives
\[
 (1_E\otimes L\Delta^*\gamma)AH_E
     =AH_E[1]\chi_E(\gamma).
\]
Take perfect trace with coefficients and apply HKR as in
(5.20)--(5.24). This gives
$\sigma_E^1\operatorname{ob}_E^1=c_E^1$,
so the claimed rank bound follows in the invariant image.

We calculate both contraction ranks explicitly. A
$B$-field factor $e^{-B}$ does not change them:
conjugation by $e^B$ replaces the degree-one contraction
operator $\iota_v$ by $\iota_v+\iota_vB\wedge$ and fixes
wedge operators. Since $B$ has type $(1,1)$, this is an
invertible shear of $HT^1(A)=U\oplus T$ and hence of its
exterior algebra, while $e^B$ preserves each Hodge-difference
space on the target.
For $w=a\alpha+b\beta$, the degree-two calculation in S-20
has first two blocks
\[
 aD+bA,\qquad\theta\wedge(bD-adA).
\]
The determinant is $-(a^2d+b^2)\ne0$, giving
$\operatorname{rank}c_E^2=12$.
In degree one, put $L=v\mathbin{\lrcorner}\theta\in U$
for $(\eta,v)\in U\oplus T$.
The first two blocks are
\[
 a\eta+bL,\qquad\theta\wedge(b\eta-adL).
\]
The same determinant is nonzero, and wedge by $\theta$
is injective on $U$. Thus $\operatorname{rank}c_E^1=8$.

Finally, $B$ cancels from $\operatorname{ch}(E)^\vee
\operatorname{ch}(E)$. If
$c_\pm=(a\pm b/s)/2$, then
$c_+c_-=(a^2+b^2/d)/4$.
The top term of $w^\vee w$ is
$2c_+c_-(2s)^4\theta^4/24$, proving (S21.3).
Put $h=\dim V_0$. The assumed negative-Ext vanishing and
invariant Serre duality imply vanishing above degree four,
$\dim V_4=h$, and $\dim V_3=\dim V_1$.
Equations (S21.7)--(S21.8) therefore imply
\[
 12=\dim V_2
 =\frac{\chi_A(E,E)}{|G|}-2h+2\dim V_1
 \ge\frac{\chi_A(E,E)}{|G|}-2h+16.
\]
Rearranging gives (S21.2), including its contradiction
when $h=1$. The proof applies separately to every coherent
linearization; none was assumed canonical.
\end{proof}

The earlier fourfold criterion and its free-translation
version are recorded in rStation paper~37, version~4,
Theorems~A, B and $B_G$. S-21 adds the coefficient-sector
and induction argument for arbitrary effective finite
subgroups of $A\times\widehat A$, and permits simplicity
only after equivariantization. It excludes this ample
imaginary-quadratic secant character for equivariantly
simple seeds; it does not exclude arbitrary corrected
characters on the final abelian eightfold.

\begin{proposition}[S-22: The balanced real-multiplication character has no rank-one torsion-free realization]
\label{prop:semireg-rm-balanced-obstruction}
Let $A$ be an abelian fourfold with a Rosati-invariant action
of a real quadratic field $F$, of multiplicity two at each
real embedding on $H^0(T_A)$. Let $\theta$ be an $F$-compatible
polarization and write $\theta=\theta_1+\theta_2$ for its two
real eigenspace components. Put
\[
 v_q=\left(1-\frac q2\theta_1^2\right)
       \left(1-\frac q2\theta_2^2\right),
 \qquad q>0,
 \qquad D=\frac{\int_A\theta^4}{24}.
 \tag{S22.1}
\]
There is no nonzero effective surface cycle whose class
belongs to $\mathbf R\theta_1^2+\mathbf R\theta_2^2$.
Consequently there is no rank-one torsion-free coherent
sheaf $E$ with $e^B\operatorname{ch}(E)=v_q$, for any
rational algebraic divisor class $B$.

The contraction ranks of $v_q$ are
\[
 \operatorname{rank}c_{v_q}^1=8,
 \qquad \operatorname{rank}c_{v_q}^2=18,
 \qquad \chi(Nv_q,Nv_q)=4Dq^2N^2.
 \tag{S22.2}
\]
If a perfect object with normalized character $Nv_q$,
$N\in\mathbf Z_{>0}$, has a gluable, semiregular, equivariantly simple
lift for an effective coherent finite target action as in
S-21, then
\[
 \frac{4Dq^2N^2}{|G|}\in\{2,4\},
 \qquad |G|\in\{Dq^2N^2,\,2Dq^2N^2\}.
 \tag{S22.3}
\]
Here only integer group orders are possible. In particular,
when $G=1$, $q=1$, and $\theta$ is integral, a simple pure
coherent sheaf with normalized character $Nv_1$ cannot be
fully semiregular, for any $N>0$.
\end{proposition}
\begin{proof}
Let $V_1\oplus V_2=H^0(T_A)$ be the two complex
$F$-eigenspaces, each of dimension two, and put
$W_i=V_i^*$. The invariant form representing $\theta_i$
is supported on $V_i$ and has rank two. In particular
$\theta_i^3=0$.
For every mixed holomorphic two-form
$\omega\in W_1\otimes W_2\subset H^0(\Omega_A^2)$,
\[
 (a\theta_1^2+b\theta_2^2)\wedge\omega\wedge\overline\omega=0
 \tag{S22.4}
\]
for all real $a,b$: the first term has excessive holomorphic
and antiholomorphic degree in $V_1$, and the second in $V_2$.

Suppose $Z=\sum n_jZ_j$, $n_j>0$, were such an effective
cycle. On a resolution of each component $Z_j$, the integral
of $\omega\wedge\overline\omega$ is nonnegative, and vanishes
only when the pullback of $\omega$ vanishes. Equation
(S22.4) and positivity therefore force every mixed two-form
to vanish on every component. At a smooth point of a
component, its two-dimensional tangent plane $T$ consequently
has no mixed component in its Pl\"ucker bivector.
Write that bivector as $u_1+u_2$, with
$u_i\in\wedge^2V_i$. Its square is $2u_1\wedge u_2$.
Since it is decomposable, this square vanishes, so one
$u_i$ is zero. Thus $T=V_1$ or $T=V_2$.
On each irreducible component the Gauss map is therefore
constant, with one of these two values.

Here is an explicit Albanese argument for the required
consequence of that constancy. Let $f:Y\to Z_j\subset A$
be a smooth projective resolution, translated so that a
chosen point maps to zero. Every invariant one-form
annihilating $V_i$ pulls back to zero on the dense smooth
locus and hence on $Y$. The induced homomorphism
$\operatorname{Alb}(Y)\to A$ therefore has tangent image
contained in $V_i$. Its image is a positive-dimensional
abelian subvariety, because $f$ is nonconstant. But $V_i$
cannot contain the tangent space of such a subvariety.
Indeed, choose $u\in F\setminus\mathbf Q$. Its action on
$H_1(A,\mathbf Q)$ is rational, while on the underlying
real eigenspace of $V_i$ it is multiplication by the
irrational number $\sigma_i(u)$. That eigenspace contains
no nonzero rational vector: otherwise a nonzero rational
coordinate of $z$ in $uz=\sigma_i(u)z$ would make
$\sigma_i(u)$ rational. The nonzero rational homology of
an abelian subvariety gives a contradiction.
This proves the assertion about effective cycles, without
assuming that $A$ itself is simple.

For a rank-one torsion-free $E$, its double dual is a line
bundle $L$, and $L/E$ has support of codimension at least
two. Its codimension-two character is the effective
fundamental cycle of its surface support. The degree-one
part of $e^B\operatorname{ch}(E)=v_q$ gives
$c_1(L)=-B$. The degree-two part thus identifies that
effective cycle with
$q(\theta_1^2+\theta_2^2)/2$.
This class is nonzero, since its pairing with $\theta^2$
is $4Dq>0$, contradicting the first part.

For completeness, the ranks in (S22.2) can be checked on
the complex Hodge decomposition, without supposing that
$A$ is a product. It factors as the exterior tensor product
of the two two-dimensional eigenspaces. For
$v_i=1-q\theta_i^2/2$, the contraction ranks in degrees
zero, one and two are respectively $1,4,1$.
In degree one, wedge operators on $H^1(\mathcal O)$ and
contraction operators on $H^0(T)$ have independent images
of dimension two each. In degree two the only nonzero
output is the one-dimensional $H^2(\mathcal O)$ block;
the middle wedge--contraction operators vanish for degree
reasons. Consequently the product has degree-one rank
$4+4=8$ and degree-two rank $1+4\cdot4+1=18$.
The three degree-two images lie in distinct pairs of Hodge
differences on the two eigenspace factors, so their sum is
direct. The $B$-field shear of S-21 preserves these ranks.
Since $v_q^\vee=v_q$, the top term of $v_q^2$ is
$q^2\theta_1^2\theta_2^2$. Also
$\int_A\theta_1^2\theta_2^2=4D$, proving the Euler formula.

The all-group argument in S-21 applies to every character.
For a semiregular simple lift here it gives $e_2^G=18$,
$e_1^G\ge8$, and
\[
 \frac{\chi_A(E,E)}{|G|}=2-2e_1^G+18.
\]
The left side is strictly positive, so it is either $2$ or
$4$, proving (S22.3). If $G=1$, $q=1$, and $\theta$ is
integral, then $D\in\mathbf Z_{>0}$ and
$4DN^2\le4$, forcing $D=N=1$.
A pure rank-one sheaf on the smooth integral fourfold is
torsion free, so the first part rules out this remaining
case.
\end{proof}

For $q=1$, (S22.1) is precisely
$(\alpha+\widetilde\alpha)/2$ in the real-multiplication
space of Markman, arXiv:2509.23079, \S11.2.1, pp.~40--41:
\[
 \frac{\cos(\theta_1+\theta_2)+
             \cos(\theta_1-\theta_2)}2
       =\cos(\theta_1)\cos(\theta_2).
\]
Its numerical eligibility was already recorded in the
real-multiplication remark of rStation paper~37, version~4.
S-22 adds the rank-one effectivity obstruction. It does
not exclude a higher-rank weakly equivariant object, or a
rank-one perfect complex which is not a pure sheaf.

\begin{proposition}[S-23: A nonsimple semiregular control on an abelian fourfold]
\label{prop:semireg-nonsimple-control}
For every abelian fourfold $A$ and point $p\in A$, the
perfect sheaf
$E=\mathcal O_A\oplus\mathcal O_p$ is fully semiregular,
with self-Ext dimensions
\[
 (e_0,e_1,e_2,e_3,e_4)=(3,8,12,8,3),
 \qquad \chi(E,E)=2.
 \tag{S23.1}
\]
Its endomorphism algebra is triangular and contains a
nonzero nilpotent. Hence full semiregularity does not
imply simplicity or a semisimple endomorphism algebra.
\end{proposition}
\begin{proof}
The self-Ext algebras of $\mathcal O_A$ and $\mathcal O_p$
are respectively $H^*(\mathcal O_A)$ and
$\wedge^*T_pA$, the latter by a regular local Koszul
resolution. The cross groups are
\[
 \operatorname{Ext}^i(\mathcal O_A,\mathcal O_p)=
 \begin{cases}\mathbf C&i=0,\\0&i\ne0,\end{cases}
\]
and, by Serre duality, the opposite cross groups are
one-dimensional in degree four only. This proves (S23.1)
and identifies the off-diagonal endomorphism
$\mathcal O_A\to\mathcal O_p$ as a nonzero square-zero
endomorphism of $E$.

The Atiyah class of $\mathcal O_A$ is zero, so its
semiregularity map identifies its degree-two self-Ext with
$H^2(\mathcal O_A)$. For the point, the local Koszul
calculation identifies evaluation of a global bivector
with its value in $\wedge^2T_pA$. Thus evaluation is onto
its six-dimensional degree-two self-Ext.
Contraction with $\operatorname{ch}(\mathcal O_p)=[p]$
identifies this bivector space with
$H^4(\Omega_A^2)$, also of dimension six; the other
Hochschild degree-two blocks annihilate $[p]$ for degree
reasons. By 032 Proposition~5.1, equation (5.9), the full
semiregularity map of the point is therefore injective,
with image exactly this last Hodge block.
There are no cross degree-two self-Ext groups for $E$.
Perfect trace on its two diagonal summands has the two
images just described, which are disjoint Hodge blocks.
Their direct sum is injective. This proves full
semiregularity of $E$.
\end{proof}

\begin{proposition}[S-24: Iterated self-extensions do not repair the positive-Euler secant seed]
\label{prop:semireg-iterated-extensions}
Let $A$, $G$ and $\mathcal C^G$ be as in S-21, and let
$\mathcal H$ be the heart of a bounded t-structure on
$\mathcal C^G$. Let $\widetilde F\in\mathcal H$ satisfy
$\operatorname{End}_{\mathcal C^G}(\widetilde F)=\mathbf C$,
with underlying normalized imaginary-quadratic secant
character as in (S21.1), and suppose
\[
 k:=\chi_{\mathcal C^G}(\widetilde F,\widetilde F)\ge2.
 \tag{S24.1}
\]
If $\widetilde E\in\mathcal H$ admits a finite filtration
of length $m>0$ whose every successive quotient is
$\widetilde F$, then $\widetilde E$ is not semiregular.
In particular no iterated self-extension of the quotient
ideal in S-19, in its equivariant coherent heart, repairs
its semiregularity failure: there $k=8d$.
\end{proposition}
\begin{proof}
We first prove the precise endomorphism bound. If $X,Y$
have filtrations of lengths $a,b$ with factors
$\widetilde F$, then
\[
 \dim\operatorname{Hom}_{\mathcal H}(X,Y)\le ab.
 \tag{S24.2}
\]
For $X=\widetilde F$, induction on a filtration of $Y$
and left exactness of $\operatorname{Hom}(\widetilde F,-)$
give the bound $b$, since
$\operatorname{End}(\widetilde F)=\mathbf C$.
Then use a short exact sequence
$0\to X'\to X\to\widetilde F\to0$ and left exactness
of $\operatorname{Hom}(-,Y)$ to obtain
$\dim\operatorname{Hom}(X,Y)\le b+(a-1)b$.
This proves (S24.2), hence $e_0^G(\widetilde E)\le m^2$.

In K-theory $[\widetilde E]=m[\widetilde F]$.
Bilinearity of the Euler pairing gives
$\chi_G(\widetilde E,\widetilde E)=m^2k$.
Its character is multiplied by $m$, which does not
change either contraction rank, namely $8$ and $12$.
Objects of a heart have no negative self-Ext.
If $\widetilde E$ were semiregular, S-21 would therefore
give $e_2^G=12$, $e_1^G\ge8$, and invariant Serre duality
would imply
\[
 12=m^2k-2e_0^G+2e_1^G
 \ge m^2(k-2)+16\ge16,
\]
a contradiction. No splitting of the filtration, and no
injection of the original trace kernel into that of an
extension, was assumed.
\end{proof}

The hypothesis (S24.1) is essential to this estimate.
S-24 concerns a fixed factor with scalar endomorphisms and a fixed coherent
action; it does not exclude arbitrary cones, constituents
with different numerical characters, or a newly chosen
finite action not preserving such a filtration.

\begin{proposition}[S-25: Correlated finite target actions do not repair a simple secant product]
\label{prop:semireg-correlated-products}
Let $X$ be an abelian fourfold and $Y$ an abelian variety.
Let $F\in\operatorname{Perf}(X)$ and
$K\in\operatorname{Perf}(Y)$ satisfy
\[
 \operatorname{End}(F)=\operatorname{End}(K)=\mathbf C,
 \qquad \operatorname{Ext}^{j}(F,F)=
        \operatorname{Ext}^{j}(K,K)=0\quad(j<0).
 \tag{S25.1}
\]
Assume that $F$ has the normalized ample
imaginary-quadratic secant character of (S21.1).
For any derived equivalence
$\Phi:\operatorname{Perf}(X\times Y)\simeq
\operatorname{Perf}(B)$ to an abelian variety $B$, put
$E=\Phi(F\boxtimes K)$.
Then $E$ admits no semiregular lift for any effective
finite subgroup $G\subset B\times\widehat B$ with a
specified coherent translation--tensor action.
The group need not be a product of factor actions.

In particular, changing the finite target action cannot
make the actual genus-four secant-square objects
$\Phi(F_1\boxtimes F_2^\vee)$ weakly semiregular when the
factors are simple and gluable. The assertion also applies
when the second factor has a different Chern character;
only the first factor is required to lie in the ample
imaginary-quadratic secant plane.
\end{proposition}
\begin{proof}
We spell out the coherence issue, since projective factor
linearizations need not inherit the given product cocycle.

First, the identity-sector argument of S-21 works on an
abelian variety $Z$ of any dimension $g$. For an effective
finite translation--tensor action $H$, all nonidentity
Hochschild coefficient sectors vanish, and $H$ acts
trivially on the identity sector. The invariant category
is smooth proper with Serre functor $[g]$. Perry's
Hochschild-cohomology formula and the identification
$HH_j\cong HH^{g-j}$ give equality of homology dimensions.
Induction and forgetting satisfy
$\operatorname{Forg}_*\operatorname{Ind}_*=|H|$.
Consequently
\[
 HH_{-2}(\operatorname{Perf}(Z)^H)
   \longrightarrow HH_{-2}(Z)
 \quad\hbox{is an isomorphism}.
 \tag{S25.2}
\]
Thus a nonzero invariant element of the ordinary full
trace kernel also gives a nonzero equivariant trace kernel.
This uses the same primary inputs as S-21 and is
independent of any choice of scalar multiplication cocycle.

Next, a finite subgroup of $Z\times\widehat Z$ stabilizing
a simple gluable object $Q$ can be given a coherent action
for which $Q$ has a lift. Choose representative functors
$\phi_h$ and isomorphisms $\lambda_h:\phi_hQ\to Q$,
with the identity choices normalized. Natural isomorphisms
$\phi_g\phi_h\to\phi_{gh}$ form a scalar torsor,
since $HH^0(Z)=\mathbf C$. Their evaluations at $Q$ also
form a scalar torsor, since $\operatorname{End}(Q)=\mathbf C$,
and evaluation is an isomorphism between those torsors.
There is therefore a unique multiplication isomorphism
$\mu_{g,h}$ satisfying
\[
 \lambda_{gh}(\mu_{g,h})_Q
       =\lambda_g\,\phi_g(\lambda_h).
 \tag{S25.3}
\]
Both associativity composites agree after evaluation at
$Q$, by (S25.3) and naturality. They consequently agree as natural
transformations, since evaluation is faithful on scalars.
This produces a unital two-categorical action and a
compatible lift of $Q$.
There is no missing higher coherence: $HH^j(Z)=0$ for
$j<0$, and $\operatorname{Ext}^j(Q,Q)=0$ for $j<0$.
Bayer--Perry Lemma~3.3 and Corollary~3.4, pp.~15--16,
identify two-categorical lifts with enhanced lifts in
this situation. Equivalently, the spaces of natural
isomorphisms and of isomorphisms of the framed object
have no higher homotopy groups; their scalar torsors
cancel in the framed automorphism groupoid.
Changing any $\lambda_h$ or the multiplication by
scalars does not change its conjugation action on
$\operatorname{Ext}^*(Q,Q)$.

By Orlov Theorem~2.10 and Corollary~2.13, author PDF
pp.~9--12, conjugating by $\Phi$ identifies
$B\times\widehat B$ with
$(X\times Y)\times\widehat{(X\times Y)}$ and sends
translation--tensor operations to translation--tensor
operations. Transport the given action and any proposed
lift of $E$ across $\Phi$. We obtain an effective finite
subgroup
\[
 H\subset (X\times\widehat X)\times(Y\times\widehat Y)
\]
with a lift of $F\boxtimes K$. Let $H_X,H_Y$ be its
projections. Every projected functor stabilizes its
corresponding factor. Here is a proof valid even for
perfect complexes, without a common-heart assumption.
For $h\in H$ write its component functors as $\phi$ and
$\psi$, so
\[
 \phi F\boxtimes\psi K\simeq F\boxtimes K.
 \tag{S25.4}
\]
Choose a point of $Y$ at which the derived fiber of $K$
is nonzero. Restriction of (S25.4) gives an isomorphism
between finite direct sums of shifts of $\phi F$ and
finite direct sums of shifts of $F$, with both sums
nonzero. A diagonal component of the composite of this
isomorphism and its inverse is the identity of a shift
of $F$. Since its endomorphisms are scalar, at least one
summand of this composite is a nonzero scalar. It makes
that shift of $F$ a retract of a shift of $\phi F$.
The latter also has scalar endomorphisms, so the retraction
is an isomorphism. Thus $\phi F\simeq F[k]$ for some
integer $k$. Translations and degree-zero tensor products
preserve the least and greatest nonzero standard
cohomological degrees, so boundedness and nonvanishing
force $k=0$. The same argument with the factors exchanged
shows $\psi K\simeq K$.

Use (S25.3) to give $H_X$ a coherent action and a lift of
$F$. This group is effective, and the lift has scalar
endomorphisms. S-21 therefore implies that it is not
semiregular. By (S25.2), there exists
\[
 0\ne u\in
 \bigl(\ker(\sigma_F:\operatorname{Ext}^2(F,F)
                         \to HH_{-2}(X))\bigr)^{H_X}.
 \tag{S25.5}
\]
The action on this space is the canonical conjugation
action determined by the stabilizing functor classes;
it is unchanged by the scalar choices used to arrange
coherence.

The class $u\boxtimes\operatorname{id}_K$ is nonzero,
by K\"unneth. It is in the full trace kernel of
$F\boxtimes K$, by naturality under the exact enhanced
functor $(-)\boxtimes K$ (Perry 2026 Remark~2.3).
It is also invariant under the originally given $H$-lift.
Indeed, after the factor stabilizations just proved,
K\"unneth and the negative-Ext vanishing in (S25.1) give
\[
 \operatorname{Hom}(\phi F\boxtimes\psi K,F\boxtimes K)
  =\operatorname{Hom}(\phi F,F)\otimes
       \operatorname{Hom}(\psi K,K),
\]
a one-dimensional space. The given product linearization
therefore differs from the tensor product of any two
factor isomorphisms only by a scalar. Conjugation fixes
$u$ in the first factor and the identity in the second;
the scalar cancels. This addresses correlated actions
and mutually cancelling projective factor cocycles.

Finally, apply (S25.2) on $X\times Y$. The nonzero
invariant class is a nonzero kernel in the equivariant
category, so the given lift is not semiregular. Derived
equivalence preserves this conclusion by full trace
naturality. The proof treats every possible coherent
finite target action and lift separately.
\end{proof}

\begin{proposition}[S-26: An exact tangent-dimension and amplitude gate for the balanced RM character]
\label{prop:semireg-rm-tangent-gate}
In the setting of S-22, let $\widetilde E$ be an
equivariantly simple gluable lift with normalized
character $Nv_q$, and set
$k=4Dq^2N^2/|G|$.
Then full semiregularity is equivalent to
\[
 \operatorname{ext}^1_G(\widetilde E,\widetilde E)
        =10-\frac{k}{2}.
 \tag{S26.1}
\]
The positive-Euler possibilities are exactly $k=4$ with
$e_1^G=8$, or $k=2$ with $e_1^G=9$.

There is also a cohomological restriction independent of
semiregularity. A rank-one perfect object with normalized
character $v_q$ cannot have all the following properties:
$H^j(E)=0$ for $j<0$; $H^0(E)$ is rank-one torsion free;
all $H^j(E)$ for $j>0$ have support of codimension at least
two; and those with positive even $j$ have support of
codimension at least three. In particular, the natural
amplitude $[0,1]$ form with torsion-free $H^0$ and
codimension-two-or-more $H^1$ cannot realize this character.
\end{proposition}
\begin{proof}
S-22 gives contraction rank eighteen in degree two.
The restricted CY4 adjunction and the identity
$\sigma\operatorname{ob}=c^2$ imply that the full trace
is injective if and only if $e_2^G=18$: one implication
is adjunction, and for the other the composition already
has rank eighteen on an eighteen-dimensional target
of evaluation. Invariant Serre duality gives
$e_2^G=k-2+2e_1^G$, proving (S26.1).
The rank-eight degree-one image and positivity of $k$
give the two possibilities in S-22.

For the second assertion put $F=H^0(E)$ and
$L=F^{**}$. The higher cohomology does not change rank
or first Chern class, so $c_1(L)=-B$.
The degree-two part of $e^B\operatorname{ch}(F)$ is the
negative of the effective surface cycle of $L/F$.
The only other degree-two contributions are the negatives
of the effective leading cycles of the odd positive-degree
cohomology sheaves; by hypothesis the even positive-degree
terms contribute zero. Thus
$q(\theta_1^2+\theta_2^2)/2$ would be an effective
surface class, contrary to S-22.
\end{proof}

\begin{proposition}[S-27: A two-sheaf cone gives a concrete gluable derived candidate]
\label{prop:semireg-two-sheaf-cone}
Let $A,T$ be nonzero coherent sheaves on a smooth projective
variety, with
\[
 \operatorname{End}(A)=\operatorname{End}(T)=\mathbf C,
 \qquad \operatorname{Ext}^i(T,A)=0\quad(0\le i\le2).
 \tag{S27.1}
\]
For every nonzero $e\in V:=\operatorname{Ext}^3(T,A)$,
form the triangle
\[
 A\longrightarrow E_e\longrightarrow T[-2]
       \xrightarrow{e}A[1].
 \tag{S27.2}
\]
Then $E_e$ is simple and gluable, with
$H^0(E_e)=A$, $H^2(E_e)=T$, and all other cohomology zero.
Its character is $\operatorname{ch}(A)+\operatorname{ch}(T)$,
and
\[
 V/\mathbf C e\hookrightarrow
       \operatorname{Ext}^1(E_e,E_e).
 \tag{S27.3}
\]
Consequently, in the nonequivariant principal balanced-RM
case of S-26, such a cone with the required character is
fully semiregular precisely when
$\operatorname{ext}^1(E_e,E_e)=8$; in particular it requires
$\dim V\le9$.
\end{proposition}
\begin{proof}
Put $B=T[-2]$. There are no negative self-Ext groups of
$A$ or $B$. Also
$\operatorname{Hom}(A,B[j])=\operatorname{Ext}^{j-2}(A,T)=0$
for $j<0$, while
$\operatorname{Hom}(B,A[j])=\operatorname{Ext}^{j+2}(T,A)=0$
for $j<0$ by (S27.1). Taking extensions preserves these
negative-Hom vanishings, by the long exact Hom sequences
in both variables. Thus $E_e$ is gluable. The cohomology
and character assertions follow directly from (S27.2).

Moreover $\operatorname{Hom}(A,B)=
\operatorname{Hom}(A,B[-1])=0$, so
$\operatorname{Hom}(A,E_e)=\operatorname{End}(A)=\mathbf C$.
Similarly $\operatorname{Hom}(E_e,B)=\operatorname{End}(B)
=\mathbf C$. Every endomorphism of $E_e$ therefore induces
scalar endomorphisms $a,b$ on the two ends of the triangle.
The obstruction to this morphism of triangles is
$(a-b)e$, so $a=b$. The kernel of the map to the two
end endomorphisms is a quotient of
$\operatorname{Hom}(B,A)=\operatorname{Ext}^2(T,A)=0$.
Hence $\operatorname{End}(E_e)=\mathbf C$.

For (S27.3), write $i:A\to E_e$ and $p:E_e\to B$
for the triangle maps. Define
\[
 j(v)=i[1]\,v\,p\in\operatorname{Hom}(E_e,E_e[1]),
 \qquad v\in V=\operatorname{Hom}(B,A[1]).
\]
The long exact sequence obtained by applying
$\operatorname{Hom}(E_e,-)$ to the shifted triangle
contains
\[
 \operatorname{Hom}(E_e,B)\xrightarrow{e\circ(-)}
 \operatorname{Hom}(E_e,A[1])\xrightarrow{i[1]\circ(-)}
 \operatorname{Hom}(E_e,E_e[1]).
\]
Its first map is zero, since $\operatorname{Hom}(E_e,B)=
\mathbf C p$ and $ep=0$. Hence the second map is injective.
Applying $\operatorname{Hom}(-,A[1])$ to (S27.2) gives
\[
 \operatorname{End}(A[1])\xrightarrow{(-)\circ e}
 V\xrightarrow{(-)\circ p}\operatorname{Hom}(E_e,A[1]).
\]
The kernel of the latter map is exactly $\mathbf C e$.
Thus $\ker j=\mathbf C e$, proving (S27.3) directly.
The last assertion follows from S-26.
\end{proof}

S-27 is an explicit construction conditional on its two
coherent sheaves and the character equation. No such pair
with the nonsplit RM character and the required tangent
dimension has been constructed here. It isolates an
allowed amplitude $[0,2]$ and a computable extension space,
in contrast to the excluded amplitude $[0,1]$ forms.

\begin{proposition}[S-28: The two-sheaf cone has an exact Yoneda tangent criterion]
\label{prop:semireg-cone-exact-tangent}
Retain the hypotheses and notation of S-27, and put
$B=T[-2]$. Define
\[
 \delta_e:\operatorname{Ext}^1(A,A)\oplus
                 \operatorname{Ext}^1(B,B)
        \longrightarrow\operatorname{Ext}^2(B,A),
 \qquad
 \delta_e(a,b)=a[1]e+e[1]b.
 \tag{S28.1}
\]
Here the sign is the usual convention in which the shifted
triangle has first two maps $i[1],p[1]$ and last map
$-e[1]$. There is a natural exact sequence
\[
 0\longrightarrow V/\mathbf C e
   \xrightarrow{j}\operatorname{Ext}^1(E_e,E_e)
   \xrightarrow{r}\ker\delta_e\longrightarrow0.
 \tag{S28.2}
\]
In particular,
\[
 \operatorname{ext}^1(E_e,E_e)
   =\dim V-1+\dim\ker\delta_e.
 \tag{S28.3}
\]

Suppose now that the ambient variety is an abelian fourfold
in the principal nonequivariant balanced-RM case of S-26,
and that $e^{\mathfrak b}\operatorname{ch}(E_e)=v_1$,
where $\mathfrak b$ denotes the normalizing $B$-field.
Let
\[
 \rho:HH^1(X)\longrightarrow
       \operatorname{Ext}^1(A,A)\oplus
       \operatorname{Ext}^1(T,T)
 \tag{S28.4}
\]
be simultaneous Hochschild evaluation, equivalently the
infinitesimal translation--tensor action on the two sheaves,
and put $h=\dim\ker\rho$. Then
\[
 E_e\text{ is fully semiregular}
 \quad\Longleftrightarrow\quad
 \dim V=h+1\ \text{ and }\ \ker\delta_e=\operatorname{im}\rho.
 \tag{S28.5}
\]
If the simultaneous infinitesimal stabilizer is zero,
the remaining conditions are exactly
\[
 \dim\operatorname{Ext}^3(T,A)=1,
 \qquad
 \operatorname{rank}\delta_e=
       \operatorname{ext}^1(A,A)+\operatorname{ext}^1(T,T)-8.
 \tag{S28.6}
\]
No existence of sheaves satisfying these conditions and
the balanced-RM character equation is asserted.
\end{proposition}
\begin{proof}
For coherent $A,T$ one has
$\operatorname{Hom}(A,B[k])=0$ for $k\le1$.
The long exact Hom sequences therefore give isomorphisms
\[
 \operatorname{Ext}^1(A,A)
     \xrightarrow{i[1]\circ(-)}\operatorname{Hom}(A,E_e[1]),
 \qquad
 \operatorname{Ext}^1(B,B)
     \xrightarrow{(-)\circ p}\operatorname{Hom}(E_e,B[1]).
 \tag{S28.7}
\]
Consequently every $u:E_e\to E_e[1]$ determines unique
$a,b$ by $ui=i[1]a$ and $p[1]u=bp$; set $r(u)=(a,b)$.
The morphism-of-triangles axiom and uniqueness of $b$
imply $a[1]e=-e[1]b$. Conversely, if this equality holds,
apply the same axiom to the rotated triangles, starting
with the commuting square on $B\to A[1]$, to obtain a
map $u$ inducing $a,b$. Thus $\operatorname{im}r=\ker\delta_e$.

If $r(u)=0$, then $ui=0$, so exactness gives $u=wp$ for
some $w:B\to E_e[1]$. Since $p[1]u=0$ and precomposition
with $p$ is injective in the second map of (S28.7),
$p[1]w=0$. A further long exact Hom sequence gives
$w=i[1]v$ for some $v\in V$. Hence $u=j(v)$.
The converse is immediate from consecutive triangle maps
having zero composite. S-27 already proves
$\ker j=\mathbf C e$, so (S28.2) and (S28.3) follow.

For the abelian-fourfold assertion, write
$\operatorname{ob}^1_E:HH^1(X)\to\operatorname{Ext}^1(E_e,E_e)$.
The degree-one contraction rank is eight by S-22.
Since $\sigma^1\operatorname{ob}^1_E=c_E^1$ and
$\dim HH^1(X)=8$, the map $\operatorname{ob}^1_E$ is
injective. Naturality gives
$r\operatorname{ob}^1_E=\rho$; the even shift $[-2]$
does not change the identification for the $T$ term.
Thus $\operatorname{im}\rho\subseteq\ker\delta_e$,
$\dim\operatorname{im}\rho=8-h$, and
\[
 \dim\bigl(\operatorname{im}\operatorname{ob}^1_E
                 \cap\ker r\bigr)=h\le\dim V-1.
 \tag{S28.8}
\]
If $E_e$ is fully semiregular, S-26 gives
$\operatorname{ext}^1(E_e,E_e)=8$. From (S28.3) and
$\dim\ker\delta_e\ge8-h$ one gets $\dim V-1\le h$.
Together with (S28.8), this forces $\dim V=h+1$ and
$\ker\delta_e=\operatorname{im}\rho$.
Conversely those two equalities make (S28.3) equal to
eight, so S-26 gives full semiregularity. Setting $h=0$
proves (S28.6).
\end{proof}

\begin{proposition}[S-29: The transverse-surface cone requires special cohomology]
\label{prop:semireg-surface-cone-cohomology}
Let $X$ be a simple complex abelian fourfold, and let
$Z,S\subset X$ be smooth connected surfaces meeting
transversely in the finite scheme $D$. Write $i:S\hookrightarrow X$,
let $L$ be a line bundle on $X$ and $M$ a line bundle on $S$,
and set
\[
 A=I_Z\otimes L,\qquad T=i_*M,\qquad
 M'=M\otimes L^{-1}|_S,\qquad
 N=K_S\otimes(M')^{-1},\qquad d=\operatorname{length}D.
\]
There are natural identifications
\[
 \operatorname{Ext}^{j+2}_X(T,A)=H^j(S,N\otimes I_{D/S})
       \quad(0\le j\le2),
 \qquad
 \operatorname{Ext}^{0,1}_X(T,A)=0.
 \tag{S29.1}
\]
The two sheaves have scalar endomorphisms, and each
translation--tensor evaluation $HH^1(X)\to\operatorname{Ext}^1$
is injective.

Suppose that a cone of S-27 formed from these sheaves has
the principal balanced-RM normalized character $v_1$.
If it is fully semiregular, then
\[
 \begin{split}
 H^0(S,N\otimes I_D)&=0,\\
 h^1(S,N\otimes I_D)&=1,\\
 h^0(S,M')&\ge8,\\
 h^2(S,M')-h^1(S,M')&=d-1\ge1.
 \end{split}
 \tag{S29.2}
\]
In fact the only possibilities are
\[
 (h^1(S,M'),h^2(S,M'))=(0,d-1)\quad\text{or}\quad(1,d).
 \tag{S29.2a}
\]
In the first case evaluation of $N$ at $D$ has
one-dimensional cokernel; in the second it is an
isomorphism. Together with the full character equation, the first two
conditions and the rank condition (S28.6) are sufficient
as well. In particular, choosing $M'$ with
$H^{>0}(S,M')=0$ cannot give a fully semiregular cone.
More precisely, for every such acyclic choice satisfying
the character equation and every nonzero extension,
\[
 \dim\ker\bigl(\sigma^2_{E_e}\bigr)\ge d-1\ge1.
 \tag{S29.3}
\]
\end{proposition}
\begin{proof}
Transversality gives
$Li^*I_Z=I_{D/S}$: pull back the ideal triangle and
compute the other term by the Koszul complex of $Z$,
which remains regular on $S$.
For the codimension-two regular immersion $i$, duality
gives $i^!\mathcal O_X=\det N_{S/X}[-2]=K_S[-2]$ and
$i^!A=Li^*A\otimes K_S[-2]$.
The precise general statements are Stacks Project,
Lemma~48.9.7 (tag \texttt{0A9X}),
Lemma~48.10.3 (tag \texttt{0E2N}), and
Lemma~48.15.5 (tag \texttt{0BQZ}); the last is the
Koszul duality computation with shift $[-2]$.
Adjunction therefore yields
\[
 R\operatorname{Hom}_X(i_*M,I_Z\otimes L)
     =R\Gamma(S,N\otimes I_D)[-2],
 \tag{S29.4}
\]
which proves (S29.1).
Endomorphisms of $A$ extend to its line-bundle double
dual and hence are scalar; those of $T$ are
$H^0(S,\mathcal O_S)=\mathbf C$.

A translation--tensor operation fixing $A$ must preserve
$Z$, and its degree-zero tensor factor is then determined
by the double dual $L$. The translation stabilizer of
$Z$ is finite: its connected component is an abelian
subvariety of the simple variety $X$, and cannot be $X$.
For $T$, preservation of the support likewise forces the
translation into the finite stabilizer of $S$. With
translation fixed, the remaining stabilizer lies in
$\ker(\operatorname{Pic}^0(X)\to\operatorname{Pic}^0(S))$,
which has zero Lie algebra: $S$ generates $X$ by simplicity,
so its Albanese map onto $X$ is surjective and the dual
map on Picard varieties has injective differential.
Thus both infinitesimal stabilizers vanish, proving the
two rank-eight evaluation assertions.

In particular $h=0$ in S-28. Its criterion gives the
first two conditions of (S29.2) and
\[
 \operatorname{rank}\delta_e
   =\operatorname{ext}^1(A,A)+\operatorname{ext}^1(T,T)-8
   \ge8.
\]
The target of $\delta_e$ has dimension
$h^2(N\otimes I_D)=h^2(N)=h^0(M')$ by Serre duality,
so $h^0(M')\ge8$.
In the cohomology sequence of
$0\to N\otimes I_D\to N\to N|_D\to0$, the first two
conditions imply
\[
 1=d-h^0(N)+h^1(N)
   =d-h^2(M')+h^1(M').
\]
Equivalently, $h^2(M')-h^1(M')=d-1$.
More precisely that exact sequence is
\[
 0\longrightarrow H^0(N)\longrightarrow H^0(N|_D)
 \longrightarrow \mathbf C\longrightarrow H^1(N)
 \longrightarrow0.
 \tag{S29.5}
\]
It follows that $h^1(N)=h^1(M')$ is either zero or one;
the two cases and the assertions about evaluation follow.

For the lower bound on $d$, the first Chern equation
forces the normalizing $B$-field to be $-c_1(L)$.
The codimension-two character equation therefore gives
\[
 [Z]-[S]=C,\qquad C=\frac{\theta_1^2+\theta_2^2}{2}.
 \tag{S29.6}
\]
Both $\theta_i$ have invariant semipositive representatives,
so $C\cdot[S]\ge0$. On a simple complex abelian variety
every irreducible subvariety is geometrically nondegenerate;
Debarre, \emph{Fulton--Hansen and Barth--Lefschetz theorems
for subvarieties of abelian varieties} (1995),
Examples~(1), pp.~188--189, and Proposition~1.4,
pp.~189--190, imply $S-S=X$.
Thus the difference map $S\times S\to X$ is generically
finite and has positive degree $[S]^2$.
This number is even: in
$H^*(X,\mathbf Z)=\bigwedge^*H^1(X,\mathbf Z)$ the square
of every degree-four integral class has even top
coefficient, because each basis monomial squares to zero
and each pair of distinct monomials occurs twice.
Consequently
\[
 d=[Z]\cdot[S]=[S]^2+C\cdot[S]\ge2.
 \tag{S29.7}
\]

Finally, if $H^{>0}(M')=0$, Serre duality gives
$H^0(N)=H^1(N)=0$. The same restriction sequence gives
$H^0(N\otimes I_D)=0$ and
$V=H^1(N\otimes I_D)\simeq H^0(N|_D)$, of dimension $d$.
The first map in the rank-eight simultaneous orbit is
injective, so $\dim\ker\delta_e\ge8$.
S-28 implies $\operatorname{ext}^1(E_e,E_e)\ge d+7$,
and the CY4 Euler equation, with $\chi(E_e,E_e)=4$,
gives $m:=\operatorname{ext}^2(E_e,E_e)\ge2d+16$.
Put $k=\dim\ker\sigma^2_{E_e}$.
By the CY4 transpose identity,
$\operatorname{rank}\sigma^2_{E_e}
 =\operatorname{rank}\operatorname{ob}^2_{E_e}=m-k$,
while the composition has rank eighteen by S-22.
The elementary rank inequality yields
$18\ge2(m-k)-m=m-2k$, hence $k\ge d-1$.
The sufficiency statement is exactly S-28 after (S29.1)
has checked the S-27 hypotheses.
\end{proof}

\begin{theorem}[S-30: An integral lattice obstruction closes the simple rank-one balanced RM gate]
\label{thm:semireg-rm-integral-obstruction}
Let $(X,\theta)$ be a complex abelian fourfold with a
Rosati-invariant action of the real quadratic field
$F=\mathbf Q(\sqrt{\Delta_0})$, where $\Delta_0>1$ is
squarefree, and with the two complex eigenspaces both
of dimension two. Use the notation
\[
 D=\int_X\frac{\theta^4}{24},\qquad
 C=\frac{\theta_1^2+\theta_2^2}{2},\qquad
 v_q=1-qC+\frac{q^2}{4}\theta_1^2\theta_2^2,
 \qquad q\in\mathbf Q_{>0}.
\]
Here $\theta$ is an integral polarization and
$\theta=\theta_1+\theta_2$ is its orthogonal decomposition
along the two eigenspaces.
If $a\in\mathbf Q$ is nonzero and $aC$ is an integral
cohomology class, then
\[
 t=a^2D\in\mathbf Z_{>0},\qquad \Delta_0\mid t.
 \tag{S30.1}
\]
Consequently no rank-one perfect object with normalized
character $v_q$ can be both simple, gluable, and fully
semiregular. More quantitatively, every simple gluable
rank-one perfect object with that normalized character
satisfies
\[
 \dim\ker\sigma_E^2\ge2Dq^2-2\ge2\Delta_0-2>0.
 \tag{S30.2}
\]
In particular, in the principal $q=1$ case, the character
$v_1$ is not the normalized character of any rank-one
perfect object at all. This last assertion does not
require simplicity, gluability, or semiregularity.

For a positive-rank perfect object with
$e^B\operatorname{ch}(E)=Nv_q$, one always has
\[
 N^4Dq^2\in\mathbf Z_{>0},\qquad
 \Delta_0\mid N^4Dq^2.
 \tag{S30.3}
\]
If, in addition, $E$ has a simple gluable fully semiregular
lift for an effective finite coherent target action $G$
as in S-22, write $k=\chi_G\in\{2,4\}$. Then
\[
 \begin{cases}
 \Delta_0\mid N^2|G|,&k=4,\\
 2\Delta_0\mid N^2|G|,&k=2.
 \end{cases}
 \tag{S30.4}
\]
For $G=1$, in particular, $N$ must be a multiple of
$\Delta_0$, and $N$ must also be even in the $k=2$ case.
Moreover, in every rank, a nonequivariant simple gluable
fully semiregular realization requires a genuinely
fractional normalizing class: $B$ cannot be the Chern
class of a line bundle. Equivalently, its determinant
cannot admit an $N$th root in $\operatorname{Pic}(X)$.
The integer $\Delta_0$ here is the squarefree radicand,
not necessarily the discriminant of the quadratic field.
\end{theorem}
\begin{proof}
Put $L=H_1(X,\mathbf Z)$ and let $\nu$ be the primitive
positive generator of $\bigwedge^8L^*$, so
$\theta^4/24=D\nu$. The two real eigenspaces give
$L_{\mathbf R}=U\oplus W$, each of real dimension four.
Write $\nu_U=\theta_1^2/2$ and $\nu_W=\theta_2^2/2$.
Then $C=\nu_U+\nu_W$ and $\nu_U\wedge\nu_W=D\nu$.

Suppose $c=aC\in\bigwedge^4L^*$. Define the integral
endomorphism $R_c$ of $\bigwedge^2L$ by
\[
 \iota_{R_c(x)}\nu=(\iota_xc)\wedge c.
 \tag{S30.5}
\]
Contraction with the primitive volume form is a
unimodular identification between
$\bigwedge^2L$ and $\bigwedge^6L^*$, so $R_c$ is indeed
integral. A direct calculation on the three summands
\[
 \bigwedge^2L_{\mathbf R}
  =\bigwedge^2U\oplus(U\otimes W)\oplus\bigwedge^2W
\]
shows that $R_c=tP$, where $t=a^2D$ and $P$ is the
projector which is the identity on the first and third
summands and zero on the middle summand.
Moreover $c^2=2t\nu$. The square of every integral
four-form on an eight-dimensional lattice has even top
coefficient: basis monomials square to zero and their
cross terms occur in pairs. Hence $t$ is an integer.

Choose the rational endomorphism $T$ of $L_{\mathbf Q}$
given by $\sqrt{\Delta_0}\in F$. It satisfies
$T^2=\Delta_0\operatorname{id}$ and has eigenvalues
$\sqrt{\Delta_0},-\sqrt{\Delta_0}$ on $U,W$. Therefore
\[
 2P-\operatorname{id}
      =\frac{\bigwedge^2T}{\Delta_0}.
 \tag{S30.6}
\]
Let $p$ divide $\Delta_0$. If $p\nmid t$, the integrality
of $tP$ makes $P$ integral over $\mathbf Z_p$.
Thus the left side of (S30.6) is an integral involution,
and hence an automorphism of $\bigwedge^2L_{\mathbf Z_p}$.
Apply Smith normal form over $\mathbf Z_p$ to the
invertible rational matrix $T$. If its elementary-divisor
exponents are integers $m_1,\ldots,m_8$, those of its
exterior square, before reordering, are
$m_i+m_j$ for $i<j$. Equation (S30.6) and unimodularity
would force
\[
 m_i+m_j=v_p(\Delta_0)=1\qquad(i<j).
 \tag{S30.7}
\]
Using any three indices forces every $m_i=1/2$, a
contradiction. This argument works also at $p=2$.
Every prime factor of the squarefree integer $\Delta_0$
therefore divides $t$, proving (S30.1).

For rank one, the degree-two character equation gives
$B=-c_1(E)$. The determinant of a perfect complex is a
line bundle, so $E'=E\otimes\det(E)^{-1}$ is an actual
perfect complex with character $v_q$. Its second Chern
class is $qC$, which is integral. Riemann--Roch on the
abelian fourfold gives $\chi(E')=Dq^2$ and
$\chi(E,E)=4Dq^2$. By (S30.1), $Dq^2$ is a positive
multiple of $\Delta_0$.

If $E$ is simple and gluable, the rank-eight degree-one
contraction forces $e_1\ge8$, and CY4 Serre duality gives
\[
 m:=e_2=4Dq^2-2+2e_1\ge4Dq^2+14.
\]
Put $b=\dim\ker\sigma_E^2$. The full transpose identity
makes the ranks of $\sigma_E^2$ and
$\operatorname{ob}_E^2$ both $m-b$. Their composition has
rank eighteen. Thus
$18\ge2(m-b)-m=m-2b$, which proves (S30.2).
For $D=q=1$, (S30.1) already contradicts $qC$ being
integral; this proves the assertion without simplicity.

For general positive rank $N$, put $\ell=c_1(E)$.
The degree-two equation gives $B=-\ell/N$, and the
next character equation is
\[
 N^2qC=Nc_2(E)-(N-1)\frac{\ell^2}{2}.
 \tag{S30.8}
\]
The class $\ell^2/2$ is integral in the exterior algebra
of a torus. Thus the right side is integral and (S30.1)
with $a=N^2q$ proves (S30.3). In the equivariant case
S-22 gives $k=4Dq^2N^2/|G|\in\{2,4\}$, so
\[
 N^4Dq^2=\frac{kN^2|G|}{4}.
\]
Substitution into (S30.3) gives (S30.4).
For $G=1$, squarefreeness of $\Delta_0$ implies
$\Delta_0\mid N$; when $k=2$, the additional divisibility
by two forces $N$ even.

Finally suppose $G=1$ and $B=c_1(Q)$ for a line bundle
$Q$. Then $E\otimes Q$ is an actual perfect object with
character $Nv_q$, so $s:=\chi(E\otimes Q)=NDq^2$ is a
positive integer. Full semiregularity, simplicity and
gluability give
$4Ns=\chi(E,E)=2-2e_1+18\le4$.
Thus $N=s=1$, which is the excluded rank-one case.
Since $NB=-c_1(E)$, a line-bundle determinant root would
make $B$ integral, and conversely an integral $B$ is a
line-bundle class by the integral $(1,1)$ theorem.
The residual degree-zero line bundle can be absorbed
into an $N$th root because multiplication by $N$ on
$\operatorname{Pic}^0(X)$ is surjective.
\end{proof}

\begin{proposition}[S-31: Rank-two and rank-three bundles cannot realize the balanced character]
\label{prop:semireg-rm-low-rank-bundles}
In the real-multiplication setting of S-30, no vector
bundle of rank $N\in\{1,2,3\}$ has normalized character
$e^B\operatorname{ch}(E)=Nv_q$ with $q>0$.
Consequently a nonequivariant simple gluable fully
semiregular vector-bundle realization would have rank
at least four and divisible by $\Delta_0$.
For example, the smallest remaining ranks allowed by
these two tests are four for $F=\mathbf Q(\sqrt2)$,
and six for $F=\mathbf Q(\sqrt3)$.
These are necessary conditions, not constructions.
\end{proposition}
\begin{proof}
Under the splitting principle, replace the $N$ Chern
roots $x_i$ by $x_i+B$. The resulting normalized Chern
polynomial still has degree at most $N$, regardless of
whether $B$ is the Chern class of a line bundle.
The normalized character has first and third components
zero, second component $-NqC$, and fourth component
$NDq^2\nu$, with $\nu$ the primitive top class.
Newton's identities therefore give
\[
 \bar c_1=\bar c_3=0,\qquad
 \bar c_2=NqC,\qquad
 \bar c_4=\frac{\bar c_2^2-12\overline{\operatorname{ch}}_4}{2}
          =N(N-6)Dq^2\nu.
 \tag{S31.1}
\]
For $N=1$, the nonzero second Chern class contradicts
the degree-one Chern polynomial. For $N=2,3$, the
nonzero fourth class contradicts degree at most $N$.
The final rank restrictions follow from S-30.
\end{proof}

\begin{theorem}[S-32: Two-adic integrality forces an even polarization or group degree]
\label{thm:semireg-rm-two-adic-obstruction}
Let $(X,\theta)$ be a polarized abelian
fourfold with Rosati-invariant real multiplication by
$F=\mathbf Q(\sqrt{\Delta_0})$ of type $(2,2)$, where
$\Delta_0>1$ is squarefree and
$\Delta_0\equiv2$ or $3\pmod4$.
Assume that $H_1(X,\mathbf Z_2)$ is stable under the
maximal local order
$R=\mathbf Z_2[\sqrt{\Delta_0}]$, and that every rational
divisor class has the form
$\ell(x,y)=\theta(\alpha x,y)$ for some $\alpha\in F$.
Put $D=\int_X\theta^4/24$.
The latter hypothesis is the usual minimal divisor space
$\operatorname{NS}(X)_{\mathbf Q}=F\theta$.
If $D$ is odd, then for every integral divisor class
$\ell$, every positive integer $N$, and every rational
two-adic unit $u$, the class
\[
 uC-\frac{\ell^2}{2N},\qquad
 C=\frac{\theta_1^2+\theta_2^2}{2},
 \tag{S32.1}
\]
is not two-adically integral. Consequently, if $D$ is odd,
no positive-rank perfect object has normalized character
$Nv_q$ with $Nq$ a two-adic unit.
More generally, every simple gluable fully semiregular
lift of a character $Nv_q$, $q\in\mathbf Q_{>0}$, for an effective finite coherent
target action $G$ must satisfy
\[
 2\mid D|G|.
 \tag{S32.1a}
\]
In particular, for a principal polarization there is no
nonequivariant simple gluable fully semiregular balanced-RM
perfect object, in any positive rank or for any
normalizing $B$-field.
\end{theorem}
\begin{proof}
For the integral-class assertion assume $D$ is odd.
Write $\omega=\sqrt{\Delta_0}$ and $K=\mathbf Q_2(\omega)$.
The ring $R=\mathbf Z_2[\omega]$ is a discrete valuation
ring: $x^2-\Delta_0$ is Eisenstein when $\Delta_0$ is even;
when $\Delta_0\equiv3\pmod4$, the element $\omega-1$
satisfies the Eisenstein polynomial
$x^2+2x+1-\Delta_0$.
Here we use the Eisenstein DVR criterion, for example
Sutherland, \emph{18.785 Number Theory I}, Fall 2019,
Lecture~11, Lemma~11.4 and its proof, pp.~1--2.
Put $L_2=H_1(X,\mathbf Z_2)$.
It is a finite torsion-free $R$-module, hence free of
rank four (Stacks Project, Lemma~15.22.11,
tag \texttt{0AUW}).
Set $\kappa=1/(2\omega)$.
The trace pairing $(a,b)\mapsto\operatorname{Tr}_{K/\mathbf Q_2}
(\kappa ab)$ on $R$ is perfect: its matrix in the basis
$1,\omega$ is
$\left(\begin{smallmatrix}0&1\\1&0\end{smallmatrix}\right)$.

The self-adjoint $R$-action therefore gives a unique
$R$-bilinear alternating form $\psi:L_2\times L_2\to R$
with
\[
 \theta(x,y)=\operatorname{Tr}_{K/\mathbf Q_2}
                         (\kappa\psi(x,y)).
 \tag{S32.2}
\]
For completeness, define $\psi(x,y)$ by requiring
$\theta(ax,y)=\operatorname{Tr}(\kappa a\psi(x,y))$
for all $a\in R$. Perfectness of the trace pairing gives
existence and uniqueness; self-adjointness gives
bilinearity. Also $\theta(ax,x)=0$ for every $a$, since
it equals both $\theta(x,ax)$ and $-\theta(x,ax)$;
thus $\psi(x,x)=0$ even at two.
The $\psi$-dual lattice equals the $\theta$-dual lattice:
test the trace against every $a\in R$, using $RL_2=L_2$.
Since $D$ is odd, the polarization form is unimodular
over $\mathbf Z_2$ and makes the latter lattice
$L_2$ itself. Hence $\psi$ is perfect.
Elementary symplectic elimination over the local ring
now gives an $R$-basis $e_1,f_1,e_2,f_2$ with
$\psi(e_i,f_j)=\delta_{ij}$ and the other pairings zero.
Indeed choose a primitive vector, pair it to a unit by
perfectness, rescale the second vector, and split off
its orthogonal hyperbolic plane; repeat once.

By the divisor-space hypothesis, $\ell=\theta(\alpha\cdot,\cdot)$
for $\alpha\in F$. Integrality of $\ell$ and unimodularity
of $\theta$ imply $\alpha L_2\subset L_2$.
The free $R$-basis therefore implies
$\alpha=a+b\omega$ with $a,b\in\mathbf Z_2$.
On either hyperbolic plane, the three pairings of $\ell$
corresponding to products $1,\omega,\omega^2$ are
\[
 \operatorname{Tr}(\kappa\alpha)=b,\qquad
 \operatorname{Tr}(\kappa\alpha\omega)=a,\qquad
 \operatorname{Tr}(\kappa\alpha\omega^2)=\Delta_0b.
 \tag{S32.3}
\]
Over a splitting field, $C$ is the sum of the squares
of the two embedding components of $\theta$, divided
by two. Evaluating on one ordered pair from each
hyperbolic plane consequently gives
$\operatorname{Tr}(\kappa^2\omega^j)$ when the product
of the four $R$-coefficients is $\omega^j$.
In particular, on the respective ordered quadruples
\[
 (e_1,f_1,e_2,f_2),\quad
 (\omega e_1,f_1,\omega e_2,f_2),\quad
 (\omega e_1,f_1,e_2,f_2),
\]
the values of $C$ are $1/(2\Delta_0),1/2,0$,
and those of $\ell^2/2$ are $b^2,a^2,ab$.
Thus integrality of (S32.1) would require
\[
 \frac{b^2}{N}-\frac{u}{2\Delta_0}\in\mathbf Z_2,
 \qquad
 \frac{a^2}{N}-\frac{u}{2}\in\mathbf Z_2,
 \qquad
 \frac{ab}{N}\in\mathbf Z_2.
 \tag{S32.4}
\]

Because $u$ is a two-adic unit, if $\Delta_0$ is even,
the first two congruences force
$2v_2(b)-v_2(N)=-2$ and
$2v_2(a)-v_2(N)=-1$, which give incompatible parities
for $v_2(N)$. If $\Delta_0$ is odd, both force valuation
$-1$, so $v_2(a)=v_2(b)$ and
$v_2(ab/N)=-1$, contradicting the third congruence.
This proves the integral-class assertion.

For a perfect object with normalized character
$Nv_q$, put $u=Nq$. The first two Chern equations give
$B=-\ell/N$ and
\[
 c_2(E)=\frac{N-1}{2N}\ell^2+uC
       =\frac{\ell^2}{2}
          +\left(uC-\frac{\ell^2}{2N}\right).
\]
Both $c_2(E)$ and $\ell^2/2$ are integral, contradicting
(S32.1) whenever $u$ is a two-adic unit.
Finally suppose a simple gluable fully semiregular lift
exists and $D|G|$ is odd. By S-22 its invariant Euler
pairing is
$k=4D(Nq)^2/|G|\in\{2,4\}$.
The value $k=2$ would give
$v_2((Nq)^2)=-1$, impossible for a rational square.
The value $k=4$ gives $v_2(Nq)=0$, which has just been
excluded. This proves (S32.1a) and the principal
nonequivariant corollary.
\end{proof}

\begin{theorem}[S-33: Odd ramification obstructs the balanced character in every rank]
\label{thm:semireg-odd-ramification}
Let \((X,\theta)\) be a polarized complex abelian fourfold, with polarization volume \[
 D=\int_X\frac{\theta^4}{24}\in\mathbf Z_{>0}.
\] Let \(F=\mathbf Q(\sqrt d)\) be a real quadratic field, where \(d>1\) is squarefree. Assume its rational action is fixed by the Rosati involution and its two complex eigenspaces each have dimension two. Write \[
 \theta=\theta_++\theta_-,\qquad
 C=\frac{\theta_+^2+\theta_-^2}{2},\qquad
 \nu=\frac{\theta^4}{24}
      =\frac{\theta_+^2\theta_-^2}{4},\qquad \int_X\nu=D.
 \tag{S33.1}
\] For \(r\in F\), let \(\theta_r\) denote the rational alternating form \((x,y)\mapsto\theta(rx,y)\), equivalently \(\sigma_+(r)\theta_++\sigma_-(r)\theta_-\). Fix an odd prime \(p\) ramified in \(F\), so \(p\mid d\), and assume \[
 p\nmid D,\qquad
 (\mathcal O_F\otimes\mathbf Z_p)H_1(X,\mathbf Z_p)
                  \subseteq H_1(X,\mathbf Z_p).
 \tag{S33.2}
\] Thus only maximal-order action at this prime is required.

There is \textbf{no} perfect object \(\mathcal E\) of positive rank \(N\), with \(\ell=c_1(\mathcal E)\in F\theta\), and no \(u\in\mathbf Q^\times\) with \(v_p(u)=0\), satisfying \[
 e^{-\ell/N}\operatorname{ch}(\mathcal E)
       =N-uC+\frac{u^2}{N}\nu.
 \tag{S33.3}
\] The assertion allows an arbitrary fractional normalizing class \(-\ell/N\) and does not require \(\mathcal E\) to be a sheaf, simple or gluable. For the balanced character \(Nv_q\), the scalar here is \(u=Nq\).

In particular, suppose \(\operatorname{NS}(X)_{\mathbf Q}=F\theta\) and maximal-order action is integral at every ramified prime. For a simple gluable fully semiregular lift with an effective finite coherent target action \(G\) as in S-22, every odd ramified prime must divide \(D|G|\). Combining with S-32's even-ramification obstruction gives the necessary condition \[
 \operatorname{rad}(|\operatorname{Disc}(F)|)\mid D|G|.
 \tag{S33.4}
\] In the nonequivariant case, all ramified primes must divide \(D\). In particular, there is no such positive-rank realization for a principal polarization and a genuine real quadratic field. This excludes the proposed normalized character under the stated local maximal-order and determinant hypotheses; it does not assert nonexistence for arbitrary derived characters or arbitrary polarizations/actions.
\end{theorem}
\begin{proof}
\par\medskip\noindent\textbf{1. The integral endomorphism corresponding to the determinant}\quad 

Since \(\ell\in F\theta\), write \(\ell=\theta_r\) for \(r\in F\). Because \(p\nmid D\), the integral alternating form \(\theta\) is perfect over \(\mathbf Z_p\): its determinant is \(D^2\) up to a unit in any integral basis. Integrality of \(\ell\) therefore implies \(rH_1(X,\mathbf Z_p)\subseteq H_1(X,\mathbf Z_p)\). Put \(K=F\otimes\mathbf Q_p\), \(O=\mathcal O_F\otimes\mathbf Z_p\) and \(L=H_1(X,\mathbf Z_p)\). The ring \(O\) is the valuation ring of a ramified quadratic extension. The \(O\)-module \(L\) is torsion-free of rank four, hence free (Stacks Project, Lemma~15.22.11, tag \texttt{0AUW}). The scalar \(r\) preserves \(O^4\), and consequently \(r\in O\).

Set \(s=\sqrt d\). The Eisenstein DVR criterion (Sutherland, Fall 2019 Lecture~11, Lemma~11.4, pp.~1--2) applies to \(x^2-d\). Because \(p\) is odd, \(O=\mathbf Z_p[s]\); the different is the principal ideal \((s)\), and \(s\) is a uniformizer up to a unit. Normalize its valuation by \(v_\pi(s)=1\), so \(v_\pi(p)=2\). We can write \[
 r=a+bs,\qquad a,b\in\mathbf Z_p.
 \tag{S33.5}
\]

\par\medskip\noindent\textbf{2. The self-dual trace-symplectic normal form}\quad 

For completeness, the precise normal form used in the coefficient calculation is justified here. The Rosati condition says \(\theta(ax,y)=\theta(x,ay)\) for \(a\in K\). Nondegeneracy of the field trace determines a unique alternating \(K\)-bilinear form \(\beta\) on \(L\otimes\mathbf Q_p\) by \[
 \operatorname{Tr}_{K/\mathbf Q_p}(a\beta(x,y))=\theta(ax,y)
 \quad(a\in K).
\] The balance identity proves bilinearity in both variables. Alternation follows from \(\theta(ax,x)=0\), using the balance identity and skew symmetry in characteristic zero.

For \(x,y\in L\), integrality after multiplying \(x\) by every \(a\in O\) gives \(\beta(x,y)\in O^\vee=s^{-1}O\), where \(O^\vee\) is the trace-dual lattice. The trace pairing identifies \(\operatorname{Hom}_O(L,s^{-1}O)\) with \(\operatorname{Hom}_{\mathbf Z_p}(L,\mathbf Z_p)\). Since the local form \(\theta\) is perfect, \(\beta\) is a perfect \(s^{-1}O\)-valued alternating form. Thus \(\psi=s\beta\) is a unimodular alternating \(O\)-bilinear form.

Choose a symplectic \(O\)-basis \((e_1,f_1,e_2,f_2)\) for \(\psi\). One obtains it by choosing a primitive vector, a second vector pairing to a unit, rescaling that vector to make the pairing one, and splitting off the orthogonal hyperbolic plane; the residual rank-two summand is handled in the same way. No division by two is involved. In this basis \[
 L=U_1\oplus U_2,
 \quad U_i=Oe_i\oplus Of_i,
 \quad \psi(e_i,f_j)=\delta_{ij},
 \quad \theta=\operatorname{Tr}_{K/\mathbf Q_p}(\psi/s).
 \tag{S33.6}
\] In particular each \(U_i\) is a unimodular symplectic \(\mathbf Z_p\)-lattice of rank four for the restricted trace form.

\par\medskip\noindent\textbf{3. The mixed second-Chern-class coefficients}\quad 

The degree-four part of (S33.3) gives \[
 \operatorname{ch}_2(\mathcal E)=\frac{\ell^2}{2N}-uC.
\] Since \(\ell^2/2\) is integral on a torus and \(c_2(\mathcal E)\) is integral, the identity \(\operatorname{ch}_2=\ell^2/2-c_2\) shows \[
 uC-\frac{\ell^2}{2N}\in H^4(X,\mathbf Z_p).
 \tag{S33.7}
\] For \(x,y\in O\), evaluate this four-form on \((xe_1,f_1,ye_2,f_2)\). The two embedding components of the trace form give \[
 C(xe_1,f_1,ye_2,f_2)=\operatorname{Tr}_{K/\mathbf Q_p}(xy/d),
\] while the orthogonality of \(U_1,U_2\) gives \[
 (\ell^2/2)(xe_1,f_1,ye_2,f_2)
 =\operatorname{Tr}(rx/s)\operatorname{Tr}(ry/s).
 \tag{S33.8}
\] Indeed each embedding component on one \(O\)-hyperbolic plane has rank two, so its square vanishes there; the mixed term in \(\theta_\sigma^2/2\) is the product of its two plane restrictions, with coefficient one.

In the basis \((1,s)\) of \(O\), the trace Gram matrix in (S33.8) for \(C\) is \(\operatorname{diag}(2/d,2)\). The corresponding linear functional for \(\ell\) has values \(2b,2a\). Therefore (S33.7) implies \[
 \frac{b^2}{N}-\frac{u}{2d}\in\mathbf Z_p,
 \qquad \frac{a^2}{N}\in\mathbf Z_p,
 \qquad \frac{ab}{N}\in\mathbf Z_p.
 \tag{S33.9}
\] Here division by four is legitimate because \(p\) is odd. These displayed factors are tied to the convention \(\theta=\operatorname{Tr}(\psi/s)\); no \(1/(2s)\) normalization has been substituted.

Let \(m=v_p(N)\) and \(k=v_p(b)\). The first congruence forces \(b\ne0\) and \[
 m=2k+1>0,
 \qquad v_p(a)\ge k+1,
 \qquad v_\pi(r)=2k+1=m.
 \tag{S33.10}
\] The summands \(a\) and \(bs\) have distinct \(\pi\)-valuations, at least \(2k+2\) and exactly \(2k+1\), respectively, proving the last equality. Moreover \[
 \frac{r^2}{N}
 =\frac{a^2}{N}+d\frac{b^2}{N}
             +2\frac{ab}{N}s
 \equiv\frac u2\pmod\pi.
 \tag{S33.11}
\] The first summand has positive \(p\)-valuation, the second has the asserted residue by (S33.9), and the third is divisible by \(\pi\). In particular \(u-r^2/N\) is a unit of \(O\).

\par\medskip\noindent\textbf{4. The integral Euler characteristic gives the contradiction}\quad 

The tangent bundle of an abelian variety is trivial, so Riemann--Roch (Stacks Project, Section~42.66, tag \texttt{02UO}) gives \(\chi(\mathcal E)=\int_X\operatorname{ch}_4(\mathcal E)\in\mathbf Z\). The character (S33.3) yields \[
 \operatorname{ch}_4(\mathcal E)
 =\frac{\ell^4}{24N^3}
       -\frac{u\ell^2C}{2N^2}+\frac{u^2}{N}\nu.
\] Using (S33.1) and writing \(r_\pm=\sigma_\pm(r)\), one computes \[
 \int_X\frac{\ell^4}{24}=D(r_+r_-)^2,
 \qquad
 \int_X\ell^2C=2D(r_+^2+r_-^2).
\] It follows that \[
 \chi(\mathcal E)
 =D\frac{\operatorname{Nm}(r)^2-Nu\operatorname{Tr}(r^2)+N^2u^2}{N^3}
 =D\frac{\operatorname{Nm}_{F/\mathbf Q}(Nu-r^2)}{N^3}.
 \tag{S33.12}
\] By (S33.11), \(Nu-r^2=N(u-r^2/N)\) has \(\pi\)-valuation \(2m\). Its local field norm has \(p\)-valuation \(2m\), since this ramified quadratic extension has residue degree one. Therefore \[
 v_p(\chi(\mathcal E))=2m-3m=-m<0,
 \tag{S33.13}
\] contradicting the integer Euler characteristic. This proves the main assertion.

For the semiregularity corollary, S-22 gives \(4D(Nq)^2/|G|=k\in\{2,4\}\). If an odd ramified prime does not divide \(D|G|\), then \(u=Nq\) is a \(p\)-adic unit. The degree-two character equation forces \(B=-\ell/N\), so the proposed normalized character is exactly (S33.3). The Néron--Severi hypothesis gives \(\ell\in F\theta\). The contradiction already proved shows that every odd ramified prime divides \(D|G|\). At \(2\), when \(2\nmid D|G|\), the value \(k=2\) would give the rational square \(u^2\) odd \(2\)-adic valuation \(-1\), which is impossible; \(k=4\) gives a \(2\)-adic unit, to which S-32 applies. This proves the combined condition (S33.4).
\end{proof}
\begin{remark}[An independent third-character certificate]
There is also a local contradiction before using the top-degree Euler calculation. This records the exact coefficient and its sign.

Let \(\theta_{\sigma,i}\) be the embedding component of \(\theta|_{U_i}\). Put \[
 \eta_1=\theta_{+,1}\theta_{-,1}
       =\theta|_{U_1}^2/2.
\] This is a primitive \(\mathbf Z_p\)-volume form on \(U_1\). In the particular ordered basis \((e_1,se_1,f_1,sf_1)\) its coefficient is \(4\) up to sign, which is harmless at odd \(p\). The degree \((4,2)\) component for the direct sum \(U_1\oplus U_2\) satisfies \[
 \begin{aligned}
 (\ell^3/6)_{(4,2)}
    &=\operatorname{Nm}(r)\eta_1\ell|_{U_2},\\
 (\ell C)_{(4,2)}
    &=\eta_1\theta_{\bar r}|_{U_2}.
 \end{aligned}
\] The second formula follows by writing \(C=\theta_{+,1}\theta_{+,2}+\theta_{-,1}\theta_{-,2}\): wedging its first summand with \(\ell|_{U_1}\) picks out \(r_-\), and wedging its second picks out \(r_+\). Thus \[
 \bigl(\operatorname{ch}_3(\mathcal E)\bigr)_{(4,2)}
 =\eta_1\theta_S|_{U_2},
 \qquad
 S=\frac{\operatorname{Nm}(r)r}{N^2}-\frac{u\bar r}{N}
   =\frac{\bar r(r^2-Nu)}{N^2}.
 \tag{S33.14}
\] For any integral degree-two form \(\ell\) on a torus, all divided powers \(\ell^j/j!\) are integral: in an exterior basis the coefficient of each ordered \(2j\)-monomial is a Pfaffian with integral entries. Hence \[
 \operatorname{ch}_3(\mathcal E)
 =\frac{\ell^3}{6}-\frac{\ell c_2(\mathcal E)}2
                         +\frac{c_3(\mathcal E)}2
 \in H^6(X,\mathbf Z_p)
\] for every odd \(p\), including \(p=3\). Since \(\eta_1\) is primitive and the restricted local trace form is perfect, (S33.14) would imply \(S\in O\). But (S33.10)--(S33.11) imply \[
 v_\pi(S)=v_\pi(\bar r)-v_\pi(N)
                  +v_\pi(r^2/N-u)=m-2m+0=-m<0.
\] This independently contradicts integrality. The argument uses only the usual integral Chern classes of a perfect object and the exterior algebra of a torus; no general integral topological Chern-character theorem is needed.
\end{remark}

\begin{theorem}[S-34: Determinant descent excludes every rank-one equivariant balanced seed]
\label{thm:semireg-rm-rank-one-equivariant}
Let $(X,\theta)$ be a polarized complex abelian fourfold
with Rosati-invariant real multiplication by a genuine
real quadratic field $F$, of type $(2,2)$. Write
$\theta=\theta_1+\theta_2$ and
\[
 v_q=\left(1-\frac q2\theta_1^2\right)
       \left(1-\frac q2\theta_2^2\right),\qquad q\in\mathbf Q_{>0}.
 \tag{S34.1}
\]
Let $G\subset X\times\widehat X$ be a finite effective
subgroup with a specified coherent action on
$\operatorname{Perf}(X)$ by translation and tensor product
with degree-zero line bundles. Let $E$ be a perfect complex
of virtual rank one satisfying
$e^B\operatorname{ch}(E)=v_q$ for a rational algebraic
divisor class $B$. No lift $\widetilde E$ to
$\operatorname{Perf}(X)^G$ which is simple and gluable
in that category is fully semiregular.

Here simplicity means
$\operatorname{End}_{\operatorname{Perf}(X)^G}
(\widetilde E)=\mathbf C$, and gluability means the
vanishing of its negative self-Ext groups. Neither condition
is required of the underlying object $E$. No maximal-order
or minimal-divisor-space hypothesis is imposed.
\end{theorem}
\begin{proof}
Write the action functors and their coherent multiplication
as
\[
 \phi_g(K)=t_{a_g}^*K\otimes P_g,
 \qquad
 \mu_{g,h}:\phi_g\phi_h\xrightarrow{\sim}\phi_{gh}.
\]
A lift consists of maps
$\lambda_g:\phi_g(E)\xrightarrow{\sim}E$ satisfying
\[
 \lambda_{gh}\circ\mu_{g,h,E}
       =\lambda_g\circ\phi_g(\lambda_h).
 \tag{S34.2}
\]
All natural identifications below use these specified
multiplication maps; replacing them by arbitrary
isomorphisms of the same functors would lose the argument.

Put $L=\det(E)$. The determinant of a perfect complex and
its compatibility with pullback are supplied by
Knudsen--Mumford, \emph{The projectivity of the moduli
space of stable curves I}, Theorem~2, p.~42, with the
extended determinant construction on pp.~40--42.
For a line bundle $P$ and a perfect complex of virtual
rank $r$, the local alternating determinant construction
gives
\[
 \det(E\otimes P)=\det(E)\otimes P^r,
 \qquad \det(c\operatorname{id}_E)=c^r\operatorname{id}_{\det E}.
 \tag{S34.3}
\]
Since $r=1$, taking determinants in (S34.2) gives maps
\[
 \delta_g:t_{a_g}^*L\otimes P_g\xrightarrow{\sim}L
 \tag{S34.4}
\]
with exactly the original scalar coherence, not a nontrivial
power of it. More explicitly, composition of (S34.4) for
$h$ and then $g$ agrees with (S34.4) for $gh$ after the
line-bundle identification induced by $\mu_{g,h}$.
This follows by applying the determinant functor to
(S34.2); the scalar part of $\mu_{g,h}$ is unchanged by
(S34.3). The formula can also be checked on any local
bounded complex of vector bundles, where the exponent
is the alternating sum of its ranks, namely one.

Conjugate the action by $K\mapsto K\otimes L^{-1}$.
Its functor for $g$ becomes
\[
 K\longmapsto t_{a_g}^*K\otimes R_g,
 \qquad
 R_g=t_{a_g}^*L\otimes P_g\otimes L^{-1}.
 \tag{S34.5}
\]
The maps $\delta_g$ trivialize $R_g$, and the coherence
just proved says that these trivializations identify the
entire action with the ordinary coherent translation action
of $G$ through $g\mapsto a_g$. They cancel any scalar
two-cocycle as well as the tensor factors. In particular
$E'=E\otimes L^{-1}$ acquires an ordinary translation
linearization. If $a_g=0$, (S34.4) implies $P_g\cong
\mathcal O_X$. Effectivity then gives $g=1$. Thus
$G\to X$ is injective. Its image $H$ acts freely, and
\[
 \pi:X\longrightarrow Y=X/H
 \tag{S34.6}
\]
is a finite etale isogeny of degree $|G|$.

For clarity, perfect descent here can be realized directly,
without assuming that descent in an arbitrary derived
homotopy category is effective. On $\pi_*E'$ the ordinary
linearization gives a genuine finite-group action. Form its
averaging idempotent and its image
\[
 E_Y=\operatorname{im}\left(
  \frac1{|G|}\sum_{g\in G}g:\pi_*E'\longrightarrow\pi_*E'
                         \right).
 \tag{S34.7}
\]
The pushforward is perfect: etale locally on $Y$ the
cover is a disjoint union of copies of $Y$, and the
assertion is then the preservation of perfection by a
finite direct sum. Perfection is fpqc local by Stacks
Project, Lemma~36.12.4, tag \texttt{09UG}.
The perfect category is closed under direct summands, so
the idempotent has a perfect image. After pulling back to
$X$, flat base change identifies $\pi^*\pi_*E'$ with
$\bigoplus_{h\in H}t_h^*E'$. The linearization identifies
the invariant summand with the diagonal copy of $E'$.
Consequently $\pi^*E_Y\cong E'$, compatibly with descent.
The same construction on morphism complexes, or the
adjunction and exactness of finite-group invariants, gives
for every integer $j$
\[
 \operatorname{Ext}^j_Y(E_Y,E_Y)
      \cong \operatorname{Ext}^j_X(E',E')^G
      \cong
      \operatorname{Ext}^j_{\operatorname{Perf}(X)^G}
                   (\widetilde E,\widetilde E).
 \tag{S34.8}
\]
Thus the descended object is simple and gluable. The
underlying sheaf descent used by this construction is
Stacks Project, Proposition~35.5.2, tag \texttt{023T};
the explicit finite averaging construction supplies the
perfect-complex statement and the Ext comparison.

Full semiregularity is also preserved in this descent.
One may use the enhanced equivalence
$\operatorname{Perf}(Y)\simeq
\operatorname{Perf}(X)^H$ just constructed and naturality
of the full Hochschild trace, as in S-10 and Perry 2026,
Remark~2.3. Equivalently, suppose
$u\in\operatorname{Ext}^2_Y(E_Y,E_Y)$ has zero full trace.
Its pullback has zero full trace on $X$. Under (S34.8)
it is an invariant class. The forgetting map on
$HH_{-2}$ is injective for this translation action by
the identity-sector argument of S-21. Naturality
therefore makes the equivariant full trace of $u$ zero;
full semiregularity of the lift forces $u=0$.
Tensoring by $L^{-1}$ is an equivalence and preserves
this property. Hence a fully semiregular lift would
give a fully semiregular rank-one perfect $E_Y$ on $Y$.

The character and the real-multiplication data descend
with the required normalization. The degree-two component
of $e^B\operatorname{ch}(E)=v_q$ gives $B=-c_1(E)$,
so $\operatorname{ch}(E')=v_q$.
Conjugating the rational $F$-action by the isogeny
defines an $F$-action on $Y$, and
$\theta_Y=(\pi^*)^{-1}\theta$ is a rational ample class
for which that action is Rosati invariant. Choose a
positive integer $m$ such that $\theta'=m\theta_Y$
is an integral polarization. With the corresponding
two real components of $\theta'$, the equality of
characters reads
\[
 \operatorname{ch}(E_Y)
   =\left(1-\frac{q/m^2}{2}(\theta'_1)^2\right)
      \left(1-\frac{q/m^2}{2}(\theta'_2)^2\right).
 \tag{S34.9}
\]
All equalities follow by pulling back along $\pi$,
which is an isomorphism on rational cohomology. The
polarization volume changes by
$\int_Y(\theta')^4/24=m^4D/|G|$, and $q/m^2>0$.
S-30 excludes a simple gluable fully semiregular
rank-one object with (S34.9), independently of the
polarization degree or an integral maximal-order action.
This is the desired contradiction.
\end{proof}

\section{Bounded division in families}

Run: DEG\_L, 2026-11-21. Author: \texttt{/root}, GPT-6 Astra Pro; effort
setting not exposed. These are new formulations and proofs for this run.
All five have received an independent whole-proof check, recorded in
CHECK\_EVENTS.tsv; root reread the amended proofs in full on 21 November
2026. Equality of a cycle with a rational cohomology class always means
equality under the cycle-class map, not equality in a Chow group whose
right-hand side is undefined.

\subsection{DIV1. A proper parameter space for a bounded signed
representative}\label{bounded-division-div1.-a-proper-parameter-space-for-a-bounded-signed-representative}

\textbf{Setup.} Let \texttt{f:X\ -\textgreater{}\ S} be a smooth
projective morphism over a smooth irreducible complex algebraic variety
of finite type. Its fibres have dimension \texttt{g}. Fix a relatively
very ample line bundle \texttt{H}, an integer
\texttt{0\ \textless{}=\ p\ \textless{}=\ g}, and a global flat section
\texttt{t} of \texttt{R\^{}\{2p\}f\_*\ Q(p)}. Put \texttt{r=g-p}. For a
complex point \texttt{s}, let \texttt{c\_t(s)} be the least integer
\texttt{C\ \textgreater{}=\ 1} for which there are an integer
\texttt{1\ \textless{}=\ N\ \textless{}=\ C} and effective integral
\texttt{r}-cycles \texttt{Z+}, \texttt{Z-} on \texttt{X\_s}, each of
\texttt{H}-degree at most \texttt{C}, such that

\texttt{N\ t\_s\ =\ cl(Z+)\ -\ cl(Z-)}.

Set \texttt{c\_t(s)=infinity} if no such representative exists. Zero
cycles in the sense of the zero element of the cycle group are allowed
among the effective cycles, including when \texttt{r\textgreater{}0}.

\textbf{Statement.} For each integer \texttt{C}, the subset
\texttt{S\_C=\{s\ :\ c\_t(s)\textless{}=C\}} is Zariski closed.
Consequently:

\begin{enumerate}
\def\labelenumi{\arabic{enumi}.}
\tightlist
\item
  A bound on any Zariski-dense subset gives the same bound on all of
  \texttt{S}.
\item
  If \texttt{t\_s} is algebraic at every complex point, then
  \texttt{S\_C=S} for some \texttt{C}.
\item
  The same conclusion holds if \texttt{t\_s} is algebraic outside a
  countable union of proper Zariski-closed subsets.
\end{enumerate}

Thus algebraicity on the entire family is equivalent to bounded signed
integral complexity on a dense subset. This conclusion supplies
existence of a bound, not an algorithm for computing it.

\textbf{Proof.} We spell out the finiteness which the informal reference
to a bounded-degree Chow scheme suppresses. Grothendieck,
\emph{Techniques de construction et théorèmes d'existence en géométrie
algébrique IV: Les schémas de Hilbert}, Séminaire Bourbaki, exposé 221
(1960/61), Lemme 2.4, printed p.255, states that the reduced,
pure-dimensional closed subschemes of fibres of a projective family with
bounded degree form a bounded family. Applying it in dimension
\texttt{r} and degree at most \texttt{C}, and stratifying a finite-type
parameter space by flatness, gives a finite list of Hilbert polynomials
\texttt{P\_1,...,P\_k} containing those of all reduced irreducible
subvarieties which can occur as components of these cycles. Each
corresponding relative Hilbert scheme is projective over \texttt{S}, by
the same exposé, Théorème 3.2, printed p.260.

Every nonzero effective cycle of degree at most \texttt{C} is a sum
\texttt{sum\_i\ m\_i{[}V\_i{]}}, where \texttt{V\_i} is reduced and
irreducible of dimension \texttt{r}, \texttt{m\_i} is a positive
integer, there are at most \texttt{C} summands, and
\texttt{m\_i\textless{}=C}. Indeed every component has positive integral
degree. There are only finitely many lists of the polynomials and
multiplicities with total degree at most \texttt{C}. For each list take
the corresponding fibre product over \texttt{S} of Hilbert schemes.
Include \texttt{S} itself to parameterize the zero cycle. This gives
finitely many projective parameter schemes whose universal
top-dimensional fundamental cycles include every desired effective
cycle. Additional points with nonreduced schemes or lower-dimensional
embedded components cause no problem: the top-dimensional fundamental
cycle is still effective and its degree is the leading coefficient of
the fixed Hilbert polynomial, multiplied by \texttt{r!}.

Take pairs of these parameter schemes for the two signs. On each such
finite-type projective parameter scheme \texttt{T\ -\textgreater{}\ S},
the class of the difference of the two universal fundamental cycles is a
flat section of the pulled-back cohomology local system. The following
argument verifies this also when \texttt{T} is singular or nonreduced.
For each universal subscheme \texttt{Z\ -\textgreater{}\ T}, its
structure sheaf is coherent and \texttt{T}-flat, and the ambient
morphism \texttt{X\_T\ -\textgreater{}\ T} is smooth. Thus \texttt{O\_Z}
is an absolutely perfect complex on \texttt{X\_T}: locally the ambient
ring is flat over the base with regular fibres of bounded dimension, and
a base-flat coherent module has finite Tor amplitude over that ambient
ring. Precisely this assertion is Stacks Project, Lemma 15.79.5, tag
\texttt{068X}, with base Tor amplitude \texttt{{[}0,0{]}}. It is
absolute perfection on \texttt{X\_T}, rather than only relative
perfection over \texttt{T}, that is used here.

One may work over affine opens of \texttt{S}; then \texttt{X\_T} is
quasi-projective, so the perfect complex admits a bounded locally free
resolution and defines a topological Chern character on
\texttt{X\_T(C)}. Derived restriction to any fibre is the ordinary
\texttt{O\_\{Z\_t\}} because of \texttt{T}-flatness. Its component
\texttt{ch\_p(O\_\{Z\_t\})} equals the Betti class of the
top-dimensional fundamental cycle \texttt{{[}Z\_t{]}\_r}, including its
multiplicities. An exact algebraic locator is Fulton, \emph{Intersection
Theory}, second edition, Theorem 18.3(3),(5), printed p.353, together
with Example 18.3.11, printed p.363: the latter identifies the top term
of the Riemann--Roch class with the associated cycle of any coherent
sheaf; on a smooth ambient fibre the Riemann--Roch class is
\texttt{ch(O\_\{Z\_t\})\ td(TX\_t)\ cap\ {[}X\_t{]}}, and the Todd
factor has constant term one. Lower-dimensional support components do
not contribute in codimension \texttt{p}.

Therefore all the fibre classes in question are restrictions of one
topological class \texttt{ch\_p(O\_Z)} on \texttt{X\_T(C)}. The smooth
proper family \texttt{X\ -\textgreater{}\ S} is locally a differentiable
fibre bundle; its pullback to \texttt{T(C)} is topologically locally
trivial. Restricting a global cohomology class in such a family gives a
flat section, since the fibre inclusions over any path in a trivializing
neighborhood are homotopic. Complex algebraic spaces are locally path
connected, and their finitely many connected components agree in the
analytic and algebraic topologies. The argument also treats nilpotents
because passing to the reduction does not alter the underlying analytic
topological space. Taking sums and differences gives the asserted flat
section for the universal signed cycle.

For any fixed \texttt{N} between 1 and \texttt{C}, equality of this
section with \texttt{N\ t} therefore selects a union of connected
components of \texttt{T}. There are finitely many such components, each
open and closed. The selected union is projective over \texttt{S}. Its
image is Zariski closed. A finite union over the parameter schemes and
\texttt{N} is precisely \texttt{S\_C}, proving the first assertion. A
closed subset containing a Zariski-dense subset is all of \texttt{S},
which proves (1).

Algebraicity of a rational class at a point is equivalent to finite
\texttt{c\_t(s)}: clear all coefficients of an algebraic rational cycle
and split the resulting integral cycle by sign. Hence algebraicity at
all points says \texttt{S(C)=union\_\{C\textgreater{}=1\}\ S\_C(C)}. If
every \texttt{S\_C} were proper, choose a small closed analytic ball
with nonempty interior in \texttt{S(C)}. Each \texttt{S\_C} meets that
ball in a closed nowhere-dense set, since a proper algebraic subset of
an irreducible smooth complex variety contains no analytic open subset.
The Baire category theorem rules out a countable covering of the ball by
these sets. Thus some \texttt{S\_C=S}. The same argument, adding the
given countably many exceptional proper closed subsets to the purported
covering, proves (3). QED.

\textbf{Scope and Grothendieck locator.} This is a statement about
relative algebraic cycles and their Betti realization. It uses FGA 221,
Lemme 2.4 and Théorème 3.2 for boundedness and properness. The analogous
cycle classes in any compatible Weil realization are transported by the
same algebraic correspondences, but a new assertion of integral Tate
surjectivity is not part of the statement. It proves the precise form
needed from sketch Route 2, statement 2.1; it does not use André--Oort
or any asserted density of a particular Hecke orbit.

\subsection{DIV2. Bounded witnesses on a dense set suffice in a Weil
family}\label{bounded-division-div2.-bounded-witnesses-on-a-dense-set-suffice-in-a-weil-family}

\textbf{Statement.} Let \texttt{S} be a smooth irreducible complex
algebraic variety of finite type, as in DIV1. Let
\texttt{A\ -\textgreater{}\ S} be a smooth projective family of
polarized abelian varieties of dimension \texttt{2n} and Weil type for a
fixed imaginary quadratic order, with globally defined order action,
fixed polarization type and a fixed rational flat Weil class \texttt{t}.
A finite base change may be made to define the action and the class.
Assume there is a constant \texttt{D\_0} and a Zariski-dense subset
\texttt{T} such that for every \texttt{s\ in\ T} there is an effective
integral \texttt{n}-cycle \texttt{Z\_s} of degree at most \texttt{D\_0}
whose Weil component is nonzero. Then \texttt{c\_t} is bounded on all of
\texttt{S}. An explicit bound on \texttt{T}, and therefore on
\texttt{S}, is the ceiling of the maximum of the two quadratic
expressions in B8 evaluated at \texttt{D=D\_0}, together with 1.

\textbf{Proof.} Apply B6--B8 with the fixed generic quadratic action.
Their projector denominator, signed-degree multiplier, polarization
volume and Weil-lattice discriminant are constant on the family; so is
the norm of a fixed flat multiple of \texttt{t}. B8 produces, for each
witness, an integral multiple of \texttt{t\_s} and signed effective
representatives with the displayed uniform quadratic bounds. This gives
\texttt{T\ subset\ S\_C} for one integer \texttt{C}. DIV1 gives
\texttt{S\_C=S}. If the polarization is only relatively ample, replace
it by a fixed relatively very ample power; degrees change by the fixed
factor equal to that power raised to \texttt{n}, so boundedness is
equivalent. QED.

The conclusion removes the need to control denominators of a projector
defined using each point's larger CM field. It leaves a geometric
requirement: uniformly bounded effective witnesses with a nonzero
component. It does not prove that requirement from the existence of
rational cycles on each CM fibre.

\subsection{DIV3. What a successful semiregular cycle construction would
prove}\label{bounded-division-div3.-what-a-successful-semiregular-cycle-construction-would-prove}

\textbf{Statement.} In the setting of DIV1, with
\texttt{1\ \textless{}=\ p\ \textless{}=\ g}, suppose at one point
\texttt{s\_0} there are smooth regularly embedded subvarieties
\texttt{Y\_1,...,Y\_k} of codimension \texttt{p} and fixed integers
\texttt{a\_j,N}, \texttt{N\textgreater{}=1}, with
\texttt{N\ t\_\{s\_0\}=sum\_j\ a\_j\ cl(Y\_j)}. Suppose each embedding
is Bloch semiregular and the flat continuation of \textbf{each full
class} \texttt{cl(Y\_j)} is of Hodge type on an analytic neighborhood of
\texttt{s\_0}. Then \texttt{c\_t} is bounded on all of \texttt{S}.

\textbf{Proof.} Apply Buchweitz--Flenner, \emph{A semiregularity map for
modules and applications to deformations}, Compositio Math.137 (2003),
Theorem 5.2, p.175, with the singleton set of indices \texttt{I=\{p\}}.
The hypotheses are a deformation of a compact complex algebraic central
manifold over a smooth complex germ, a horizontal Hodge section
\texttt{alpha\_p}, and an \texttt{I}-semiregular subspace with
\texttt{alpha\_p(s\_0)=ch\_p(O\_Y)}. Our smooth projective family over
the analytic germ of \texttt{S} meets the ambient hypotheses. For a
regular codimension-\texttt{p} immersion,
\texttt{T\^{}2\_\{Y/X\}(O\_Y)=H\^{}1(Y,N\_\{Y/X\})},
\texttt{ch\_p(O\_Y)={[}Y{]}}, and Proposition 8.2, p.200, identifies the
component \texttt{tau\_\{p-1\}} with Bloch's map. Thus Bloch
semiregularity is exactly the required \texttt{I}-semiregularity, and
preservation of this full fundamental class is sufficient; no other
component of \texttt{ch(O\_Y)} is required.

The proof of Theorem 5.2, p.180, invokes the same lifting argument as
the proof of Theorem 5.1, p.179, using the subspace obstruction maps of
§§4.11 and 4.13 in place of the sheaf maps. That argument constructs
formal lifts over the infinitesimal neighborhoods of the smooth base
germ, uses a versal deformation, and applies Artin approximation to
obtain an actual local flat deformation. For subspaces this is a local
section of the relative Douady space, equivalently of the
analytification of the relative Hilbert scheme in the present projective
setting. Apply it separately to every \texttt{Y\_j} and intersect the
finitely many parameter neighborhoods. The local extensions have the
original Hilbert polynomials.

The degrees of the fundamental cycles of these extensions are constant,
because their Hilbert polynomials are fixed. Their cohomology classes
are the flat continuations of the original classes, so the displayed
signed equality with \texttt{N\ t} persists. Splitting the fixed
integers \texttt{a\_j} by sign gives effective representatives of
degrees at most \texttt{sum\_j\ \textbar{}a\_j\textbar{}\ deg\_H(Y\_j)},
uniformly on a nonempty analytic open set. It is Zariski dense in
\texttt{S}. DIV1 therefore extends the bound to every fibre.
Equivalently, the finite product of the relevant projective Hilbert
schemes has a component dominating that analytic open set; its proper
image is all of \texttt{S}. QED.

\textbf{Necessary distinctions.} Preserving only
\texttt{sum\_j\ a\_j\ cl(Y\_j)} does not meet the stated hypotheses.
Semiregularity of a perfect complex is likewise a statement about its
Atiyah-trace map and the Chern components detected by that map; it is
not supplied by its \texttt{K\_0} class. A projective bounded parameter
space or an independent local-to-cycle argument is needed when a sheaf
or complex replaces the embeddings. S-01--S-07 show why the original
rigid 032 ball quotient does not directly satisfy these conditions on a
varying generated abelian family.

\subsection{DIV4. An explicit upper bound for the natural Hecke
representative}\label{bounded-division-div4.-an-explicit-upper-bound-for-the-natural-hecke-representative}

\textbf{Statement.} Let \texttt{psi:A\_x\ -\textgreater{}\ A\_0} be an
isogeny of polarized abelian \texttt{g}-folds, with equal polarization
type and \texttt{psi\^{}*L\_0=m\ L\_x} at the level of first Chern
classes, where \texttt{m} is a positive integer. Suppose
\texttt{psi\^{}*\ t\_0=(a/b)t\_x}, with coprime integers \texttt{a!=0},
\texttt{b\textgreater{}0}. If
\texttt{N\_0t\_0=cl(Z\_0\^{}+)-cl(Z\_0\^{}-)} and both degrees are at
most \texttt{D\_0}, then

\texttt{c\_t(x)\ \textless{}=\ max(N\_0\ \textbar{}a\textbar{},\ b\ m\^{}p\ D\_0)}.

\textbf{Proof.} Isogenies in characteristic zero are finite flat, so
flat pullback of each effective integral cycle is effective integral.
Pulling back the given equality and multiplying by \texttt{b} gives
\texttt{N\_0\ a\ t\_x=b(psi\^{}*cl(Z\_0\^{}+)-psi\^{}*cl(Z\_0\^{}-))}.
Reverse the signs if \texttt{a\textless{}0}. B4 computes the degree of
each pullback as \texttt{m\^{}p} times the old degree. This gives
exactly the bound asserted. QED.

This upper bound is for a specified representative. It is not a lower
bound for the minimum \texttt{c\_t(x)}. When \texttt{a/b=m\^{}p},
normalizing the rational pullback keeps its numerical degree fixed while
its displayed denominator grows. T6.4--T6.5 cancel auxiliary source
covers and target multiplication for specially normalized correspondence
constructions; they do not cancel an arbitrary Hecke transport between
different target fibres.

\subsection{DIV5. What the known fourfold theorem already
supplies}\label{bounded-division-div5.-what-the-known-fourfold-theorem-already-supplies}

\textbf{Statement.} On any fixed finite-type family of complex abelian
fourfolds and any global rational flat Hodge class, the signed integral
complexity is uniformly bounded, provided one admits the supplied
Markman theorem of the rational Hodge conjecture for all abelian
fourfolds. The same conclusion applies to any higher-dimensional family
for which the supplied theorem proves algebraicity of the particular
flat class on every fibre.

\textbf{Proof.} The admitted theorem supplies algebraicity pointwise. If
the base is smooth and irreducible, DIV1(2) supplies a uniform bound.
For a general finite-type base, choose a finite stratification of its
reduction into smooth irreducible locally closed complex subvarieties,
apply DIV1(2) to the restricted family and flat class on each stratum,
and take the maximum of the finitely many resulting bounds. Fibres over
the original base and its reduction are identical at complex points.
QED.

This bound is a posteriori and ineffective in this argument. It does not
calculate a theta degree, identify a Rallis factor, or extend the known
fourfold theorem to the open Weil families in sketch T8. The exact
supplied Markman locators are Theorem 1.5.1 (\texttt{thm-algebraicity}),
proving algebraicity of the Hodge--Weil classes on polarized abelian
sixfolds of Weil type for an arbitrary imaginary quadratic field and
discriminant \texttt{-1}, and Corollary 1.6.1, asserting the Hodge
conjecture for all abelian fourfolds. These numbers were checked from
the source's subsection-based shared theorem counter after removing
comments and hidden draft blocks; the sixfold assertion must not be
broadened to arbitrary discriminant or arbitrary Hodge classes.

\section{Polarization constraints and syzygy candidates}

Run \texttt{DEG\_L\_20261121}, author \texttt{/root}, GPT-6 Astra Pro,
effort not exposed. These statements attack the replacement-object gate
after S-09 and S-14. They do not assert that the required weakly
equivariant object already exists.

\subsection{DER1. A strict polarization-discriminant test for a
Hermite--Einstein
candidate}\label{derived-candidates-der1.-a-strict-polarization-discriminant-test-for-a-hermiteeinstein-candidate}

Let \texttt{(A,H)} be an abelian variety of Weil type of dimension
\texttt{2n}, \texttt{n\textgreater{}=2}, with \texttt{theta=c1(H)}. Let
\texttt{0!=t} be a rational Weil class, so \texttt{theta\ t=0}. Let
\texttt{E} be a holomorphic vector bundle of positive rank \texttt{r},
let \texttt{B} be a rational algebraic divisor class, and assume

\[
 w=e^B\operatorname{ch}(E)=P(\theta)+\lambda t,\qquad
 P(x)=r+p_1x+p_2x^2+\cdots,\quad \lambda\ne0.
\]

If \texttt{E} admits a Hermite--Einstein metric for the polarization,
then necessarily

\[
 p_1^2-2rp_2>0.
\tag{DER1}
\]

In particular a single exponential \texttt{P(x)=r\ e\^{}\{ax\}} is
impossible. More generally any finite sum
\texttt{P(x)=sum\_i\ r\_i\ e\^{}\{a\_i\ x\}} with positive rational
\texttt{r\_i} is impossible. These are obstructions to this
Hermite--Einstein route, not claims that every semiregular bundle is
Hermite--Einstein.

\textbf{Proof.} Write
\texttt{Delta(E)=c1(E)\^{}2-2r\ ch2(E)=2r\ c2(E)-(r-1)c1(E)\^{}2}.
Multiplication of a Chern character by \texttt{e\^{}B} leaves this
expression invariant: its first two components become \texttt{c1+rB} and
\texttt{ch2+B\ c1+rB\^{}2/2}, and the extra terms cancel. If
\texttt{n=2}, the Weil contribution to \texttt{ch2} is
\texttt{lambda\ t}, whose product with \texttt{theta\^{}\{2n-2\}}
vanishes because \texttt{theta\ t=0}. If \texttt{n\textgreater{}2},
there is no Weil contribution to \texttt{ch2}. Therefore

\[
 \int_A\Delta(E)\theta^{2n-2}
 =(p_1^2-2rp_2)\int_A\theta^{2n}.
\]

Kobayashi, \emph{Differential Geometry of Complex Vector Bundles},
Theorem 4.4.7, printed p.106 in the 2013 retypesetting of the 1987 book,
says that this number is nonnegative and that equality for an
Einstein--Hermitian bundle forces projective flatness. The full
statement and its curvature proof were read in the official Mathematical
Society of Japan PDF. Projective flatness gives
\texttt{ch(E)=r\ exp(c1(E)/r)}, by the same book, equation (2.3.3),
printed p.41. Thus in the equality case

\[
 w=r\exp(B+c_1(E)/r)
   =r\exp((p_1/r)\theta),
\]

where the second identity follows by comparing degree-two components of
\texttt{w}. Its Weil component is zero, contradicting
\texttt{lambda\ t!=0}. Hence the inequality is strict.

For a sum of exponentials the same coefficient is

\[
 p_1^2-2rp_2
 =-\sum_{i<j}r_i r_j(a_i-a_j)^2\le0.
\]

This contradicts the strict inequality, including the case in which all
slopes coincide. QED.

\textbf{Exact primary source.}
\url{https://www.mathsoc.jp/assets/pdf/publications/pubmsj/Vol15.pdf};
Theorem4.4.7, PDF page116 / printed p106, and equation2.3.3, PDF page51
/ printed p41. This is a cited curvature theorem applied to the
particular character, not a new proof of the Kobayashi--Lübke
inequality.

\subsection{DER2. The globally generated syzygy construction preserves a
full trace
kernel}\label{derived-candidates-der2.-the-globally-generated-syzygy-construction-preserves-a-full-trace-kernel}

Let \texttt{A} be an abelian variety of dimension
\texttt{g\textgreater{}=3}. Let \texttt{E} be a globally generated
vector bundle with \texttt{H\^{}i(A,E)=0} for all
\texttt{i\textgreater{}0}. Put \texttt{V=H\^{}0(A,E)} and let

\[
 0\longrightarrow M_E\longrightarrow V\otimes\mathcal O_A
 \longrightarrow E\longrightarrow0
\tag{DER2a}
\]

be the evaluation sequence. Then the Fourier--Mukai functor with kernel
the ideal of the diagonal induces an injection

\[
 \operatorname{Ext}^2_A(E,E)
 \hookrightarrow \operatorname{Ext}^2_A(M_E,M_E).
\tag{DER2b}
\]

Under this injection, the kernel of the full Buchweitz--Flenner
semiregularity map goes into the corresponding kernel for \texttt{M\_E}.
Consequently, if \texttt{E} is not fully semiregular, neither is
\texttt{M\_E}.

\textbf{Proof.} Let \texttt{Psi=Phi\_\{I\_Delta\}} on \texttt{Dperf(A)}.
The exact sequence of diagonal kernels induces the functorial
distinguished triangle

\[
 \Psi(E)\longrightarrow R\Gamma(A,E)\otimes\mathcal O_A
 \longrightarrow E\xrightarrow{\eta_E}\Psi(E)[1].
\]

The vanishing hypothesis identifies \texttt{RΓ(E)} with \texttt{V} in
degree zero; surjectivity of evaluation identifies \texttt{Psi(E)} with
the vector bundle \texttt{M\_E} in degree zero. Serre duality and
\texttt{K\_A=O\_A} give

\[
 \operatorname{Ext}^2(E,\mathcal O_A)
 =H^2(E^\vee)
 \simeq H^{g-2}(E)^*=0,
\]

since \texttt{g-2\textgreater{}0}. Apply \texttt{Hom(E,-)} to the
triangle shifted by two. Because
\texttt{Hom(E,V\ tensor\ O\_A{[}2{]})=0}, composition with
\texttt{eta\_E{[}2{]}} injects \texttt{Ext²(E,E)} into
\texttt{Hom(E,Psi(E){[}3{]})}.

If \texttt{u:E\ -\textgreater{}\ E{[}2{]}} has \texttt{Psi(u)=0},
naturality of \texttt{eta} gives

\[
 \eta_E[2]\circ u=\Psi(u)[1]\circ\eta_E=0.
\]

The preceding injection implies \texttt{u=0}. This proves (DER2b).

The full categorical semiregularity morphism is natural for a functor
between smooth proper categories: Perry, \emph{The semiregularity
theorem for equivariant noncommutative varieties}, arXiv2604.00511v2,
Remark2.3, printed p.5. Remark2.5 on that page identifies it with the
Buchweitz--Flenner map under HKR. Hence

\[
 \sigma_{M_E}(\Psi(u))=\Psi_*\sigma_E(u).
\]

A nonzero element in \texttt{ker\ sigma\_E} therefore has a nonzero
image in \texttt{ker\ sigma\_\{M\_E\}}. No assertion that \texttt{Psi}
is an equivalence is used or needed. QED.

\textbf{Scope.} This concerns the full object trace map. A map-level
obstruction kernel \texttt{ker\ tau\_f} cannot be inserted unless its
nonzero image in \texttt{Ext²(Rf\_*O,Rf\_*O)} has separately been
proved. It also does not rule out unrelated equivariant structures which
eliminate the original invariant kernel. For a finite translation action
preserving all displayed data, the same argument is equivariant and
preserves a kernel already present in the invariant category; no
assertion is made for arbitrary tensor actions, which need not commute
with the evaluation functor.

\subsection{DER3. Opposite exponential coefficients fix the numerical
test but not this trace
kernel}\label{derived-candidates-der3.-opposite-exponential-coefficients-fix-the-numerical-test-but-not-this-trace-kernel}

Suppose an integral \texttt{K\_0} class with character \texttt{N\ t} has
been realized, as in S-16, by a vector bundle \texttt{E\_0} with

\[
 \operatorname{ch}(E_0)=R e^{a\theta}+N t,
 \qquad R>0,\quad a\in\mathbf Z.
\]

Assume \texttt{E\_0} has an established nonzero full semiregularity
kernel. This hypothesis holds in the precise Q(i) fourfold situation of
S-13; it is not asserted for every higher-dimensional Weil family.

For every sufficiently large positive integer \texttt{q}, write
\texttt{b=a+q}, set \texttt{E=E\_0\ tensor\ H\^{}q}, and let
\texttt{M=h\^{}0(A,E)}. Then the evaluation kernel is a vector bundle of
positive rank \texttt{M-R}, with

\[
 \operatorname{ch}(M_E)
 =M-R e^{b\theta}-N t.
\tag{DER3a}
\]

Its polynomial character has strictly positive discriminant coefficient

\[
 p_1^2-2(M-R)p_2=MRb^2>0,
\tag{DER3b}
\]

but \texttt{M\_E} is not fully semiregular.

\textbf{Proof.} Choose \texttt{q} so that \texttt{E} is globally
generated, has no higher cohomology, \texttt{b!=0}, and
\texttt{M\textgreater{}R}. Global generation and higher-cohomology
vanishing hold for large twists by Stacks Project, Proposition
28.27.13(7), tag \texttt{01Q3}, and Lemma 30.16.1(4), tag \texttt{0B5T},
respectively; the Hilbert polynomial has positive leading coefficient,
giving \texttt{M\textgreater{}R}. Since \texttt{theta\ t=0},
multiplication by \texttt{e\^{}\{q\ theta\}} fixes \texttt{t}, so
\texttt{ch(E)=R\ e\^{}\{b\ theta\}+Nt}. Taking the character of the
evaluation sequence gives (DER3a). Here \texttt{p1=-Rb} and
\texttt{p2=-Rb²/2}; direct substitution gives (DER3b).

Tensoring by a line bundle is an equivalence, so S-10, or Perry
Remarks2.3 and2.5 applied in both directions, preserves the nonzero full
trace kernel of \texttt{E\_0}. DER2 then preserves a nonzero kernel in
\texttt{M\_E}. Thus this canonical two-exponential construction passes
the necessary numerical inequality while failing the desired trace
injectivity. QED.

\textbf{Consequences for the next approach.} The polarization polynomial
cannot be chosen solely to satisfy a positivity inequality. A successful
replacement must change the relevant obstruction representation or its
invariant trace kernel. The tensor--determinant construction after a
successful deformation is different: it is used to produce effective
cycles and does not require the resulting bundle to be semiregular.

\subsection{DER4. Adams untwisting removes the virtual-rank
condition}\label{derived-candidates-der4.-adams-untwisting-removes-the-virtual-rank-condition}

Let \texttt{A\ -\textgreater{}\ S} be a smooth projective family of
complex abelian varieties of dimension \texttt{2n}, with \texttt{S}
smooth and irreducible and a relatively very ample line bundle
\texttt{H}. Let \texttt{theta=c1(H)} be the relative polarization class
and \texttt{0!=t} a global flat Weil class, with \texttt{theta\ t=0}.
Suppose a point \texttt{0} has an object and actual finite target-group
action satisfying Perry Theorem 1.2, with

\[
 e^{B_0}\operatorname{ch}(E_0)=P(\theta_0)+\lambda t_0,
 \qquad \lambda\ne0.
\tag{DER4a}
\]

The complete displayed character is required to remain Hodge. There is
\textbf{no restriction on the virtual rank of \texttt{E\_0}}. Then the
conclusion can be strengthened from existence of algebraic classes to a
concrete construction of a fixed signed-cycle bound, as follows.

After passage to an etale neighborhood \texttt{U\ -\textgreater{}\ S},
let \texttt{mu\_M} be the band of a Kummer gerbe used in Perry's proof,
enlarged as there if necessary. There is an ordinary integral perfect
\texttt{K\_0} class \texttt{xi} on \texttt{A\_U} whose central-fibre
character is

\[
 \operatorname{ch}(\xi_0)=P(M\theta_0)+M^n\lambda t_0.
\tag{DER4b}
\]

Choose a finite bundle presentation and integers \texttt{R,m,a,k} by the
construction below. Write

\[
 Q(x)=e^{kx}\bigl(P(Mx)+R e^{-mx}+a\bigr),
 \qquad \mu=M^n\lambda,
\]

and define

\[
 C_n=[u^n]\exp\!\left(
 \sum_{j=1}^{n}(-1)^{j-1}(j-1)![x^j]Q(x)u^j\right),
 \qquad \beta=(-1)^{n-1}(n-1)!\mu.
\tag{DER4c}
\]

Let \texttt{v=integral\_A\ theta\^{}\{2n\}}. For an integer
\texttt{d\textgreater{}=1} clearing the denominators of \texttt{C\_n}
and \texttt{beta}, put

\[
 N=|d\beta|,\qquad D=dC_n v.
\tag{DER4d}
\]

The construction gives \texttt{C\_n\textgreater{}0}, \texttt{N} a
positive integer, and effective integral cycles of degree at most
\texttt{D} representing \texttt{N\ t} as their difference on every fibre
of \texttt{S}. Thus

\[
 c_t(s)\le \max\{1,N,\lceil D\rceil\}\qquad(s\in S(\mathbf C)).
\tag{DER4e}
\]

These numbers are explicit once the successful local deformation and its
finite bundle presentation and generation twists are specified. The
assertion does not bound those choices in terms of the original 032
level.

\textbf{Proof: the gerbe and the Adams operation.} Perry,
arXiv2604.00511v2, section6.2, proofs of Theorems1.1 and1.2, printed
pp26--27, constructs the deformation using a \texttt{mu\_M} gerbe
\texttt{pi:mathcal\ A\ -\textgreater{}\ A\_U} whose central class is the
reduction of \texttt{M\ B\_0} modulo \texttt{M}. The integer may be
increased by a multiple when extending the group action, as specified in
Corollary5.17 and the proof on p27. Replacing the initial integer by
this multiple does not change \texttt{B\_0}. The category of twisted
perfect complexes is explicitly the category of weight-one perfect
complexes on this gerbe (section3.2, pp7--8). Hence the resulting object
\texttt{F} is of weight one. Under the central trivialization specified
by the algebraic B-field, the gerbe is the gerbe of \texttt{M}th roots
of a line bundle \texttt{L\_0} with \texttt{c\_1(L\_0)=M\ B\_0}; its
tautological root \texttt{T} has \texttt{T\^{}M=pi\^{}*L\_0}, and

\[
 F_0\simeq \pi^*E_0\otimes T.
\]

The sign is fixed by the convention that the weight-one Chern character
under the rational cohomology identification is
\texttt{exp(B\_0)\ ch(E\_0)}, as in Perry Theorem3.2(2), equation(3.3),
p9.

For a perfect complex in characteristic zero, each derived exterior
power is again perfect. One can form it as the alternating idempotent
summand of a derived tensor power; these summands exist in the
idempotent-complete perfect category. Exterior powers of a weight-one
object have weight equal to their exterior exponent. Newton's recursion

\[
 \psi^j(x)=\lambda^1(x)\psi^{j-1}(x)
 -\lambda^2(x)\psi^{j-2}(x)+\cdots
 +(-1)^{j-2}\lambda^{j-1}(x)\psi^1(x)
 +(-1)^{j-1}j\lambda^j(x)
\tag{DER4f}
\]

therefore writes \texttt{psi\^{}M({[}F{]})} as an integral finite
virtual sum of perfect objects, every term of weight \texttt{M}, hence
weight zero modulo \texttt{M}. A weight-zero perfect complex on a tame
\texttt{mu\_M} gerbe descends uniquely to the base: this is checked
etale-locally on the neutral gerbe, where a weight-zero representation
is the trivial representation tensored with a complex on the base; the
identifications descend. Applying this to each term defines the ordinary
integral class \texttt{xi} on \texttt{A\_U}.

The Newton operations are additive and multiplicative on \texttt{K\_0},
and on a line bundle satisfy \texttt{psi\^{}M(L)=L\^{}M}. Their formal
source is SGA6, ExposeV, Appendix7.1, equation(7.1.1), printed
pp360--361, and Proposition7.5, pp362--364 (SLMath original scan, PDF
pages367--371). The weighted homogeneity in (DER4f) is also explicitly
noted in7.1. In the present characteristic-zero perfect setting, the
exterior-power filtration of a triangle supplies the same lambda
identities, so the recursion is compatible with passing to the perfect
\texttt{K\_0} group. On a split bundle with Chern roots \texttt{x\_i},
the Chern character of its Adams operation is
\texttt{sum\_i\ exp(Mx\_i)}. The splitting principle, and additivity for
a difference of bundles, consequently give
\texttt{ch\_j(psi\^{}M\ x)=M\^{}j\ ch\_j(x)}.

At the central fibre we may compute before descending:

\[
 \psi^M([F_0])
 =\pi^*\psi^M([E_0])\,[T]^M
 =\pi^*\bigl(\psi^M([E_0])\,[L_0]\bigr).
\]

Thus \(\xi_0=\psi^M(E_0)\otimes L_0\). Its Chern character is
\(\exp(MB_0)\operatorname{ch}(\psi^M(E_0))\), which is the degreewise
operation multiplying the degree-\(2j\) component of
\(\exp(B_0)\operatorname{ch}(E_0)\) by \(M^j\). Applying it to (DER4a)
gives (DER4b). All objects and Chern classes are defined over \(U\);
their fibre classes are flat, so this character formula persists on the
connected neighborhood.

\textbf{Proof: an ordinary bundle and its effective Chern cycle.} Shrink
\texttt{U} to a smooth affine neighborhood. The total space
\texttt{A\_U} is smooth and quasi-projective and has the resolution
property. Stacks Project Lemma36.37.2, tag0F8F, gives a bounded finite
locally free presentation of every perfect object. Grouping its even and
odd terms writes \texttt{xi={[}V\^{}+{]}-{[}V\^{}-{]}} for actual vector
bundles. Choose \texttt{m} so that \texttt{V\^{}-\ tensor\ H\^{}m} is
globally generated, and a finite surjection
\texttt{O\^{}R\ -\textgreater{}\ V\^{}-\ tensor\ H\^{}m}. Its kernel
\texttt{K} is locally free because its quotient is locally free. Hence

\[
 [V^+\oplus (K\otimes H^{-m})]
 =\xi+R[H^{-m}].
\]

Add \texttt{a\textgreater{}=0} trivial line bundles so that the
resulting bundle has rank at least \texttt{n}, then choose \texttt{k} so
that its tensor product \texttt{V} with \texttt{H\^{}k} is globally
generated and has no higher fibre cohomology on a neighborhood. The
relevant standard facts are Stacks Project Proposition28.27.13(7),
tag01Q3, and Lemma30.16.1(4), tag0B5T, together with cohomology and base
change; all twists and ranks are finite fixed integers. Since
\texttt{theta\ t=0}, the resulting character is

\[
 \operatorname{ch}(V_s)=Q(\theta_s)+\mu t_s.
\]

The identity between total Chern class and Chern character is

\[
 c_u(V_s)=\exp\left(\sum_{j\ge1}(-1)^{j-1}(j-1)!
 \operatorname{ch}_j(V_s)u^j\right).
\]

Only the linear occurrence of \texttt{t\_s} contributes in codimension
\texttt{n}; therefore

\[
 c_n(V_s)=C_n\theta_s^n+\beta t_s.
\tag{DER4g}
\]

For a globally generated bundle of rank \texttt{q\textgreater{}=n}, the
degeneracy cycle of \texttt{q-n+1} sufficiently general sections is
effective of codimension \texttt{n} and represents \texttt{c\_n(V\_s)};
if empty its class is zero. This is Fulton, Intersection Theory, second
edition, Theorem14.4, p254, and Examples14.4.2--3, pp255--256. Call that
cycle \texttt{Z\_s}. Its degree is \texttt{C\_n\ v}, because
\texttt{theta\_s\ t\_s=0}. Its Weil component is the nonzero class
\texttt{beta\ t\_s}, so it is not the zero cycle and
\texttt{C\_n\textgreater{}0}. Let \texttt{W\_s} be an intersection of
\texttt{n} sufficiently general polarization hyperplanes, representing
\texttt{theta\_s\^{}n}, of degree \texttt{v}. Then

\[
 d\beta t_s=\operatorname{cl}(dZ_s)
             -\operatorname{cl}(dC_n W_s),
\]

and both effective cycles have degree \texttt{dC\_n\ v}. Reversing signs
if \texttt{beta\textless{}0} gives (DER4d)--(DER4e) on the nonempty
etale image of this neighborhood. That image is Zariski open and dense
in \texttt{S}. Use the fixed-multiplier version of the closed-locus
argument in the proof of DIV1: retain this one \texttt{N} and degree
cutoff \texttt{D} in its finite projective parameter spaces, instead of
taking the finite union over multipliers. Its closed image contains the
dense open set and is all of \texttt{S}, preserving these same
\texttt{N,D}. This completes the proof.

\textbf{What was removed.} Tensor--determinant untwisting requires
nonzero virtual rank and gives a factor \texttt{r\^{}r}. Adams
untwisting uses the torsion order of the explicitly constructed gerbe,
preserves the first nonzero Weil component even in virtual rank zero,
and gives the factor \texttt{M\^{}n}. It takes place after deformation;
no semiregularity assertion about the new bundle is needed. The
remaining hypothesis is existence of the original weakly equivariant
semiregular object with its complete Hodge-compatible character.

\subsection{Check state}\label{derived-candidates-check-state}

Root read DER1--DER4 in full after writing. The semiregularity agent
independently read all four proofs and checked the central gerbe sign,
Adams weight/descent and effective Chern cycle. Its requests for the
exact bundle-resolution locator and fixed-N specialization argument are
incorporated above. Recorded events distinguish these readings from
formal verification. No assertion here proves existence of a new
semiregular object.

\section{Faithful operations and the remaining replacement gate}

Run \texttt{DEG\_L\_20261121}, author \texttt{/root}, GPT-6 Astra Pro,
effort not exposed. These findings strengthen DER2 and test broader
attempts to repair an established full semiregularity kernel. They do
not construct the successful replacement required by S-14.

\subsection{DER5. Nonzero-rank tensor and nonzero-Euler convolution
preserve every full trace
kernel}\label{derived-functors-der5.-nonzero-rank-tensor-and-nonzero-euler-convolution-preserve-every-full-trace-kernel}

Let \(A\) be a positive-dimensional complex abelian variety. For perfect
complexes define derived Pontryagin convolution by \[
 E*F=Rm_*(p_1^*E\otimes^{\mathbf L}p_2^*F),
\] where \(m\) is addition. If \(\chi(A,F)\ne0\), this functor is
injective on every morphism space: \[
 \operatorname{Hom}(E,E'[j])\hookrightarrow
 \operatorname{Hom}(E*F,(E'*F)[j])\qquad(j\in\mathbf Z).
\tag{DER5a}
\] It therefore injects \(\ker\sigma_E\) into \(\ker\sigma_{E*F}\). More
generally, tensoring by a perfect complex of nonzero virtual rank has
these same properties on every smooth connected projective complex
variety.

In particular, for the ideal of the diagonal, \[
 \operatorname{Ext}^j(E,E)\hookrightarrow
 \operatorname{Ext}^j(\Phi_{\mathcal I_\Delta}E,
                         \Phi_{\mathcal I_\Delta}E)
\tag{DER5b}
\] for \textbf{every} perfect complex \(E\) and every \(j\). The
evaluation-cone construction cannot remove a full semiregularity kernel
even without the global generation, cohomology vanishing, or
dimension-at-least-three hypotheses of DER2.

\textbf{Proof.} Let \(K\) be perfect of constant virtual rank \(r\ne0\)
on a smooth connected complex variety \(X\). Dualizability gives
coevaluation and graded trace maps \[
 \mathcal O_X\longrightarrow K\otimes^{\mathbf L}K^\vee
 \longrightarrow\mathcal O_X
\] whose composite is \(r\operatorname{id}\). Indeed, a bounded locally
free presentation computes the trace of the identity as the alternating
sum of its ranks. This local computation defines the global categorical
trace. For \(u:E\to E'[j]\), naturality identifies the composite through
\(u\otimes\operatorname{id}_K\otimes\operatorname{id}_{K^\vee}\) with
\(ru\). If \(u\otimes\operatorname{id}_K=0\), that composite vanishes;
since \(r\) is invertible over \(\mathbf C\), \(u=0\). This argument
applies to complexes of either sign of virtual rank and in every
morphism degree.

Let \(\mathcal F:D_{\mathrm{perf}}(A)\to D_{\mathrm{perf}}(A^\vee)\) be
the normalized Poincare Fourier--Mukai equivalence. Mukai, \emph{Duality
between \(D(X)\) and \(D(\widehat X)\) with its application to Picard
sheaves}, Nagoya Math.J.81(1981), Theorem2.2 p156 and equation(3.7)
p160, gives \[
 \mathcal F(E*F)\simeq
 \mathcal F(E)\otimes^{\mathbf L}\mathcal F(F).
\tag{DER5c}
\] There is no shift in this direction of the identity. The virtual rank
of \(\mathcal F(F)\) is \(\chi(A,F)\): derived base change identifies
its fibre at a degree-zero line bundle \(P\) with
\(R\Gamma(A,F\otimes P)\). The Euler characteristic of this family is
constant and at \(P=\mathcal O_A\) equals \(\chi(A,F)\). Apply the
preceding tensor argument to \(K=\mathcal F(F)\) and conjugate by the
equivalence, proving (DER5a).

For the evaluation cone, use the flat difference map \(d(x,y)=y-x\).
Then \(d^*\mathcal I_0=\mathcal I_\Delta\). The change of variables
\((x,z)\mapsto(x,x+z)\) in the integral transform gives \[
 \Phi_{\mathcal I_\Delta}(E)\simeq E*\mathcal I_0.
\] Since \(\dim A\ge1\), \(\chi(\mathcal O_A)=0\) and the exact sequence
of the origin gives \(\chi(\mathcal I_0)=-1\). One can also see the sign
directly from Mukai Theorem2.2 and Example2.6 p157: applying
\(\mathcal F\) to this sequence gives the triangle \[
 \mathcal F(\mathcal I_0)\longrightarrow
 \mathcal O_{\widehat0}[-\dim A]\longrightarrow
 \mathcal O_{A^\vee}\longrightarrow.
\] Its first term has virtual rank \(-1\). Thus (DER5a) applies and
proves (DER5b). Positive dimension is necessary, because the ideal of
the origin is zero in dimension zero.

For each exact functor \(\Psi\) under consideration, Perry,
arXiv2604.00511v2, Remarks2.3 and2.5 p5, gives \[
 \sigma_{\Psi(E)}(\Psi(u))=\Psi_*\sigma_E(u)
\] for the full semiregularity map. Faithfulness makes the image of a
nonzero kernel element nonzero, and the displayed identity keeps it in
the kernel. This conclusion persists under composition with other such
functors and with derived equivalences. No equivalence of the
convolution functor itself is asserted. QED.

\textbf{Scope.} This concerns full object semiregularity. A map-level
obstruction killed before reaching self-Ext is not covered. A different
equivariant structure on a new object, or a functor not satisfying the
displayed faithfulness criterion, is not excluded. The result closes the
proposed repair by allowing arbitrary cohomology in the evaluation cone;
DER2 is retained as its simpler bundle case.

\textbf{Primary source and Grothendieck placement.} Mukai's PDF is
\url{https://www.cambridge.org/core/services/aop-cambridge-core/content/view/BDBEBAC584BE15236C2D62C383A34245/S002776300001922Xa.pdf/duality-between-dx-and-with-its-application-to-picard-sheaves.pdf}.
The whole short proofs of Theorem2.2 and equation(3.7) were read. The
convolution identity uses the Poincare biextension additivity of SGA7I,
ExposeVII, section2.4 pp164--165 and Example2.9.5 p175. Virtual rank is
the perfect-complex K-theory invariant used in SGA6, Expose0 section1
p1. This placement does not attribute the later Fourier--Mukai
equivalence or semiregularity theorem to those earlier texts.

\subsection{DER6. A polarization-only tensor repair retains the kernel
or loses the Weil
component}\label{derived-functors-der6.-a-polarization-only-tensor-repair-retains-the-kernel-or-loses-the-weil-component}

Let \(A\) be polarized of Weil type, of dimension \(2n\) with \(n\ge2\).
Suppose a perfect \(E\) has a nonzero full semiregularity kernel and \[
 e^B\operatorname{ch}(E)=P(\theta)+\lambda t,
 \qquad\lambda\ne0,\quad\theta t=0.
\] Let \(K\) be a perfect complex whose entire Chern character is
\(Q(\theta)\). If \(Q(0)\ne0\), the tensor product
\(E\otimes^{\mathbf L}K\) retains a nonzero full kernel. If \(Q(0)=0\),
its \(B\)-normalized character is a polarization polynomial and its Weil
component is zero. Thus none of these tensor operations repairs full
semiregularity while retaining the specified nonzero Weil component.

\textbf{Proof.} The constant coefficient \(Q(0)\) is the virtual rank of
\(K\). If nonzero, DER5 gives the kernel injection. For either value of
the coefficient, \[
 e^B\operatorname{ch}(E\otimes^{\mathbf L}K)
 =P(\theta)Q(\theta)+\lambda Q(0)t,
\tag{DER6}
\] since all positive powers of \(\theta\) annihilate \(t\). The
polynomial part has zero Weil projection by B6. Hence when \(Q(0)=0\) no
Weil component remains. QED.

The same alternative holds on the Fourier--Mukai dual if the transformed
input has the displayed polynomial-plus-Weil form there and the
transformed convolution kernel has polarization-polynomial character.
This qualification is necessary; no general assertion about an arbitrary
\(B\)-field under Fourier transform is made. A correction using
different non-Lefschetz data or a new equivariant structure is not
excluded. Adding a direct summand likewise retains the original full
kernel by the block-diagonal trace calculation underlying S-08.

The Chern-character product belongs to SGA6, Expose0 section1,
equations(1.1)--(1.3), p1; the precise generic Weil projector is B6/B16.
The conclusion is an obstruction for these specified constructions, not
an assertion that all objects with the corrected K-class fail
semiregularity.

\subsection{DER7. Adams--determinant untwisting gives a multiplier
independent of the gerbe order in nonzero
rank}\label{derived-functors-der7.-adamsdeterminant-untwisting-gives-a-multiplier-independent-of-the-gerbe-order-in-nonzero-rank}

Keep the successful-deformation hypotheses of DER4, and write \[
 w=e^{B_0}\operatorname{ch}(E_0)=P(\theta_0)+\lambda t_0,
 \qquad P(x)=r+p_1x+\cdots.
\] Assume the virtual rank \(r\) is nonzero; it may be negative. Put
\(d=|r|\), \(a=\operatorname{sgn}(r)\), and \[
 \beta_0=(-1)^{n-1}(n-1)!\lambda,\qquad
 v=\int_{A_s}\theta_s^{2n}.
\] Choose an integer \(b>0\) with \(b\beta_0\in\mathbf Z\). Then the
single positive integer \[
 N_*=v\,|b\beta_0|\,|r|^n
\tag{DER7a}
\] works as a multiplier for signed integral representatives of \(t_s\)
on every fibre. A common degree bound is obtained by the
bundle-generation construction below. Unlike the degree bound, this
multiplier depends only on \(r,\lambda,n\) and the fixed polarization
volume, and not on the gerbe order, its enlargement for the group
action, a bundle resolution, or its generation twists.

\textbf{Proof: exact weight cancellation.} On the local finite gerbe
\(\pi:\mathcal A\to A_U\) in DER4, let \(F\) be the deformed weight-one
perfect object. Its determinant is an invertible object of weight \(r\):
for a bounded vector-bundle presentation, the determinant is the
alternating tensor product of determinants and its weight is the
alternating sum of the ranks. The determinant construction is invariant
under quasi-isomorphism and descends; an exact primary locator is
Knudsen--Mumford, \emph{The projectivity of the moduli space of stable
curves I}, Mathematica Scandinavica39(1976), Theorem2 p42, as used in
S-17.

The Newton expression for \(\psi^d([F])\) is an integral virtual sum
whose terms all have weight \(d\), by DER4. Tensoring every term by
\(\det(F)^{-a}\) makes its weight \[
 d-ar=|r|-\operatorname{sgn}(r)r=0.
\] Thus it descends term by term to an ordinary integral perfect K-class
\(\xi\) on \(A_U\). This is exact weight zero, not just weight zero
modulo a particular band order.

On the central root gerbe write \(F_0=\pi^*E_0\otimes T\). The factors
of the tautological root cancel, so \[
 \xi_0=\psi^d(E_0)\otimes\det(E_0)^{-a}.
\] Comparing the degree-two component of \(w\) gives
\(c_1(E_0)+rB_0=p_1\theta_0\). Consequently \[
 \begin{aligned}
 \operatorname{ch}(\xi_0)
 &=e^{-a c_1(E_0)}e^{-dB_0}
       \bigl(P(d\theta_0)+d^n\lambda t_0\bigr)\\
 &=e^{-a p_1\theta_0}P(d\theta_0)+d^n\lambda t_0.
 \end{aligned}
\tag{DER7b}
\] Here \(ar=d\) cancels the B-field and \(\theta_0t_0=0\) removes the
exponential on the Weil term. The character formula is flat on the
connected neighborhood, exactly as in DER4.

\textbf{Proof: integral cycles with the stated fixed multiplier.} Apply
DER4's finite bundle resolution, addition of a negative polarization
exponential and trivial summands, and positive generation twist to
\(\xi\). None of these changes the Weil coefficient. We obtain a
globally generated ordinary vector bundle \(V\) of rank at least \(n\)
with \[
 c_n(V_s)=C\theta_s^n+d^n\beta_0t_s.
\] Its general degeneracy cycle \(Z_s\) is effective, has this class,
and has degree \(Cv\). The nonzero Weil component implies \(Z_s\ne0\),
hence \(C>0\). Since \(Z_s\) is an integral cycle, its degree is an
integer: \[
 Cv\in\mathbf Z_{>0}.
\tag{DER7c}
\] This elementary integrality, rather than a denominator estimate on
the entire polarization polynomial, suffices. Set \(c=vb\). Then
\(cC=b(Cv)\) is a positive integer and \(c d^n\beta_0\) is a nonzero
integer. Let \(W_s\) be a general intersection of \(n\) polarization
hyperplanes. We have \[
 c d^n\beta_0t_s
   =\operatorname{cl}(cZ_s)-\operatorname{cl}(cC W_s).
\] Both cycles on the right are effective integral and have the same
degree \[
 D=vb\,\deg Z_s=vbCv.
\tag{DER7d}
\] Reverse the signs if necessary. The multiplier is exactly (DER7a).
The generation construction fixes \(C\) on the local neighborhood, hence
fixes \(D\). The fixed-\(N_*\), fixed-\(D\) projective closed-locus
argument of DIV1 propagates these same integers from its dense etale
image to every fibre. QED.

\textbf{Optional gcd refinement.} If the gerbe band is \(\mu_M\), put
\(d_0=\gcd(|r|,M)\) and choose \(a_0\in\mathbf Z\) with
\(a_0r\equiv d_0\pmod M\). Then \(\psi^{d_0}(F)\otimes\det(F)^{-a_0}\)
has weight zero modulo \(M\) and descends. If \(k=(d_0-a_0r)/M\), its
central class is \[
 \psi^{d_0}(E_0)\otimes\det(E_0)^{-a_0}\otimes L_0^k,
\] where \(T^M=\pi^*L_0\). Its character is \[
 e^{-a_0p_1\theta_0}P(d_0\theta_0)+d_0^n\lambda t_0.
\] Indeed the B-field exponent is \(-d_0B_0+a_0rB_0+(d_0-a_0r)B_0=0\).
The preceding cycle proof gives multiplier \(v|b\beta_0|d_0^n\).
Multiplying both effective cycles by \((|r|/d_0)^n\) recovers the same
\(N_*\) if desired. For virtual rank zero, take \(d_0=M,a_0=0\); this is
DER4, and no bound independent of \(M\) is asserted in that case.

\textbf{Relation to S-17 and the remaining gate.} The Adams degree can
be the absolute virtual rank, in place of the tensor exponent used by
the tensor--determinant construction. The Weil coefficient grows by
\(|r|^n\), irrespective of the larger powers arising from a tensor
expansion. All these operations occur after the successful deformation;
their output is used for cycles and is not claimed semiregular. The
required initial weakly equivariant semiregular object remains
unconstructed in the open cases. The operation uses SGA6, ExposeV,
Appendix7.1 pp360--361 and Proposition7.5 pp362--364, together with its
determinant formalism and the FGA221 closed parameter spaces of DIV1.

\subsection{Check status}\label{derived-functors-check-status}

Root and the semiregularity agent have read all three complete proofs.
DER5--DER6 passed the graded-trace, virtual-rank, no-shift convolution,
origin-ideal and dimension-zero checks. DER7 passed the exact-weight
cancellation for either sign of rank, central B-field calculation,
integral coefficient argument, fixed-multiplier closure and gcd
refinement. The events record actual whole-proof readings; no
proof-assistant check is claimed.

\clearpage
\part{Open questions and unproved implications}
\section{The remaining main questions}
No statement in this part is asserted as a theorem without an explicit
conditional antecedent. The exact raw degree is not a function of source
level alone. A complete normalization specifies the primitive integral
maps, their polarization, theta inputs, adelic measures and local factors.
The finite spectral formula does not automatically reduce every rank-two
standard $L$-function to Hecke characters.

The uniform geometric question is the existence, on dense fibres of each
open Weil family, of effective middle-dimensional cycles of bounded degree
whose Weil component is nonzero. B6--B8 and DIV1--DIV2 remove the generic
projector and bounded-division extraction from that question. They do not
prove the bounded witness assertion. The integral Hom covolume, its rank,
and the generic mapping degree must be retained when converting theta
norms into a bound for the reduced image.
T6.20 cancels all varying input-isogeny factors in the normalized Chow
class. T6.21 computes the remaining full-CM orbit saturation exactly and
proves a uniform integral division bound when one output factor has
maximal-order action. This does not identify that orbit with a generic
quadratic Weil plane or bound the degrees of its signed representatives.
B25 shows that an unbounded-stabilizer escape must already import
target-sized core complexity into the intermediate quotient.

For a specified smooth reduced 032 image the Bloch question remains the
surjectivity in S-02. The universal assertion for all allowed deeper-level
map presentations is refuted by S-05--S-06. A deformation of a projected
class requires a different semiregular object, or another argument; the
full unprojected rigid surface class need not stay Hodge. A projected
$K_0$ class is not a proof that a representative object is semiregular.
S-14 and DER4 reduce the sufficient deformation route to one gluable
weakly target-equivariant object with complete character
$e^{B_0}\operatorname{ch}(E_0)=P(\theta_0)+\lambda t_0$.
DER4 constructs fixed signed-cycle bounds from its local deformation
without a virtual-rank assumption; DER7 makes the multiplier independent
of the gerbe order when the virtual rank is nonzero. DER1--DER6 and
S-18 exclude several specified canonical attempts to construct that object.
S-30 also excludes every nonequivariant simple gluable rank-one balanced
real-multiplication object, and every positive-rank such realization with
an integral normalizing class. S-31 excludes vector-bundle ranks at most
three. Fractional normalization, genuinely equivariant constructions and
nonsimple objects require the further tests recorded in the later findings.

Every approach tried, its outcome, and the next precise gate are recorded
in the component reports, QUESTIONS.tsv, and LEDGER.md. Scheduled progress
reviews and their actual timestamps are recorded in the dated ledger;
the session remains active until the prompt's stopping rule is met.
\section{Consequences and exact framework locators}

This file is a live research deliverable. The session remains in
progress; open questions below have not been promoted to theorems. All
findings use run ID \texttt{DEG\_L\_20261121}; their check events are
collected in \texttt{CHECK\_EVENTS.tsv}.

\subsection{The degree and L-value
question}\label{IMPLICATIONS-the-degree-and-l-value-question}

T6.1 identifies the effective surface degree with a positive sum of
norms of all pulled-back polarization-coordinate wedges. T6.2 separates
it from the pure four-factor codimension-two output, whose degree for
the product polarization is zero. T6.3 refutes a raw degree determined
by source level alone, because the construction permits multiplication
of its denominator-clearing maps at the same level. The meaningful
normalized Chow cycle is unchanged under auxiliary covers and these map
multiplications once its divisor kernels are chosen as biextensions
(T6.4--T6.5).

T6.7 removes a proposed analytic obstruction: algebraic norm and trace
clear the nonzero extracted period without requiring a lower bound for
its individual complex absolute value. T6.8 bounds the effective degrees
created by fixed divisor kernels. These statements retain the
coefficient field, integral generator lattice and kernel constants
explicitly.

T6.9 uses the supplied human proof and the checked theta spanning
literature to cover every pulled-back holomorphic coordinate. The
Hodge-index estimate B5 therefore reduces the effective surface degree
to a finite character-theta Gram form. The global character factor
occurs at 3/2 and has uniform elementary Euler-product bounds for fixed
field degree (B1). B1a also computes the archimedean factor in the
actual cohomological normalization. B14 proves that unit input norm and
fixed target level do not bound a split local factor; it supplies sharp
shell-support and weighted-seminorm bounds instead. B17 then removes
arbitrary radial support from a specified existing split theta vector by
an exact one-shell section, using the global norm identity to prove
preservation of the vector itself. It also gives uniformly bounded
analytic test inputs, with actual compact level and averaging
denominators retained. The remaining issue is the size of the output in
its integral Hom normalization. B18 gives a uniform positive lower bound
for the normalized Weil period once a generic middle-dimensional
integral witness is already realized; it does not supply that preceding
realization or degree bound. T6.12--T6.13 now prove the exact finite
wedge-norm and degree formulas through rank-two sources, including
actual multiplicity copies, local forms and unitary intertwiners. T6.14
verifies the dependence on finite adelic level. T6.15 now constructs an
actual balanced theta wedge and CM-Albanese surface map whose degree
formula contains a rank-two source with nonzero Steinberg monodromy.
This refutes universal character-only reduction of the source
parameters. Its scope is a decomposable balanced relation; it does not
decide a preselected indecomposable switch relation or compare isolated
numerical L-values.

\subsection{Route 2 and the Hodge
conjecture}\label{IMPLICATIONS-route-2-and-the-hodge-conjecture}

B6 constructs the Weil projector using only a fixed generic imaginary
quadratic action, with an explicit resultant denominator. B7 bounds the
projected component of an effective cycle by its degree in a positive,
flat Weil metric. B8 turns any nonzero effective Weil witness into a
representative of every fixed Weil class, with both denominator and
signed effective degree bounded quadratically in the witness degree.
DIV1 supplies the proper parameter-space argument; DIV2 carries a
uniform witness bound on a dense set to every fibre.

T6.16 extends the source-parameter counterexample to an actual 032
witness for four distinct labeled CM types satisfying an indecomposable
switch. Prescribed mixed real discrete series, fixed-input continuity, a
finite Hecke projector and sequential extraction preserve both a nonzero
four-map period and a degree term with Steinberg monodromy. This
resolves the existential distinct-switch question left open by T6.15. It
does not claim the same occurrence in every prescribed relation or
identify numerical special values with geometric Galois constituents.

T6.17 gives a positive analytic conclusion for this actual fixed witness
and its four CM isogeny classes. Short integral isogenies control each
polarization eigenvalue, including for imprimitive CM types. Further
actual integer multiplications balance the four factors. A single common
metric-period normalization then bounds the energy, the ratio of raw
degree to metric period, and all coefficients in fixed finite theta
presentations. This removes the fixed-source analytic normalization gate
in that setting. The metric period and the algebraically normalized
Betti coefficient are different data; the construction does not bound
integral cycle degree or the denominator of the Betti division.

T6.18 computes the actual integral Betti coefficient eigenlines. Their
full pullback image is the core/hull ideal, with additive index J\_i
H\_i in each factor, where H\_i is the maximal-order hull index; a
compatible principal polarization gives H\_i=J\_i. The specified
four-factor period ideal therefore has index product(J\_i H\_i), and
product(J\_i squared) in the principal case. The final period remains in
the fixed Galois CM coefficient field. This closes the missing
eigenline-normalization calculation behind T6.17. It does not infer an
integer trace denominator from ideal index alone: a fully computed
Gaussian-prime example shows why that inference fails.

This is a strict reduction of the denominator and projector problem. It
is not a proof of uniformly bounded witnesses in the open Weil families.
No new general Hodge-conjecture case is asserted from the presently
established estimates. The supplied Markman theorem yields a posteriori
bounded complexity for abelian fourfold families by DIV5. That existence
argument gives no numerical constant or explicit theta representative.
B15 gives a different unconditional quantitative result: from one fixed
representative, an O-linear polarized isogeny supplies a bound linear in
its degree, even if it rotates the Weil plane. When it transports the
prescribed flat class by a scalar, the bound improves to the square root
of the isogeny degree. Neither is a bound independent of the isogeny.

B19 removes an additional comparison in a fixed simple CM isogeny class.
The positive metric covolume of the integral Hom lattice is exactly
proportional to the index J of the maximal-order core of the target
lattice. Its smallest nonzero squared norm has order J\^{}(1/g). A small
genuine isogeny therefore yields effective witness degree O(sqrt J) and
signed complexity O(J), with explicit fixed-field constants and a
conductor-norm upper bound. On a fixed source, normalizing the generic
Weil period gives uniformly bounded analytic energy and finite theta
coordinates; that real normalization changes the integral map lattice.
B20 proves that taking reduced images of a fixed finite list of
reference witnesses does not remove the index: finite stabilizers leave
their degrees bounded below by a positive constant times sqrt J. B21
strengthens this from a finite list to the entire family of
bounded-degree source subvarieties: a universal Hilbert stabilizer is
finite, giving a uniform stabilizer-order bound. B22 makes the bound
explicit by intersecting at most n+1 translates and applying the
generalized Bezout inequality. A uniformly bounded target degree would
require source degree at least c J\^{}(1/(2(n+1))); for surfaces the
exponent is1/6. Thus even target-dependent selection among source
witnesses of bounded degree cannot remove square-root index growth.
These results identify the arithmetic cost of the fixed-reference
construction while leaving unrelated target-dependent bounded cycles
open.

B22a strengthens that necessary source-degree growth to J\^{}(1/4) for
middle-dimensional subvarieties of a simple source. The difference map
has positive degree equal to the self-intersection; quotienting by the
full stabilizer leaves a positive even integral square. A self-dual
middle-form comass estimate then gives \textbar Stab Z\textbar{} at most
D squared divided by 4(n!) squared V0. The sharper explicit integer
bound also reduces the lcm needed for reduced images below.

B23 eliminates cancellation among arbitrarily many integral reference
transports as a way to remove the arithmetic index. In the generic Weil
determinant lattice the sum of all isogeny pushforward images is exactly
the conjugate determinant-core ideal times the target lattice, of index
J, including ramified primes. A relation Nt made from arbitrary integral
pushes must satisfy J dividing R squared N squared c\_t. With a bounded
source degree, the same assertion holds for reduced images after
multiplying N by the explicit lcm of all possible mapping
multiplicities. A fixed nonmaximal quadratic order costs only a
controlled conductor factor. This is an obstruction to the specified
construction, not to arbitrary target cycles; reduced-image
constructions with unbounded stabilizers remain to be attacked.

B24 rules out a bounded-index escape on an infinite fixed-field,
fixed-type, fixed-volume polarized CM set. It counts bounded-index
targets using ideal classes, bounded-index overlattices and polarization
image lattices modulo real-unit squares. Such an infinite set
necessarily has unbounded J. No Hecke-density theorem is asserted.
Together B23 and B24 identify a common integral-content obstruction
throughout every infinite fixed-data target collection, beyond an
estimate for a single preferred isogeny.

T6.19 computes both scalar trace content and the full unscaled
tensor-presentation index on actual principally polarized CM families.
The scalar formula uses the minimum local valuation; the full tensor
formula uses a maximum. T6.20 then removes a possible
overinterpretation: all input-isogeny factors cancel exactly in the
normalized Chow class at fixed output, by birigidification and the
projection formula. Fixed integral input eigenvectors give a common Chow
denominator N0 equal to the absolute norm of their nonzero period, for
every factorwise integral output tensor.

T6.21 identifies the remaining primitive-lattice comparison exactly. If
M is the additive product of the four output homology lattices, with the
prescribed conjugations, the integral orbit is its trace dual and the
factorwise trace span is the inverse maximal-order hull. Their index is
\textbar Disc(E)\textbar{[}Mplus:M{]}. If one output factor is stable
under the maximal order, M is already an ideal and the saturation
exponent is the fixed different exponent d\_E. Thus N0*d\_E represents
every integral orbit class by an integral algebraic cycle uniformly over
the other three lattices. When all four factors vary, actual principal
families have saturation exponent p\^{}r. That is a limitation of the
factorized construction, not an intrinsic obstruction to other algebraic
cycles. A comparison with the generic quadratic Weil plane still
requires the actual integral correspondence, and no signed-cycle degree
bound is inferred from this division result.

B25 now treats the unbounded-stabilizer escape left by B23. Saturating
the kernel intersection with the source stabilizer under the quadratic
order produces an intermediate quotient on which the residual map is
birational on the image. Bounded target degree forces bounded residual
degree and an intermediate core index at least a fixed positive multiple
of the target core index. If this saturated kernel is stable under the
full CM order, the intermediate core is one and the original obstruction
remains. B25a shows the alternative is real when simplicity is dropped:
explicit product CM fourfolds have degree-two source and target surfaces
but target core m\^{}4, with multiplicity m\^{}2 removed by the
quotient. The quotient has already imported the target arithmetic
complexity.

B26 imports the known principal-angle Weil calibration and equality
rigidity from rStation4/v1, Theorem1(b) in Section2 and Theorem3(a),(c)
in Section4. Its new contribution here is the exact q/V normalization
and integral gcd margin. It improves the residual-degree constants in
B25 and makes the simple-ambient inequality strict, without changing the
growth exponent or constructing new witnesses.

\subsection{Semiregularity and
deformation}\label{IMPLICATIONS-semiregularity-and-deformation}

The 032 paper supplies a morphism with two-dimensional image, not a
smooth embedded surface. S-01 shows that its permitted covers do not
even preserve injectivity of that morphism. S-02 is the exact
conormal-cohomology test for an actual smooth embedding. S-03 and S-11
refute general shortcuts from a nonzero period or from an ample complete
intersection to semiregularity.

For the rigid ball-quotient source, S-04 forces the generated quotient
to remain constant in any embedded semiregular deformation of its entire
class. S-05--S-06 refute universal semiregularity of all deeper-level
morphism presentations even for the full generalized map: the
obstruction space grows beyond its fixed Hodge target. S-07 removes that
cover growth by equivariant descent but returns to the base-level
obstruction.

S-09 produces a perfect-complex K-class whose surviving Chern character
is the desired projected class plus a chosen polarization polynomial. It
removes extraneous Hodge conditions at the K-class level. It does not
construct a semiregular object. S-10 explains why an equivalence alone
does not remove an existing full trace kernel. S-12 shows that
preserving a pure Weil character without polarization data can force a
contradiction with Voisin's nonprojective torus family for coherent
sheaves. S-13 excludes a single exponential repair in that precise
coherent fourfold setting. S-14, using Perry Theorem1.2, reduces the
construction to one gluable weakly target-equivariant perfect object
whose full B-normalized character is P(theta)+lambda t. It already
implies bounded division. S-15 shows that canonical source-cover lifts
do not provide such a repair, and S-16 separates bundle realization from
trace injectivity. DER1 gives a strict discriminant constraint for a
Hermite--Einstein attempt; DER2--DER3 show that the canonical globally
generated syzygy can pass that constraint while preserving a full trace
kernel. DER5 removes their positivity restrictions: nonzero-Euler
convolution and nonzero-virtual-rank tensor functors are faithful on all
morphisms, so the ideal-diagonal evaluation cone preserves the kernel
for every perfect object on a positive-dimensional abelian variety. DER6
shows that a polarization-polynomial tensor correction either has
nonzero virtual rank and retains the kernel, or has rank zero and
removes the desired Weil component.

DER4 makes the successful-object implication constructive without
requiring nonzero virtual rank. An Adams operation whose degree equals
the band order of the Kummer gerbe descends to ordinary perfect K-theory
and multiplies the Weil component by M\^{}n.~A relative bundle
presentation, a generation twist, and a Chern degeneracy locus yield
fixed explicit N and D once those finite choices are supplied. DIV1
carries that very bound from the local etale image to every fibre. This
step requires no proper moduli space of complexes and no Baire-category
extraction of an unspecified constant. The remaining geometric problem
is existence of the original weakly equivariant semiregular object with
the complete compatible character. DIV3 is the separate embedded-cycle
implication.

DER7 further makes the multiplier independent of the gerbe order when
the virtual rank r is nonzero. The Adams operation of degree
\textbar r\textbar, tensored by the appropriate determinant power, has
exactly zero inertia weight and retains \textbar r\textbar\^{}n times
the Weil coefficient. Integrality of the resulting Chern-cycle degree
gives N*=v\textbar b beta0\textbar\textbar r\textbar\^{}n without
estimating the denominators of the whole polarization polynomial. Its
common degree bound still depends on the finite generation choices. S-17
is the tensor--determinant alternative, and S-18 excludes canonical
categorical induction as a way to remove an existing full kernel. These
operations are used at different stages: untwisting happens after the
successful deformation and is not claimed to preserve semiregularity.

S-19 and S-20 quantify the failure of the actual genus-four
isolated-point and partial-normalization secant seeds. Their qualitative
failure and the CY4 adjunction criterion were already recorded in
rStation paper37/v4; the present bounds are respectively ext2 at
least24d+6 and full-kernel dimension at least4d-3. S-21 extends the
obstruction from the prior free-translation setting to every effective
coherent finite translation/tensor action. Vanishing of nonidentity
Hochschild coefficient sectors, induction/forgetting and categorical
Riemann--Roch force an equivariantly semiregular imaginary-secant seed
to have endomorphism dimension strictly larger than2. Equivariantly
simple seeds are therefore excluded even when the underlying object is
nonsimple. The new argument does not impose a free-translation
Fourier--Mukai model.

S-22 tests the balanced real-multiplication character which passes the
earlier numerical secant test. On a genuine real-quadratic eigenspace
decomposition, positivity would force an effective surface in its
prescribed second-Chern class to have an irrational constant tangent
space; its Albanese map rules that out. The character thus has no
rank-one torsion-free realization. The degree-one and degree-two
contraction ranks are8 and18, and a simple weakly equivariant
semiregular lift must have one of two precisely specified finite-group
orders. S-23 provides a necessary control: O\_A plus O\_p is fully
semiregular, with endomorphism dimension3 and a nilpotent. It prevents
promoting the simple-seed exclusion to all nonsimple objects. S-24 does
exclude arbitrary iterated self-extensions of a fixed
scalar-endomorphism imaginary-secant factor of Euler characteristic at
least2; a heart filtration bounds End by its squared length,
contradicting the semiregularity rank equality. Different constituents,
different compatible characters, and genuinely new derived constructions
remain open.

S-25 removes a further concrete candidate family: an external product
containing a simple gluable ample imaginary-secant factor cannot become
fully semiregular under any effective finite coherent translation/tensor
action, even if that action correlates the factors and even after any
derived equivalence. Framing the factor actions at their
scalar-endomorphism objects supplies coherence; a derived-fibre
cancellation proves that every stabilizing product functor stabilizes
each factor without a shift. The nonzero factor kernel therefore remains
invariant for the full correlated group. This applies to the actual
genus-four secant-square construction. It leaves genuinely nonsplit
real-multiplication derived objects and nonsimple constructions as
concrete next tests; it asserts no exclusion for all perfect objects.

S-26 reduces the balanced real-multiplication simple-object test to an
exact tangent dimension: equivariant Euler characteristic4 requires Ext1
dimension8, while Euler characteristic2 requires dimension9. It also
excludes the usual amplitude{[}0,1{]} rank-one form with torsion-free H0
and higher cohomology in codimension at least two. S-27 supplies an
actual conditional construction template in amplitude{[}0,2{]}: a
nonzero Ext3(T,A) class between two scalar-endomorphism sheaves with the
specified lower cross-Ext vanishing produces a simple gluable cone. Its
projective extension parameters inject into Ext1, so in the principal
required-character case the extension space has dimension at most9 and
full semiregularity is exactly Ext1 dimension8. No pair of sheaves with
the required nonsplit real-multiplication character has yet been
constructed.

S-28 identifies the two-sheaf cone tangent space by an exact Yoneda
compatibility sequence. S-29 computes that sequence for transverse
smooth surface sheaves: it gives a unique-extension/evaluation test and
proves that acyclic positive twists retain a kernel of dimension at
least d-1. These remain valid conditional tools, but the previously
proposed principal rank-one balanced character is now excluded
arithmetically.

S-30 closes the entire nonequivariant simple gluable rank-one balanced
real-multiplication branch, including the non-pure two-sheaf cone
attempt. The proof uses the actual integral cohomology lattice:
integrality of aC forces the squarefree real-quadratic radicand to
divide a squared times the polarization volume. Exterior-square Smith
exponents prove this even at two. The required Euler pairing is then too
large for the fixed rank-eighteen full semiregularity map. The principal
q=1 character is not the character of any rank-one perfect object, with
no simplicity assumption. In higher rank the proof gives exact
rank/group divisibility and rules out every nonequivariant simple
gluable realization with an integral B-normalizer. S-31 independently
excludes vector-bundle ranks one, two and three by Newton identities.
Fractional determinant normalization and a genuinely new higher-rank or
nonsimple construction are therefore necessary in that route.

S-32 strengthens the higher-rank test at two. With maximal local real
multiplication and the minimal rational divisor space, an odd
polarization volume and a two-adic unit Nq make the second Chern class
nonintegral. Every equivariantly simple gluable fully semiregular
balanced object in this setting must therefore have even D times the
group order. The proof includes the local trace-dual symplectic normal
form, and applies to all perfect objects rather than only bundles. Its
principal nonequivariant corollary excludes every positive rank when the
real quadratic field is ramified at two.

S-33 supplies the odd-prime counterpart, and its proof works at every
odd ramified prime where the polarization is locally unimodular and the
real multiplication order is maximal. The second-Chern congruences force
the integer Euler characteristic to have negative valuation; an
independent third-Chern calculation confirms the same obstruction,
including at3. Combined with S-32, every prime ramified in the real
quadratic field must divide D times the effective group order.
Consequently every nonequivariant positive rank is excluded on the
principally polarized maximal-order locus with minimal rational divisor
space. The remaining hypotheses that can fail are now concrete: local
maximal-order action, unramified polarization volume, a determinant
lying in the RM divisor plane, or absence of a sufficiently ramified
effective group. Nonsimple equivariant objects remain separately outside
the simple-lift Euler argument.

T6.22 extends the construction obstruction beyond factorwise pure
kernels. On the actual indecomposable principal-shear family, arbitrary
integral divisor polynomials on a fixed CM source times the outputs have
the same local orbit content, even for source cycles of unbounded
degree. It also prevents returning a maximal-order auxiliary factor by a
joint integral cohomological correspondence of bounded denominator on
any proper block of one, two or three output factors, even when the
class on the anchored product is an arbitrary integral Hodge class. The
proof uses pure Hodge rows and the uniqueness of the allowed vertex in a
conjugate-coordinate sign cube. Genuinely joint all-four
correspondences, orbit-mixing maps with specified integral content, and
intrinsic Chow divisibility remain unexcluded.

S-34 removes the rank-one equivariant exception entirely. The
determinant of a rank-one perfect object carries exactly the same action
coherence, so tensoring by its inverse turns the whole
translation--tensor action into an ordinary free translation action.
Perfect descent produces a rank-one object on the quotient abelian
variety with the same normalized real-multiplication character after
rational polarization rescaling. S-30 then excludes it. The theorem
requires simplicity and gluability only in the equivariant category and
imposes no maximal-order or minimal-divisor-space hypothesis.
Higher-rank or nonsimple lifts are now necessary for this
balanced-character route.

\subsection{GRH hypotheses}\label{IMPLICATIONS-grh-hypotheses}

The task's Alpöge Dirichlet GRH theorem is admitted, with the supplied
Lean declarations as its exact formal statement. Consequently the cited
conditional Dirichlet edge estimates may be used without a further GRH
assumption in this run. The full Lean dependency archive was not
supplied or compiled.

General Hecke edge estimates are conditional on the Masullo claim
identified by the prompt as §1.2. The supplied index does not contain
the complete claim or proof and does not identify a numbered Hecke-GRH
theorem. Its numbered trace formulas have not been substituted for such
a statement. No Hecke GRH hypothesis is needed for the absolutely
convergent 3/2 factor. A nonabelian rank-two standard L-function is
outside the stated Dirichlet/Hecke scope unless a factorization is
established.

\subsection{Tate and standard
conjectures}\label{IMPLICATIONS-tate-and-standard-conjectures}

The supplied 032 and Milne consequences remain the background theorems
requested by the task: CM Hodge, Tate for abelian varieties over finite
fields in the cited formulation, and the Hodge standard conjecture for
abelian varieties. The new degree estimates do not enlarge their scope.
The correspondences in B6 exist algebraically and define projectors in
homological realizations and numerical motives. B16 adds the canonical
Deninger--Murre degree projector and Moonen's endomorphism polynomial
law to obtain an exact idempotent modulo rational equivalence, with
denominator dividing R times the explicit interpolation denominator.
This upgrade uses a motivic theorem; it is not inferred from
cohomological eigenvalues. It still does not identify every Weil
realization with an algebraic cycle. A positive cohomological norm does
not by itself construct an effective cycle or control an integral
saturation index.

\subsection{Grothendieck locators and the previous
runs}\label{IMPLICATIONS-grothendieck-locators-and-the-previous-runs}

\begin{itemize}
\tightlist
\item
  \textbf{FGA, Bourbaki exposé 221, Lemme 2.4 p.255 and Théorème 3.2
  p.260:} bounded reduced subvarieties and projective
  fixed-Hilbert-polynomial parameter spaces give DIV1--DIV3 their
  precise algebraic-geometric mechanism.
\item
  \textbf{SGA 7 I, Exposé VII, §2.4 pp164--165, Example2.9.5 p175 and
  Remark2.9.6 p176:} the original scans have been visually verified.
  Additivity in each argument, and the equivalence with birigidified
  line bundles, give the exact cancellation T6.5. Equation2.4.2 itself
  displays additivity in the coefficient group; the surrounding text
  states all three arguments.
\item
  \textbf{SGA6, ExposéV, Appendix7.1 pp360--361 and Proposition7.5
  pp362--364:} weighted Newton polynomials and multiplicative Adams
  operations give DER4 its exact untwisting mechanism. Original scans
  PDF367--371 were read, including the whole proof of Proposition7.5.
\item
  \textbf{SGA6, Exposé0 §1, equations(1.1)--(1.3), p1:} the
  Chern-character and relative tangent-class formulation places S-09,
  S-16 and DER1--DER4 in the K-theoretic intersection framework. The
  original scan PDF8 was read. The later precise semiregularity maps are
  supplied by Bloch, Buchweitz--Flenner and Perry; no handwritten
  positivity statement is substituted for them.
\item
  \textbf{Cote 16, the Coates letter, PDF page 45 / archive footer 44,
  and PDF page 46 / footer 45:} the displayed equivalence involving
  algebraic Lefschetz inverse and cycle-rank symmetry is consistent with
  retaining the algebraic projector as a separate datum. The current
  results do not infer missing cycles from positivity alone.
\item
  \textbf{Cote 44, PDF page 2 / archive footer 1, Th.4; PDF page 4 /
  footer 3, Th.5:} polarization and the intersection/height form remain
  distinct from the unnormalized degree of the 032 surface. No new claim
  about the arithmetic height formula is made in this run.
\end{itemize}

The sketch uses PDF-page numbering, including the cover, for several of
these scan references. We record both PDF and archive-footer numbers;
this is a locator convention, not a mathematical transcription
correction.

The rows directly treated are \texttt{HODGE\_BRIDGE\_SKETCH\_20261121}
T6, T7, Route 2 §2.1, and the generic-projector/witness part of §2.4. T8
is affected only by the reduction DIV2; its higher-dimensional cases
remain open here. Other runs named SA06, SA07, H7, H8 and H9 were not
supplied beyond quotations in the sketch; no unseen row is declared
proved or changed.

\section{Per-finding Grothendieck framework placement}

The index \texttt{FRAMEWORK\_LOCATORS.tsv} has one row for every
numbered finding, including the explicit B19 substatements. The grouped
descriptions below enumerate the same findings. A locator is a framework
placement unless the proof expressly invokes a result on that page.
Later analytic, categorical and number-theoretic theorems are attributed
to their actual authors; they are not retroactively attributed to
Grothendieck. Provisional identifiers, if present in this index, acquire
no proved status from being indexed.

The archival scan convention includes the cover; the corresponding
archive footer is given as well. No mathematical transcription error has
been established. In particular the Tate and standard-conjecture
consequences admitted in the prompt are not enlarged by the analytic
estimates or by these framework assignments.

\subsection{Polarized cycle degree and finite
maps}\label{FRAMEWORK-polarized-cycle-degree-and-finite-maps}

\textbf{Findings:} T6.1, T6.2, T6.3, T6.4, T6.6, B3, B4, B5.

\textbf{Exact placement:} Cote44, PDF scan2 / archive footer1, Theorem4;
PDF scan4 / footer3, Theorem5; FGA, Bourbaki exposé221, Lemme2.4 p255
and Théorème3.2 p260.

Intersection with a polarization in homological and numerical cycle
groups; proper pushforward retains generic multiplicity. Analytic norms
are evaluated in the Betti realization, not treated as integral cycle
classes.

\textbf{Later inputs:} Projection formula, surface Hodge index theorem,
supplied032 §7; exact normalization proved in these findings.

\subsection{Correspondences, Picard duality and
biextensions}\label{FRAMEWORK-correspondences-picard-duality-and-biextensions}

\textbf{Findings:} T6.5, T6.7, T6.8, T6.9, T6.20, DER5.

\textbf{Exact placement:} SGA7 I, ExposéVII, §2.4 pp164--165,
Example2.9.5 p175 and Remark2.9.6 p176; Cote16, PDF scans45--46 /
archive footers44--45, Coates letter on algebraic Lefschetz inverse and
cycle-rank symmetry.

Divisor kernels are algebraic correspondences from birigidified line
bundles; norm/trace and normalized correspondence cancellation take
place in cycle groups. The Albanese and Fourier--Mukai comparisons
require the separately cited later results.

\textbf{Later inputs:} Supplied032 Lemma2.2/Proposition3.1; BMM
Theorem9.1; Mukai Theorem2.2 and equation3.7; Perry trace naturality.

\subsection{Theta norms and their arithmetic
normalization}\label{FRAMEWORK-theta-norms-and-their-arithmetic-normalization}

\textbf{Findings:} T6.10, T6.11, T6.12, T6.13, T6.14, T6.15, T6.16,
T6.17, T6.18, T6.19, T6.21, T6.22, B1, B1a, B2, B14, B14a, B17, LN7.

\textbf{Exact placement:} Cote44, PDF scan2 / archive footer1, Theorem4;
PDF scan4 / footer3, Theorem5; SGA7 I, ExposéVII, §2.4 pp164--165,
Example2.9.5 p175 and Remark2.9.6 p176.

The theta identities calculate polarization pairings in Betti
realization and constrain integral correspondence inputs. These are
framework placements only: no Rallis formula, GRH estimate, automorphic
classification, or CM eigenline-content theorem is attributed to these
Grothendieck pages.

\textbf{Later inputs:} Haan Theorem3.4; Gan--Qiu--Takeda Theorem11.4;
BMM; Shin4.8/5.13; Labesse5.3/5.9; Alpöge supplied declarations; Masullo
conditional §1.2; exact locators in the findings.

\subsection{Weil projectors and integral realization
lattices}\label{FRAMEWORK-weil-projectors-and-integral-realization-lattices}

\textbf{Findings:} B6, B7, B8, B9, B16, B18, B23, B26.

\textbf{Exact placement:} Cote16, PDF scans45--46 / archive
footers44--45, Coates letter on algebraic Lefschetz inverse and
cycle-rank symmetry; SGA6, Exposé0, §1, equations(1.1)--(1.3), p1 (PDF
scan8).

The generic quadratic action defines correspondences on the middle
homological motive; B16 gives an exact Chow idempotent using a later
motivic theorem. Integral Betti intersection lattices and their positive
middle pairing control division. The algebraic correspondences also
admit étale realizations after spreading out; no new integral Tate or
positivity theorem is inferred.

\textbf{Later inputs:} Moonen §2.1, Proposition2.5 and Corollary2.6;
Deninger--Murre degree projectors as imported there; MilneANT4.1--4.2
for B23.

\subsection{CM isogeny transport and bounded source
geometry}\label{FRAMEWORK-cm-isogeny-transport-and-bounded-source-geometry}

\textbf{Findings:} B10, B11, B12, B13, B15, B19, B20, B21, B22, B22a,
B24, B25, B25a.

\textbf{Exact placement:} Cote44, PDF scan2 / archive footer1, Theorem4;
PDF scan4 / footer3, Theorem5; SGA7 I, ExposéVII, §2.4 pp164--165,
Example2.9.5 p175 and Remark2.9.6 p176; FGA, Bourbaki exposé221,
Lemme2.4 p255 and Théorème3.2 p260.

Integral CM Hom modules define genuine isogeny correspondences; their
determinant indices and polarization metrics remain distinct. Hilbert
boundedness and stabilizers control reduced-image multiplicities. B24
counts polarized CM points with fixed field and type; this is not a
theorem about integral Shimura models.

\textbf{Later inputs:} MilneANT4.19/4.26 and4.4/5.1; Ullmo--Yafaev
Proposition4.1(1); Fulton Examples8.4.5--8.4.7; Stacks02LS; B22a has its
separate sum-dimension source.

\subsection{Embedded and map-level deformation; bounded
division}\label{FRAMEWORK-embedded-and-map-level-deformation-bounded-division}

\textbf{Findings:} S-01, S-02, S-03, S-04, S-05, S-06, S-07, S-11, DIV1,
DIV2, DIV3, DIV4, DIV5.

\textbf{Exact placement:} FGA, Bourbaki exposé221, Lemme2.4 p255 and
Théorème3.2 p260; SGA6, Exposé0, §1, equations(1.1)--(1.3), p1 (PDF
scan8); Cote16, PDF scans45--46 / archive footers44--45, Coates letter
on algebraic Lefschetz inverse and cycle-rank symmetry.

Hilbert spaces parameterize bounded integral cycles, and their proper
images propagate a fixed multiplier and signed degree. A morphism,
reduced image, embedding and perfect pushforward are different
deformation objects. Full Hodge compatibility is retained before
invoking a deformation theorem.

\textbf{Later inputs:} Bloch semiregularity; Buchweitz--Flenner
Theorem5.2 and Proposition8.2; Stacks068X; Fulton18.3/18.3.11;
suppliedMarkman Theorem1.5.1 and Corollary1.6.1.

\subsection{Perfect objects, Chern characters and full trace
kernels}\label{FRAMEWORK-perfect-objects-chern-characters-and-full-trace-kernels}

\textbf{Findings:} S-08, S-09, S-10, S-12, S-13, S-14, S-15, S-16, S-18,
S-19, S-20, S-21, S-22, S-23, S-24, S-25, S-26, S-27, S-28, S-29, S-30,
S-31, S-32, S-33, S-34, DER1, DER2, DER3, DER6.

\textbf{Exact placement:} SGA6, Exposé0, §1, equations(1.1)--(1.3), p1
(PDF scan8); SGA7 I, ExposéVII, §2.4 pp164--165, Example2.9.5 p175 and
Remark2.9.6 p176.

Perfect K-theory and Chern character provide the algebraic intersection
framework. Hochschild evaluation, equivariant categories and full
semiregularity are later refinements, not statements located in these
SGA pages. A numerical K-class is never identified with a semiregular
object without proof.

\textbf{Later inputs:} Buchweitz--Flenner; Perry2026 Theorem1.2 and
trace naturality; Perry2018 Theorem4.4; Bayer--Perry3.3/3.4/3.15;
Shklyarov3.2--3.5; Orlov2.10/2.13; prior rStation37/v4 criterion
explicitly attributed.

\subsection{Integral untwisting and effective Chern
cycles}\label{FRAMEWORK-integral-untwisting-and-effective-chern-cycles}

\textbf{Findings:} S-17, DER4, DER7.

\textbf{Exact placement:} SGA6, ExposéV, Appendix7.1 pp360--361 and
Proposition7.5 pp362--364 (PDF scans367--371); SGA6, Exposé0, §1,
equations(1.1)--(1.3), p1 (PDF scan8); FGA, Bourbaki exposé221, Lemme2.4
p255 and Théorème3.2 p260.

Homogeneous Adams operations and determinant weights descend an integral
K-class from a finite gerbe after successful deformation. Effective
Chern degeneracy cycles and projective parameter spaces then give fixed
signed-cycle bounds. These operations do not construct the initial
semiregular object.

\textbf{Later inputs:} Knudsen--Mumford Definition4 pp40--42 and
Theorem2 p42; Perry2026 Kummer-gerbe construction; Stacks0F8F; the
generation and fixed-multiplier proofs in DER4/DER7.


\end{document}
